% Sensibilisation sur la nécessité de préserver le fonctionnement des neurones — SVT Tle D — Cours signé OnBuch+
% Programme officiel MINESEC. Compiler avec : tectonic N08-fonctionnement-des-neurones.tex
\newcommand{\DOCMATIERE}{SVT}
\newcommand{\DOCNIVEAU}{Tle D}
\newcommand{\DOCMODULE}{Module I : Le Monde Vivant — Fonctionnement des neurones}
\newcommand{\DOCLECON}{N08}
\newcommand{\DOCTITRE}{Sensibilisation sur la nécessité de préserver le fonctionnement des neurones}
% OnBuch+ — préambule « cours signés » ENRICHI (compilé avec tectonic / XeLaTeX)
% Système visuel « Pop » d'OnBuch+ : fond crème (jamais blanc), bordures encre
% épaisses, ombres « dures » décalées, titres Archivo Black en capitales.
% Version MATHÉMATIQUES : graphes (pgfplots/tikz), boîtes pédagogiques
% (définition, propriété, méthode, attention, le savais-tu, exercice résolu,
% exercice, TP, bilan), page de garde et signature OnBuch+ ancrée en bas.
%
% Macros attendues dans main.tex AVANT \input{preamble} :
%   \DOCMATIERE  \DOCNIVEAU  \DOCMODULE  \DOCLECON (numéro)  \DOCTITRE
\documentclass[11pt]{article}

\usepackage{fontspec}
\usepackage[french]{babel}
\usepackage[a4paper,margin=2cm,top=3.4cm,bottom=2.8cm,headheight=1.4cm,footskip=1.3cm]{geometry}
\usepackage{amsmath,amssymb,amsfonts}
\usepackage{mathtools} % \xmapsto, \coloneqq... souvent utilisés par les modèles
\usepackage{cancel} % simplifications barrées (bilans de forces)
% \abs{z} (module, valeur absolue) et \norm{u} (norme), utilisés par les modèles
\providecommand{\abs}{}\let\abs\relax\DeclarePairedDelimiter{\abs}{\lvert}{\rvert}
\providecommand{\norm}{}\let\norm\relax\DeclarePairedDelimiter{\norm}{\lVert}{\rVert}
\usepackage[table]{xcolor} % [table] : fournit aussi \rowcolors (alternance de lignes)
\usepackage{pagecolor}
\usepackage{tikz}
\usetikzlibrary{arrows.meta,positioning,calc,shapes.geometric,shapes.misc,shapes.arrows,decorations.pathmorphing,decorations.pathreplacing,decorations.markings,patterns,shadows,mindmap,trees,fit,backgrounds,matrix,intersections,angles,quotes}
\usepackage{pgfplots}
\usepackage[version=4]{mhchem} % équations nucléaires \ce{^{235}_{92}U}
\usepackage[european,siunitx]{circuitikz} % schémas électriques
\pgfplotsset{compat=1.18}
\usepgfplotslibrary{fillbetween,groupplots} % aires entre courbes (calcul intégral)
\usepackage{enumitem}
\usepackage{fancyhdr}
% [most] charge toutes les bibliothèques tcolorbox (listings, minted, raster,
% documentation...) et gonfle inutilement la mémoire TeX sur un document long
% (~300 boîtes) : on ne charge que ce qui est réellement utilisé.
\usepackage[skins,breakable]{tcolorbox}
\usepackage{graphicx}
\usepackage{colortbl}
\usepackage{booktabs}
\usepackage{tabularx}
\usepackage{diagbox} % case d'en-tête diagonale des tableaux à double entrée
% Colonnes extensibles centrée / gauche / droite (C, L, R), souvent utilisées par les modèles
\newcolumntype{C}{>{\centering\arraybackslash}X}
\newcolumntype{L}{>{\raggedright\arraybackslash}X}
\newcolumntype{R}{>{\raggedleft\arraybackslash}X}
\usepackage{array}
\usepackage{multicol}
\usepackage{siunitx}
% parse-numbers=false : accepte aussi \SI{2,5 \times 10^{-3}}{...}
\sisetup{locale=FR,output-decimal-marker={,},group-minimum-digits=4,parse-numbers=false}
\DeclareSIUnit\ohm{\ensuremath{\symup{\Omega}}} % U+2126 absent de la police
\DeclareSIUnit\division{div} % réglages d'oscilloscope (ms/div, V/div)
\DeclareSIUnit\tr{tr} % tours (vitesse de rotation, ex. \SI{1500}{\tr\per\minute})
\DeclareSIUnit\molar{M} % concentration molaire (ex. \SI{3}{\molar}, \SI{1}{\micro\molar})
\DeclareSIUnit\an{an} % année (le modèle confond parfois avec \year, un compteur TeX) — ex. \SI{5}{\kilo\meter\per\an}
\DeclareSIUnit\FCFA{FCFA} % franc CFA (ex. \SI{500}{\FCFA\per\kilo\gram})
\DeclareSIUnit\degreeBrix{\textdegree Brix} % teneur en sucre d'un sirop/jus (réfractométrie)
\DeclareSIUnit\month{mois} % le modèle utilise parfois \month, un compteur TeX, comme unité de durée
\usepackage{qrcode}
% \degree : commande fréquemment utilisée par les modèles pour un angle ou une
% température en dehors de \SI{}{} (non fournie par les paquets chargés).
\providecommand{\degree}{^{\circ}}
% \qed (fin de démonstration), utilisé par les modèles sans amsthm
\providecommand{\qed}{\hfill$\square$}
% \barre : double barre des valeurs interdites dans un tableau de variations (tkz-tab)
\providecommand{\barre}{\ensuremath{\|}}
% environnement proof (amsthm non chargé) utilisé par les modèles
\ifdefined\proof\else\newenvironment{proof}[1][Démonstration]{\par\noindent\textit{#1.}\ }{\qed\par}\fi
\usepackage{lastpage}
\usepackage{hyperref}
\hypersetup{hidelinks}

% ============================================================
% Palette « Pop » OnBuch+ — mêmes hex que la classe P (plus_ui.dart)
% ============================================================
\definecolor{poporange}{HTML}{F59321}
\definecolor{popdark}{HTML}{E07A0C}
\definecolor{popcream}{HTML}{FFF6E8}
\definecolor{popink}{HTML}{1C1714}
\definecolor{poppurple}{HTML}{7A5AE0}
\definecolor{popgreen}{HTML}{2E9E6A}
\definecolor{poppink}{HTML}{E0457B}
\definecolor{popblue}{HTML}{2F7DE1}
\definecolor{popgold}{HTML}{C98A12}
\definecolor{popmuted}{HTML}{8A7F76}
\definecolor{popline}{HTML}{EADFD2}
\definecolor{popgray}{HTML}{8A7F76}
\colorlet{popgrey}{popgray}
% Alias tolérants (le modèle nomme parfois les couleurs différemment)
\colorlet{popwhite}{white}
\colorlet{popblack}{popink}
\colorlet{popred}{poppink}
\colorlet{popyellow}{popgold}
% Teintes claires (fonds de tableaux/encadrés) — le modèle nomme parfois
% les couleurs avec un suffixe L (ex. popblueL) sans les avoir définies.
\colorlet{poporangeL}{poporange!20}
\colorlet{popdarkL}{popdark!20}
\colorlet{poppurpleL}{poppurple!20}
\colorlet{popgreenL}{popgreen!20}
\colorlet{poppinkL}{poppink!20}
\colorlet{popblueL}{popblue!20}
\colorlet{popgoldL}{popgold!20}
\colorlet{popredL}{popred!20}
\colorlet{popgrayL}{popgray!20}
% Teintes claires pour fonds de tableaux / zones
\colorlet{popblueL}{popblue!10!white}
\colorlet{poporangeL}{poporange!14!white}
\colorlet{popgreenL}{popgreen!12!white}
\colorlet{poppurpleL}{poppurple!12!white}
\colorlet{poppinkL}{poppink!12!white}
\colorlet{popgoldL}{popgold!14!white}

\pagecolor{popcream}

% ============================================================
% Typographie
% ============================================================
\defaultfontfeatures{Ligatures=TeX}
\setmainfont{PlusJakartaSans-Medium.ttf}[Path=../../../fonts/,BoldFont=PlusJakartaSans-Bold.ttf]
% Maths en sans-serif (Fira Math) avec chiffres et lettres droites de Plus
% Jakarta, pour que les formules se fondent dans le texte.
\usepackage[math-style=ISO,bold-style=ISO]{unicode-math}
\setmathfont{FiraMath-Regular.otf}[Path=../../../fonts/]
\setmathfont{PlusJakartaSans-Medium.ttf}[Path=../../../fonts/,range={up/num,up/{Latin,latin},\mathup}]
\usepackage{icomma} % virgule décimale sans espace en mode math
% Caractères mathématiques Unicode absents des polices de texte (Plus Jakarta,
% Archivo Black) : rendus par leur équivalent mathématique, en texte comme en
% formule — sinon ils s'affichent en carré vide (ex. « Arithmétique dans ℤ »).
\usepackage{newunicodechar}
\newunicodechar{ℝ}{\ensuremath{\mathbb{R}}}
\newunicodechar{ℤ}{\ensuremath{\mathbb{Z}}}
\newunicodechar{ℂ}{\ensuremath{\mathbb{C}}}
\newunicodechar{ℕ}{\ensuremath{\mathbb{N}}}
\newunicodechar{ℚ}{\ensuremath{\mathbb{Q}}}
\newunicodechar{𝔻}{\ensuremath{\mathbb{D}}}
\newunicodechar{α}{\ensuremath{\alpha}}
\newunicodechar{β}{\ensuremath{\beta}}
\newunicodechar{γ}{\ensuremath{\gamma}}
\newunicodechar{δ}{\ensuremath{\delta}}
\newunicodechar{ε}{\ensuremath{\varepsilon}}
\newunicodechar{θ}{\ensuremath{\theta}}
\newunicodechar{λ}{\ensuremath{\lambda}}
\newunicodechar{μ}{\ensuremath{\mu}}
\newunicodechar{σ}{\ensuremath{\sigma}}
\newunicodechar{τ}{\ensuremath{\tau}}
\newunicodechar{φ}{\ensuremath{\varphi}}
\newunicodechar{ψ}{\ensuremath{\psi}}
\newunicodechar{ω}{\ensuremath{\omega}}
\newunicodechar{Σ}{\ensuremath{\Sigma}}
% majuscules produites par \MakeUppercase dans les titres (α-aminés -> Α)
\newunicodechar{Α}{\ensuremath{\alpha}}
\newunicodechar{Β}{\ensuremath{\beta}}
\newunicodechar{Γ}{\ensuremath{\Gamma}}
\newunicodechar{Δ}{\ensuremath{\Delta}}
\newunicodechar{Ε}{\ensuremath{\varepsilon}}
\newunicodechar{Θ}{\ensuremath{\Theta}}
\newunicodechar{Λ}{\ensuremath{\Lambda}}
\newunicodechar{Μ}{\ensuremath{\mu}}
\newunicodechar{Τ}{\ensuremath{\tau}}
\newunicodechar{Φ}{\ensuremath{\Phi}}
\newunicodechar{Ψ}{\ensuremath{\Psi}}
\newunicodechar{Ω}{\ensuremath{\Omega}}
\newunicodechar{↦}{\ensuremath{\mapsto}}
\newunicodechar{∈}{\ensuremath{\in}}
\newunicodechar{∉}{\ensuremath{\notin}}
\newunicodechar{∧}{\ensuremath{\wedge}}
\newunicodechar{⇔}{\ensuremath{\Leftrightarrow}}
\newunicodechar{⟺}{\ensuremath{\Longleftrightarrow}}
\newunicodechar{⇒}{\ensuremath{\Rightarrow}}
\newunicodechar{⟹}{\ensuremath{\Longrightarrow}}
\newunicodechar{∘}{\ensuremath{\circ}}
\newunicodechar{∪}{\ensuremath{\cup}}
\newunicodechar{∩}{\ensuremath{\cap}}
\newunicodechar{≡}{\ensuremath{\equiv}}
\newunicodechar{⊂}{\ensuremath{\subset}}
\newunicodechar{⊥}{\ensuremath{\perp}}
\newunicodechar{∃}{\ensuremath{\exists}}
\newunicodechar{∀}{\ensuremath{\forall}}
\newunicodechar{∛}{\ensuremath{\sqrt[3]{\,}}}
\newunicodechar{∜}{\ensuremath{\sqrt[4]{\,}}}
\newunicodechar{⃗}{\ensuremath{{}^{\to}}}
\newunicodechar{ⁿ}{\ifmmode^{n}\else\textsuperscript{\ensuremath{n}}\fi}
\newunicodechar{ˣ}{\ifmmode^{x}\else\textsuperscript{\ensuremath{x}}\fi}
\newunicodechar{ᵇ}{\ifmmode^{b}\else\textsuperscript{\ensuremath{b}}\fi}
\newunicodechar{ʳ}{\ifmmode^{r}\else\textsuperscript{\ensuremath{r}}\fi}
\newunicodechar{ᵘ}{\ifmmode^{u}\else\textsuperscript{\ensuremath{u}}\fi}
\newunicodechar{ʸ}{\ifmmode^{y}\else\textsuperscript{\ensuremath{y}}\fi}
\newunicodechar{ᵃ}{\ifmmode^{a}\else\textsuperscript{\ensuremath{a}}\fi}
\newunicodechar{ᵉ}{\ifmmode^{e}\else\textsuperscript{\ensuremath{e}}\fi}
\newunicodechar{ᵏ}{\ifmmode^{k}\else\textsuperscript{\ensuremath{k}}\fi}
\newunicodechar{ᵖ}{\ifmmode^{p}\else\textsuperscript{\ensuremath{p}}\fi}
\newunicodechar{ᴺ}{\ifmmode^{N}\else\textsuperscript{\ensuremath{N}}\fi}
\newunicodechar{ᶜ}{\ifmmode^{c}\else\textsuperscript{\ensuremath{c}}\fi}
\newunicodechar{ʰ}{\ifmmode^{h}\else\textsuperscript{\ensuremath{h}}\fi}
\newunicodechar{ᵠ}{\ifmmode^{q}\else\textsuperscript{\ensuremath{q}}\fi}
\newunicodechar{ₙ}{\ifmmode_{n}\else\textsubscript{\ensuremath{n}}\fi}
\newunicodechar{ₐ}{\ifmmode_{a}\else\textsubscript{\ensuremath{a}}\fi}
\newunicodechar{ᵢ}{\ifmmode_{i}\else\textsubscript{\ensuremath{i}}\fi}
\newunicodechar{ᵣ}{\ifmmode_{r}\else\textsubscript{\ensuremath{r}}\fi}
\newunicodechar{ₖ}{\ifmmode_{k}\else\textsubscript{\ensuremath{k}}\fi}
\newunicodechar{ᵦ}{\ifmmode_{\beta}\else\textsubscript{\ensuremath{\beta}}\fi}
\newunicodechar{ₓ}{\ifmmode_{x}\else\textsubscript{\ensuremath{x}}\fi}
\newfontfamily\popdisplay{ArchivoBlack-Regular.ttf}[Path=../../../fonts/]
\newfontfamily\poplogo{SpaceGrotesk-Bold.ttf}[Path=../../../fonts/]
\newfontfamily\popsemi{PlusJakartaSans-SemiBold.ttf}[Path=../../../fonts/]
\newfontfamily\popxbold{PlusJakartaSans-ExtraBold.ttf}[Path=../../../fonts/]
\newfontfamily\popxboldls{PlusJakartaSans-ExtraBold.ttf}[Path=../../../fonts/,LetterSpace=3.0]

\setlist[enumerate]{leftmargin=1.7em,itemsep=5pt,topsep=4pt}
\setlist[itemize]{leftmargin=1.7em,itemsep=4pt,topsep=4pt,label={\color{poporange}\textbullet}}
\setlength{\parindent}{0pt}
\setlength{\parskip}{5pt}
\renewcommand{\arraystretch}{1.25}

% pgfplots : style Pop par défaut
\pgfplotsset{
  every axis/.append style={axis line style={popink,line width=1pt},tick style={popink},
    grid style={popline,line width=0.5pt},label style={font=\small},tick label style={font=\footnotesize}},
  popcurve/.style={poporange,line width=1.8pt,no markers,smooth},
}
% Même style disponible dans un \draw TikZ simple (hors pgfplots)
\tikzset{popcurve/.style={poporange,line width=1.8pt}}
\tikzset{popgrid/.style={popline,very thin}}
\colorlet{popcurve}{poporange} % le nom du style est parfois utilisé comme couleur

% ============================================================
% Pill / squiggle
% ============================================================
\newcommand{\poppill}[3]{%
  \tikz[baseline=-0.55ex]\node[draw=popink,line width=1.2pt,rounded corners=2.6pt,%
    fill=#2,inner xsep=6.5pt,inner ysep=3pt]{%
    \popxboldls\fontsize{7}{7}\selectfont\color{#3}\MakeUppercase{#1}};%
}
\newcommand{\popsquiggle}[1]{%
  \tikz[baseline=-0.4ex]{
    \draw[#1,line width=2.6pt,line cap=round]
      plot [smooth] coordinates {(0,0.09) (0.34,0.01) (0.68,0.21) (1.02,0.01) (1.36,0.10)};
  }%
}
% Mot-clé surligné (marqueur orange sous le texte)
\newcommand{\cle}[1]{{\popxbold\color{popdark}#1}}

% ============================================================
% En-tête / pied
% ============================================================
\pagestyle{fancy}
\fancyhf{}
\renewcommand{\headrulewidth}{0pt}
\renewcommand{\footrulewidth}{0pt}
\lhead{\raisebox{-0.15ex}{\poplogo\fontsize{13}{13}\selectfont\color{popink}OnBuch\color{poporange}+}}
\rhead{\poppill{\DOCMATIERE\ \textbullet\ \DOCNIVEAU\ \textbullet\ Leçon \DOCLECON}{white}{popink}}
\lfoot{\popxbold\fontsize{7.6}{7.6}\selectfont\color{popmuted}\MakeUppercase{Cours signé OnBuch+}\ \textbullet\ {\popsemi\DOCTITRE}}
\rfoot{\tikz[baseline=-0.6ex]\node[rounded corners=8.5pt,fill=popink,minimum height=17pt,minimum width=17pt,inner xsep=5pt,inner sep=0]{\popxbold\fontsize{8}{8}\selectfont\color{popcream}\thepage\,/\,\pageref*{LastPage}};}

% ============================================================
% Titres
% ============================================================
\newcommand{\chaptitre}[2]{%
  \begin{tcolorbox}[enhanced,colback=poporange,colframe=popink,boxrule=2.6pt,arc=11pt,
      drop shadow=popink,left=15pt,right=15pt,top=12pt,bottom=11pt,before skip=3pt,after skip=18pt]
    \poppill{#2}{popink}{popcream}\par\vspace{9pt}
    {\raggedright\hyphenpenalty=10000\exhyphenpenalty=10000\popdisplay\fontsize{22}{24}\selectfont\color{popcream}\MakeUppercase{#1}\par}
  \end{tcolorbox}
}
% Section numérotée : pastille orange avec le numéro + titre Archivo
\newcounter{popsec}
\newcommand{\coursec}[1]{%
  \stepcounter{popsec}\setcounter{popsub}{0}%
  \par\vspace{16pt}\noindent
  \begin{minipage}{\linewidth}
  \tikz[baseline=-0.7ex]\node[draw=popink,line width=1.6pt,fill=poporange,rounded corners=4pt,minimum size=22pt,inner sep=0]{\popdisplay\fontsize{12}{12}\selectfont\color{popcream}\thepopsec};%
  \hspace{8pt}{\popdisplay\fontsize{15}{17}\selectfont\color{popink}\MakeUppercase{#1}}\par
  \vspace{4pt}\noindent{\color{poporange}\rule{36pt}{3.6pt}}
  \end{minipage}\par\nopagebreak\vspace{8pt}
}
\newcounter{popsub}
\newcommand{\courssub}[1]{%
  \stepcounter{popsub}%
  \par\vspace{9pt}\noindent{\popxbold\fontsize{11.5}{13}\selectfont\color{popink}{\color{poporange}\thepopsec.\thepopsub}\hspace{6pt}#1}\par\nopagebreak\vspace{4pt}
}

% ============================================================
% Boîtes pédagogiques — même recette : fond blanc, bordure épaisse colorée,
% ombre dure encre, badge plein. Titre optionnel [#1].
% ============================================================
\tcbset{popbase/.style={breakable,enhanced,colback=white,boxrule=2.3pt,arc=9pt,
  shadow={3pt}{-3pt}{0pt}{popink},left=14pt,right=14pt,top=11pt,bottom=11pt,before skip=11pt,after skip=15pt,
  fontupper=\popsemi\fontsize{10.6}{14.5}\selectfont\color{popink},
  fontlower=\popsemi\fontsize{10.6}{14.5}\selectfont\color{popink}}}
% Coche dessinée (le glyphe ✓ n'existe pas dans les polices du document)
\newcommand{\popcheck}[1]{\tikz[baseline=-0.5ex]\draw[#1,line width=1.6pt,line cap=round,line join=round](-0.13,0.0)--(-0.04,-0.09)--(0.13,0.1);}
\providecommand{\coche}{\popcheck{popgreen}}
\providecommand{\resultat}[1]{\par\smallskip\noindent\fcolorbox{popgreen}{popgreenL}{\parbox{\dimexpr\linewidth-2\fboxsep-2\fboxrule\relax}{\textbf{Résultat.} #1}}\par\smallskip}
\newcommand{\popboxhead}[3]{\poppill{#1}{#2}{white}\ifx\relax#3\relax\else\hspace{6pt}{\popxbold\fontsize{10.6}{12}\selectfont\color{popink}#3}\fi\par\vspace{7pt}}

\newtcolorbox{aretenir}[1][]{popbase,colframe=popgreen,before upper={\popboxhead{À retenir}{popgreen}{#1}}}
\newtcolorbox{exemplebox}[1][]{popbase,colframe=poporange,before upper={\popboxhead{Exemple}{poporange}{#1}}}
\newtcolorbox{definition}[1][]{popbase,colframe=popblue,before upper={\popboxhead{Définition}{popblue}{#1}}}
\newtcolorbox{propriete}[1][]{popbase,colframe=poppurple,before upper={\popboxhead{Propriété}{poppurple}{#1}}}
\newtcolorbox{methode}[1][]{popbase,colframe=popgold,before upper={\popboxhead{Méthode}{popgold}{#1}}}
\newtcolorbox{attention}[1][]{popbase,colframe=poppink,before upper={\popboxhead{Attention}{poppink}{#1}}}
\newtcolorbox{experience}[1][]{popbase,colframe=popblue,colback=popblueL,before upper={\popboxhead{Expérience}{popblue}{#1}}}
% « Le savais-tu ? » : carte violette pleine (contexte, culture, Cameroun)
\newtcolorbox{savaistu}[1][]{popbase,colback=poppurple,colframe=popink,
  fontupper=\popsemi\fontsize{10.4}{14.2}\selectfont\color{white},
  before upper={\poppill{Le savais-tu\,?}{white}{poppurple}\ifx\relax#1\relax\else\hspace{6pt}{\popxbold\color{white}#1}\fi\par\vspace{7pt}}}
% Exercice résolu : énoncé en haut, \tcblower puis la solution sur fond crème
\newcounter{popexo}\newcounter{popexr}
\newtcolorbox{exoresolu}[1][]{popbase,colframe=poporange,colbacklower=popcream,
  segmentation style={popink,line width=1.2pt,dashed},
  before upper={\stepcounter{popexr}\popboxhead{Exercice résolu \thepopexr}{poporange}{#1}},
  before lower={\poppill{Solution}{popink}{popcream}\par\vspace{6pt}}}
% Exercice (non résolu en place) — la correction est dans la partie Corrigés
\newtcolorbox{exercice}[1][]{popbase,colframe=popink,
  before upper={\stepcounter{popexo}\popboxhead{Exercice \thepopexo}{popink}{#1}}}
% Corrigé
\newtcolorbox{corrige}[1][]{popbase,colframe=popgreen,colback=popgreenL,
  before upper={\popboxhead{Corrigé}{popgreen}{#1}}}
% Objectifs / prérequis / bilan
\newtcolorbox{objectifs}{popbase,colframe=popink,colback=poporangeL,
  before upper={\popboxhead{Objectifs de la leçon}{popink}{}}}
\newtcolorbox{prerequis}{popbase,colframe=popink,colback=popblueL,
  before upper={\popboxhead{Prérequis}{popblue}{}}}
\newtcolorbox{bilan}{popbase,colframe=popink,colback=white,boxrule=2.8pt,
  before upper={\popboxhead{Fiche bilan}{poporange}{L'essentiel à connaître pour l'examen}}}
% Figure encadrée avec légende
\newtcolorbox{popfigure}[1][]{enhanced,breakable=false,colback=white,colframe=popline,boxrule=1.4pt,arc=8pt,
  left=10pt,right=10pt,top=10pt,bottom=8pt,before skip=10pt,after skip=12pt,halign=center,
  lower separated=false,halign lower=center,
  fontlower=\popsemi\fontsize{9}{11.5}\selectfont\color{popmuted},
  #1}
\newcommand{\legende}[1]{\par\vspace{6pt}{\popsemi\fontsize{9}{11.5}\selectfont\color{popmuted}#1}}

% ============================================================
% Signature OnBuch+
% ============================================================
\newcommand{\poplogoinline}[1]{{\poplogo\fontsize{#1}{#1}\selectfont\color{popink}OnBuch\color{poporange}+}}
\newcommand{\signatureonbuch}{%
  \begin{tcolorbox}[enhanced,colback=popink,colframe=popink,boxrule=2.2pt,arc=10pt,
      drop shadow=poporange,left=14pt,right=14pt,top=10pt,bottom=10pt,before skip=0pt,after skip=0pt]
    \tikz[baseline=-0.7ex]\node[circle,fill=popgreen,draw=popcream,line width=1.3pt,minimum size=22pt,inner sep=0]{\popcheck{white}};%
    \hspace{9pt}\begin{minipage}[c]{0.62\linewidth}
      {\popxbold\fontsize{12}{13}\selectfont\color{popcream}Signé\ \textbullet\ {\poplogo OnBuch\color{poporange}+}}\par\vspace{2pt}
      {\raggedright\popsemi\fontsize{8}{10.5}\selectfont\color{popline}\MakeUppercase{Équipe pédagogique OnBuch+}\\\MakeUppercase{Conforme au programme officiel MINESEC}\par}
    \end{minipage}\hfill
    {\poplogo\fontsize{22}{22}\selectfont\color{popcream}OnBuch\color{poporange}+}
  \end{tcolorbox}%
}

% ============================================================
% Page de garde
% #1 titre · #2 matière · #3 niveau · #4 module · #5 n° leçon · #6 durée ·
% #7 description / objectifs (texte ou itemize)
% ============================================================
\newcommand{\pagegarde}[7]{%
  \thispagestyle{empty}
  \newgeometry{a4paper,margin=1.7cm,top=1.6cm,bottom=1.4cm}%
  \begin{tikzpicture}[remember picture,overlay]
    \fill[poporange!18] ([xshift=-3.2cm,yshift=-3.6cm]current page.north east) circle (4.6cm);
    \fill[poppurple!14] ([xshift=2.4cm,yshift=7.6cm]current page.south west) circle (3.8cm);
  \end{tikzpicture}%
  {\poplogo\fontsize{30}{30}\selectfont\color{popink}OnBuch\color{poporange}+}\hfill
  \raisebox{0.6ex}{\poppill{Cours signé \textbullet\ Premium}{popink}{popcream}}\par
  \vspace{1pt}\popsquiggle{poppurple}\par
  \vspace{8pt}
  \begin{tcolorbox}[enhanced,colback=poporange,colframe=popink,boxrule=2.8pt,arc=13pt,
      drop shadow=popink,left=18pt,right=18pt,top=12pt,bottom=13pt,after skip=10pt]
    \poppill{#2 \textbullet\ #3}{popink}{popcream}\hspace{5pt}\poppill{#4}{white}{popink}\par\vspace{8pt}
    {\popdisplay\fontsize{14}{14}\selectfont\color{popink}LEÇON #5}\par\vspace{4pt}
    {\raggedright\hyphenpenalty=10000\exhyphenpenalty=10000\popdisplay\fontsize{25}{27}\selectfont\color{popcream}\MakeUppercase{#1}\par}
    \vspace{8pt}
    {\popxbold\fontsize{9.5}{9.5}\selectfont\color{popink}\MakeUppercase{Durée conseillée : #6}}
  \end{tcolorbox}
  \begin{tcolorbox}[enhanced,colback=white,colframe=popink,boxrule=2.2pt,arc=10pt,
      drop shadow=popink,left=16pt,right=16pt,top=10pt,bottom=10pt,after skip=8pt]
    \poppill{Dans cette leçon}{poppurple}{white}\par\vspace{5pt}
    \popsemi\fontsize{9.3}{12.6}\selectfont\color{popink}#7
  \end{tcolorbox}
  \noindent\begin{minipage}[t]{0.487\linewidth}
    \begin{tcolorbox}[enhanced,colback=popgreen,colframe=popink,boxrule=2.2pt,arc=10pt,
        drop shadow=popink,left=11pt,right=11pt,top=10pt,bottom=10pt,height=3.1cm,valign=center]
      \tcbox[colback=white,colframe=popink,boxrule=1.3pt,arc=5pt,boxsep=0pt,nobeforeafter,
        left=3pt,right=3pt,top=3pt,bottom=3pt,valign=center]{\qrcode[height=1.9cm]{https://play.google.com/store/apps/details?id=com.luvvix.onbuch&hl=fr}}%
      \hspace{8pt}\begin{minipage}[c]{0.5\linewidth}
        {\popdisplay\fontsize{11}{12.5}\selectfont\color{white}\MakeUppercase{Télécharge OnBuch}}\par\vspace{4pt}
        {\raggedright\popsemi\fontsize{8.4}{11}\selectfont\color{white}Gratuit \textbullet\ Google Play\\Tuteur IA, annales, concours, résultats\par}
      \end{minipage}
    \end{tcolorbox}
  \end{minipage}\hfill
  \begin{minipage}[t]{0.487\linewidth}
    \begin{tcolorbox}[enhanced,colback=poppurple,colframe=popink,boxrule=2.2pt,arc=10pt,
        drop shadow=popink,left=11pt,right=11pt,top=10pt,bottom=10pt,height=3.1cm,valign=center]
      \tikz[baseline=-0.7ex]\node[circle,fill=white,draw=popink,line width=1.3pt,minimum size=1.6cm,inner sep=0]{\popdisplay\fontsize{24}{24}\selectfont\color{poppurple}+};%
      \hspace{8pt}\begin{minipage}[c]{0.6\linewidth}
        {\popdisplay\fontsize{11}{12.5}\selectfont\color{white}\MakeUppercase{Passe à OnBuch+}}\par\vspace{4pt}
        {\raggedright\popsemi\fontsize{8.4}{11}\selectfont\color{white}Tous les cours signés, Tuteur IA illimité, fiches PDF — onglet OnBuch+ de l'app\par}
      \end{minipage}
    \end{tcolorbox}
  \end{minipage}
  \vspace{8pt}
  \popcontenu
  \vfill
  \signatureonbuch
  \restoregeometry
  \newpage
}

% Rangée « Ce cours contient » (page de garde)
\newcommand{\poptile}[3]{% #1 couleur #2 gros texte #3 légende
  \begin{tcolorbox}[enhanced,colback=#1,colframe=popink,boxrule=2pt,arc=9pt,shadow={2.5pt}{-2.5pt}{0pt}{popink},
      left=6pt,right=6pt,top=9pt,bottom=9pt,halign=center,valign=center,height=1.75cm,width=\linewidth,nobeforeafter]
    {\popdisplay\fontsize{12.5}{13}\selectfont\color{white}\MakeUppercase{#2}}\par\vspace{3pt}
    {\centering\popsemi\fontsize{7.8}{9.5}\selectfont\color{white}#3\par}
  \end{tcolorbox}}
\newcommand{\popcontenu}{%
  \par\noindent{\popxboldls\fontsize{7.5}{7.5}\selectfont\color{popmuted}CE COURS SIGNÉ CONTIENT}\par\vspace{7pt}
  \noindent\begin{minipage}[t]{0.235\linewidth}\poptile{popblue}{Cours}{complet, illustré, pas à pas}\end{minipage}\hfill
  \begin{minipage}[t]{0.235\linewidth}\poptile{poporange}{Méthodes}{et exercices résolus}\end{minipage}\hfill
  \begin{minipage}[t]{0.235\linewidth}\poptile{poppink}{Exercices}{type examen + corrigés}\end{minipage}\hfill
  \begin{minipage}[t]{0.235\linewidth}\poptile{popgreen}{Fiche bilan}{l'essentiel à retenir}\end{minipage}\par}

% Sommaire
\newtcolorbox{sommaire}{enhanced,colback=white,colframe=popink,boxrule=2.2pt,arc=10pt,
  drop shadow=popink,left=18pt,right=18pt,top=13pt,bottom=13pt,before skip=6pt,after skip=20pt,
  before upper={\poppill{Sommaire}{poppurple}{white}\par\vspace{9pt}\popsemi\fontsize{10.6}{14.5}\selectfont\color{popink}}}

% Signature de fin de cours — poussée tout en bas de la dernière page
\newcommand{\findecours}{%
  \par\vfill\nopagebreak
  \begin{center}\popsquiggle{poporange}\end{center}\vspace{4pt}
  \signatureonbuch
}
\AtBeginDocument{\def\llbracket{\mathopen{[\mkern-3.5mu[}}\def\rrbracket{\mathclose{]\mkern-3.5mu]}}} % intervalles d'entiers (glyphe ⟦ absent de la police)
% Glyphes absents de Fira Math : redéfinis à partir de glyphes disponibles
\AtBeginDocument{%
  \def\setminus{\mathbin{\backslash}}\let\smallsetminus\setminus
  \def\vdots{\mathord{\vbox{\baselineskip3.2pt\lineskiplimit0pt\kern4pt\hbox{.}\hbox{.}\hbox{.}}}}%
  \def\ddots{\mathinner{\mkern1mu\raise6.5pt\hbox{.}\mkern2mu\raise3.6pt\hbox{.}\mkern2mu\raise0.7pt\hbox{.}\mkern1mu}}%
  \def\triangle{\mathord{\scalebox{0.9}{$\symup{\Delta}$}}}%
  \def\nabla{\mathord{\raisebox{\depth}{\scalebox{1}[-1]{$\symup{\Delta}$}}}}%
  \def\checkmark{\popcheck{popdark}}%
  \def\star{\mathbin{\ast}}\def\bigstar{\mathord{\ast}}%
  \def\diamond{\mathbin{\scalebox{0.6}{$\Diamond$}}}%
  \def\bigcup{\mathop{\mathchoice{\raisebox{-0.3em}{\scalebox{1.7}{$\cup$}}}{\scalebox{1.25}{$\cup$}}{\cup}{\cup}}\displaylimits}%
  \def\bigcap{\mathop{\mathchoice{\raisebox{-0.3em}{\scalebox{1.7}{$\cap$}}}{\scalebox{1.25}{$\cap$}}{\cap}{\cap}}\displaylimits}%
  \def\lesseqqgtr{\mathrel{\lessgtr}}\def\bigoplus{\mathop{\oplus}\displaylimits}%
}
\providecommand{\ssi}{si et seulement si} % abréviation fréquente
\AtBeginDocument{\providecommand{\wideparen}{\overparen}} % arc de cercle

%% FIGFIX-BEGIN v5 — correctifs globaux des figures (docs/onbuch-generation/tools/figfix)
\usepackage{adjustbox}\usepackage{etoolbox}\tcbuselibrary{raster}
% toute figure TikZ/pgfplots plus large que la ligne est réduite à la largeur disponible
\BeforeBeginEnvironment{tikzpicture}{\begin{adjustbox}{max width=\linewidth}}
\AfterEndEnvironment{tikzpicture}{\end{adjustbox}}
% légendes pgfplots lisibles par-dessus les courbes
\pgfplotsset{every axis/.append style={legend style={fill=white,fill opacity=0.92,text opacity=1,draw=black!15,inner sep=3pt}}}
% étiquettes d'arêtes (syntaxe edge["..."]) avec fond blanc
\tikzset{every edge quotes/.append style={fill=white,inner sep=1.2pt}}
%% FIGFIX-END
\begin{document}

\pagegarde{Sensibilisation sur la nécessité de préserver le fonctionnement des neurones}{SVT}{Tle D}{Module I : Le Monde Vivant — Fonctionnement des neurones}{N08}{10 h}{Cette leçon étudie, à partir d'expériences historiques et de courbes d'enregistrement, les mécanismes ioniques du potentiel de repos et du potentiel d'action, ainsi que les conditions de naissance et de propagation du message nerveux. Elle fournit les outils d'interprétation des tracés électrophysiologiques indispensables à l'épreuve de SVT du Baccalauréat.
\vspace{4pt}
\begin{multicols}{2}\small
\begin{itemize}
\item À la fin de cette leçon, je sais réaliser et interpréter le schéma du montage mettant en évidence le potentiel de repos
\item À la fin de cette leçon, je sais interpréter ioniquement le potentiel de repos à partir des gradients de concentration en Na+ et K+
\item À la fin de cette leçon, je sais décrire la courbe du potentiel d'action et l'expliquer par les variations de perméabilité membranaire
\item À la fin de cette leçon, je sais interpréter les tracés du potentiel de récepteur et le codage du message nerveux
\item À la fin de cette leçon, je sais interpréter la réponse d'une fibre et d'un nerf à des stimulations d'intensités croissantes
\item À la fin de cette leçon, je sais exploiter la courbe intensité-durée (rheobase, temps utile)
\end{itemize}\end{multicols}}

\begin{sommaire}
\begin{multicols}{2}
\begin{enumerate}
\item Le potentiel de repos : mise en évidence expérimentale et interprétation ionique
\item Le potentiel d'action : enregistrement, phases et mécanismes ioniques
\item Naissance du message nerveux : récepteurs sensoriels et seuil d'excitation
\item Propagation du potentiel d'action le long de la fibre nerveuse
\item Synthèse et sensibilisation : préserver le fonctionnement des neurones
\item Activité d'intégration
\item Méthodes et exercices résolus
\item Exercices
\item Corrigés des exercices
\item Fiche bilan
\end{enumerate}
\end{multicols}
\end{sommaire}

\begin{objectifs}
\begin{itemize}
\item À la fin de cette leçon, je sais réaliser et interpréter le schéma du montage mettant en évidence le potentiel de repos
\item À la fin de cette leçon, je sais interpréter ioniquement le potentiel de repos à partir des gradients de concentration en Na+ et K+
\item À la fin de cette leçon, je sais décrire la courbe du potentiel d'action et l'expliquer par les variations de perméabilité membranaire
\item À la fin de cette leçon, je sais interpréter les tracés du potentiel de récepteur et le codage du message nerveux
\item À la fin de cette leçon, je sais interpréter la réponse d'une fibre et d'un nerf à des stimulations d'intensités croissantes
\item À la fin de cette leçon, je sais exploiter la courbe intensité-durée (rheobase, temps utile)
\item À la fin de cette leçon, je sais expliquer le sens et le mode de propagation du potentiel d'action et le rôle de la période réfractaire
\item À la fin de cette leçon, je sais calculer la vitesse de propagation d'un message nerveux et citer les facteurs qui l'influencent
\end{itemize}
\end{objectifs}

\begin{prerequis}
\begin{itemize}
\item Structure du neurone et de la fibre nerveuse (corps cellulaire, dendrites, axone, gaine de myéline, nœuds de Ranvier) vues en classe de Première
\item Notion de membrane plasmique, de perméabilité sélective et de transport membranaire (diffusion, transport actif)
\item Notions d'électricité de base : différence de potentiel, électrode, millivolt, oscilloscope/galvanomètre
\item Organisation générale du système nerveux et notion de message nerveux (Première)
\item Lecture et exploitation de courbes expérimentales (axes, unités, variations)
\end{itemize}
\end{prerequis}

\newpage

\coursec{Le potentiel de repos : mise en évidence expérimentale et interprétation ionique}

Imagine un instant la scène à l'hôpital de district de Bafoussam. Un jeune motard, victime d'un accident de la circulation, est admis en urgence : il ne sent plus ses jambes, il ne peut plus les bouger. Dans la salle voisine, un paysan arrive, la main gonflée par une piqûre de scorpion ; une douleur fulgurante lui remonte le long du bras à une vitesse stupéfiante. Quelques kilomètres plus loin, chez un dentiste de Yaoundé, une injection d'anesthésique local rend une dent insensible en quelques minutes, le temps d'une extraction. Pourquoi l'information douloureuse voyage-t-elle si vite le long des nerfs ? Pourquoi une lésion de la moelle épinière prive-t-elle définitivement le cerveau des messages venant des membres ? Et comment un simple produit chimique bloque-t-il temporairement la douleur ? Toutes ces questions trouvent leur réponse dans la physique et la chimie de la membrane neuronale. Commençons par le fondement : comment un neurone, au repos, entretient-il une différence de potentiel électrique à travers sa membrane ?

\textbf{Mise en évidence expérimentale du potentiel de repos}\par

\textbf{Le montage de référence : l'axone géant du calmar.}
C'est sur l'axone géant du calmar (\emph{Dosidicus gigas} ou \emph{Loligo}), dont le diamètre peut atteindre \SI{1}{\milli\metre} (visible à l'œil nu), que les physiologistes Alan Hodgkin et Andrew Huxley ont réalisé, dans les années 1940-1950, les expériences fondatrices qui leur vaudront le prix Nobel en 1963. La taille exceptionnelle de cette fibre permet d'introduire une microélectrode à l'intérieur même de l'axoplasme sans la détruire.

\begin{popfigure}
\begin{tikzpicture}[scale=0.9, every node/.style={font=\small}]
    % Couleurs
    \colorlet{axoplasm}{popblue!10}
    \colorlet{membrane}{popdark}
    \colorlet{electrode}{poporange}
    \colorlet{oscillo}{poppurple!10}
    
    % --- Situation A : Deux électrodes extracellulaires ---
    \begin{scope}[xshift=0cm]
        \node[align=center, font=\bfseries\small] at (2.5, 4.5) {Situation A : Électrodes en surface};
        % Bain
        \draw[fill=popblue!5, draw=popblue] (0.5, 0.5) rectangle (4.5, 3.5);
        % Axone
        \draw[fill=axoplasm, draw=membrane, thick] (1.5, 1.2) rectangle (3.5, 2.8);
        \node[align=center] at (2.5, 2.0) {Axone\\(axoplasme)};
        % Électrodes
        \draw[electrode, thick] (1.0, 3.5) -- (1.0, 4.2) node[above] {E1};
        \draw[electrode, thick] (4.0, 3.5) -- (4.0, 4.2) node[above] {E2};
        % Fil vers oscilloscope
        \draw[popline, thick] (1.0, 4.2) -- (0.5, 4.7) -- (0.5, 5.5);
        \draw[popline, thick] (4.0, 4.2) -- (4.5, 4.7) -- (4.5, 5.5);
        % Oscilloscope A
        \draw[fill=oscillo, draw=poppurple, thick, rounded corners] (0, 5.5) rectangle (5, 7.5);
        \node[font=\bfseries] at (2.5, 7.2) {Oscilloscope / Galvanomètre};
        % Trace A : ligne plate
        \draw[popcurve, thick, domain=0.3:4.7, samples=2] plot (\x, {6.5});
        \draw[popmuted, thin] (0.3, 6.5) -- (4.7, 6.5);
        \node[popmuted] at (2.5, 6.0) {Trace : ligne plate (0 mV)};
        \node[align=center, font=\small] at (2.5, 0.0) {Milieu extracellulaire\\(électrodes E1 et E2 à l'extérieur)};
    \end{scope}

    % --- Situation B : Une électrode intracellulaire ---
    \begin{scope}[xshift=8.5cm]
        \node[align=center, font=\bfseries\small] at (2.5, 4.5) {Situation B : Une électrode intracellulaire};
        % Bain
        \draw[fill=popblue!5, draw=popblue] (0.5, 0.5) rectangle (4.5, 3.5);
        % Axone
        \draw[fill=axoplasm, draw=membrane, thick] (1.5, 1.2) rectangle (3.5, 2.8);
        \node[align=center] at (2.5, 2.0) {Axone};
        % Électrodes
        \draw[electrode, thick] (2.5, 3.5) -- (2.5, 2.8); % tige
        \draw[electrode, thick, ->] (2.5, 2.8) -- (2.5, 2.0) node[left, pos=0.5] {Pénétration};
        \draw[electrode, thick] (4.0, 3.5) -- (4.0, 4.2) node[above] {E2 (ref)};
        % Fil vers oscilloscope
        \draw[popline, thick] (2.5, 3.5) -- (2.5, 4.2) -- (0.5, 4.7) -- (0.5, 5.5);
        \draw[popline, thick] (4.0, 4.2) -- (4.5, 4.7) -- (4.5, 5.5);
        % Oscilloscope B
        \draw[fill=oscillo, draw=poppurple, thick, rounded corners] (0, 5.5) rectangle (5, 7.5);
        \node[font=\bfseries] at (2.5, 7.2) {Oscilloscope / Galvanomètre};
        % Trace B : déviation brutale
        \draw[popmuted, thin] (0.3, 6.5) -- (1.5, 6.5);
        \draw[popcurve, thick] (1.5, 6.5) -- (1.5, 5.0); % front descendant
        \draw[popcurve, thick, domain=1.5:4.7, samples=2] plot (\x, {5.0});
        \node[popcurve] at (3.0, 4.5) {Déflection brutale};
        \node[popcurve] at (3.0, 4.1) {$\approx -70$ mV};
        \node[align=center, font=\small] at (2.5, 0.0) {E1 dans l'axoplasme ; E2 à l'extérieur (référence)};
    \end{scope}
\end{tikzpicture}
\legende{Montage électrophysiologique classique sur axone géant de calmar. \textbf{A :} Les deux électrodes (E1, E2) plongées dans le milieu extracellulaire ne détectent aucune différence de potentiel (ddp = 0 mV). \textbf{B :} L'électrode E1 pénètre l'axoplasme ; l'oscilloscope enregistre une déflexion brutale vers le négatif : le potentiel de repos $V_m \approx -70$ mV (intérieur négatif par rapport à l'extérieur).}
\end{popfigure}

\begin{methode}[Interpréter un montage électrophysiologique]
Pour analyser ou schématiser une expérience de mesure du potentiel membranaire, suis ces étapes rigoureuses :
\begin{enumerate}
    \item \textbf{Préciser le montage :} Nature de la préparation (axone géant de calmar, fibre musculaire, etc.), type d'électrodes (microélectrodes en verre remplies de KCl 3 mol/L), appareil de mesure (oscilloscope à balayage lent ou voltmètre à haute impédance).
    \item \textbf{Indiquer la position des électrodes :} Électrode de référence (indifférente) toujours dans le milieu extracellulaire ; électrode d'exploration (active) soit extracellulaire, soit intracellulaire.
    \item \textbf{Noter le résultat chiffré :} Valeur de la ddp (en mV), signe (négatif = intérieur négatif), stabilité du tracé.
    \item \textbf{Conclure :} Déduire l'état de polarisation de la membrane (polarisée au repos, dépolarisée, hyperpolarisée).
\end{enumerate}
\end{methode}

\begin{definition}[Potentiel de repos (PR)]
Le \cle{potentiel de repos} est la différence de potentiel électrique transmembranaire permanente, d'environ \textbf{-70 mV} (valeur variant de -50 à -90 mV selon les types cellulaires), qui existe entre les deux faces de la membrane plasmique d'une cellule vivante non excitée (au repos). L'intérieur de la cellule est \textbf{négatif} par rapport à l'extérieur. La membrane est alors dite \cle{polarisée}.
\end{definition}

\textbf{Composition ionique des milieux intra- et extracellulaires : les gradients de concentration}\par

L'expérience montre que le potentiel de repos n'est pas un hasard : il est le reflet direct de la répartition inégale des ions de part et d'autre de la membrane. Des analyses chimiques (par spectrophotométrie d'absorption atomique ou électrodes sélectives d'ions) sur l'axoplasme et le milieu extracellulaire du calmar donnent les concentrations moyennes suivantes :

\begin{exemplebox}[Concentrations ioniques types dans l'axone géant de calmar (mM)]
\begin{center}
\begin{tabularx}{\linewidth}{|l|X|X|X|}\hline
\rowcolor{popblueL}
\textbf{Ion} & \textbf{Intracellulaire (axoplasme)} & \textbf{Extracellulaire (eau de mer)} & \textbf{Gradient (Int/Ext)} \\ \hline
$\mathbf{K^+}$ & $\approx 400$ & $\approx 20$ & \textbf{Sortie} favorisée (gradient chimique vers l'extérieur) \\ \hline
$\mathbf{Na^+}$ & $\approx 50$ & $\approx 440$ & \textbf{Entrée} favorisée (gradient chimique vers l'intérieur) \\ \hline
$\mathbf{Cl^-}$ & $\approx 40\text{--}150$ & $\approx 560$ & Entrée favorisée \\ \hline
$\mathbf{Anions\ organiques^-}$ (protéines, acides aminés, $\text{ATP}^{4-}$) & $\approx 300\text{--}400$ & $\approx 0$ & \textbf{Non diffusibles} (piégés à l'intérieur) \\ \hline
\end{tabularx}
\end{center}
\end{exemplebox}

\begin{savaistu}[L'axone géant du calmar : une aubaine pour la science]
Pourquoi le calmar ? Parce que son axone géant (jusqu'à \SI{1}{\milli\metre} de diamètre, soit \textbf{1000 fois plus gros} qu'un axone mammifère typique) permet l'insertion d'électrodes intracellulaires et même le remplacement de l'axoplasme par des solutions artificielles. C'est sur cette préparation que Hodgkin et Huxley (Prix Nobel 1963) ont élucidé les mécanismes ioniques du potentiel d'action. Au Cameroun, les laboratoires de physiologie animale de l'Université de Yaoundé I ou de Douala utilisent souvent la préparation « nerf sciatique - muscle gastrocnémien » de grenouille (\emph{Xenopus} ou \emph{Ptychadena}), plus accessible, pour illustrer ces mêmes principes en travaux pratiques.
\end{savaistu}

\textbf{Le potassium, artisan principal du potentiel de repos : l'expérience décisive}\par

Si le $\text{Na}^+$ est bien plus concentré à l'extérieur, pourquoi le potentiel de repos est-il négatif (proche de l'équilibre du $\text{K}^+$) et non positif ? L'expérience clé consiste à varier la concentration extracellulaire en potassium $[\text{K}^+]_e$ et à mesurer la variation du potentiel de repos $V_m$.

\begin{experience}[{Variation du potentiel de repos en fonction de $[\text{K}^+]_e$}]
\textbf{Matériel :} Axone géant de calmar, montage à deux électrodes, bains de perfusion à $[\text{K}^+]_e$ contrôlée (de 10 mmol/L à 400 mmol/L), oscilloscope.
\textbf{Protocole :} On perfuse l'axone avec des solutions de $[\text{K}^+]_e$ croissantes. On attend la stabilisation du $V_m$ à chaque palier. On relève la valeur du potentiel de repos.
\textbf{Résultats attendus (données types) :}
\begin{center}
\begin{tabularx}{\linewidth}{|c|X|}\hline
\rowcolor{popblueL}
$[\text{K}^+]_e$ (mmol/L) & $V_m$ (mV) \\ \hline
10 & -95 \\ \hline
20 (physiologique) & -70 \\ \hline
40 & -50 \\ \hline
100 & -25 \\ \hline
200 & -10 \\ \hline
400 (isotonique intracellulaire) & 0 \\ \hline
\end{tabularx}
\end{center}
\textbf{Interprétation :} Lorsque $[\text{K}^+]_e$ augmente, le gradient de concentration $\text{K}^+$ (int $\to$ ext) diminue. La force chimique expulsant le $\text{K}^+$ faiblit. Moins de charges positives sortent, l'intérieur devient \textbf{moins négatif} : c'est une \cle{dépolarisation} (la valeur absolue du PR diminue). Quand $[\text{K}^+]_e = [\text{K}^+]_i$ (400 mmol/L), le gradient s'annule, il n'y a plus de force chimique pour le $\text{K}^+$, le potentiel de repos s'effondre à \textbf{0 mV}.
\end{experience}

\begin{popfigure}
\begin{tikzpicture}
\begin{axis}[
    width=\linewidth, height=8cm,
    xmode=log, log basis x=10,
    xlabel={$[\text{K}^+]_e$ (mmol/L)}, ylabel={Potentiel de repos $V_m$ (mV)},
    xmin=5, xmax=500, ymin=-110, ymax=10,
    grid=both, minor tick num=1,
    xtick={10,20,40,100,200,400},
    xticklabels={10,20,40,100,200,400},
    ytick={-100,-80,-60,-40,-20,0},
    axis lines=middle,
    enlargelimits=0.05,
    legend style={at={(0.98,0.98)}, anchor=north east, fill=white, draw=popline},
    popcurve/.style={popblue, thick, mark=*, mark size=2pt}
]
\addplot[popcurve] coordinates {
    (10, -95)
    (20, -70)
    (40, -50)
    (100, -25)
    (200, -10)
    (400, 0)
};
\addlegendentry{Données expérimentales}
% Droite de Nernst théorique pour K+ (pente 58 mV/décade à 20°C)
\addplot[poporange, dashed, thick, domain=10:400, samples=50] {58 * (log10(x/400))};
\addlegendentry{Équilibre de Nernst $E_K$ (théorique)}
\end{axis}
\end{tikzpicture}
\legende{Variation du potentiel de repos $V_m$ en fonction de la concentration extracellulaire en potassium $[\text{K}^+]_e$ (échelle logarithmique). Les points expérimentaux (bleu) suivent la droite de Nernst pour le $\text{K}^+$ (orange, pointillés), prouvant que $V_m \approx E_K$ au repos. L'augmentation de $[\text{K}^+]_e$ provoque une dépolarisation.}
\end{popfigure}

\begin{attention}[Piège classique : Confondre gradient chimique et potentiel électrique]
\begin{itemize}
    \item \textbf{Erreur fréquente :} « Le $\text{Na}^+$ entre au repos car il est concentré dehors. » \textbf{FAUX.} Au repos, la membrane est \textbf{quasi imperméable au $\text{Na}^+$} (canaux fermés). Le $\text{Na}^+$ ne traverse pratiquement pas la membrane.
    \item \textbf{Vrai :} C'est la \textbf{sortie du $\text{K}^+$} (membrane très perméable au $\text{K}^+$ via des canaux de fuite) qui laisse derrière elle les anions fixes négatifs (protéines), créant la négativité interne.
    \item \textbf{Repère :} Au repos, $P_{\text{K}^+} \gg P_{\text{Na}^+}$ (perméabilité au $\text{K}^+$ environ 20 à 100 fois supérieure à celle au $\text{Na}^+$).
\end{itemize}
\end{attention}

\textbf{Interprétation ionique : diffusion, perméabilité et anions non diffusibles}\par

Comment expliquer quantitativement ce $\approx -70$ mV ? C'est l'équilibre entre deux forces agissant sur le $\text{K}^+$ :

\begin{enumerate}
    \item \textbf{Force chimique (gradient de concentration) :} Pousse le $\text{K}^+$ vers l'extérieur (de 400 mM vers 20 mM).
    \item \textbf{Force électrique (champ transmembranaire) :} Dès que des $\text{K}^+$ sortent, l'intérieur devient négatif. Cette négativité attire les cations $\text{K}^+$ vers l'intérieur (force électrique opposée à la sortie).
\end{enumerate}
Au \cle{potentiel d'équilibre de Nernst pour le $\text{K}^+$ ($E_K$)}, ces deux forces s'annulent. À \SI{20}{\celsius}, l'équation de Nernst simplifiée donne :
\[
E_K = 58 \log_{10}\left(\frac{[\text{K}^+]_e}{[\text{K}^+]_i}\right) \approx 58 \log_{10}\left(\frac{20}{400}\right) = 58 \times (-1.30) \approx \mathbf{-75 \text{ mV}}
\]
La valeur expérimentale (-70 mV) est très proche de $E_K$. La légère différence (le $V_m$ est un peu moins négatif que $E_K$) s'explique par une \textbf{faible perméabilité résiduelle au $\text{Na}^+$} qui laisse entrer quelques $\text{Na}^+$, dépolarisant légèrement la membrane.

\begin{popfigure}
\begin{tikzpicture}[scale=1.1, every node/.style={font=\small}]
    % Membrane
    \draw[popdark, line width=2pt] (0,0) -- (10,0) node[midway, above=2pt] {Face extracellulaire};
    \draw[popdark, line width=2pt] (0,-2.5) -- (10,-2.5) node[midway, below=2pt] {Face intracellulaire (axoplasme)};
    \draw[popblue!20, fill=popblue!5] (0,0) rectangle (10,-2.5);
    \node[popdark, rotate=90] at (-0.6, -1.25) {Bicouche lipidique};
    
    % Ions extracellulaires
    \foreach \x in {1.5, 3.5, 5.5, 7.5, 9} {
        \node[poporange, circle, draw, thick, inner sep=2pt] at (\x, 0.6) {$\text{Na}^+$};
        \node[popblue, circle, draw, thick, inner sep=2pt] at (\x+0.4, 1.1) {$\text{K}^+$};
        \node[poppurple, circle, draw, thick, inner sep=2pt] at (\x-0.4, 1.1) {$\text{Cl}^-$};
    }
    \node[poporange] at (1, 1.5) {$[\text{Na}^+] \approx 440$ mM};
    \node[popblue] at (9, 1.5) {$[\text{K}^+] \approx 20$ mM};
    
    % Ions intracellulaires
    \foreach \x in {1.5, 3.5, 5.5, 7.5, 9} {
        \node[popblue, circle, draw, thick, inner sep=2pt] at (\x, -3.1) {$\text{K}^+$};
        \node[poporange, circle, draw, thick, inner sep=2pt] at (\x+0.4, -3.6) {$\text{Na}^+$};
        \node[poppurple, circle, draw, thick, inner sep=2pt] at (\x-0.4, -3.6) {$\text{Cl}^-$};
        \node[popgreen, draw, thick, rounded corners, inner sep=1pt] at (\x, -1.8) {$\text{A}^-$};
    }
    \node[popblue] at (1, -4.2) {$[\text{K}^+] \approx 400$ mM};
    \node[poporange] at (9, -4.2) {$[\text{Na}^+] \approx 50$ mM};
    \node[popgreen] at (5, -4.5) {Anions organiques non diffusibles ($\text{Pr}^-$, $\text{ATP}^{4-}$...) $\approx 350$ mM};
    
    % Canaux de fuite K+
    \draw[popblue, thick, fill=popblue!20] (2, -0.2) rectangle (3, -2.3);
    \node[popblue, align=center] at (2.5, -1.25) {Canal $\text{K}^+$\\de fuite (ouvert)};
    \draw[popblue, ->, thick] (2.5, -1.25) -- (2.5, -0.5) node[above] {Sortie $\text{K}^+$};
    
    % Canaux Na+ (fermés au repos)
    \draw[poporange, thick, fill=poporange!20, dashed] (7, -0.2) rectangle (8, -2.3);
    \node[poporange, align=center] at (7.5, -1.25) {Canal $\text{Na}^+$\\voltage-dépendant (FERMÉ)};
    \draw[poporange, ->, thick, dashed] (7.5, -1.25) -- (7.5, -2.0);
    \node[popmuted, font=\tiny] at (7.5, -2.2) {Perméabilité quasi nulle};
    
    % Pompe Na+/K+ ATPase
    \draw[popgold, thick, fill=popgold!20, rounded corners] (5, -0.2) rectangle (6, -2.3);
    \node[popgold, align=center] at (5.5, -1.25) {Pompe $\text{Na}^+/\text{K}^+$\\ATPase};
    \draw[poporange, ->, thick] (5.5, -1.25) -- (5.5, -0.5) node[above] {3 $\text{Na}^+$ out};
    \draw[popblue, ->, thick] (5.5, -1.25) -- (5.5, -2.0) node[below] {2 $\text{K}^+$ in};
    \node[popmuted, font=\tiny] at (5.5, 0.4) {Consomme 1 ATP};
    
    % Charges de surface
    \foreach \x in {0.5, 1.5, ..., 9.5} {
        \node[poporange, font=\Large] at (\x, 0.1) {$+$};
        \node[popblue, font=\Large] at (\x, -2.6) {$-$};
    }
    \node[align=center, font=\bfseries\small] at (5, -5.5) {Séparation des charges : Excès de $+$ dehors, excès de $-$ dedans $\Rightarrow$ $V_m < 0$};
\end{tikzpicture}
\legende{Schéma interprétatif de la membrane au repos. \textbf{1.} Canaux $\text{K}^+$ de fuite ouverts $\to$ sortie de $\text{K}^+$ down son gradient. \textbf{2.} Canaux $\text{Na}^+$ voltage-dépendants fermés $\to$ entrée de $\text{Na}^+$ quasi nulle. \textbf{3.} Anions organiques ($\text{A}^-$) piégés à l'intérieur. \textbf{4.} Pompe $\text{Na}^+/\text{K}^+$ ATPase active (3 $\text{Na}^+$ out / 2 $\text{K}^+$ in) qui entretient les gradients à long terme. L'accumulation de charges positives à l'extérieur et négatives à l'intérieur crée le champ électrique transmembranaire.}
\end{popfigure}

\textbf{Le rôle indispensable de la pompe $\text{Na}^+/\text{K}^+$ ATPase}\par

Si la membrane était seulement perméable au $\text{K}^+$, les ions $\text{K}^+$ sortiraient jusqu'à atteindre l'équilibre de Nernst ($E_K$), mais les concentrations finiraient par s'égaliser à très long terme (perte des gradients). La vie cellulaire nécessite le maintien actif des gradients.

\begin{propriete}[Pompe $\text{Na}^+/\text{K}^+$ ATPase (Savoir admis, complément utile au Bac)]
La \cle{pompe $\text{Na}^+/\text{K}^+$ ATPase} est une protéine transmembranaire à activité enzymatique (ATPase) qui utilise l'énergie de l'hydrolyse d' \textbf{1 ATP} pour transporter activement :
\begin{itemize}
    \item \textbf{3 ions $\text{Na}^+$} vers l'extérieur (contre leur gradient électrochimique),
    \item \textbf{2 ions $\text{K}^+$} vers l'intérieur (contre leur gradient chimique).
\end{itemize}
C'est un transport \textbf{actif primaire} (électrogène car il déplace 3 charges positives vers l'extérieur pour 2 vers l'intérieur $\to$ contribution directe d'environ -3 à -5 mV au $V_m$). Elle est inhibée par l'ouabaïne (glycoside cardiaque) et nécessite $\text{Mg}^{2+}$.
\end{propriete}

\begin{aretenir}[Synthèse : Origine du potentiel de repos]
Le potentiel de repos ($\approx -70$ mV) résulte de la \textbf{conjonction de trois facteurs} :
\begin{enumerate}
    \item \textbf{Gradients de concentration} maintenus par la pompe $\text{Na}^+/\text{K}^+$ ATPase ($[\text{K}^+]_i \gg [\text{K}^+]_e$ ; $[\text{Na}^+]_e \gg [\text{Na}^+]_i$).
    \item \textbf{Perméabilité différentielle} de la membrane au repos : \textbf{Élevée pour $\text{K}^+$} (canaux de fuite ouverts), \textbf{Très faible pour $\text{Na}^+$} (canaux voltage-dépendants fermés), intermédiaire pour $\text{Cl}^-$.
    \item \textbf{Présence d'anions non diffusibles} ($\text{A}^-$) à l'intérieur qui contrebalancent la sortie du $\text{K}^+$ et fixent la négativité interne.
\end{enumerate}
$\Rightarrow$ Le $V_m$ se stabilise proche du potentiel d'équilibre du $\text{K}^+$ ($E_K$), légèrement dépolarisé par la fuite entrante de $\text{Na}^+$.
\end{aretenir}

\textbf{Exercice résolu : Analyse d'une courbe $V_m = f([\text{K}^+]_e)$}\par

\begin{exoresolu}[Détermination de la concentration intracellulaire en $\text{K}^+$]
\textbf{Énoncé :} On a réalisé l'expérience de variation de $[\text{K}^+]_e$ sur un axone géant de calmar à \SI{20}{\celsius}. On a tracé la courbe théorique du potentiel d'équilibre du potassium $E_K = 58 \log_{10}([\text{K}^+]_e / [\text{K}^+]_i)$ pour trois valeurs hypothétiques de $[\text{K}^+]_i$ : 100 mmol/L, 200 mmol/L et 400 mmol/L. Les points expérimentaux (cercles) sont superposés.
\begin{center}
\begin{tikzpicture}
\begin{axis}[
    width=0.9\linewidth, height=7cm,
    xmode=log, log basis x=10,
    xlabel={$[\text{K}^+]_e$ (mmol/L)}, ylabel={$V_m$ (mV)},
    xmin=5, xmax=500, ymin=-100, ymax=10,
    grid=both, minor tick num=1,
    xtick={10,20,40,100,200,400},
    xticklabels={10,20,40,100,200,400},
    ytick={-100,-80,-60,-40,-20,0},
    axis lines=middle,
    legend style={at={(0.98,0.98)}, anchor=north east, fill=white, draw=popline, font=\small},
]
% Courbes théoriques
\addplot[popmuted, dashed, domain=10:400, samples=50] {58 * (log10(x/100))}; \addlegendentry{$[\text{K}^+]_i = 100$ mM}
\addplot[popmuted, dashed, domain=10:400, samples=50] {58 * (log10(x/200))}; \addlegendentry{$[\text{K}^+]_i = 200$ mM}
\addplot[popmuted, dashed, domain=10:400, samples=50] {58 * (log10(x/400))}; \addlegendentry{$[\text{K}^+]_i = 400$ mM}
% Points exp
\addplot[popcurve, only marks, mark=*, mark size=3pt] coordinates {
    (10, -95) (20, -70) (40, -50) (100, -25) (200, -10) (400, 0)
};
\addlegendentry{Expérimental}
\end{axis}
\end{tikzpicture}
\end{center}
\textbf{Question :} Quelle est la concentration intracellulaire en $\text{K}^+$ ($[\text{K}^+]_i$) de cet axone ? Justifier.

\tcblower
\textbf{Solution détaillée :}
\begin{enumerate}
    \item \textbf{Observation :} Les points expérimentaux (bleu) se superposent parfaitement à la courbe théorique tracée en pointillés pour $[\text{K}^+]_i = \mathbf{400 \text{ mM}}$ (courbe du bas).
    \item \textbf{Interprétation :} Le potentiel de repos $V_m$ suit la loi de Nernst pour le $\text{K}^+$ ($V_m \approx E_K$) sur toute la gamme de $[\text{K}^+]_e$ testée. Cela confirme que la membrane au repos se comporte comme une électrode sélective au $\text{K}^+$ (perméabilité quasi exclusive au $\text{K}^+$).
    \item \textbf{Vérification numérique (point physiologique) :} Pour $[\text{K}^+]_e = 20$ mM (milieu physiologique), la courbe $[\text{K}^+]_i = 400$ mM donne $E_K = 58 \log_{10}(20/400) = 58 \times (-1.301) \approx -75.5$ mV. La mesure expérimentale donne $-70$ mV. La légère différence (environ $+5$ mV) est due à la faible perméabilité résiduelle au $\text{Na}^+$ (équation de Goldman-Hodgkin-Katz).
    \item \textbf{Conclusion :} La concentration intracellulaire en potassium de cet axone de calmar est de \textbf{$\mathbf{[\text{K}^+]_i \approx 400 \text{ mM}}$}.
\end{enumerate}
\end{exoresolu}

\begin{exercice}[Application : Hyperkaliémie et excitabilité]
Un patient arrive aux urgences de l'Hôpital Central de Yaoundé avec une \textbf{hyperkaliémie sévère} ($[\text{K}^+]_e = 8,0 mmol/L$, normale $\approx 4,0 mmol/L$) due à une insuffisance rénale aiguë.
\begin{enumerate}
    \item En utilisant l'équation de Nernst ($E_K = 58 \log([\text{K}^+]_e/[\text{K}^+]_i)$) et en supposant $[\text{K}^+]_i = 150 mmol/L$ (valeur typique cellule mammifère), calculer $E_K$ pour $[\text{K}^+]_e = 4$ mM et pour $[\text{K}^+]_e = 8$ mM.
    \item Quel est l'effet de cette hyperkaliémie sur le potentiel de repos des cellules excitables (neurones, cardiomyocytes) ? Utiliser le terme exact : \textbf{dépolarisation} ou \textbf{hyperpolarisation} ?
    \item Expliquer pourquoi cette situation est \textbf{grave} pour la conduction cardiaque (risque d'arythmie, fibrillation ventriculaire) et la transmission neuromusculaire (paralysie flasque). Relier à la notion de \textbf{seuil d'excitation} et d'\textbf{inactivation des canaux $\text{Na}^+$}.
\end{enumerate}
\end{exercice}

\begin{corrige}[Exercice : Hyperkaliémie et excitabilité]
\begin{enumerate}
    \item \textbf{Calculs :}
    \begin{align*}
    E_K([\text{K}^+]_e=4\text{ mM}) &= 58 \log_{10}(4/150) = 58 \log_{10}(0.0267) = 58 \times (-1.574) \approx \mathbf{-91.3 \text{ mV}} \\
    E_K([\text{K}^+]_e=8\text{ mM}) &= 58 \log_{10}(8/150) = 58 \log_{10}(0.0533) = 58 \times (-1.273) \approx \mathbf{-73.8 \text{ mV}}
    \end{align*}
    \item \textbf{Effet :} Le potentiel de repos $V_m$ (qui suit $E_K$) passe d'environ \textbf{-91 mV à -74 mV}. La valeur absolue diminue : l'intérieur devient \textbf{moins négatif}. C'est une \textbf{\cle{dépolarisation}} de la membrane au repos (d'environ $+17$ mV).
    \item \textbf{Gravité physiopathologique :}
    \begin{itemize}
        \item \textbf{Cœur :} La dépolarisation au repos rapproche $V_m$ du seuil d'activation des canaux $\text{Na}^+$ voltage-dépendants (seuil $\approx -55$ mV). Cela augmente l'excitabilité de manière anarchique (extrasystoles). Mais surtout, une dépolarisation \textbf{prolongée} maintient les canaux $\text{Na}^+$ dans un état \textbf{inactivé} (fermés de façon réfractaire), empêchant la naissance de potentiels d'action normaux $\to$ bloc de conduction, fibrillation ventriculaire, arrêt cardiaque.
        \item \textbf{Jonction neuromusculaire / Muscle :} La dépolarisation au repos inactivera les canaux $\text{Na}^+$ des fibres musculaires. Le potentiel d'action nerveux arrive, libère l'acétylcholine, ouvre les canaux $\text{Na}^+$ post-synaptiques, mais ceux-ci sont inactivés $\to$ pas de potentiel d'action musculaire $\to$ \textbf{paralysie flasque} (le motard de Bafoussam ne sent plus et ne bouge plus ses jambes si sa moelle est comprimée, mais ici c'est une cause métabolique).
    \end{itemize}
    \textbf{Traitement d'urgence :} Injection de gluconate de calcium (stabilise la membrane en augmentant le seuil), insuline + glucose (fait rentrer le $\text{K}^+$ dans les cellules), dialyse.
\end{enumerate}
\end{corrige}

\coursec{Le potentiel d'action : enregistrement, phases et mécanismes ioniques}

Dans la section précédente, nous avons établi que le neurone au repos maintient une différence de potentiel stable d'environ \SI{-70}{mV} entre l'intérieur et l'extérieur de sa membrane, grâce à la pompe \ce{Na+}/\ce{K+}-ATPase et à la perméabilité sélective de la membrane au \ce{K+}. Mais ce potentiel de repos n'est pas une fin en soi : il est la condition préalable à l'activité nerveuse. Comment le neurone utilise-t-il cette réserve d'énergie électrochimique pour produire un signal ? C'est ce que nous allons découvrir : le \cle{potentiel d'action}, véritable « étincelle électrique » du système nerveux.

\courssub{Mise en évidence expérimentale du potentiel d'action}

\textbf{Observation et question.} Un muscle de grenouille se contracte lorsqu'on stimule électriquement le nerf qui l'innerve. Que se passe-t-il, au niveau électrique, dans la fibre nerveuse stimulée ? Le potentiel de repos est-il modifié ?

\begin{experience}[Mise en évidence du potentiel d'action]
\textbf{Matériel :} une fibre nerveuse géante (axone géant de calmar, diamètre $\approx$ \SI{0,5}{mm}, ou nerf sciatique de grenouille), une microélectrode intracellulaire reliée à un oscilloscope (ou à un voltmètre à déviation rapide), une électrode de référence dans le milieu extracellulaire, et deux électrodes de stimulation délivrant des chocs électriques brefs (\SI{0,1}{ms}) d'intensité réglable.

\textbf{Protocole :}
\begin{enumerate}
\item On enregistre d'abord le potentiel de repos : l'oscilloscope affiche une valeur stable de \SI{-70}{mV}.
\item On applique ensuite une stimulation électrique brève sur l'axone, à quelques millimètres de l'électrode d'enregistrement.
\item On fait varier l'intensité de la stimulation et on observe la réponse enregistrée.
\end{enumerate}

\textbf{Observations :}
\begin{itemize}
\item Si la stimulation est trop faible (\cle{stimulation subliminale}), on n'enregistre qu'une petite dépolarisation locale qui s'éteint : pas de signal propagé.
\item À partir d'une certaine intensité (\cle{seuil}), apparaît une variation brève, transitoire et de grande amplitude du potentiel membranaire : le potentiel passe de \SI{-70}{mV} à environ \SI{+30}{mV} puis revient à sa valeur initiale en 1 à \SI{2}{ms}. C'est le \cle{potentiel d'action} (PA).
\item Si l'on augmente encore l'intensité de la stimulation au-delà du seuil, l'amplitude du PA \textbf{ne change pas} : seule sa fréquence peut varier.
\end{itemize}

\textbf{Interprétation :} le PA est une réponse stéréotypée de la membrane excitable, déclenchée uniquement si la stimulation atteint le seuil. Son amplitude constante suggère un mécanisme « tout ou rien ».
\end{experience}

\begin{definition}[Potentiel d'action]
Le \cle{potentiel d'action} (PA) est une variation brève et transitoire du potentiel membranaire, caractérisée par une inversion de polarité (l'intérieur de la fibre devient momentanément positif par rapport à l'extérieur), qui se propage sans s'affaiblir le long de la fibre nerveuse. C'est le \cle{signal élémentaire du message nerveux}.
\end{definition}

\courssub{Description de la courbe du potentiel d'action}

L'enregistrement oscillographique du PA fait apparaître une courbe caractéristique, dont on distingue cinq phases :

\begin{enumerate}
\item \textbf{Le potentiel de repos} : avant la stimulation, le potentiel vaut environ \SI{-70}{mV}.
\item \textbf{La dépolarisation} : le potentiel s'élève brutalement, franchit \SI{0}{mV} et atteint un pic de $+30$ à \SI{+40}{mV}. Il y a \cle{inversion de polarité} : l'intérieur devient positif. Cette phase dure environ \SI{0,5}{ms}.
\item \textbf{La repolarisation} : le potentiel redescend rapidement vers des valeurs négatives.
\item \textbf{L'hyperpolarisation} : le potentiel devient transitoirement plus négatif que le potentiel de repos (environ \SI{-80}{mV}).
\item \textbf{Le retour au potentiel de repos} : la membrane retrouve sa valeur de \SI{-70}{mV}.
\end{enumerate}

La durée totale du phénomène est de 1 à \SI{2}{ms} : c'est un événement extrêmement bref.

\begin{popfigure}
\begin{center}
\begin{tikzpicture}
\begin{axis}[
width=0.95\linewidth, height=7cm,
xlabel={Temps (ms)}, ylabel={Potentiel membranaire (mV)},
xmin=0, xmax=6, ymin=-100, ymax=60,
xtick={0,1,2,3,4,5,6},
ytick={-80,-70,0,30},
grid=both, grid style={popline!40},
axis line style={popdark},
legend style={draw=none, fill=none}
]
% Repos
\addplot[popcurve, domain=0:1, samples=20]{-70};
% Dépolarisation
\addplot[popcurve, domain=1:1.5, samples=60]{-70 + 105*(1 - exp(-(x-1)*12))};
% Repolarisation + Hyperpolarisation
\addplot[popcurve, domain=1.5:2.6, samples=80]{35 - 115*(1 - exp(-(x-1.5)*3.2))};
% Retour hyperpolarisation vers repos
\addplot[popcurve, domain=2.6:4.2, samples=60]{-80 + 10*(1 - exp(-(x-2.6)*2.5))};
% Repos final
\addplot[popcurve, domain=4.2:6, samples=20]{-70};
% Annotations
\node[popdark, font=\small, align=center] at (axis cs:0.5,-60) {Repos \SI{-70}{mV}};
\node[poporange, font=\small\bfseries] at (axis cs:1.35,10) {Dépolarisation};
\draw[poporange, ->, thick] (axis cs:1.25,-20) -- (axis cs:1.4,25);
\node[popblue, font=\small\bfseries] at (axis cs:1.75,48) {Pic $\approx$ \SI{+30}{mV}};
\node[popgreen, font=\small\bfseries] at (axis cs:2.3,15) {Repolarisation};
\draw[popgreen, ->, thick] (axis cs:2.05,20) -- (axis cs:2.3,-40);
\node[poppurple, font=\small\bfseries] at (axis cs:3.6,-92) {Hyperpol. \SI{-80}{mV}};
\draw[poppurple, ->, thick] (axis cs:3.4,-75) -- (axis cs:3.2,-80);
% Mouvements ioniques
\node[popink, font=\footnotesize, align=center] at (axis cs:1.3,-45) {Entrée massive\\ de \ce{Na+}};
\node[popink, font=\footnotesize, align=center] at (axis cs:2.5,-55) {Sortie de \ce{K+}};
\end{axis}
\end{tikzpicture}
\end{center}
\legende{Courbe du potentiel d'action enregistrée par microélectrode intracellulaire. La dépolarisation (entrée de \ce{Na+}) porte le potentiel de \SI{-70}{mV} à environ \SI{+30}{mV} ; la repolarisation (sortie de \ce{K+}) le ramène vers le potentiel de repos, avec une hyperpolarisation transitoire. Durée totale : 1 à \SI{2}{ms}.}
\end{popfigure}

\begin{methode}[Décrire une courbe de potentiel d'action au Bac]
Face à une courbe de PA, procède toujours ainsi :
\begin{enumerate}
\item \textbf{Nomme} chaque phase dans l'ordre chronologique : repos, dépolarisation, pic, repolarisation, hyperpolarisation, retour au repos.
\item \textbf{Donne les valeurs} : potentiel de repos $\approx$ \SI{-70}{mV}, pic $\approx$ $+30$ à \SI{+40}{mV}, hyperpolarisation $\approx$ \SI{-80}{mV}, durée totale $\approx$ 1 à \SI{2}{ms}.
\item \textbf{Associe} chaque phase à un mouvement ionique précis : dépolarisation = entrée de \ce{Na+} ; repolarisation = sortie de \ce{K+} ; hyperpolarisation = fermeture tardive des canaux \ce{K+}.
\item \textbf{Conclus} : le PA est une inversion brève de polarité due à des variations séquentielles de perméabilité membranaire.
\end{enumerate}
\end{methode}

\courssub{Le seuil d'excitation, la loi du tout ou rien et la relation intensité-durée}

Reprenons l'expérience de stimulation en faisant varier l'intensité du choc électrique. Les résultats sont les suivants :

\begin{center}
\begin{tabularx}{\linewidth}{|l|X|X|}
\hline
\rowcolor{popblueL} \textbf{Stimulation} & \textbf{Réponse observée (fibre unique)} & \textbf{Conclusion} \\
\hline
Subliminaire (sous le seuil) & Dépolarisation locale faible, non propagée, proportionnelle à la stimulation & Pas de PA \\
\hline
Liminale (au seuil) & PA d'amplitude maximale ($\approx$ \SI{100}{mV} de variation) & Déclenchement du PA \\
\hline
Supraliminaire (au-dessus du seuil) & PA d'amplitude identique, quelle que soit l'intensité & Amplitude constante \\
\hline
\end{tabularx}
\end{center}

\begin{propriete}[Loi du tout ou rien (fibre unique)]
Le potentiel d'action obéit à la \cle{loi du tout ou rien} : si la stimulation n'atteint pas le seuil, il n'y a \textbf{aucun} potentiel d'action ; si la stimulation atteint ou dépasse le seuil, le PA apparaît avec son amplitude \textbf{maximale et constante}, indépendante de l'intensité de la stimulation supraliminaire. L'information sur l'intensité d'un stimulus n'est donc pas codée par l'amplitude du PA, mais par la \textbf{fréquence} des PA (codage fréquentiel) et, pour un nerf entier, par le \textbf{nombre de fibres recrutées} (codage spatial).
\end{propriete}

\begin{attention}[Fibre unique vs Nerf entier : ne pas confondre !]
Au Bac, distingue clairement :
\begin{itemize}
\item \textbf{Fibre unique (axone) :} loi du tout ou rien stricte. Amplitude constante. Seuil unique.
\item \textbf{Nerf (faisceau de fibres) :} pas de loi du tout ou rien globale. L'amplitude de la réponse globale (potentiel d'action composé) augmente avec l'intensité du stimulus car on recrute davantage de fibres (fibres de gros diamètre = seuil bas ; fibres de petit diamètre = seuil haut). Le « seuil du nerf » est l'intensité minimale pour recruter les fibres les plus excitables.
\end{itemize}
\end{attention}

\begin{exemplebox}[Une analogie pour comprendre]
Le PA ressemble au tir d'un pistolet : il faut appuyer sur la gâchette avec une force minimale (le seuil) pour que le coup parte ; mais une fois le coup parti, appuyer plus fort ne rend pas la balle plus rapide. Le coup est « tout ou rien ». De même, quand tu touches délicatement puis fermement la braise d'un feu de bois au village, tes fibres sensorielles émettent des PA de même amplitude, mais de fréquence croissante : c'est la fréquence qui informe ton cerveau de l'intensité de la douleur.
\end{exemplebox}

\paragraph{Relation intensité-durée (Chronaxie et Rhéobase).}
Le seuil d'excitation dépend non seulement de l'intensité $I$ du courant, mais aussi de sa durée $t$. Pour déclencher un PA, la charge électrique $Q = I \times t$ doit dépasser une valeur critique.

\begin{experience}[Détermination de la chronaxie et de la rhéobase]
On stimule une fibre nerveuse par des impulsions carrées de durée variable $t$ et on mesure l'intensité minimale $I$ nécessaire pour déclencher un PA (seuil).
\begin{itemize}
\item Pour des durées longues ($t > 1$ ms), l'intensité seuil tend vers une valeur minimale : la \cle{rhéobase} ($I_{rh}$), intensité minimale pour exciter avec une impulsion de durée « infinie ».
\item Pour des durées très courtes, l'intensité seuil augmente fortement.
\item La \cle{chronaxie} ($t_{ch}$) est la durée d'impulsion nécessaire pour exciter avec une intensité égale à \textbf{deux fois la rhéobase} ($2 I_{rh}$). C'est une caractéristique de l'excitabilité de la fibre (chronaxie courte = fibre très excitable, ex. fibre motrice $\approx$ \SI{0,1}{ms} ; chronaxie longue = fibre peu excitable, ex. fibre cardiaque $\approx$ \SI{3}{ms}).
\end{itemize}
\end{experience}

\begin{popfigure}
\begin{center}
\begin{tikzpicture}
\begin{axis}[
width=0.9\linewidth, height=6cm,
xlabel={Durée de l'impulsion $t$ (ms)}, ylabel={Intensité seuil $I$ (mA)},
xmin=0, xmax=3, ymin=0, ymax=5,
xtick={0,0.1,0.5,1,2,3},
ytick={0,1,2,3,4},
grid=both, grid style={popline!40},
axis line style={popdark},
legend style={at={(0.97,0.97)}, anchor=north east, draw=popline, font=\small}
]
% Courbe hyperbolique I = I_rh * (1 + t_ch / t)  approx
\addplot[popcurve, poporange, domain=0.05:3, samples=200]{1*(1 + 0.1/x)};
\addlegendentry{Courbe $I = f(t)$}
% Rheobase
\draw[popblue, dashed] (axis cs:0,1) -- (axis cs:3,1) node[right, popblue, font=\small]{Rhéobase $I_{rh}$};
% Chronaxie
\draw[popgreen, dashed] (axis cs:0.1,0) -- (axis cs:0.1,2) node[above, popgreen, font=\small]{Chronaxie $t_{ch}$};
\draw[popgreen, dashed] (axis cs:0,2) -- (axis cs:0.1,2);
\draw[popdark, <->] (axis cs:0,1.2) -- (axis cs:0.1,1.2) node[midway, above, font=\small]{$t_{ch}$};
\draw[popdark, <->] (axis cs:-0.1,1) -- (axis cs:-0.1,2) node[midway, left, font=\small]{$I_{rh}$};
\end{axis}
\end{tikzpicture}
\end{center}
\legende{Courbe de la relation intensité-durée (courbe de Weiss-Lapicque). La rhéobase $I_{rh}$ est l'asymptote horizontale. La chronaxie $t_{ch}$ est l'abscisse du point d'ordonnée $2 I_{rh}$.}
\end{popfigure}

\begin{methode}[Exploiter une courbe Intensité-Durée]
\begin{enumerate}
\item Identifie l'asymptote horizontale : c'est la \textbf{rhéobase} $I_{rh}$ (unité : mA ou $\mu$A).
\item Repère l'ordonnée $2 I_{rh}$, trace l'horizontale jusqu'à la courbe, puis la verticale vers l'axe des abscisses : c'est la \textbf{chronaxie} $t_{ch}$ (unité : ms).
\item Compare : une fibre de chronaxie plus courte est \textbf{plus excitable} (nécessite une impulsion plus brève pour être excitée à $2 I_{rh}$).
\item Application clinique : le stimulateur cardiaque (pacemaker) utilise des impulsions de durée adaptée à la chronaxie du myocarde ($\approx$ \SI{1}{ms}) pour économiser la batterie.
\end{enumerate}
\end{methode}

\courssub{Les mécanismes ioniques du potentiel d'action}

\textbf{Question scientifique.} Quels mouvements ioniques expliquent une variation aussi rapide et aussi ample du potentiel membranaire ?

\begin{experience}[Les expériences de Hodgkin et Huxley : mesure des conductances (Voltage-clamp)]
Entre 1949 et 1952, Alan Hodgkin et Andrew Huxley (Prix Nobel 1963) ont utilisé la technique du \cle{voltage-clamp} (« potentiel imposé ») sur l'axone géant du calmar. Le principe : on maintient artificiellement le potentiel membranaire à une valeur fixée ($V_m$ constant) et l'on mesure le courant ionique $I_m$ que l'appareil doit injecter pour compenser les courants membranaires. On en déduit la \cle{conductance} membranaire ($g = I / (V_m - E_{eq})$), qui reflète le degré d'ouverture des canaux ioniques.

\textbf{Résultats majeurs :}
\begin{itemize}
\item Lors d'une dépolarisation imposée, la conductance au sodium $g_{\ce{Na}}$ augmente \textbf{d'abord}, très rapidement (en moins de \SI{0,5}{ms}), puis retombe vite : elle est \textbf{transitoire} (pic puis inactivation).
\item La conductance au potassium $g_{\ce{K}}$ augmente \textbf{ensuite}, plus lentement, et de façon \textbf{durable} (plusieurs millisecondes) : pas d'inactivation rapide.
\end{itemize}

\textbf{Interprétation :} la membrane possède deux types de canaux voltage-dépendants distincts : des canaux \ce{Na+} à ouverture rapide suivie d'inactivation (état « fermé-inactivé »), et des canaux \ce{K+} à ouverture retardée sans inactivation rapide.
\end{experience}

\begin{popfigure}
\begin{center}
\begin{tikzpicture}
\begin{axis}[
width=0.95\linewidth, height=6.5cm,
xlabel={Temps (ms)}, ylabel={Conductance relative (a.u.)},
xmin=0, xmax=5, ymin=0, ymax=1.3,
xtick={0,1,2,3,4,5},
grid=both, grid style={popline!40},
axis line style={popdark},
legend style={at={(0.97,0.97)}, anchor=north east, draw=popline, font=\small}
]
% g_Na : gaussienne centrée vers 1.35, largeur ~0.35
\addplot[popcurve, poporange, domain=0.5:5, samples=200]{1.2*exp(-((x-1.35)*(x-1.35))/0.12)};
\addlegendentry{$g_{\ce{Na}}$ (conductance \ce{Na+})}
% g_K : monté sigmoïde puis descente lente
\addplot[popcurve, popgreen, domain=0.5:5, samples=200]{0.9*(1/(1+exp(-(x-1.2)*5)))*exp(-(x-1.2)*0.3)};
\addlegendentry{$g_{\ce{K}}$ (conductance \ce{K+})}
\draw[popline, dashed] (axis cs:1,0) -- (axis cs:1,1.25) node[above, font=\footnotesize, popdark]{Stimulation};
\end{axis}
\end{tikzpicture}
\end{center}
\legende{Variations des conductances membranaires $g_{\ce{Na}}$ et $g_{\ce{K}}$ au cours du potentiel d'action (d'après Hodgkin et Huxley). La conductance \ce{Na+} (orange) augmente la première, rapidement et transitoirement (inactivation) ; la conductance \ce{K+} (verte) augmente ensuite, plus lentement et durablement.}
\end{popfigure}

\textbf{Interprétation ionique détaillée des phases du PA.}

\begin{itemize}
\item \textbf{Dépolarisation (montée) :} La stimulation seuil (dépolarisation locale $\approx$ \SI{-55}{mV}) provoque l'ouverture des \cle{canaux \ce{Na+} voltage-dépendants} (canaux « rapides »). Le gradient électrochimique étant très favorable (concentration \ce{Na+}$_{\text{ext}}$ $\approx$ 10-15 $\times$ \ce{Na+}$_{\text{int}}$ ; $V_m$ négatif), les ions \ce{Na+} \textbf{entrent massivement} dans la fibre (courant entrant). L'intérieur devient positif : c'est l'inversion de polarité, jusqu'au pic ($+30$ à \SI{+40}{mV}), proche du potentiel d'équilibre de Nernst pour le sodium ($E_{\ce{Na}}$).
\item \textbf{Inactivation des canaux \ce{Na+} :} Au pic du PA, les canaux \ce{Na+} s'\textbf{inactivent} spontanément : une « boule » cytoplasmique (boucle d'inactivation) bouche le pore. Ils deviennent \textbf{fermés et inouvrables} (état réfractaire absolu). L'entrée de \ce{Na+} cesse brutalement.
\item \textbf{Repolarisation (descente) :} Simultanément, les \cle{canaux \ce{K+} voltage-dépendants} (canaux « lents » ou « retardés »), sensibles à la dépolarisation, s'ouvrent. Les ions \ce{K+}, très concentrés à l'intérieur, \textbf{sortent} de la fibre (courant sortant), entraînant des charges positives vers l'extérieur : le potentiel redevient négatif.
\item \textbf{Hyperpolarisation (potentiel de suite) :} Les canaux \ce{K+} se ferment \textbf{tardivement} (pas d'inactivation rapide) ; la sortie de \ce{K+} se poursuit donc un peu trop longtemps après le retour au potentiel de repos, et le potentiel descend transitoirement en dessous de \SI{-70}{mV} (environ \SI{-80}{mV}), se rapprochant du potentiel d'équilibre du potassium ($E_{\ce{K}}$).
\item \textbf{Retour au repos :} Tous les canaux voltage-dépendants se referment (canaux \ce{K+} fermés, canaux \ce{Na+} revenus à l'état « fermé-ouvrable »). La pompe \ce{Na+}/\ce{K+}-ATPase et les fuites passives de \ce{K+} (canaux de fuite) rétablissent et maintiennent le potentiel de repos à \SI{-70}{mV}.
\end{itemize}

\begin{attention}[Erreurs fatales au Bac : mouvements ioniques et rôle de la pompe]
\begin{itemize}
\item \textbf{Dépolarisation = ENTRÉE de \ce{Na+}} (charges $+$ entrent $\Rightarrow$ $V_m$ augmente). JAMAIS « sortie de \ce{Na+} ».
\item \textbf{Repolarisation = SORTIE de \ce{K+}} (charges $+$ sortent $\Rightarrow$ $V_m$ diminue). JAMAIS « entrée de \ce{K+} » ni « sortie de \ce{Na+} ».
\item La pompe \ce{Na+}/\ce{K+}-ATPase \textbf{ne crée pas le PA}. Elle consomme de l'ATP pour maintenir les gradients de concentration à long terme (3 \ce{Na+} sortent, 2 \ce{K+} entrent). Les mouvements du PA sont \textbf{passifs} (diffusion le long des gradients électrochimiques) à travers des canaux voltage-dépendants.
\item L'inactivation des canaux \ce{Na+} (boule) est distincte de leur fermeture (porte d'activation). Inactivé = fermé \textbf{et} inouvrable tant que $V_m$ n'est pas revenu au repos.
\end{itemize}
\end{attention}

\begin{popfigure}
\begin{center}
\begin{tikzpicture}[font=\small, >=Stealth, scale=0.9]
% Canal Na+ fermé (repos)
\begin{scope}[shift={(0,0)}]
\draw[fill=popcream, rounded corners] (-1.1,-0.7) rectangle (1.1,0.7);
\draw[fill=poporange, rounded corners=2pt] (-0.7,-0.7) rectangle (-0.3,0.3);
\draw[fill=poporange, rounded corners=2pt] (0.3,-0.7) rectangle (0.7,0.3);
\node[poporange, font=\footnotesize\bfseries] at (0,-1.0) {Canal \ce{Na+} fermé};
\node[font=\footnotesize, popmuted] at (0,0.55) {Repos (\SI{-70}{mV})};
\draw[->, popdark, thick] (0,0.9) -- (0,1.3) node[above, font=\tiny]{Dépolarisation seuil};
\end{scope}
% Canal Na+ ouvert (dépolarisation)
\begin{scope}[shift={(4,0)}]
\draw[fill=popcream, rounded corners] (-1.1,-0.7) rectangle (1.1,0.7);
\draw[fill=poporange, rounded corners=2pt] (-0.7,-0.7) rectangle (-0.3,0.6);
\draw[fill=poporange, rounded corners=2pt] (0.3,-0.7) rectangle (0.7,0.6);
\draw[->, very thick, popblue] (0,0.9) -- (0,-1.0);
\node[popblue, font=\footnotesize] at (0.85,0.1) {\ce{Na+}};
\node[poporange, font=\footnotesize\bfseries] at (0,-1.3) {Canal \ce{Na+} OUVERT};
\node[font=\footnotesize, popmuted] at (0,0.8) {Dépolarisation};
\end{scope}
% Canal Na+ inactivé (pic)
\begin{scope}[shift={(8,0)}]
\draw[fill=popcream, rounded corners] (-1.1,-0.7) rectangle (1.1,0.7);
\draw[fill=poporange, rounded corners=2pt] (-0.7,-0.7) rectangle (-0.3,0.6);
\draw[fill=poporange, rounded corners=2pt] (0.3,-0.7) rectangle (0.7,0.6);
\draw[fill=poppink, rounded corners=2pt] (-0.2,-0.7) rectangle (0.2,-0.15);
\node[poppink, font=\footnotesize] at (0.85,-0.4) {Boule d'inactivation};
\node[poporange, font=\footnotesize\bfseries] at (0,-1.0) {Canal \ce{Na+} INACTIVÉ};
\node[font=\footnotesize, popmuted] at (0,0.8) {Pic du PA};
\end{scope}
% Canal K+ ouvert (repolarisation)
\begin{scope}[shift={(12,0)}]
\draw[fill=popcream, rounded corners] (-1.1,-0.7) rectangle (1.1,0.7);
\draw[fill=popgreen, rounded corners=2pt] (-0.7,-0.7) rectangle (-0.3,0.6);
\draw[fill=popgreen, rounded corners=2pt] (0.3,-0.7) rectangle (0.7,0.6);
\draw[->, very thick, poppurple] (0,-1.0) -- (0,0.9);
\node[poppurple, font=\footnotesize] at (0.85,0.1) {\ce{K+}};
\node[popgreen, font=\footnotesize\bfseries] at (0,-1.3) {Canal \ce{K+} OUVERT};
\node[font=\footnotesize, popmuted] at (0,0.8) {Repolarisation};
\end{scope}
\draw[->, thick, popdark] (1.3,0) -- (2.7,0);
\draw[->, thick, popdark] (5.3,0) -- (6.7,0);
\draw[->, thick, popdark] (9.3,0) -- (10.7,0);
\end{tikzpicture}
\end{center}
\legende{États conformationnels des canaux ioniques voltage-dépendants au cours du PA. Le canal \ce{Na+} : fermé-ouvrable (repos) $\to$ ouvert (dépolarisation, entrée \ce{Na+}) $\to$ inactivé (pic, boule, inouvrable). Le canal \ce{K+} : fermé (repos) $\to$ ouvert retardé (repolarisation, sortie \ce{K+}) $\to$ fermé (retour repos).}
\end{popfigure}

\courssub{Preuves pharmacologiques du rôle des canaux ioniques}

Pour valider l'interprétation ionique, on utilise des inhibiteurs sélectifs des canaux :

\begin{center}
\begin{tabularx}{\linewidth}{|l|X|X|}
\hline
\rowcolor{popblueL} \textbf{Substance} & \textbf{Cible moléculaire} & \textbf{Effet observé sur le PA} \\
\hline
Tétrodotoxine (TTX) & Bloque le pore des canaux \ce{Na+} voltage-dépendants (site extracellulaire) & Suppression totale de la dépolarisation et du PA (seuil infini) \\
\hline
Tétraéthylammonium (TEA) & Bloque le pore des canaux \ce{K+} voltage-dépendants (site intracellulaire) & Dépolarisation normale, mais abolition de la repolarisation : PA très prolongé (plateau) \\
\hline
Anesthésiques locaux (Lidocaïne, Novocaïne) & Bloquent les canaux \ce{Na+} voltage-dépendants (site intracellulaire, état-dépendants) & Blocage réversible de la conduction nerveuse (analgésie locale) \\
\hline
\end{tabularx}
\end{center}

Ces expériences confirment que le PA repose sur l'ouverture séquentielle de deux populations distinctes de canaux : les canaux \ce{Na+} pour la dépolarisation, les canaux \ce{K+} pour la repolarisation.

\begin{savaistu}[La tétrodotoxine, poison du poisson-globe et anesthésie locale]
La \cle{tétrodotoxine} (TTX) est un poison extrêmement puissant (LD50 $\approx$ \SI{10}{\micro g/kg}) présent dans le poisson-globe (\emph{Takifugu}), mets raffiné (fugu) mais dangereux de la cuisine japonaise. Quelques milligrammes suffisent à bloquer tous les canaux \ce{Na+} de l'organisme : plus aucun PA ne peut naître, les muscles — y compris le diaphragme — se paralysent, et la mort survient par arrêt respiratoire. Il n'existe pas d'antidote connu, seule la ventilation artificielle maintient la vie le temps de l'élimination.
À l'inverse, certains anesthésiques locaux utilisés quotidiennement dans les hôpitaux et cabinets dentaires camerounais, comme la \textbf{lidocaïne} ou la \textbf{novocaïne}, bloquent \textbf{transitoirement} et \textbf{réversiblement} les canaux \ce{Na+} voltage-dépendants (ils se fixent sur le site intracellulaire quand le canal est ouvert/inactivé). C'est ainsi qu'un dentiste à Yaoundé ou Douala supprime la douleur : en empêchant la naissance et la propagation des messages nerveux douloureux dans les fibres sensitives de la dent !
\end{savaistu}

\courssub{Bilan et exercice d'application}

\begin{aretenir}[L'essentiel sur le potentiel d'action]
\begin{itemize}
\item Le PA est une \textbf{inversion brève de la polarité membranaire} (de \SI{-70}{mV} à $\approx$ \SI{+30}{mV}), durant 1 à \SI{2}{ms}.
\item Il obéit à la \textbf{loi du tout ou rien} pour une fibre unique : amplitude constante dès que le seuil est atteint. Codage de l'intensité = fréquence des PA + recrutement spatial (nerf).
\item Le seuil dépend de la durée de stimulation : relation \textbf{intensité-durée} caractérisée par la \textbf{rhéobase} ($I_{rh}$) et la \textbf{chronaxie} ($t_{ch}$). Chronaxie courte = forte excitabilité.
\item Il résulte de \textbf{variations séquentielles de perméabilité} : ouverture rapide et transitoire des canaux \ce{Na+} (entrée de \ce{Na+} $\Rightarrow$ dépolarisation), inactivation des canaux \ce{Na+}, puis ouverture retardée et durable des canaux \ce{K+} (sortie de \ce{K+} $\Rightarrow$ repolarisation + hyperpolarisation).
\item Preuves expérimentales : voltage-clamp (Hodgkin \& Huxley, conductances) et pharmacologie (TTX, TEA, anesthésiques locaux).
\end{itemize}
\end{aretenir}

\begin{exoresolu}[Interprétation d'une expérience pharmacologique]
\textbf{Énoncé.} On enregistre le potentiel membranaire d'un axone géant de calmar stimulé par un choc électrique supraliminaire. Sur le témoin, on obtient un PA normal (pic à \SI{+30}{mV}, durée \SI{1,5}{ms}). On ajoute ensuite de la TTX dans le milieu extracellulaire : la même stimulation ne provoque plus aucune réponse. Enfin, sur un autre axone, on ajoute du TEA : la dépolarisation est normale, mais le potentiel reste bloqué autour de \SI{+20}{mV} pendant plusieurs dizaines de millisecondes avant de redescendre très lentement. Interprète ces résultats.
\tcblower
\textbf{Solution détaillée.}
\begin{enumerate}
\item \textbf{Témoin :} La stimulation supraliminaire dépolarise la membrane au-delà du seuil. Les canaux \ce{Na+} voltage-dépendants s'ouvrent rapidement $\Rightarrow$ entrée massive de \ce{Na+} $\Rightarrow$ dépolarisation jusqu'au pic (\SI{+30}{mV}). Au pic, les canaux \ce{Na+} s'inactivent et les canaux \ce{K+} voltage-dépendants s'ouvrent $\Rightarrow$ sortie de \ce{K+} $\Rightarrow$ repolarisation rapide, hyperpolarisation, retour au repos. PA normal.
\item \textbf{Avec TTX :} La tétrodotoxine se fixe sur le site extracellulaire des canaux \ce{Na+} voltage-dépendants et les bloque physiquement. Malgré la stimulation supraliminaire, les canaux \ce{Na+} ne peuvent pas s'ouvrir. Il n'y a \textbf{aucune entrée de \ce{Na+}}, donc \textbf{aucune dépolarisation} regenerative. Le PA est aboli. Cela prouve formellement que la phase montante du PA est due à l'entrée de \ce{Na+} par ces canaux spécifiques.
\item \textbf{Avec TEA :} Le TEA bloque les canaux \ce{K+} voltage-dépendants (par l'intérieur). La dépolarisation (ouverture canaux \ce{Na+}, entrée \ce{Na+}) se déroule normalement car les canaux \ce{Na+} sont intacts. Mais l'ouverture des canaux \ce{K+} est empêchée : \textbf{pas de sortie de \ce{K+}}, donc \textbf{pas de repolarisation} active. Le potentiel reste positif (plateau) jusqu'à ce que des mécanismes lents (fuites, pompe, inactivation complète \ce{Na+}) le fassent redescendre. Cela prouve que la repolarisation est due à la sortie de \ce{K+} par les canaux \ce{K+} voltage-dépendants.
\item \textbf{Conclusion globale :} Le potentiel d'action résulte bien de l'ouverture séquentielle et indépendante de deux populations de canaux voltage-dépendants : canaux \ce{Na+} (dépolarisation) puis canaux \ce{K+} (repolarisation).
\end{enumerate}
\end{exoresolu}

\begin{exercice}[Application : Chronaxie et excitabilité]
On étudie l'excitabilité de deux types de fibres nerveuses, A et B. On trace leurs courbes intensité-durée ($I$ en mA, $t$ en ms).
\begin{itemize}
\item Fibre A : Rhéobase $I_{rh} = \SI{0,5}{mA}$ ; Chronaxie $t_{ch} = \SI{0,1}{ms}$.
\item Fibre B : Rhéobase $I_{rh} = \SI{0,5}{mA}$ ; Chronaxie $t_{ch} = \SI{2,0}{ms}$.
\end{itemize}
\textbf{Questions :}
\begin{enumerate}
\item Laquelle de ces deux fibres est la plus excitable ? Justifie en te basant sur la définition de la chronaxie.
\item Pour une impulsion de durée \SI{0,5}{ms}, quelle fibre nécessite l'intensité de stimulation la plus faible ? Explique.
\item Quel type de fibre (myélinisée de gros diamètre ou amyélinisée de petit diamètre) correspond probablement à la fibre A ? Pourquoi ?
\end{enumerate}
\end{exercice}

\begin{corrige}[Exercice : Chronaxie et excitabilité]
\begin{enumerate}
\item \textbf{Fibre A est la plus excitable.} La chronaxie est la durée minimale pour exciter à $2 I_{rh}$. Une chronaxie courte (\SI{0,1}{ms} pour A vs \SI{2,0}{ms} pour B) signifie qu'une impulsion très brève suffit à déclencher le PA à intensité modérée. La fibre A répond donc à des stimulations plus brèves : elle est plus excitable.
\item \textbf{Fibre A.} Pour $t = \SI{0,5}{ms}$ ($> t_{ch,A}$ mais $< t_{ch,B}$), la fibre A est proche de sa rhéobase (intensité seuil $\approx I_{rh} = \SI{0,5}{mA}$). La fibre B, elle, nécessite encore une intensité $> 2 I_{rh} = \SI{1,0}{mA}$ car $t < t_{ch,B}$. Donc $I_{seuil,A} < I_{seuil,B}$.
\item \textbf{Fibre A = fibre myélinisée de gros diamètre (ex. fibre motrice $\alpha$).} Les fibres myélinisées de gros diamètre ont une constante de temps membranaire courte, une faible capacitance, et une densité élevée de canaux \ce{Na+} aux nœuds de Ranvier : elles sont très excitables (chronaxie courte, \SI{0,05}{ms} à \SI{0,2}{ms}). La fibre B correspond à une fibre amyélinisée de petit diamètre (ex. fibre C douleur, chronaxie \SI{0,5}{ms} à \SI{2}{ms}).
\end{enumerate}
\end{corrige}

Nous savons désormais comment naît et ce qu'est un potentiel d'action au niveau d'une fibre unique, et comment le seuil varie avec la durée de stimulation. Reste à comprendre comment ce signal élémentaire prend naissance dans les récepteurs sensoriels (potentiel de récepteur) et au niveau de la zone de déclenchement du neurone (cône d'émergence), puis comment il voyage le long de la fibre nerveuse sans jamais s'affaiblir : ce sera l'objet des sections suivantes.

\coursec{Naissance du message nerveux : récepteurs sensoriels et seuil d'excitation}

Dans la section précédente, nous avons établi que le potentiel d'action (PA) est un phénomène \textbf{tout ou rien} : une fois le seuil de dépolarisation atteint au niveau de la zone de déclenchement (cône d'émergence ou premier nœud de Ranvier), le PA se propage invariablement le long de l'axone sans déformation. Mais comment ce seuil est-il atteint au départ ? Comment une information physique venue du monde extérieur (lumière, pression, chaleur, molécule odorante) se transforme-t-elle en signal électrique codé en fréquence de PA ? C'est l'objet de cette section : nous allons suivre le cheminement du message nerveux depuis sa \textbf{genèse au niveau du récepteur sensoriel} jusqu'au \textbf{recrutement des fibres nerveuses} au sein d'un nerf, en passant par les lois fondamentales de l'excitabilité (seuil, relation intensité-durée).

\courssub{Du stimulus à l'influx nerveux : le rôle du récepteur sensoriel}

\textbf{Dispositif expérimental de référence : le corpuscule de Pacini.}\par
Pour comprendre la transduction sensorielle, les neurophysiologistes étudient souvent le \textbf{corpuscule de Pacini}, un mécanorécepteur de pression profonde et de vibrations hautes fréquences (présent dans la peau, les méninges, les articulations). C'est un modèle idéal car sa structure (lamelles de tissu conjonctif entourant la terminaison nerveuse) permet d'isoler l'effet du stimulus mécanique sur la membrane du neurone sensoriel.

\begin{experience}[Mise en évidence du potentiel de récepteur et du codage fréquentiel]
\textbf{Matériel :} Préparation isolée de corpuscule de Pacini (ex. chat, grenouille) maintenue in vitro ; microélectrode intracellulaire enfoncée dans la terminaison sensorielle ; électrode de référence extracellulaire ; système de stimulation mécanique contrôlée (sonde appliquant des pressions d'intensités croissantes) ; oscilloscope ou acquisition numérique.
\textbf{Protocole :}
\begin{enumerate}
    \item Enregistrer le potentiel de membrane au repos (environ \SI{-70}{mV}).
    \item Appliquer une stimulation mécanique de faible intensité (pression légère) pendant une durée fixe (ex. \SI{100}{ms}).
    \item Augmenter progressivement l'intensité de la pression (paliers de \SI{10}{g}, \SI{50}{g}, \SI{200}{g}, etc.) tout en enregistrant la réponse électrique.
    \item Mesurer l'amplitude de la dépolarisation initiale et, le cas échéant, compter le nombre de PA déclenchés par seconde (fréquence instantanée).
\end{enumerate}
\textbf{Observations (tracés typiques) :}
\begin{itemize}
    \item \textbf{Stimulus faible :} Une dépolarisation locale, \textbf{graduée} (amplitude proportionnelle à l'intensité du stimulus), \textbf{non propagée} (elle s'éteint spatialement), sans PA. C'est le \textbf{potentiel de récepteur} (ou potentiel générateur).
    \item \textbf{Stimulus modéré :} Le potentiel de récepteur atteint un niveau critique (\textbf{seuil}, env. \SI{-50}{mV}) au niveau du premier nœud de Ranvier. Un PA naît, puis une salve de PA à fréquence modérée.
    \item \textbf{Stimulus fort :} Potentiel de récepteur plus ample, dépassant largement le seuil. Salve de PA à \textbf{fréquence élevée}.
\end{itemize}
\textbf{Interprétation :} Le récepteur sensoriel est un \textbf{transducteur}. Il convertit une énergie physique (mécanique ici) en signal électrique (variation de perméabilité ionique $\rightarrow$ courant dépolarisant entrant $I_{Na^+}$ ou $I_{Ca^{2+}}$). Cette dépolarisation est locale et graduée car elle ne fait pas intervenir de canaux voltage-dépendants à ouverture coopérative (pas de boucle de rétroaction positive) mais des canaux \textbf{activés par le stimulus} (canaux mécanosensibles type Piezo). Ce n'est qu'au niveau de la zone de déclenchement (riche en canaux $Na^+$ voltage-dépendants) que le tout ou rien s'applique.
\legende{Expérience classique d'Adrian et Zotterman (1926) sur le nerf vague de grenouille ou préparation de Pacini : relation entre intensité du stimulus, amplitude du potentiel de récepteur et fréquence des PA.}
\end{experience}

\begin{popfigure}
\begin{tikzpicture}[scale=0.9]
\begin{axis}[
    width=\linewidth, height=0.6\linewidth,
    xlabel={Temps (ms)}, ylabel={Potentiel de membrane (mV)},
    xmin=0, xmax=300, ymin=-80, ymax=50,
    axis lines=middle, grid=major, grid style={popmuted},
    xtick={0,50,100,150,200,250,300},
    ytick={-70,-50,0,30},
    legend style={at={(0.02,0.98)}, anchor=north west, fill=popcream, draw=popline, font=\footnotesize},
    clip=false
]
% Stimulus faible : Potentiel de récepteur seul
\addplot[popblue, thick, domain=0:300, samples=300] 
    ({x}, {-70 + 25*exp(-(x-50)*(x-50)/2000) * (x>50 && x<250)});
\addlegendentry{Stimulus faible : Pot. récepteur seul}

% Stimulus modéré : Pot. récepteur + 2 PA
\addplot[poporange, thick, domain=0:300, samples=600] 
    ({x}, {-70 + 35*exp(-(x-50)*(x-50)/2000) * (x>50 && x<250) 
        + 100*exp(-(x-80)*(x-80)/50) * (x>75 && x<85)
        + 100*exp(-(x-180)*(x-180)/50) * (x>175 && x<185)});
\addlegendentry{Stimulus modéré : Pot. récepteur + 2 PA}

% Stimulus fort : Pot. récepteur + 5 PA
\addplot[popgreen, thick, domain=0:300, samples=900] 
    ({x}, {-70 + 45*exp(-(x-50)*(x-50)/2000) * (x>50 && x<250) 
        + 100*exp(-(x-70)*(x-70)/50) * (x>65 && x<75)
        + 100*exp(-(x-110)*(x-110)/50) * (x>105 && x<115)
        + 100*exp(-(x-150)*(x-150)/50) * (x>145 && x<155)
        + 100*exp(-(x-190)*(x-190)/50) * (x>185 && x<195)
        + 100*exp(-(x-230)*(x-230)/50) * (x>225 && x<235)});
\addlegendentry{Stimulus fort : Pot. récepteur + 5 PA}

% Seuil line
\draw[popred, dashed, thick] (0,-50) -- (300,-50) node[pos=0.95, right, font=\footnotesize] {Seuil ($\approx -50$ mV)};
% Repos line
\draw[popmuted, dashed] (0,-70) -- (300,-70) node[pos=0.95, right, font=\footnotesize] {Repos ($-70$ mV)};
% Stimulus bar
\draw[popdark, very thick] (50,-75) -- (250,-75) node[midway, below=2pt, font=\footnotesize] {Durée du stimulus mécanique};
\end{axis}
\end{tikzpicture}
\legende{Enregistrement intracellulaire au niveau d'un corpuscule de Pacini soumis à trois intensités de stimulation mécanique croissantes. On observe le potentiel de récepteur gradué (dépolarisation locale) qui, dès qu'il atteint le seuil au premier nœud de Ranvier, déclenche une salve de potentiels d'action dont la fréquence code l'intensité du stimulus.}
\end{popfigure}

\begin{definition}[Potentiel de récepteur (ou potentiel générateur)]
C'est une \textbf{dépolarisation locale, graduée et non propagée}, générée au niveau de la membrane d'un récepteur sensoriel soumis à un stimulus adéquat. Son amplitude est proportionnelle à l'intensité du stimulus (dans une certaine plage). Il résulte de l'ouverture de canaux ioniques \textbf{spécifiques à la modalité du stimulus} (canaux mécanosensibles Piezo pour le toucher, canaux TRP pour la température/chimique, cascade rhodopsine/GMPc pour la vision), entraînant un courant entrant dépolarisant (principalement $Na^+$, $Ca^{2+}$). Il ne suit pas la loi du tout ou rien.
\end{definition}

\begin{attention}[Piège majeur : Potentiel de récepteur $\neq$ Potentiel d'action]
Ne confondez jamais ces deux entités :
\begin{itemize}
    \item \textbf{Potentiel de récepteur :} Local (ne se propage pas), gradué (amplitude variable), pas de période réfractaire, somme algébrique possible, généré par canaux \emph{stimulus-dépendants}.
    \item \textbf{Potentiel d'action :} Propagé (le long de l'axone), tout ou rien (amplitude fixe), période réfractaire absolue et relative, non sommable, généré par canaux \emph{voltage-dépendants} ($Na^+_v$, $K^+_v$).
\end{itemize}
Le potentiel de récepteur n'est que le \textbf{générateur} de courant qui va charger la capacité membranaire du premier nœud de Ranvier jusqu'au seuil de déclenchement des PA.
\end{attention}

\courssub{Codage du message nerveux : la fréquence, langage universel}

Une fois le seuil franchi, l'information n'est plus portée par l'amplitude (fixe) mais par la \textbf{fréquence} des PA. C'est le principe fondamental du codage neural établi par Edgar Adrian (Prix Nobel 1932).

\begin{propriete}[Codage fréquentiel de l'intensité du stimulus]
Dans une fibre afferente donnée, \textbf{l'intensité du stimulus est codée par la fréquence des potentiels d'action} (nombre de PA par seconde, en Hz). L'amplitude du PA ne transporte \textbf{aucune information} sur l'intensité (loi du tout ou rien). La durée du stimulus est codée par la durée de la décharge (adaptation du récepteur). La nature (qualité) de la sensation (toucher, douleur, chaud) dépend de \textbf{l'identité de la voie nerveuse} (ligne étiquetée : loi de la projection spécifique de Müller) et de l'aire corticale qui la reçoit.
\end{propriete}

\begin{exemplebox}[Exemple concret : discrimination tactile au Cameroun]
Imaginez un cultivateur de cacao à \textbf{Ebolowa} qui manipule une cabosse.
\begin{itemize}
    \item \textbf{Toucher léger (effleurement) :} Stimule faiblement les corpuscules de Meissner (bas seuil). Fréquence de décharge : \textbf{$\approx 10$ à $20$ PA/s}.
    \item \textbf{Pression forte (soulever la cabosse) :} Recrute aussi les corpuscules de Pacini et de Ruffini (seuil plus élevé). Fréquence de décharge : \textbf{$\approx 80$ à $150$ PA/s} sur les mêmes fibres + recrutement de fibres à seuil plus élevé.
\end{itemize}
Le cortex somesthésique « lit » cette fréquence : $20$ Hz $\rightarrow$ « contact léger » ; $120$ Hz $\rightarrow$ « pression forte ». C'est ce qui permet d'ajuster la force de préhension sans écraser le fruit.
\end{exemplebox}

\courssub{Seuil d'excitation d'une fibre et d'un nerf : la sommation spatiale}

L'expérience sur une \textbf{fibre unique} (axone isolé ou fibre unique disséquée dans un nerf) montre une réponse \textbf{tout ou rien} en fonction de l'intensité du stimulus électrique appliqué.

\begin{experience}[Courbe réponse/intensité sur fibre unique et sur nerf entier]
\textbf{Dispositif :} Nerf sciatique de grenouille (ou rat) isolé, maintenu humide (serum physiologique \SI{0,9}{\%} NaCl, \SI{20}{\celsius}). Deux électrodes de stimulation (cathode/anode) à une extrémité, deux électrodes d'enregistrement (différentielles) à l'autre extrémité. Stimulateur délivrant des impulsions carrées de durée fixe (ex. \SI{1}{ms}) d'intensité croissante ($I$ en mA ou V). Oscilloscope mesurant l'amplitude de la réponse composée (somme des PA de toutes les fibres excitées).
\textbf{Observations :}
\begin{enumerate}
    \item \textbf{Fibre unique (théorique ou microélectrode) :} Réponse nulle pour $I < I_{seuil}$. À $I = I_{seuil}$, apparition brutale d'un PA d'amplitude maximale ($A_{max}$). Pour $I > I_{seuil}$, amplitude constante $A_{max}$. \textbf{Seuil net, tout ou rien.}
    \item \textbf{Nerf entier (réponse composée) :} Amplitude nulle pour $I < I_{min}$ (seuil de la fibre la plus excitable). Puis augmentation \textbf{par paliers} (escaliers) à mesure que $I$ augmente : chaque palier correspond au recrutement d'un nouveau groupe de fibres de seuil identique. Plateau final à $I_{max}$ : toutes les fibres sont recrutées (réponse maximale).
\end{enumerate}
\legende{Mise en évidence du seuil d'excitation et du recrutement des fibres dans un nerf. La réponse composée (somme vectorielle des PA) augmente par paliers car les fibres ont des diamètres (donc des seuils) différents.}
\end{experience}

\begin{popfigure}
\begin{tikzpicture}
\begin{axis}[
    width=\linewidth, height=0.55\linewidth,
    xlabel={Intensité du stimulus $I$ (mA)}, ylabel={Amplitude réponse composée (mV)},
    xmin=0, xmax=10, ymin=0, ymax=5.5,
    axis lines=left, grid=major, grid style={popmuted},
    xtick={0,1,2,3,4,5,6,7,8,9,10},
    ytick={0,1,2,3,4,5},
    clip=false
]
% Escalier idéalisé
\addplot[popblue, very thick, const plot] coordinates {
    (0,0) (1.5,0) (1.5,1.2) (3.0,1.2) (3.0,2.5) (4.5,2.5) (4.5,3.8) (6.0,3.8) (6.0,4.8) (8.0,4.8) (8.0,5.2) (10,5.2)
};
% Fibre unique (inset conceptuel ou trait fin)
\addplot[poporange, thick, dashed, const plot] coordinates {
    (0,0) (2.5,0) (2.5,5.2) (10,5.2)
};
\addlegendentry{Nerf entier (recrutement spatial)}
\addlegendentry{Fibre unique (tout ou rien, seuil $\approx 2.5$ mA)}

% Annotations
\node[popblue, anchor=south west, font=\footnotesize] at (1.5,1.2) {Seuil fibre 1};
\node[popblue, anchor=south west, font=\footnotesize] at (3.0,2.5) {Seuil fibre 2};
\node[popblue, anchor=south west, font=\footnotesize] at (4.5,3.8) {Seuil fibre 3};
\node[popblue, anchor=south west, font=\footnotesize] at (6.0,4.8) {Seuil fibre 4};
\node[popblue, anchor=south west, font=\footnotesize] at (8.0,5.2) {Seuil fibre 5 ($I_{max}$)};
\end{axis}
\end{tikzpicture}
\legende{Variation de l'amplitude de la réponse composée d'un nerf (tronc nerveux) en fonction de l'intensité du stimulus électrique. La courbe en escalier traduit le \textbf{recrutement progressif des fibres} (sommation spatiale) : les fibres de gros diamètre (faible résistance axiale, forte densité de canaux $Na^+$) ont un seuil bas et sont recrutées les premières ; les fibres fines (fort seuil) sont recrutées en dernier. Une fibre unique, elle, répond selon la loi du tout ou rien (courbe orange en pointillés).}
\end{popfigure}

\begin{definition}[Seuil d'excitation (rhéobase) d'une fibre]
C'est l'\textbf{intensité minimale} d'un stimulus électrique (de durée supraliminaire, ex. \SI{1}{ms}) capable de déclencher un potentiel d'action dans cette fibre. Il dépend du diamètre de la fibre (fibres myélinisées grosses $\rightarrow$ seuil bas), de la densité de canaux $Na^+_v$ au nœud de Ranvier, et de la constante de temps membranaire.
\end{definition}

\begin{aretenir}[Sommation spatiale dans un nerf]
Un nerf est un faisceau de fibres de diamètres et de seuils différents. L'augmentation de l'intensité du stimulus recrute progressivement les fibres : d'abord les plus grosses (A$\alpha$, A$\beta$ : proprioception, toucher fin, conduction rapide), puis les moyennes (A$\delta$ : douleur rapide, froid), enfin les fines non myélinisées (C : douleur lente, chaud, système nerveux autonome). C'est la base physiologique de la \textbf{graduation de la force musculaire} (recrutement des unités motrices selon le principe de taille de Henneman) et de la discrimination sensorielle.
\end{aretenir}

\courssub{Relation intensité-durée : rhéobase et temps utile}

Le seuil d'excitation n'est pas une valeur absolue : il dépend de la \textbf{durée} de l'impulsion stimulatrice. C'est la loi de Weiss (1901), vérifiée expérimentalement sur nerf ou muscle.

\begin{experience}[Détermination de la courbe intensité-durée (chronaxie)]
\textbf{Protocole :} Même montage nerf sciatique. On fixe la durée de l'impulsion $t$ (ex. \SI{0,05}{ms}, \SI{0,1}{ms}, \SI{0,2}{ms}, \SI{0,5}{ms}, \SI{1}{ms}, \SI{2}{ms}, \SI{5}{ms}, \SI{10}{ms}). Pour chaque $t$, on cherche l'intensité minimale $I_{seuil}$ déclenchant une réponse (PA). On trace $I_{seuil} = f(t)$.
\textbf{Observation :} Courbe hyperbolique décroissante. Pour une durée très longue, $I_{seuil}$ tend vers une asymptote horizontale. Pour une durée très courte, $I_{seuil}$ tend vers l'infini.
\textbf{Interprétation :} La membrane doit accumuler une \textbf{charge critique} $Q_{critique}$ pour atteindre le seuil. $Q = I \times t$. Si $t$ est long, un faible $I$ suffit (charge accumulée). Si $t$ est court, il faut un $I$ fort pour apporter la charge avant que la membrane ne se décharge (courant de fuite $K^+$).
\legende{Expérience classique de chronaxie (Lapicque, 1909) : mesure du seuil d'excitation en fonction de la durée de l'impulsion carrée.}
\end{experience}

\begin{popfigure}
\begin{tikzpicture}
\begin{axis}[
    width=\linewidth, height=0.55\linewidth,
    xlabel={Durée du stimulus $t$ (ms)}, ylabel={Intensité seuil $I$ (mA)},
    xmin=0, xmax=10, ymin=0, ymax=6,
    axis lines=left, grid=major, grid style={popmuted},
    xtick={0,0.5,1,2,3,4,5,6,7,8,9,10},
    ytick={0,1,2,3,4,5,6},
    domain=0.05:10, samples=200,
    clip=false
]
% Courbe hyperbolique I = Rh * (1 + Tc/t)  -> I = Rh + Rh*Tc/t
% Rh = 1 mA, Tc = 0.5 ms (exemple typique nerf myélinisé)
\def\Rh{1}
\def\Tc{0.5}
\addplot[popblue, very thick, domain=0.1:10] {\Rh + \Rh*\Tc/x};
% Asymptote Rhéobase
\draw[popred, dashed, thick] (0,\Rh) -- (10,\Rh) node[pos=0.95, right, font=\footnotesize] {Rhéobase ($I_{rh} = \SI{1}{mA}$)};
% Temps utile (Chronaxie) : t = Tc quand I = 2*Rh
\draw[poporange, dashed, thick] (\Tc,0) -- (\Tc,2*\Rh) node[pos=0.5, right, font=\footnotesize] {Temps utile / Chronaxie ($t_u = \SI{0,5}{ms}$)};
\draw[poporange, dashed, thick] (0,2*\Rh) -- (\Tc,2*\Rh);
% Point Rheobase
\fill[popred] (0.1,\Rh) circle (2pt);
% Point Chronaxie
\fill[poporange] (\Tc,2*\Rh) circle (2pt);
\end{axis}
\end{tikzpicture}
\legende{Courbe intensité-durée (hyperbole équilatère) pour une fibre nerveuse myélinisée. La \textbf{rhéobase} ($I_{rh}$) est l'intensité seuil pour une durée infinie (pratiquement $> 10 \times t_u$). Le \textbf{temps utile} (ou \textbf{chronaxie} $t_u$) est la durée minimale pour qu'un stimulus d'intensité $2 \times I_{rh}$ soit efficace. Ici $I_{rh}=\SI{1}{mA}$, $t_u=\SI{0,5}{ms}$. Plus la fibre est myélinisée (gros diamètre), plus la chronaxie est courte (excitabilité grande pour les stimuli brefs).}
\end{popfigure}

\begin{definition}[Rhéobase et Temps utile (Chronaxie)]
\begin{itemize}
    \item \textbf{Rhéobase ($I_{rh}$) :} Intensité minimale d'un stimulus capable de déclencher un PA lorsqu'il est appliqué pendant une \textbf{durée suffisamment longue} (théoriquement infinie, praktisch $\ge 10 \times t_u$). Unité : \textbf{mA} (ou V). Elle caractérise l'\textbf{excitabilité de base} (inverse de la rhéobase).
    \item \textbf{Temps utile ($t_u$) ou Chronaxie :} Durée minimale d'application d'un stimulus d'intensité \textbf{double de la rhéobase} ($2 I_{rh}$) pour déclencher un PA. Unité : \textbf{ms} (ou $\mu$s). Elle caractérise l'\textbf{aptitude à répondre aux stimuli brefs}.
\end{itemize}
\textbf{Relation mathématique (Loi de Weiss) :} $I = I_{rh} \left(1 + \frac{t_u}{t}\right)$ pour $t > 0$. La charge seuil $Q_{seuil} = I_{rh} \times t_u$ est constante pour une fibre donnée.
\end{definition}

\begin{savaistu}[Applications cliniques et technologiques au Cameroun]
\begin{itemize}
    \item \textbf{Électrodiagnostic (EMG/ENG) :} Au CHU de Yaoundé ou Douala, on mesure la chronaxie des nerfs moteurs pour diagnostiquer une \textbf{dénervation} (maladie du motoneurone, traumatisme du plexus brachial suite à un accident de moto fréquent sur l'axe Yaoundé-Douala). Une chronaxie augmentée ($> \SI{1}{ms}$) signe une dénervation (fibre démyélinisée ou atrophiée).
    \item \textbf{Stimulation cardiaque (Pacemaker) :} Les sondes d'endocavité stimulent le myocarde avec des impulsions de durée $\approx \SI{0,4}{ms}$ (chronaxie myocardique $\approx \SI{0,3}{ms}$) à une intensité $\approx 2 \times$ seuil pour garantir la capture tout en économisant la batterie (durée de vie $\approx 10$ ans).
    \item \textbf{TENS (Neurostimulation transcutanée) :} Utilisé en kinésithérapie (ex. Centre de rééducation de Mbalmayo) pour soulager la douleur chronique (lombalgie). On utilise des impulsions courtes ($\SI{0,1}{ms}$) à haute fréquence ($100$ Hz) pour recruter sélectivement les fibres A$\beta$ (toucher, chronaxie courte $\approx \SI{0,05}{ms}$) et inhiber la douleur (théorie du \emph{Gate Control}), sans recruter les fibres C (douleur, chronaxie longue $\approx \SI{0,4}{ms}$).
\end{itemize}
\end{savaistu}

\begin{exoresolu}[Détermination de la rhéobase et du temps utile à partir de données expérimentales]
\textbf{Énoncé :} On stimule une fibre nerveuse isolée par des impulsions carrées de durées variables. On relève l'intensité seuil $I$ (en mA) nécessaire pour obtenir un PA.
\begin{center}
\begin{tabular}{|c|c|c|c|c|c|}\hline
Durée $t$ (ms) & 0,10 & 0,20 & 0,50 & 1,00 & 2,00 \\ \hline
Intensité seuil $I$ (mA) & 6,0 & 3,5 & 1,5 & 1,1 & 1,0 \\ \hline
\end{tabular}
\end{center}
\begin{enumerate}
    \item Tracer la courbe $I = f(t)$ (qualitativement) et identifier la forme.
    \item Déterminer graphiquement (ou par calcul) la rhéobase $I_{rh}$ et le temps utile $t_u$.
    \item Calculer la charge seuil $Q_{seuil}$. Quelle est son unité ?
    \item Prédire l'intensité seuil pour une durée $t = \SI{0,05}{ms}$.
\end{enumerate}
\tcblower
\textbf{Solution détaillée :}
\begin{enumerate}
    \item La courbe est une \textbf{branche d'hyperbole} décroissante (loi de Weiss). Asymptote horizontale pour $t$ grand.
    \item \textbf{Rhéobase $I_{rh}$ :} Valeur asymptotique pour $t \to \infty$. D'après le tableau, pour $t=\SI{2,00}{ms}$, $I=\SI{1,0}{mA}$. On estime $I_{rh} \approx \SI{1,0}{mA}$.
    \item \textbf{Temps utile $t_u$ :} Durée pour laquelle $I = 2 \times I_{rh} = \SI{2,0}{mA}$. Par interpolation entre $t=\SI{0,20}{ms}$ ($I=\SI{3,5}{mA}$) et $t=\SI{0,50}{ms}$ ($I=\SI{1,5}{mA}$) :
    \[ t_u \approx 0,20 + \frac{3,5 - 2,0}{3,5 - 1,5} \times (0,50 - 0,20) = 0,20 + 0,75 \times 0,30 = \SI{0,425}{ms} \approx \SI{0,43}{ms}. \]
    (Vérification avec la formule $I = I_{rh}(1 + t_u/t)$ pour $t=0,5$ : $1,5 = 1,0(1 + t_u/0,5) \Rightarrow 0,5 = t_u/0,5 \Rightarrow t_u = 0,25$ ms. Les données ne sont pas parfaitement hyperboliques, on retient l'ordre de grandeur : $t_u \in [\SI{0,25}{ms}; \SI{0,43}{ms}]$).
    \item \textbf{Charge seuil $Q_{seuil} = I_{rh} \times t_u \approx \SI{1,0}{mA} \times \SI{0,35}{ms} = \SI{0,35}{\mu C}$ (microcoulombs).} Unité : $\text{A} \cdot \text{s} = \text{C}$.
    \item \textbf{Prédiction pour $t = \SI{0,05}{ms}$ :} $I = 1,0 \times (1 + 0,35 / 0,05) = 1,0 \times 8 = \SI{8,0}{mA}$. (Très fort courant nécessaire pour stimulus très bref).
\end{enumerate}
\textbf{Conclusion :} Cette fibre a une rhéobase faible ($\SI{1}{mA}$) et une chronaxie courte ($\approx \SI{0,35}{ms}$), typique d'une \textbf{fibre myélinisée de gros diamètre} (type A$\alpha$ ou A$\beta$).
\end{exoresolu}

\begin{methode}[Analyser une courbe Intensité-Durée (Chronaxie) au Bac]
\begin{enumerate}
    \item \textbf{Identifier la courbe :} Hyperbole décroissante $\rightarrow$ Loi de Weiss.
    \item \textbf{Lire la Rhéobase ($I_{rh}$) :} Projeter la branche horizontale (grand $t$) sur l'axe des ordonnées. Unité : mA. \emph{Ne pas confondre avec l'intensité à $t=0$ (infinie).}
    \item \textbf{Lire le Temps utile ($t_u$) :} Repérer $I = 2 \times I_{rh}$ sur l'axe des ordonnées, projeter sur la courbe, puis sur l'axe des abscisses. Unité : ms (ou $\mu$s).
    \item \textbf{Interpréter :} $I_{rh}$ faible + $t_u$ court = Fibre myélinisée, gros diamètre, conduction rapide, excitabilité élevée. $I_{rh}$ forte + $t_u$ long = Fibre amyélinisée, fin diamètre, conduction lente (ex. fibre C douleur).
    \item \textbf{Calculer la charge seuil :} $Q = I_{rh} \times t_u$ (en $\mu$C). Constante pour une fibre donnée.
\end{enumerate}
\end{methode}

\courssub{Synthèse de la naissance à la propagation}

\begin{aretenir}[Résumé intégré : Du stimulus à l'influx propagé]
\begin{enumerate}
    \item \textbf{Transduction (Récepteur) :} Stimulus adéquat $\rightarrow$ Ouverture canaux stimulus-dépendants $\rightarrow$ Courant dépolarisant entrant $\rightarrow$ \textbf{Potentiel de récepteur} (local, gradué, $\propto I_{stimulus}$).
    \item \textbf{Déclenchement (Zone de déclenchement / 1er nœud Ranvier) :} Courant local $\rightarrow$ Dépolarisation membrane $\rightarrow$ Si $V_m \ge V_{seuil}$ ($\approx \SI{-50}{mV}$) $\rightarrow$ Ouverture canaux $Na^+_v$ $\rightarrow$ \textbf{Potentiel d'action} (Tout ou rien).
    \item \textbf{Codage (Fibre afferente) :} $I_{stimulus} \uparrow \Rightarrow$ Amplitude Pot. Récepteur $\uparrow \Rightarrow$ Fréquence PA $\uparrow$ (Codage fréquentiel). Durée stimulus $\Rightarrow$ Durée décharge (adaptation).
    \item \textbf{Recrutement (Nerf) :} $I_{stimulus} \uparrow \Rightarrow$ Nombre fibres excitées $\uparrow$ (Sommation spatiale) $\Rightarrow$ Amplitude réponse composée $\uparrow$ (paliers).
    \item \textbf{Excitabilité (Seuil) :} Définie par le couple ($I_{rh}$, $t_u$). $I = I_{rh}(1 + t_u/t)$. Charge seuil $Q = I_{rh}t_u$ constante.
\end{enumerate}
\end{aretenir}

\begin{exercice}[Application : Codage sensoriel et recrutement]
Un nerf sciatique de grenouille contient des fibres A (myélinisées, gros diamètre) et C (amyélinisées, fin diamètre). On applique un stimulus électrique d'intensité croissante.
\begin{enumerate}
    \item Préciser l'ordre de recrutement des fibres A et C. Justifier par la rhéobase et la chronaxie.
    \item On enregistre l'activité d'une fibre A unique en réponse à un étirement musculaire (récepteur : fuseau neuromusculaire). L'étirement passe de faible à fort. Décrire l'évolution de l'enregistrement (potentiel de récepteur, PA, fréquence).
    \item Pourquoi la douleur vive (piqûre d'aiguille) est-elle perçue avant la douleur sourde (brûlure) bien que le stimulus soit appliqué au même endroit ? (Indice : vitesses de conduction et seuils).
\end{enumerate}
\end{exercice}

\begin{corrige}[Exercice : Codage sensoriel et recrutement]
\begin{enumerate}
    \item \textbf{Ordre :} Fibres A recrutées \textbf{en premier}, puis fibres C. \textbf{Justification :} Fibres A (gros diamètre, myélinisées) $\rightarrow$ Faible résistance axiale, forte densité canaux $Na^+$ aux nœuds $\rightarrow$ \textbf{Rhéobase basse}, \textbf{Chronaxie courte} ($\approx \SI{0,1}{ms}$). Fibres C (fin, amyélinisées) $\rightarrow$ Forte résistance, canaux $Na^+$ diffus $\rightarrow$ \textbf{Rhéobase haute}, \textbf{Chronaxie longue} ($\approx \SI{0,4}{ms}$). À intensité faible, seules les fibres A atteignent le seuil.
    \item \textbf{Enregistrement :}
        \begin{itemize}
            \item Étirement faible : Potentiel de récepteur gradué (amplitude faible), \textbf{pas de PA} (seuil non atteint).
            \item Étirement modéré : Potentiel de récepteur plus ample, atteint le seuil $\rightarrow$ Apparition de \textbf{PA isolés} ou salve à \textbf{faible fréquence} (ex. $10-20$ Hz).
            \item Étirement fort : Potentiel de récepteur maximal $\rightarrow$ Salve de PA à \textbf{fréquence élevée} (ex. $100-200$ Hz). Amplitude PA constante.
        \end{itemize}
    \item \textbf{Douleur vive (A$\delta$) vs Douleur sourde (C) :} La piqûre active les nocicepteurs A$\delta$ (myélinisées, conduction rapide $\approx \SI{10-30}{m/s}$, seuil bas/moyen). La brûlure active surtout les nocicepteurs C (amyélinisées, conduction lente $\approx \SI{0,5-2}{m/s}$, seuil élevé). L'influx A$\delta$ arrive au cortex \textbf{en premier} (douleur vive, précise, localisée, réflexe de retrait immédiat). L'influx C arrive \textbf{secondairement} (douleur sourde, diffuse, persistante, composante affective). C'est la base du \textbf{double pic douloureux}.
\end{enumerate}
\end{corrige}

\begin{attention}[Erreurs fréquentes à ne pas commettre à l'écrit]
\begin{itemize}
    \item \textbf{Dire « l'amplitude du PA code l'intensité ».} \textbf{FAUX.} Amplitude fixe (tout ou rien). C'est la \textbf{fréquence}.
    \item \textbf{Confondre « Seuil de la fibre » et « Rhéobase ».} Le seuil dépend de la durée du stimulus. La rhéobase est le seuil \emph{pour une durée longue}. Préciser : « Seuil pour une impulsion de \SI{1}{ms} » ou « Rhéobase ».
    \item \textbf{Oublier l'unité de la chronaxie.} C'est un \textbf{temps} : ms ou $\mu$s. Pas de mA !
    \item \textbf{Dire « le potentiel de récepteur se propage ».} \textbf{FAUX.} Il est \textbf{local} (décroissance exponentielle avec la constante de longueur $\lambda$). Seuls les PA se propagent.
    \item \textbf{Inverser recrutement A et C.} Retenir : « Gros diamètre = Seuil bas = Recruté en premier = Conduction rapide ». Mnémotechnique : « Les gros passent en premier ».
\end{itemize}
\end{attention}

\coursec{Propagation du potentiel d'action le long de la fibre nerveuse}

Après avoir analysé comment naît le message nerveux, soit au niveau d'un récepteur sensoriel (transduction d'une énergie en potentiel de récepteur puis potentiel d'action), soit au niveau du cône d'émergence d'un neurone (sommation spatiale et temporelle des PSP atteignant le \cle{seuil d'excitation}), nous devons maintenant comprendre comment ce signal électrique voyage le long de l'axone jusqu'aux terminaisons synaptiques. La propagation n'est pas un simple déplacement de charge : c'est une \textbf{réautorégénération locale successive} du potentiel d'action, soumise à des contraintes temporelles strictes (période réfractaire) et dont la vitesse dépend de propriétés physiques de la fibre (diamètre, myélinisation).

\courssub{Mise en évidence et interprétation de la période réfractaire}

\paragraph{Expérience clé : la double stimulation.}
Pour comprendre pourquoi les potentiels d'action (PA) restent des événements discrets et ne se somment pas, réalisons une expérience classique de neurophysiologie sur une fibre nerveuse isolée (ex. : nerf sciatique de grenouille, préparation standard en TP).

\begin{experience}[Mise en évidence de la période réfractaire par double stimulation]
\textbf{Matériel :} Fibre nerveuse isolée maintenue vivante dans un bain de Ringer ; deux électrodes stimulateurs (S1, S2) placées au même endroit (ou stimulation unique par deux impulsions successives) ; deux électrodes réceptrices (R1, R2) espacées d'une distance connue $d$ ; oscilloscope ou acquisition numérique.
\textbf{Protocole :}
\begin{enumerate}
    \item On délivre une première impulsion de stimulation (Stim 1) d'intensité supraseuil. On enregistre un PA au niveau de R1 puis R2.
    \item On délivre une seconde impulsion (Stim 2) d'intensité identique, après un intervalle de temps $\Delta t$ variable (de 0,1 ms à 10 ms).
    \item On observe la réponse (présence ou absence de second PA) en fonction de $\Delta t$.
\end{enumerate}
\textbf{Observations :}
\begin{itemize}
    \item Si $\Delta t < 1$ ms (environ) : Stim 2 ne déclenche \textbf{aucun} PA, quelle que soit son intensité (même très forte).
    \item Si $1 \text{ ms} < \Delta t < 4 \text{ ms}$ (valeurs variables selon la fibre) : Stim 2 ne déclenche un PA que si son intensité est \textbf{supérieure} au seuil normal.
    \item Si $\Delta t > 4 \text{ ms}$ : Stim 2 déclenche un PA normal au seuil habituel.
\end{itemize}
\textbf{Interprétation :} La membrane, après un PA, traverse un état d'inexcitabilité transitoire : la \cle{période réfractaire}.
\end{experience}

\begin{popfigure}
\begin{tikzpicture}[scale=0.9, font=\small]
    % Axon representation - Time axis
    \draw[thick, popdark] (0,0) -- (14,0) node[midway, below=0.4cm] {Axe du temps (ms)};
    
    % Stim 1
    \draw[popblue, thick] (1,1.5) -- (1,0.5) node[above] {Stim 1};
    % AP 1 at R1 (schematic)
    \begin{scope}[shift={(1,0)}]
        \draw[popblue, thick, domain=0:3, samples=100] plot (\x, {0.8*exp(-(\x-0.5)^2/0.2) - 0.3*exp(-(\x-1.5)^2/0.1)});
        \draw[dashed, popmuted] (0,-0.5) -- (3,-0.5) node[right] {Potentiel de repos};
    \end{scope}
    \node[popblue] at (2.5, 1.2) {PA 1};
    
    % Refractory period bars (manual brackets)
    % Absolute
    \draw[poporange, thick] (1,-1.2) -- (4.5,-1.2);
    \draw[poporange, thick] (1,-1.3) -- (1,-1.1);
    \draw[poporange, thick] (4.5,-1.3) -- (4.5,-1.1);
    \node[poporange, below=0.25cm] at (2.75, -1.2) {Période réfractaire absolue};
    % Relative
    \draw[poporange!70!popdark, thick] (4.5,-1.2) -- (7,-1.2);
    \draw[poporange!70!popdark, thick] (4.5,-1.3) -- (4.5,-1.1);
    \draw[poporange!70!popdark, thick] (7,-1.3) -- (7,-1.1);
    \node[poporange!70!popdark, below=0.25cm] at (5.75, -1.2) {Période réfractaire relative};
    
    % Stim 2 during absolute refractory (fails)
    \draw[popred, thick] (2,2.5) -- (2,1.8) node[above, popred] {Stim 2 (échec)};
    \draw[popred, thick] (2,0) -- (2,-0.8) node[below, popred] {Pas de PA};
    
    % Stim 2 during relative refractory (success with strong stim)
    \draw[popgreen, thick] (5.5,2.5) -- (5.5,1.8) node[above, popgreen] {Stim 2 (fort)};
    \begin{scope}[shift={(5.5,0)}]
        \draw[popgreen, thick, domain=0:3, samples=100] plot (\x, {0.6*exp(-(\x-0.5)^2/0.3) - 0.2*exp(-(\x-1.5)^2/0.15)});
    \end{scope}
    \node[popgreen] at (7, 1) {PA 2 (amplitude normale)};
    
    % Stim 2 after refractory (normal)
    \draw[popblue, thick] (8,2.5) -- (8,1.8) node[above, popblue] {Stim 2 (normal)};
    \begin{scope}[shift={(8,0)}]
        \draw[popblue, thick, domain=0:3, samples=100] plot (\x, {0.8*exp(-(\x-0.5)^2/0.2) - 0.3*exp(-(\x-1.5)^2/0.1)});
    \end{scope}
    \node[popblue] at (9.5, 1.2) {PA 2 normal};
    
    % Voltage scale
    \draw[->, popmuted] (0.2,-0.5) -- (0.2,1.5) node[left] {$V_m$ (mV)};
    \node[popmuted] at (0.2,0) {-70};
    \node[popmuted] at (0.2,1) {+30};
\end{tikzpicture}
\legende{Schéma de principe de l'expérience de double stimulation. En haut : chronologie des stimulations et réponses. En bas : barres représentant la période réfractaire absolue (aucun PA possible) puis relative (PA possible seulement si stimulus $>$ seuil).}
\end{popfigure}

\begin{definition}[Période réfractaire]
La \textbf{période réfractaire} est l'intervalle de temps qui suit un potentiel d'action pendant lequel la fibre nerveuse a une excitabilité diminuée ou nulle. On distingue deux phases :
\begin{itemize}
    \item \textbf{Période réfractaire absolue (PRA) :} Durée $\approx 0,5$ à $1,5$ ms. La fibre est \textbf{totalement inexcitable}. Aucun stimulus, si intense soit-il, ne peut déclencher un nouveau PA.
    \item \textbf{Période réfractaire relative (PRR) :} Durée $\approx 2$ à $4$ ms (variable). La fibre est \textbf{difficilement excitable} : un stimulus d'intensité \textbf{supérieure au seuil normal} peut déclencher un PA.
\end{itemize}
\end{definition}

\begin{propriete}[Base ionique de la période réfractaire]
\begin{itemize}
    \item \textbf{PRA :} Correspond à l'\textbf{inactivation} des canaux $\text{Na}^+$ voltage-dépendants (porte $h$ fermée). Ces canaux ne peuvent pas s'ouvrir avant d'être revenus à l'état de repos (dé-inactivation), ce qui nécessite une repolarisation complète. Les canaux $\text{K}^+$ sont ouverts (repolarisation).
    \item \textbf{PRR :} Correspond à la phase de \textbf{surpolarisation} (potentiel plus négatif que $V_{\text{repos}}$) due à la fermeture lente des canaux $\text{K}^+$. Pour atteindre le seuil, il faut une dépolarisation plus forte (stimulus plus intense) pour compenser ce "frein" hyperpolarisant et recruter suffisamment de canaux $\text{Na}^+$ dé-inactivés.
\end{itemize}
\end{propriete}

\begin{aretenir}[Conséquences fonctionnelles majeures]
\begin{enumerate}
    \item \textbf{Limitation de la fréquence maximale :} La durée totale de la période réfractaire ($\approx 1$ à $5$ ms) impose une fréquence maximale de décharge : $f_{\text{max}} \approx 1 / T_{\text{ref}} \approx \mathbf{200 \text{ à } 1000 \text{ Hz}}$. Le message nerveux reste un signal \textbf{discontinu} (suite de PA distincts), codé en \textbf{fréquence}, pas en amplitude.
    \item \textbf{Unidirectionnalité de la propagation (in vivo) :} La zone qui vient de générer le PA est en période réfractaire. Le PA ne peut donc pas revenir en arrière ; il se propage nécessairement vers la zone au repos (direction orthodromique : corps cellulaire $\to$ terminaisons).
\end{enumerate}
\end{aretenir}

\courssub{Mécanisme de propagation : courants locaux et autorégénération}

\begin{experience}[Preuve de la propagation par courants locaux (modèle du câble)]
\textbf{Observation :} Si l'on stimule une fibre au milieu (in vitro), deux PA apparaissent aux deux extrémités. Si l'on bloque localement les canaux $\text{Na}^+$ (ex. tétrodotoxine TTX) sur un segment, le PA s'arrête net à l'entrée du segment bloqué.
\textbf{Interprétation :} Le PA en un point $A$ (dépolarisation à $+30$ mV) crée une différence de potentiel avec le point voisin $B$ au repos ($-70$ mV). Des \textbf{courants locaux} (flux d'ions $\text{Na}^+$ intracellulaire vers $B$, flux de $\text{K}^+$ extracellulaire vers $A$) se ferment en boucle. Ce courant dépolarise $B$. Si cette dépolarisation atteint le \cle{seuil}, les canaux $\text{Na}^+$ de $B$ s'ouvrent : un \textbf{nouveau PA naît intégralement en $B$}. Le signal ne s'atténue pas : c'est une \textbf{propagation autorégénératrice de proche en proche}.
\end{experience}

\begin{popfigure}
\begin{tikzpicture}[scale=0.85, font=\small]
    % Axon membrane
    \draw[popdark, thick] (0.5, 2.5) -- (13.5, 2.5) node[midway, above=0.5cm] {Membrane axonale (extracellulaire au-dessus)};
    \draw[popdark, thick] (0.5, 0.5) -- (13.5, 0.5) node[midway, below=0.5cm] {Intracellulaire};
    
    % Zone 1: Refractary (just passed)
    \fill[poporange!20] (0.5, 0.5) rectangle (3.5, 2.5);
    \node[poporange!80!popdark, align=center] at (2, 1.5) {Zone 1 :\\ Période réfractaire\\ (Canaux $\text{Na}^+$ inactivés)};
    \draw[->, thick, popred] (2, 1.5) -- (2, 0.8); 
    \node[popred, left] at (1.7, 1) {$\text{K}^+$ sort};
    
    % Zone 2: Active (AP peak)
    \fill[popblue!20] (3.5, 0.5) rectangle (6.5, 2.5);
    \node[popblue, align=center] at (5, 1.5) {Zone 2 :\\ Pic du PA\\ ($V_m \approx +30$ mV)};
    \draw[->, thick, popblue] (5, 2.0) -- (5, 2.3); 
    \node[popblue, right] at (5.3, 2.15) {$\text{Na}^+$ entre};
    % Local currents loops
    \draw[->, thick, popgreen!70!popdark] (4.5, 1.5) .. controls (4.5, 1.0) and (3.0, 1.0) .. (3.0, 1.5);
    \draw[->, thick, popgreen!70!popdark] (3.0, 1.5) .. controls (3.0, 2.0) and (4.5, 2.0) .. (4.5, 1.5);
    \node[popgreen!70!popdark, align=center] at (3.7, 0.8) {Courant local\\ intracellulaire};
    
    % Zone 3: Threshold / Rising phase
    \fill[popgreen!20] (6.5, 0.5) rectangle (9.5, 2.5);
    \node[popgreen!70!popdark, align=center] at (8, 1.5) {Zone 3 :\\ Dépolarisation\\ seuil atteinte};
    \draw[->, thick, popblue] (8, 2.0) -- (8, 2.3);
    \node[popblue, right, align=center] at (8.3, 2.15){Ouverture\\ canaux $\text{Na}^+$};
    
    % Zone 4: Rest
    \fill[popgray!20] (9.5, 0.5) rectangle (13.5, 2.5);
    \node[popdark, align=center] at (11.5, 1.5) {Zone 4 :\\ Repos ($V_m = -70$ mV)\\ Canaux $\text{Na}^+$ fermés,\\ prêts à s'ouvrir};
    
    % Propagation arrow
    \draw[->, very thick, popdark] (7, 3.2) -- (11, 3.2) node[midway, above] {Sens de propagation};
    
    % Ion channels representation (simplified)
    \foreach \x in {4, 5, 6} {
        \draw[popblue, thick] (\x, 2.5) -- (\x, 2.7) node[above, scale=0.7] {$\text{Na}^+$};
    }
    \foreach \x in {7, 8, 9} {
        \draw[popblue, thick, dashed] (\x, 2.5) -- (\x, 2.7) node[above, scale=0.7] {$\text{Na}^+$};
    }
    \foreach \x in {10, 11, 12} {
        \draw[popblue, thick] (\x, 2.5) -- (\x, 2.7) node[above, scale=0.7] {$\text{Na}^+$};
    }
    \node[popmuted, scale=0.7] at (11.5, 2.9) {Canaux fermés (repos)};
    \node[popmuted, scale=0.7] at (5, 2.9) {Canaux ouverts};
    \node[popmuted, scale=0.7] at (2, 2.9) {Canaux inactivés};
\end{tikzpicture}
\legende{Mécanisme de propagation par courants locaux. Le PA en Zone 2 génère un courant intracellulaire vers l'avant (Zone 3) qui la dépolarise jusqu'au seuil, déclenchant un nouveau PA. La Zone 1, en période réfractaire, empêche la propagation en arrière.}
\end{popfigure}

\begin{attention}[Piège conceptuel : le PA ne « voyage » pas]
\textbf{Erreur fréquente :} « Le potentiel d'action se déplace le long de l'axone comme une vague d'eau. »
\textbf{Correction rigoureuse :} Le PA est un \textbf{événement local} qui se \textbf{reconstruit à neuf} en chaque point de la membrane. L'amplitude du PA est \textbf{constante} tout au long de la fibre (loi du tout-ou-rien). Seule la \textbf{zone de dépolarisation active} se déplace. C'est pourquoi la myélinisation (voir plus loin) accélère la propagation sans changer l'amplitude du PA aux nœuds.
\end{attention}

\courssub{Sens de la propagation : orthodromie et antidromie}

\begin{propriete}[Sens de propagation]
\begin{itemize}
    \item \textbf{In vivo (physiologique) :} La propagation est \textbf{unidirectionnelle (orthodromique)} : du corps cellulaire (ou zone de déclenchement : cône d'émergence / récepteur) vers les terminaisons synaptiques. \textbf{Cause :} La zone parcourue est en période réfractaire absolue, bloquant toute propagation rétrograde.
    \item \textbf{In vitro (expérimental) :} Si l'on stimule \textbf{artificiellement au milieu} de la fibre, le PA se propage \textbf{dans les deux sens} (bidirectionnel) : orthodromique vers les terminaisons, et \textbf{antidromique} vers le corps cellulaire. Cela prouve que l'axone est physiquement capable de conduire dans les deux sens, mais que la réfractarité impose l'unidirectionnalité naturelle.
\end{itemize}
\end{propriete}

\courssub{Mesure de la vitesse de propagation}

\begin{methode}[Calcul de la vitesse de propagation $v$]
\textbf{Dispositif :} Deux électrodes réceptrices $R_1$ et $R_2$ (non invasives, type ventouses ou aiguilles) placées sur le nerf, espacées d'une distance connue $d$ (en mètres). Une électrode stimulateur $S$ en amont de $R_1$. Enregistrement simultané sur deux voies d'un oscilloscope.
\begin{enumerate}
    \item Mesurer la distance $d$ entre $R_1$ et $R_2$ (ex. : $d = 0,10 \text{ m}$).
    \item Déclencher un PA unique par stimulation supraseuil.
    \item Mesurer sur l'écran le décalage temporel $\Delta t$ entre l'arrivée du PA sur $R_1$ et son arrivée sur $R_2$ (ex. : $\Delta t = 2,0 \text{ ms} = 2,0 \times 10^{-3} \text{ s}$).
    \item Calculer : $v = \dfrac{d}{\Delta t}$.
\end{enumerate}
\textbf{Vérification des unités :} $v$ en $\text{m/s}$ si $d$ en $\text{m}$ et $\Delta t$ en $\text{s}$.
\end{methode}

\begin{popfigure}
\begin{tikzpicture}[scale=0.9, font=\small]
    % Nerve fiber
    \draw[popdark, very thick] (0.5, 2) -- (13.5, 2);
    \draw[popdark, very thick] (0.5, 1) -- (13.5, 1);
    \node[popdark, left] at (0, 1.5) {Nerf};
    
    % Stimulating electrode
    \draw[popred, thick] (-0.5, 1.5) -- (0.5, 1.5);
    \node[popred, left] at (-0.5, 1.5) {S};
    
    % Recording electrodes R1 and R2
    \draw[popblue, thick] (3, 2.5) -- (3, 0.5);
    \node[popblue, above] at (3, 2.5) {R1};
    \draw[popblue, thick] (9, 2.5) -- (9, 0.5);
    \node[popblue, above] at (9, 2.5) {R2};
    
    % Distance d
    \draw[<->, popdark] (3, 2.7) -- (9, 2.7) node[midway, above] {$d$};
    
    % Oscilloscope traces
    \begin{scope}[shift={(0.5, -2.5)}, scale=0.8]
        \draw[popmuted] (0,0) grid[step=1] (8,4);
        \draw[->] (0,0) -- (8.5,0) node[right] {$t$ (ms)};
        \draw[->] (0,0) -- (0,4.5) node[above] {Tension};
        % Trace R1
        \draw[popblue, thick, domain=0.5:7.5, samples=200] plot (\x, {2 + 1.5*exp(-(\x-2)^2/0.3) - 0.5*exp(-(\x-3)^2/0.2)});
        \node[popblue] at (2, 3.8) {R1};
        % Trace R2 delayed
        \draw[poporange, thick, domain=0.5:7.5, samples=200] plot (\x, {0.5 + 1.5*exp(-(\x-4)^2/0.3) - 0.5*exp(-(\x-5)^2/0.2)});
        \node[poporange] at (4, 2.3) {R2};
        % Delta t arrow
        \draw[<->, popred, thick] (2, 3.3) -- (4, 3.3) node[midway, above] {$\Delta t$};
    \end{scope}
    
    % Formula box
    \node[draw, popblueL, fill=popblue!5, rounded corners, align=center, anchor=north west] at (10, -1.5) {
        \textbf{Calcul de la vitesse :}\\
        $v = \dfrac{d}{\Delta t}$\\
        Exemple : $d = 0,10 \text{ m}$, $\Delta t = 2,0 \text{ ms}$\\
        $v = \dfrac{0,10}{2,0 \times 10^{-3}} = 50 \text{ m/s}$
    };
\end{tikzpicture}
\legende{Montage classique à deux électrodes réceptrices pour mesurer la vitesse de conduction. Le décalage $\Delta t$ entre les deux tracés (R1 et R2) permet le calcul $v = d/\Delta t$.}
\end{popfigure}

\begin{exoresolu}[Calcul de vitesse de conduction]
\textbf{Énoncé :} Sur une fibre nerveuse myélinisée de grenouille, deux électrodes réceptrices sont espacées de $d = 15 \text{ cm}$. Un potentiel d'action est enregistré. Le décalage temporel entre les deux pics est $\Delta t = 1,25 \text{ ms}$. Calculer la vitesse de propagation en $\text{m/s}$ et en $\text{km/h}$.
\textbf{Solution détaillée :}
\begin{enumerate}
    \item Conversion des unités SI :
    $d = 15 \text{ cm} = 0,15 \text{ m}$.
    $\Delta t = 1,25 \text{ ms} = 1,25 \times 10^{-3} \text{ s}$.
    \item Application de la formule :
    $v = \dfrac{d}{\Delta t} = \dfrac{0,15}{1,25 \times 10^{-3}} = \dfrac{150}{1,25} = 120 \text{ m/s}$.
    \item Conversion en $\text{km/h}$ :
    $v = 120 \times 3,6 = 432 \text{ km/h}$.
\end{enumerate}
\textbf{Conclusion :} La vitesse est de $\mathbf{120 \text{ m/s}}$ ($\mathbf{432 \text{ km/h}}$), valeur typique des fibres myélinisées de gros diamètre (type A$\alpha$).
\tcblower
\textbf{Conseil de rédaction :} Toujours écrire la formule, faire les conversions d'unités \textbf{avant} le calcul, et donner le résultat avec son unité. Vérifier l'ordre de grandeur : une fibre amyélinique ne dépasse pas $2 \text{ m/s}$.
\end{exoresolu}

\courssub{Facteurs influençant la vitesse : diamètre et myélinisation}

La vitesse de propagation $v$ dépend de deux paramètres structurels majeurs de la fibre nerveuse.

\begin{propriete}[Loi de la vitesse de conduction]
\begin{enumerate}
    \item \textbf{Diamètre de la fibre (axone) :} $v \propto \sqrt{\text{diamètre}}$ (pour fibres amyélinisées) ou $v \propto \text{diamètre}$ (pour fibres myélinisées).
    \emph{Explication physique :} Un axone plus large offre une \textbf{résistance axoplasmique ($R_a$) plus faible} au passage du courant local intracellulaire. Le courant se propage plus loin avant de fuir par la membrane, dépolarisant une plus grande longueur à la fois $\to$ le PA « saute » plus loin $\to$ vitesse plus grande.
    \item \textbf{Myélinisation :} Gain de vitesse \textbf{spectaculaire} (facteur 50 à 100). La gaine de myéline (lipides) agit comme un \textbf{isolant électrique} (augmente la résistance membranaire $R_m$, diminue la capacité $C_m$).
\end{enumerate}
\end{propriete}

\begin{popfigure}
\begin{tikzpicture}[scale=0.8, font=\small]
    % Unmyelinated fiber (Top Left)
    \node[popdark, align=center] at (3, 5.5) {Fibre amyélinique (ex: fibre C, douleur diffuse)};
    \draw[popdark, thick] (0.5, 4.8) -- (5.5, 4.8);
    \draw[popdark, thick] (0.5, 4.2) -- (5.5, 4.2);
    \foreach \x in {1, 1.5, 2, 2.5, 3, 3.5, 4, 4.5, 5} {
        \draw[popblue, thick] (\x, 4.2) -- (\x, 4.8);
        \node[popblue, scale=0.6, above] at (\x, 4.8) {$\text{Na}^+$};
    }
    \draw[->, very thick, popgreen] (1, 5.2) -- (5, 5.2) node[midway, above] {Propagation continue (lente)};
    \node[popmuted, align=center] at (3, 3.8) {Courants locaux à chaque $\mu$m\\ $v \approx \mathbf{0,5 \text{ à } 2 \text{ m/s}}$};
    
    % Myelinated fiber (Top Right)
    \node[popdark, align=center] at (10.5, 5.5) {Fibre myélinisée (ex: fibre A$\alpha$, motricité)};
    \draw[popdark, thick] (7.5, 4.8) -- (13.5, 4.8);
    \draw[popdark, thick] (7.5, 4.2) -- (13.5, 4.2);
    % Myelin sheaths
    \foreach \x in {8, 9.5, 11, 12.5} {
        \fill[poporange!30] (\x, 4.2) rectangle (\x+0.8, 4.8);
        \draw[poporange!80!popdark, thick] (\x, 4.2) -- (\x, 4.8);
        \draw[poporange!80!popdark, thick] (\x+0.8, 4.2) -- (\x+0.8, 4.8);
    }
    % Nodes of Ranvier
    \foreach \x in {8.8, 10.3, 11.8} {
        \draw[popblue, very thick] (\x, 4.2) -- (\x, 4.8);
        \node[popblue, scale=0.7, above] at (\x, 4.8) {Nœud};
        \foreach \y in {4.3, 4.45, 4.6, 4.75} {
            \draw[popblue, thick] (\x, \y) -- (\x+0.05, \y); 
        }
    }
    \draw[->, very thick, popgreen] (8, 5.2) -- (13, 5.2) node[midway, above] {Conduction saltatoire (rapide)};
    \node[popmuted, align=center] at (10.5, 3.8) {PA seulement aux nœuds de Ranvier\\ $v \approx \mathbf{80 \text{ à } 120 \text{ m/s}}$};
    
    % Zoom on Node (Centered Below)
    \begin{scope}[shift={(7, 1.5)}, scale=1.3]
        \node[popdark] at (0, 2.2) {Zoom sur un nœud de Ranvier};
        \draw[popdark, thick] (-1.5, 1.5) -- (1.5, 1.5);
        \draw[popdark, thick] (-1.5, 0.5) -- (1.5, 0.5);
        \fill[poporange!30] (-1.5, 0.5) rectangle (-0.5, 1.5);
        \fill[poporange!30] (0.5, 0.5) rectangle (1.5, 1.5);
        \node[poporange!80!popdark, rotate=90, scale=0.7] at (-1.0, 1.0) {Myéline};
        \node[poporange!80!popdark, rotate=90, scale=0.7] at (1.0, 1.0) {Myéline};
        \draw[popblue, very thick] (-0.5, 0.5) -- (-0.5, 1.5);
        \draw[popblue, very thick] (0.5, 0.5) -- (0.5, 1.5);
        \node[popblue, scale=0.7] at (0, 1.8) {Nœud de Ranvier};
        \node[popblue, scale=0.7] at (0, 0.2) {Forte densité canaux $\text{Na}^+/\text{K}^+$};
        \draw[<->, popmuted] (-0.5, 0.0) -- (0.5, 0.0) node[midway, below] {$\approx 1 \mu\text{m}$};
        \draw[<->, popmuted] (-1.7, 0.5) -- (-1.7, 1.5) node[midway, left] {Internœud $\approx 1 \text{ mm}$};
    \end{scope}
\end{tikzpicture}
\legende{Comparaison structurelle et fonctionnelle : fibre amyélinique (propagation continue, lente) vs fibre myélinisée (conduction saltatoire aux nœuds de Ranvier, rapide). La myéline isole l'internœud, forçant le courant local à sauter de nœud en nœud.}
\end{popfigure}

\begin{experience}[Mise en évidence de la conduction saltatoire (Expérience de Tasaki, 1939 / Stämpfli, 1952)]
\textbf{Procédé :} Enregistrement intra-axoplasmique (microélectrode) ou extracellulaire à haute résolution spatiale sur un nerf myélinisé (ex. nerf vague de grenouille).
\textbf{Observation :} Le potentiel d'action n'est détectable \textbf{qu'aux nœuds de Ranvier}. Dans l'internœud (sous la myéline), la membrane est dépolarisée passivement (potentiel électrotonique) mais \textbf{aucun PA ne s'y régénère} (absence de canaux $\text{Na}^+$ voltage-dépendants en forte densité).
\textbf{Interprétation :} La myéline augmente la résistance membranaire ($R_m \uparrow$) et diminue la capacité ($C_m \downarrow$). La constante de longueur $\lambda = \sqrt{R_m / R_a}$ devient très grande ($\approx 1 \text{ mm}$). Le courant local généré au nœud $N$ se propage passivement sans perte significative jusqu'au nœud $N+1$, où il déclenche un nouveau PA. C'est la \textbf{conduction saltatoire} (du latin \emph{saltare}, sauter).
\end{experience}

\begin{center}
\begin{tabularx}{\linewidth}{|l|X|X|X|}
\hline
\rowcolor{popblueL}
\textbf{Type de fibre} & \textbf{Diamètre ($\mu$m)} & \textbf{Myélinisation} & \textbf{Vitesse typique (m/s)} \\
\hline
Fibre C (douleur chronique, végétatif) & 0,2 -- 1,5 & Non & \textbf{0,5 -- 2} \\
Fibre A$\delta$ (douleur vive, froid) & 1,5 -- 5 & Faible / Fine & \textbf{5 -- 30} \\
Fibre A$\beta$ (toucher, pression) & 5 -- 13 & Oui & \textbf{30 -- 70} \\
Fibre A$\alpha$ (motricité, proprioception) & 13 -- 22 & Oui (épaisse) & \textbf{80 -- 120} \\
\hline
\end{tabularx}
\end{center}

\begin{aretenir}[Synthèse comparative des vitesses]
\begin{itemize}
    \item \textbf{Fibres amyélinisées fines :} $v \approx \mathbf{1 \text{ m/s}}$. Propagation continue. Codent des signaux lents, non urgents (douleur sourde, autonomique).
    \item \textbf{Fibres myélinisées gros diamètre :} $v \approx \mathbf{100 \text{ m/s}}$ ($360 \text{ km/h}$). Conduction saltatoire. Codent l'urgence (réflexes ostéotendineux, motricité fine, toucher discriminant).
    \item \textbf{Intérêt adaptatif :} La myélinisation permet une conduction rapide \textbf{sans augmenter excessivement le volume} du nerf (gain de place et d'énergie métabolique : moins de pompes $\text{Na}^+/\text{K}^+$ à activer par PA car seule la membrane nodale dépense de l'ATP).
\end{itemize}
\end{aretenir}

\begin{savaistu}[Pathologie : La sclérose en plaques (SEP)]
La \textbf{sclérose en plaques} est une maladie auto-immune où le système immunitaire attaque la \textbf{myéline} du système nerveux central (cerveau, moelle épinière). La démyélinisation entraîne :
\begin{itemize}
    \item Une \textbf{chute drastique de la vitesse de conduction} (conduction saltatoire impossible $\to$ retour à une conduction continue lente ou blocage total).
    \item Une \textbf{augmentation de la période réfractaire} et une fatigabilité accrue (bloc de conduction à haute fréquence).
    \item \textbf{Symptômes cliniques :} Troubles moteurs (faiblesse, spasticité), sensitifs (fourmillements, perte de proprioception), visuels (névrite optique), fatigabilité majeure à l'effort (phénomène d'Uhthoff : chaleur $\to$ ralentissement conduction).
\end{itemize}
Au Cameroun, bien que moins fréquente que dans les pays tempérés, la SEP est diagnostiquée (service de neurologie de l'Hôpital Central de Yaoundé, CHU de Douala). La prise en charge repose sur les immunomodulateurs et la rééducation fonctionnelle pour préserver l'autonomie.
\end{savaistu}

\courssub{Synthèse intégrée : du stimulus à la réponse effectrice}

\begin{exemplebox}[Parcours complet d'un réflexe myotatique (ex: réflexe rotulien au CHU de Yaoundé)]
\begin{enumerate}
    \item \textbf{Réception :} Étirement du tendon rotulien $\to$ déformation du fuseau neuromusculaire $\to$ \textbf{potentiel de récepteur} (générateur) dans les fibres Ia $\to$ \textbf{PA} si seuil atteint (codage fréquence/intensité étirement).
    \item \textbf{Conduction afferente :} Fibre \textbf{A$\alpha$ myélinisée, gros diamètre} ($v \approx 100 \text{ m/s}$). PA orthodromique vers moelle épinière. Période réfractaire assure fidélité du codage fréquentiel.
    \item \textbf{Traitement central :} Synapse monosynaptique sur motoneurone $\alpha$ (EPSP). Sommation spatiale/temporelle $\to$ PA au cône d'émergence.
    \item \textbf{Conduction efferente :} Fibre \textbf{A$\alpha$ myélinisée} ($v \approx 100 \text{ m/s}$) vers muscle extenseur.
    \item \textbf{Effet :} Contraction musculaire (extension jambe) en $\approx 30 \text{ ms}$ total. Vitesse cruciale pour la posture.
\end{enumerate}
\textbf{Si fibre démyélinisée (SEP) :} $v$ chute à $< 10 \text{ m/s}$ $\to$ latence réflexe augmentée $\to$ désynchronisation motrice $\to$ troubles de la marche, spasticité.
\end{exemplebox}

\begin{exercice}[Analyse d'un tracé de période réfractaire]
On enregistre l'activité d'une fibre nerveuse soumise à deux stimulations successives S1 et S2 d'intensité identique (supraseuil). L'intervalle S1-S2 est varié. On obtient les résultats suivants :
\begin{center}
\begin{tabular}{|c|c|}
\hline
Intervalle S1-S2 (ms) & Réponse à S2 \\ \hline
0,5 & Aucun PA \\ \hline
1,5 & PA présent (amplitude normale) \\ \hline
3,0 & PA présent (amplitude normale) \\ \hline
\end{tabular}
\end{center}
\textbf{Questions :}
\begin{enumerate}
    \item Déduire la durée approximative de la période réfractaire absolue (PRA) de cette fibre.
    \item Expliquer pourquoi à 0,5 ms, aucun PA n'apparaît, même avec un stimulus supraseuil.
    \item Prédire le résultat si l'on augmentait l'intensité de S2 à 0,5 ms d'intervalle. Justifier.
\end{enumerate}
\end{exercice}

\begin{corrige}[Exercice n]
\begin{enumerate}
    \item \textbf{PRA $\approx 1,0 \text{ ms}$ (entre 0,5 et 1,5 ms).} À 0,5 ms : dans PRA (échec). À 1,5 ms : hors PRA (succès au seuil normal).
    \item \textbf{Cause ionique :} À 0,5 ms, les canaux $\text{Na}^+$ voltage-dépendants sont en conformation \textbf{inactivée} (porte $h$ fermée). Ils ne peuvent pas s'ouvrir, quelle que soit la dépolarisation apportée par le stimulus. La perméabilité $\text{Na}^+$ ne peut pas augmenter $\to$ pas de boucle de rétroaction positive $\to$ pas de PA.
    \item \textbf{Prédiction :} \textbf{Toujours aucun PA.} Par définition, pendant la PRA, \textbf{aucune intensité de stimulus} ne peut déclencher de PA. L'inexcitabilité est absolue. (Ce n'est qu'en période réfractaire relative qu'un stimulus plus fort fonctionne).
\end{enumerate}
\end{corrige}

\begin{attention}[Erreurs à ne pas commettre au Bac]
\begin{itemize}
    \item Confondre \textbf{vitesse de propagation} ($v$ en m/s, propriété de la fibre) et \textbf{fréquence de décharge} ($f$ en Hz, codage de l'intensité du stimulus). $v$ est constante pour une fibre donnée ; $f$ varie.
    \item Oublier de convertir $\Delta t$ en \textbf{secondes} et $d$ en \textbf{mètres} avant d'appliquer $v = d/\Delta t$.
    \item Dire que « le PA saute la myéline ». \textbf{Précision :} Le PA \textbf{ne se produit pas} sous la myéline (pas de canaux). Le \textbf{courant local} traverse l'internœud passivement ; le PA \textbf{renaît} au nœud suivant.
    \item Écrire « période réfractaire = temps où le neurone se repose ». \textbf{Non :} C'est un temps d'\textbf{inexcitabilité forcée} par l'état conformationnel des canaux ioniques, indispensable pour la discrétisation du signal et l'unidirectionnalité.
\end{itemize}
\end{attention}

\coursec{Synthèse et sensibilisation : préserver le fonctionnement des neurones}

Dans les sections précédentes, nous avons analysé le potentiel de repos, le potentiel d'action, sa naissance au niveau des récepteurs et des neurones, ainsi que sa propagation le long de la fibre nerveuse. Il est maintenant temps de rassembler toutes les pièces du puzzle pour comprendre comment une information sensorielle devient une sensation consciente, et surtout — c'est le sens du titre de cette leçon — d'identifier les facteurs qui perturbent cette chaîne et les moyens de la protéger.

\courssub{La chaîne complète du message nerveux : du stimulus au centre nerveux}

\textbf{Observation de départ.} Approche ta main d'une flamme : tu la retires avant même d'avoir pensé « ça brûle ». En une fraction de seconde, l'information « chaleur excessive » a été captée, transformée, codée, transportée et traitée. Reconstituons cette chaîne étape par étape.

\begin{enumerate}
\item \textbf{Le stimulus} (chaleur, pression, lumière, substance chimique) excite un \cle{récepteur sensoriel} : terminaison nerveuse libre, corpuscule de Pacini, cellule photoréceptrice, etc.
\item \textbf{Naissance d'un potentiel de récepteur} : le stimulus ouvre des canaux ioniques spécifiques de la membrane du récepteur, ce qui modifie localement le potentiel de membrane. Cette variation est \cle{graduée} : son amplitude est proportionnelle à l'intensité du stimulus. C'est le premier \cle{codage} du message : codage en amplitude.
\item \textbf{Passage au potentiel d'action} : le potentiel de récepteur se propage électrotoniquement (avec décrément) jusqu'à la zone d'intégration du neurone (le cône d'émergence de l'axone, riche en canaux Na$^+$ voltage-dépendants). Si la dépolarisation y atteint le \cle{seuil} (environ \SI{-55}{mV}), un PA se déclenche selon la loi du \cle{tout ou rien}.
\item \textbf{Codage en fréquence} : l'amplitude du PA étant constante, l'information sur l'intensité du stimulus est convertie en \cle{fréquence} des PA : plus le stimulus est intense, plus le potentiel de récepteur est ample, plus les PA se succèdent rapidement (la période réfractaire fixant toutefois une fréquence maximale, de l'ordre de 500 à 1000 PA par seconde).
\item \textbf{Propagation} : le train de PA se propage sans décrément le long de la fibre afférente, à une vitesse dépendant du diamètre et de la myélinisation.
\item \textbf{Traitement par les centres nerveux} : les PA parviennent à la moelle épinière (réflexes) et/ou au cerveau (sensation consciente), où l'information est décodée : la nature de la sensation dépend de la voie empruntée et de l'aire cérébrale activée, son intensité de la fréquence des PA et du nombre de fibres recrutées.
\end{enumerate}

\begin{popfigure}
\begin{center}
\begin{tikzpicture}[font=\small, node distance=0.35cm,
etape/.style={draw=popblue, thick, fill=popblueL, rounded corners=2pt, align=center, minimum height=0.9cm, text width=2.6cm},
code/.style={draw=poporange, thick, fill=poporangeL, rounded corners=2pt, align=center, minimum height=0.7cm, text width=2.6cm},
fleche/.style={-{Stealth[length=2.5mm]}, thick, popdark}]
\node[etape, align=center] (stim){\textbf{Stimulus}\\ (chaleur, pression, lumière)};
\node[etape, right=0.8cm of stim, align=center] (recept){\textbf{Récepteur sensoriel}\\ potentiel de récepteur};
\node[etape, right=0.8cm of recept, align=center] (pa){\textbf{Zone d'intégration}\\ PA si seuil atteint};
\node[etape, right=0.8cm of pa, align=center] (prop){\textbf{Fibre nerveuse}\\ propagation sans décrément};
\node[etape, right=0.8cm of prop, align=center] (centre){\textbf{Centres nerveux}\\ moelle, cerveau};
\draw[fleche] (stim) -- (recept);
\draw[fleche] (recept) -- (pa);
\draw[fleche] (pa) -- (prop);
\draw[fleche] (prop) -- (centre);
\node[code, below=0.5cm of recept, align=center] (c1){Codage en \textbf{amplitude}\\ (gradué)};
\node[code, below=0.5cm of pa, align=center] (c2){Loi du \textbf{tout ou rien}\\ (seuil $\approx \SI{-55}{mV}$)};
\node[code, below=0.5cm of prop, align=center] (c3){Codage en \textbf{fréquence}\\ des PA};
\node[code, below=0.5cm of centre] (c4) {\textbf{Décodage} : nature et intensité de la sensation};
\draw[-{Stealth[length=2mm]}, poporange, dashed] (recept) -- (c1);
\draw[-{Stealth[length=2mm]}, poporange, dashed] (pa) -- (c2);
\draw[-{Stealth[length=2mm]}, poporange, dashed] (prop) -- (c3);
\draw[-{Stealth[length=2mm]}, poporange, dashed] (centre) -- (c4);
\end{tikzpicture}
\end{center}
\legende{Schéma-bilan de la chaîne du message nerveux. Le stimulus engendre un potentiel de récepteur gradué (codage en amplitude) ; si le seuil est atteint, des PA tout ou rien naissent et se propagent (codage en fréquence) jusqu'aux centres nerveux qui décodent l'information.}
\end{popfigure}

\begin{aretenir}[L'essentiel de la chaîne]
Le message nerveux subit deux conversions successives : \textbf{codage en amplitude} au niveau du récepteur (potentiel gradué), puis \textbf{codage en fréquence} au niveau de la fibre (trains de PA d'amplitude constante). Aucune information n'est portée par l'amplitude du PA, toujours identique : c'est la \emph{fréquence} qui renseigne sur l'intensité du stimulus.
\end{aretenir}

\courssub{Tableau comparatif de synthèse : les trois potentiels}

Les sujets de Bac adorent te demander de distinguer ces trois potentiels. Voici le tableau de référence à connaître par cœur.

\begin{center}
\begin{tabularx}{\linewidth}{|l|X|X|X|}
\hline
\rowcolor{popblueL} \textbf{Critère} & \textbf{Potentiel de repos} & \textbf{Potentiel de récepteur} & \textbf{Potentiel d'action} \\
\hline
Valeur / amplitude & \SI{-70}{mV} (constant) & Graduée, proportionnelle au stimulus & Pic à \SI{+30}{mV} \textasciitilde \SI{+40}{mV} (amplitude \SI{\approx 100-110}{mV}, constante, tout ou rien) \\
\hline
Condition d'apparition & État permanent de la membrane au repos & Stimulus d'un récepteur sensoriel & Dépolarisation atteignant le seuil ($\approx \SI{-55}{mV}$) \\
\hline
Ions impliqués & Fuites de K$^+$ sortants ; pompe Na$^+$/K$^+$-ATPase & Ouverture de canaux spécifiques au stimulus (entrées Na$^+$ le plus souvent) & Entrée massive de Na$^+$ (dépolarisation), sortie de K$^+$ (repolarisation) \\
\hline
Propagation & Nulle (état, pas un signal) & Locale, passive, avec décrément & Active, sans décrément, le long de toute la fibre \\
\hline
Codage de l'intensité & Aucun & En amplitude & En fréquence des PA \\
\hline
Période réfractaire & Sans objet & Non (sommation possible) & Oui (absolue puis relative) \\
\hline
\end{tabularx}
\end{center}

\begin{attention}[Ne confonds pas les trois potentiels]
Erreurs fréquentes au Bac : (1) dire que le potentiel de récepteur obéit à la loi du tout ou rien — c'est \textbf{faux}, il est gradué ; (2) dire que le PA se propage avec décrément — c'est le potentiel de récepteur qui s'atténue avec la distance ; (3) affirmer que l'intensité du stimulus modifie l'amplitude du PA — elle n'en modifie que la \textbf{fréquence}. Chaque ligne du tableau ci-dessus est un piège potentiel : vérifie toujours de quel potentiel parle l'énoncé.
\end{attention}

\courssub{Facteurs perturbant le fonctionnement des neurones}

Comprendre les mécanismes ioniques du message nerveux, c'est aussi comprendre comment ils peuvent être bloqués ou déréglés. Classons les perturbations selon l'étape de la chaîne qu'elles affectent.

\textbf{1. Blocage des canaux Na$^+$ voltage-dépendants.} La \cle{tétrodotoxine} (TTX), poison du poisson-globe (fugu), bloque spécifiquement les canaux Na$^+$ : le seuil ne peut plus déclencher d'influx Na$^+$, donc \textbf{aucun PA ne naît ni ne se propage} — paralysie puis arrêt respiratoire. Les \cle{anesthésiques locaux} (lidocaïne, procaïne) exploitent exactement le même mécanisme, mais de façon localisée, transitoire et contrôlée.

\textbf{2. Toxines modifiant la cinétique des canaux.} Certains venins maintiennent les canaux Na$^+$ ouverts (hyperexcitabilité : décharges anarchiques, douleurs intenses, crampes) ou bloquent les canaux K$^+$ (repolarisation impossible, PA prolongé).

\textbf{3. Lésions de la myéline.} Dans la \cle{sclérose en plaques}, la gaine de myéline est détruite par endroits : la conduction saltatoire devient impossible, la vitesse de propagation chute brutalement, voire le PA s'éteint — d'où troubles moteurs, sensitifs et visuels fluctuants.

\textbf{4. Section des fibres.} Un traumatisme (accident de moto, plaie) peut sectionner des axones : la continuité de la voie est rompue, le message ne passe plus — paralysie et anesthésie des territoires en aval.

\textbf{5. Carences et déséquilibres ioniques.} Le potentiel de repos dépend des gradients Na$^+$/K$^+$ : une hypokaliémie (carence en K$^+$) ou des troubles rénaux modifient le potentiel de membrane et perturbent l'excitabilité. Les carences en vitamines du groupe B (B$_1$, B$_{12}$) altèrent la myéline et la survie des neurones (neuropathies, fréquentes en cas de dénutrition ou d'alcoolisme chronique).

\textbf{6. Drogues et alcool.} L'alcool, le cannabis et de nombreuses drogues modifient la libération ou la captation des neurotransmetteurs et la fluidité membranaire : ralentissement des réflexes, troubles de l'équilibre, altération du jugement — autant de causes d'accidents. L'alcoolisme chronique détruit irréversiblement des neurones.

\begin{exemplebox}[La lidocaïne chez le dentiste : un blocage utile]
Tu as mal à une molaire. Le dentiste injecte de la lidocaïne près du nerf alvéolaire : quelques minutes plus tard, il peut soigner sans douleur. Mécanisme : la lidocaïne \textbf{bloque les canaux Na$^+$ voltage-dépendants} des fibres douloureuses locales. Le stimulus (fraise, pression) existe bien et engendre localement une dépolarisation, mais le seuil ne peut plus déclencher de PA : \textbf{le message n'atteint jamais le cerveau}, donc aucune douleur n'est perçue. L'effet disparaît en 1 à 3 heures, le temps que la molécule soit éliminée : les canaux Na$^+$ redeviennent fonctionnels, les neurones sont intacts.
\end{exemplebox}

\begin{attention}[Un anesthésique ne « tue » pas les neurones]
Erreur fréquente de rédaction : écrire que l'anesthésique local « détruit » ou « endort » les neurones. C'est faux : il \textbf{bloque transitoirement les canaux Na$^+$}, empêchant la naissance et la propagation du PA, \textbf{sans détruire la cellule}. Le potentiel de repos est maintenu, la pompe Na$^+$/K$^+$ fonctionne, le neurone reste vivant. À la fin de l'action du produit, la conduction redevient normale. Emploie toujours le vocabulaire exact : « blocage réversible des canaux Na$^+$ voltage-dépendants ».
\end{attention}

\begin{savaistu}[Venins et canaux ioniques : des armes moléculaires redoutables]
Certains venins de scorpions et de serpents (présents dans les savanes et forêts du Cameroun) contiennent des toxines qui agissent directement sur les canaux ioniques : les \emph{$\alpha$-toxines} de scorpion ralentissent l'inactivation des canaux Na$^+$ (décharges répétées, douleurs intenses, sueurs, troubles cardiaques), tandis que la \emph{dendrotoxine} des mambas bloque les canaux K$^+$ (repolarisation impossible, hyperexcitabilité puis paralysie). D'où l'importance capitale de la sensibilisation : en cas de piqûre ou morsure, immobiliser le membre, ne pas inciser, et \textbf{consulter en urgence} un centre de santé disposant de sérum antivenimeux — la rapidité de la prise en charge conditionne le pronostic. Ces toxines sont aussi devenues des outils précieux des chercheurs : c'est grâce à la TTX et à la tétraéthylammonium (bloquant K$^+$) que les rôles respectifs des ions Na$^+$ et K$^+$ dans le PA ont été démontrés.
\end{savaistu}

\courssub{Sensibilisation : protéger son système nerveux au quotidien}

La biologie que tu viens d'apprendre a des conséquences directes sur ta vie. Le neurone est une cellule fragile, peu renouvelable, dont le fonctionnement repose sur des structures fines (canaux, myéline) et des équilibres ioniques délicats. Quelques réflexes simples le protègent :

\begin{itemize}
\item \textbf{Port du casque à moto} : les traumatismes crâniens sont la première cause de lésions cérébrales chez les jeunes ; au Cameroun, les accidents de motos-taxis en sont une cause majeure. Le casque absorbe l'énergie du choc et réduit drastiquement le risque de lésion irréversible.
\item \textbf{Éviter alcool et drogues} : ils perturbent la transmission nerveuse, altèrent les réflexes (temps de réaction allongé, d'où les accidents) et, à terme, détruisent des neurones et démyélinisent des fibres.
\item \textbf{Alimentation équilibrée} : apports suffisants en vitamines B (poissons, légumineuses, céréales) pour l'entretien de la myéline ; éviter les carences qui provoquent des neuropathies.
\item \textbf{Consulter rapidement} en cas de troubles sensitifs (fourmillements persistants, pertes de sensibilité) ou moteurs (faiblesse, tremblements) : ils peuvent révéler une neuropathie, une compression nerveuse ou une intoxication, d'autant mieux traitées qu'elles sont prises tôt.
\item \textbf{Prudence face aux animaux venimeux} et aux poissons ou champignons inconnus : beaucoup de toxines naturelles ciblent précisément les canaux ioniques étudiés dans cette leçon.
\end{itemize}

\courssub{Méthodologie Bac : interpréter un document électrophysiologique}

Les épreuves de SVT te présentent presque toujours des enregistrements de potentiels (oscillogrammes, courbes de perméabilité, tracés en fonction de l'intensité du stimulus). La réussite tient à une démarche rigoureuse, pas à la paraphrase du document.

\begin{methode}[Plan type d'interprétation d'un document électrophysiologique]
\begin{enumerate}
\item \textbf{Décrire les conditions expérimentales} : quel matériel (fibre isolée, nerf entier, récepteur), quel dispositif (microélectrodes, électrodes de surface), quel paramètre varie (intensité, durée, composition du milieu, toxine).
\item \textbf{Dégager les résultats observés, avec des valeurs} : « pour une stimulation de \SI{2}{V}, aucun PA ; à partir de \SI{3}{V}, un PA d'amplitude \SI{110}{mV} apparaît ; à \SI{5}{V}, l'amplitude reste de \SI{110}{mV} ». Cite toujours des chiffres et leurs unités.
\item \textbf{Comparer} les situations (avant/après, avec/sans, intensités croissantes) pour isoler le fait expérimental : « l'amplitude du PA est indépendante de l'intensité du stimulus supraliminaire ».
\item \textbf{Expliquer par les mouvements ioniques} : relier chaque phase aux canaux et aux flux (entrée Na$^+$, sortie K$^+$, pompe Na$^+$/K$^+$, blocage par une toxine).
\item \textbf{Conclure en répondant au problème posé} : formuler la loi ou le mécanisme mis en évidence (tout ou rien, codage en fréquence, rôle de la myéline\ldots).
\end{enumerate}
\end{methode}

\begin{exoresolu}[Calcul d'une vitesse de propagation]
Sur un nerf sciatique de grenouille maintenu en vie, on place deux électrodes réceptrices R$_1$ et R$_2$ distantes de $d = \SI{30}{mm}$ le long du nerf. On stimule le nerf en amont de R$_1$. Les PA sont enregistrés en R$_1$ à l'instant $t_1 = \SI{2,0}{ms}$ et en R$_2$ à l'instant $t_2 = \SI{3,0}{ms}$ après la stimulation. (1) Calculer la vitesse de propagation du PA dans ce nerf. (2) Sachant qu'une fibre amyélinique de même diamètre conduirait à environ \SI{2}{m/s}, proposer une explication.
\tcblower
(1) Le PA met $\Delta t = t_2 - t_1 = \SI{3,0}{ms} - \SI{2,0}{ms} = \SI{1,0}{ms}$ pour parcourir $d = \SI{30}{mm} = \SI{0,030}{m}$. La vitesse de propagation est :
\[
v = \frac{d}{\Delta t} = \frac{\SI{0,030}{m}}{\SI{1,0e-3}{s}} = \SI{30}{m/s}.
\]
(2) La vitesse calculée (\SI{30}{m/s}) est 15 fois supérieure à celle d'une fibre amyélinique de même diamètre (\SI{2}{m/s}). Or la vitesse de propagation augmente avec la myélinisation : la gaine de myéline isole l'axone et confine les échanges ioniques aux nœuds de Ranvier, où le PA est régénéré ; le message « saute » de nœud en nœud (conduction saltatoire). On conclut que le nerf étudié comporte des \textbf{fibres myélinisées}.
\end{exoresolu}

\begin{exoresolu}[Exploitation de la courbe intensité-durée]
On stimule une fibre nerveuse isolée avec des impulsions d'intensité et de durée variables, et on note pour chaque durée l'intensité minimale déclenchant un PA. Résultats : pour une durée de \SI{0,1}{ms}, il faut \SI{8}{V} ; pour \SI{0,5}{ms}, il faut \SI{2}{V} ; pour \SI{1}{ms} et au-delà, \SI{1}{V} suffit ; en dessous de \SI{1}{V}, aucune durée, même très longue, ne déclenche de PA. (1) Décrire la relation entre intensité et durée. (2) Définir la rhéobase et donner sa valeur ici. (3) Expliquer l'existence d'une intensité minimale.
\tcblower
(1) L'intensité minimale efficace \textbf{diminue quand la durée de stimulation augmente}, puis se stabilise : il existe une relation inverse entre l'intensité et la durée du stimulus (courbe intensité-durée hyperbolique). (2) La \cle{rhéobase} est l'intensité minimale capable de déclencher un PA pour une durée de stimulation prolongée ; ici, elle vaut \textbf{\SI{1}{V}}. (3) Le PA ne naît que si la membrane atteint le seuil de dépolarisation ($\approx \SI{-55}{mV}$), c'est-à-dire si suffisamment de canaux Na$^+$ voltage-dépendants s'ouvrent. Une stimulation trop faible, même longue, ne dépolarise pas assez la membrane : les entrées Na$^+$ restent compensées par les fuites K$^+$, le seuil n'est jamais atteint, aucun PA ne naît. Une stimulation brève doit être plus intense pour ouvrir assez de canaux en peu de temps : c'est le compromis intensité-durée.
\end{exoresolu}

\begin{exercice}[Pour t'entraîner]
Un élève enregistre l'activité d'une fibre issue d'un récepteur tactile. Pour une pression de \SI{10}{g}, il observe \SI{20}{PA/s} ; pour \SI{30}{g}, \SI{60}{PA/s} ; pour \SI{60}{g}, \SI{100}{PA/s}. Dans les trois cas, l'amplitude des PA est de \SI{110}{mV}. (1) Quel est le mode de codage de l'intensité du stimulus ? Justifie avec les valeurs. (2) Pourquoi l'amplitude des PA ne varie-t-elle pas ? (3) La fréquence peut-elle augmenter indéfiniment si la pression croît ? Justifie.
\end{exercice}

\begin{corrige}[Exercice — Pour t'entraîner]
(1) L'intensité du stimulus est codée en \textbf{fréquence} des PA : quand la pression passe de \SI{10}{g} à \SI{60}{g}, la fréquence passe de \SI{20}{PA/s} à \SI{100}{PA/s}, tandis que l'amplitude reste fixe (\SI{110}{mV}). (2) Le PA obéit à la loi du tout ou rien : dès que le seuil est atteint, l'entrée massive de Na$^+$ porte toujours le potentiel au même pic ; l'amplitude est donc indépendante de l'intensité du stimulus. (3) Non : la \textbf{période réfractaire absolue} (environ \SI{1}{ms} à \SI{2}{ms}), pendant laquelle les canaux Na$^+$ sont inactivés, interdit le déclenchement d'un nouveau PA immédiat ; la fréquence maximale est donc plafonnée (de l'ordre de 500 à 1000 PA/s).
\end{corrige}

\begin{aretenir}[Le bilan de la leçon en cinq phrases]
(1) Au repos, la membrane est polarisée à \SI{-70}{mV} grâce aux fuites K$^+$ et à la pompe Na$^+$/K$^+$. (2) Le stimulus engendre un potentiel de récepteur gradué. (3) Si le seuil est atteint, un PA tout ou rien naît : entrée Na$^+$, puis sortie K$^+$. (4) Le PA se propage sans décrément, plus vite sur fibre myélinisée, avec une période réfractaire qui impose le codage en fréquence. (5) Toxines, anesthésiques, traumatismes, démyélinisation, carences et drogues perturbent cette machinerie : protéger ses neurones, c'est protéger la communication de tout l'organisme.
\end{aretenir}

\newpage

\coursec{Activité d'intégration}

\coursec{Sensibilisation sur la nécessité de préserver le fonctionnement des neurones}

\courssub{Objectifs de l'activité}
Cette activité d'intégration vise à mobiliser l'ensemble des savoir-faire de la leçon pour résoudre une situation clinique authentique. À l'issue de ce travail, vous devez être capable de :
\begin{itemize}
    \item Décrire et interpréter ioniquement les tracés électrophysiologiques (potentiel de repos, potentiel d'action, réponse à la stimulation).
    \item Calculer la vitesse de propagation de l'influx nerveux et l'interpréter en fonction de l'état de la myéline.
    \item Caractériser le seuil d'excitabilité et la relation intensité-durée à partir de courbes expérimentales.
    \item Expliquer le mécanisme d'action d'un anesthésique local sur les canaux ioniques voltage-dépendants.
    \item Rédiger une conclusion argumentée reliant l'intégrité structurale et fonctionnelle du neurone à la préservation de la sensibilité et de la motricité.
\end{itemize}

\courssub{Situation : Enquête électrophysiologique au Centre Hospitalier Régional d'Ebolowa}
\textbf{Contexte.} Monsieur \textbf{Mbarga}, 45 ans, cultivateur de cacao dans la région du Sud (Ebolowa), consulte au service de neurologie du Centre Hospitalier Régional pour une \emph{perte de sensibilité progressive des membres inférieurs} (hypoesthésie), des fourmillements (paresthésies) et une difficulté à marcher sur la pointe des pieds, évoluant depuis 6 mois. Il signale une consommation chronique d'alcool artisanal (\emph{odontol}) et un diabète de type 2 mal équilibré (glycémie à jeun \SI{1,80}{g/L}). Le neurologue suspecte une \textbf{neuropathie périphérique sensitive et motrice} (polyneuropathie diabétique et alcoolique).

Pour confirmer le diagnostic topographique et étiologique, le médecin réalise un \textbf{électroneuromyogramme (ENMG)} sur le nerf sciatique poplité externe (nerf mixte sensitivo-moteur) au niveau de la jambe. Quatre documents sont obtenus.

\begin{popfigure}
\centering
\begin{tikzpicture}[scale=0.8, every node/.style={font=\small}]
% Axes for Resting Potential
\begin{scope}[xshift=0cm, yshift=0cm]
\draw[->] (-0.5,0) -- (5,0) node[right] {$t$ (ms)};
\draw[->] (0,-1) -- (0,1.5) node[above] {$V_m$ (mV)};
\draw[thick, popblue] (0, -0.72) -- (4.5, -0.72);
\draw[dashed, gray] (0, -0.72) node[left] {-72} -- (4.5, -0.72);
\foreach \x in {0,1,2,3,4} \draw (\x, -0.05) -- (\x, 0.05) node[below] {\x};
\foreach \y in {-70,-50,0,50} \draw (-0.05, \y/100) -- (0.05, \y/100) node[left] {\y};
\node[popblue] at (2.25, 0.5) {\textbf{Doc 1a : Potentiel de repos}};
\end{scope}

% Axes for Action Potential
\begin{scope}[xshift=7cm, yshift=0cm]
\draw[->] (-0.5,0) -- (5,0) node[right] {$t$ (ms)};
\draw[->] (0,-1) -- (0,1.5) node[above] {$V_m$ (mV)};
% AP shape corrected: peak at +35mV (y=0.35), hyperpol at -90mV (y=-0.9)
\draw[thick, poporange] (0, -0.72) -- (0.5, -0.72)
    .. controls (0.8, -0.72) and (0.8, 0.35) .. (1.0, 0.35) % overshoot +35mV
    .. controls (1.2, 0.35) and (1.2, -0.9) .. (1.5, -0.9) % hyperpol -90mV
    .. controls (1.8, -0.9) and (2.0, -0.72) .. (2.0, -0.72)
    -- (4.5, -0.72);
\draw[dashed, gray] (0, -0.72) node[left] {-72} -- (4.5, -0.72);
\draw[dashed, gray] (0, 0.35) node[left] {+35} -- (4.5, 0.35);
\foreach \x in {0,1,2,3,4} \draw (\x, -0.05) -- (\x, 0.05) node[below] {\x};
\foreach \y in {-70,-50,0,50} \draw (-0.05, \y/100) -- (0.05, \y/100) node[left] {\y};
\node[poporange] at (2.25, 0.5) {\textbf{Doc 1b : Potentiel d'action}};
\end{scope}

% Table concentrations
\begin{scope}[xshift=0cm, yshift=-5.5cm]
\node[anchor=north west, align=left, text width=13cm] {
    \textbf{Tableau 1 : Concentrations ioniques intracellulaires et extracellulaires (mM) au niveau de l'axone.}
    \begin{center}
    \begin{tabular}{|c|c|c|c|}\hline
    Ion & $[ion]_{in}$ (mM) & $[ion]_{out}$ (mM) & $E_{ion}$ calculé (mV) \\ \hline
    \ce{Na+} & 15 & 145 & \textbf{?} \\ \hline
    \ce{K+} & 140 & 4 & \textbf{?} \\ \hline
    \ce{Cl-} & 10 & 120 & \textbf{?} \\ \hline
    \end{tabular}
    \end{center}
};
\end{scope}
\end{tikzpicture}
\legende{Document 1 : Enregistrement intracellulaire (microélectrode) sur une fibre axonale saine du nerf sciatique de M. Mbarga (température \SI{37}{\celsius}). À gauche : potentiel de repos stable. À droite : potentiel d'action évoqué par une stimulation électrique brève supra-seuil. Le tableau donne les concentrations ioniques.}
\end{popfigure}

\begin{popfigure}
\centering
\begin{tikzpicture}[scale=0.75, every node/.style={font=\small}]
% Healthy Subject
\begin{scope}[xshift=0cm]
\draw[->] (0,0) -- (5.5,0) node[right] {Intensité stimulus $I$ (mA)};
\draw[->] (0,0) -- (0,4.5) node[above] {Amplitude réponse $A$ (mV)};
% Sigmoid curve healthy
\draw[thick, popgreen, domain=0:5, samples=100] plot (\x, {3.5/(1+exp(-3*(\x-0.8)))});
\draw[dashed, gray] (0.8, 0) node[below] {$I_{seuil}=0,3$} -- (0.8, 1.75);
\draw[dashed, gray] (0, 3.5) node[left] {$A_{max}=10$} -- (5, 3.5);
\node[popgreen] at (2.5, 4) {\textbf{Sujet témoin sain}};
\end{scope}

% Patient
\begin{scope}[xshift=7.5cm]
\draw[->] (0,0) -- (5.5,0) node[right] {Intensité stimulus $I$ (mA)};
\draw[->] (0,0) -- (0,4.5) node[above] {Amplitude réponse $A$ (mV)};
% Sigmoid curve patient - shifted right, lower max
\draw[thick, popink, domain=0:5, samples=100] plot (\x, {1.05/(1+exp(-1.5*(\x-2.5)))});
\draw[dashed, gray] (2.5, 0) node[below] {$I_{seuil}=1,2$} -- (2.5, 0.525);
\draw[dashed, gray] (0, 1.05) node[left] {$A_{max}=3$} -- (5, 1.05);
\node[popink] at (2.5, 4) {\textbf{Patient M. Mbarga}};
\end{scope}
\end{tikzpicture}
\legende{Document 2 : Courbes de recrutement (relation intensité-amplitude de la réponse motrice composée) obtenues par stimulation électrique du nerf sciatique poplité externe au niveau du pli poplité, enregistrement au muscle court fibulaire (pied). Le stimulus est une impulsion carrée de durée \SI{0,1}{ms}.}
\end{popfigure}

\begin{popfigure}
\centering
\begin{tikzpicture}[scale=0.8, every node/.style={font=\small}]
% Healthy: 2 electrodes, 8cm, delay 4ms
\begin{scope}[xshift=0cm]
\draw[->] (-0.5,0) -- (6.5,0) node[right] {$t$ (ms)};
\draw[->] (0,-0.5) -- (0,2.5) node[above] {Tension (mV)};
% Stim artifact at 0
\draw[thick, popdark] (0, 0) -- (0, 2) -- (0.1, 2) -- (0.1, 0);
% Response 1 (proximal) at 2ms
\draw[thick, popblue] (1.8, 0) -- (2.0, 1.5) -- (2.2, -0.5) -- (2.4, 0);
\node[popblue, above] at (2.0, 1.5) {$E_1$ (proximale)};
% Response 2 (distal) at 6ms (delay 4ms)
\draw[thick, poporange] (5.8, 0) -- (6.0, 1.2) -- (6.2, -0.4) -- (6.4, 0);
\node[poporange, above] at (6.0, 1.2) {$E_2$ (distale)};
\draw[<->, gray] (2.0, -0.3) -- (6.0, -0.3) node[midway, below] {$\Delta t = 4,0$ ms};
\node[align=center] at (3, 2.2) {\textbf{Sujet témoin sain}\\\textit{Distance $E_1-E_2 = 8,0$ cm}};
\end{scope}

% Patient: delay 16ms - extended x axis to 20
\begin{scope}[xshift=8.5cm]
\draw[->] (-0.5,0) -- (20,0) node[right] {$t$ (ms)};
\draw[->] (0,-0.5) -- (0,2.5) node[above] {Tension (mV)};
% Stim artifact
\draw[thick, popdark] (0, 0) -- (0, 2) -- (0.1, 2) -- (0.1, 0);
% Response 1
\draw[thick, popblue] (1.8, 0) -- (2.0, 1.0) -- (2.2, -0.3) -- (2.4, 0);
\node[popblue, above] at (2.0, 1.0) {$E_1$};
% Response 2 at 18ms (delay 16ms)
\draw[thick, poporange] (17.8, 0) -- (18.0, 0.8) -- (18.2, -0.2) -- (18.4, 0);
\node[poporange, above] at (18.0, 0.8) {$E_2$};
\draw[<->, gray] (2.0, -0.3) -- (18.0, -0.3) node[midway, below] {$\Delta t = 16,0$ ms};
\node[align=center] at (10, 2.2) {\textbf{Patient M. Mbarga}\\\textit{Distance $E_1-E_2 = 8,0$ cm}};
\end{scope}
\end{tikzpicture}
\legende{Document 3 : Montage à deux électrodes de surface (électrodes annulaires) espacées de \SI{8,0}{cm} sur le trajet du nerf sciatique poplité externe. Stimulation unique supra-seuil en amont de $E_1$. Mesure de la latence (délai) entre le passage de l'influx aux deux points d'enregistrement.}
\end{popfigure}

\begin{popfigure}
\centering
\begin{tikzpicture}[scale=0.8, every node/.style={font=\small}]
% Normal AP
\begin{scope}[xshift=0cm]
\draw[->] (-0.5,0) -- (5,0) node[right] {$t$ (ms)};
\draw[->] (0,-1) -- (0,1.5) node[above] {$V_m$ (mV)};
% Corrected peak at 0.35
\draw[thick, popgreen] (0, -0.72) -- (0.5, -0.72)
    .. controls (0.8, -0.72) and (0.8, 0.35) .. (1.0, 0.35)
    .. controls (1.2, 0.35) and (1.2, -0.9) .. (1.5, -0.9)
    .. controls (1.8, -0.9) and (2.0, -0.72) .. (2.0, -0.72)
    -- (4.5, -0.72);
\draw[dashed, gray] (0, -0.72) -- (4.5, -0.72);
\node[popgreen] at (2.25, 0.5) {\textbf{Avant Lidocaïne}};
\end{scope}

% With Lidocaine
\begin{scope}[xshift=7cm]
\draw[->] (-0.5,0) -- (5,0) node[right] {$t$ (ms)};
\draw[->] (0,-1) -- (0,1.5) node[above] {$V_m$ (mV)};
% Flat line at rest, small local response
\draw[thick, poppink] (0, -0.72) -- (4.5, -0.72);
\draw[thick, poppink, dashed] (0.5, -0.72) .. controls (0.8, -0.72) and (0.8, -0.4) .. (1.0, -0.4) .. controls (1.2, -0.4) and (1.2, -0.72) .. (1.5, -0.72);
\node[poppink] at (2.25, 0.5) {\textbf{Après Lidocaïne 1\%}};
\node[poppink, align=center] at (2.5, -0.2) {Réponse locale\\décroissante (pas de PA)};
\end{scope}
\end{tikzpicture}
\legende{Document 4 : Effet de l'application locale de \SI{1}{\milli\liter} de Lidocaïne chlorhydrate à 1\% (anesthésique local de type amide) sur la fibre axonale. Enregistrement intracellulaire dans les mêmes conditions que le Document 1. Le potentiel de repos reste stable à -72 mV.}
\end{popfigure}

\courssub{Consignes}
\begin{enumerate}
    \item \textbf{Document 1 :} 
    \begin{enumerate}
        \item Décrire les deux tracés (valeurs numériques, phases). 
        \item Calculer les potentiels d'équilibre de Nernst pour \ce{Na+}, \ce{K+} et \ce{Cl-} à \SI{37}{\celsius} (utiliser la formule simplifiée $E = \frac{60}{z} \log_{10} \frac{[ion]_{out}}{[ion]_{in}}$). 
        \item Interpréter la valeur du potentiel de repos (-72 mV) par rapport aux potentiels d'équilibre calculés (équation de Goldman-Hodgkin-Katz). 
        \item Expliquer ioniquement les phases de dépolarisation, repolarisation et hyperpolarisation du potentiel d'action.
    \end{enumerate}
    \item \textbf{Document 2 :} 
    \begin{enumerate}
        \item Définir la notion de \cle{seuil d'excitabilité} (ou seuil d'intensité) pour une fibre et pour un nerf. 
        \item Comparer les seuils et les amplitudes maximales entre le sujet sain et le patient. 
        \item Proposer une explication physiopathologique à l'augmentation du seuil et à la diminution de l'amplitude maximale chez M. Mbarga (distinction atteinte démyélinisante / axonale).
    \end{enumerate}
    \item \textbf{Document 3 :} 
    \begin{enumerate}
        \item Calculer la vitesse de propagation de l'influx nerveux ($v$) pour le sujet sain et pour le patient ($v = \frac{d}{\Delta t}$). 
        \item Comparer ces valeurs aux vitesses de référence (fibres myélinisées A$\alpha$/A$\beta$ : 50-100 m/s ; fibres amyéliniques C : 0,5-2 m/s). 
        \item En déduire la nature de l'atteinte (démyélinisation vs dégénérescence axonale) et justifier par le mécanisme de la conduction saltatoire.
    \end{enumerate}
    \item \textbf{Document 4 :} 
    \begin{enumerate}
        \item Décrire l'effet de la lidocaïne sur le potentiel d'action et sur le potentiel de repos. 
        \item Expliquer le mécanisme d'action moléculaire de la lidocaïne sur les canaux \ce{Na+} voltage-dépendants (blocage \emph{use-dependent}). 
        \item Pourquoi le potentiel de repos n'est-il pas modifié ?
    \end{enumerate}
    \item \textbf{Synthèse clinique :} Rédiger un compte-rendu argumenté (15-20 lignes) adressé au médecin traitant, concluant sur la nature de la neuropathie de M. Mbarga et insistant sur les mesures de prévention pour préserver le fonctionnement des neurones (contrôle glycémique, sevrage alcoolique, prévention des traumatismes, vitamine B12).
\end{enumerate}

\courssub{Ressources à mobiliser}
\begin{itemize}
    \item \textbf{Équation de Nernst (37°C) :} $E_{ion} = \frac{60}{z} \log_{10} \frac{[ion]_{out}}{[ion]_{in}}$ (en mV). $z$ = charge de l'ion (+1 pour \ce{Na+}, \ce{K+} ; -1 pour \ce{Cl-}).
    \item \textbf{Équation de Goldman-Hodgkin-Katz (GHK) :} $V_m = 60 \log_{10} \frac{P_K[K^+]_{out} + P_{Na}[Na^+]_{out} + P_{Cl}[Cl^-]_{in}}{P_K[K^+]_{in} + P_{Na}[Na^+]_{in} + P_{Cl}[Cl^-]_{out}}$. Au repos : $P_K \gg P_{Na}, P_{Cl}$.
    \item \textbf{Vitesse de propagation :} $v = \frac{d}{\Delta t}$ ($d$ en m, $\Delta t$ en s, $v$ en m/s).
    \item \textbf{Seuil d'excitabilité :} Intensité minimale du stimulus pour déclencher un PA (fibre) ou une réponse composée (nerf). Relation intensité-durée : $I = \frac{I_{rh}}{1 - e^{-t/\tau}}$ (courbe de Weiss-Lapicque), rhéobase $I_{rh}$, chronaxie $\tau$.
    \item \textbf{Période réfractaire :} Absolue (canaux \ce{Na+} inactivés) et relative (canaux \ce{K+} ouverts, hyperpolarisation).
    \item \textbf{Anesthésiques locaux :} Blocage intracellulaire des canaux \ce{Na+} voltage-dépendants (état ouvert/inactivé), empêchant la dépolarisation. N'agissent pas sur la pompe \ce{Na+/K+} ATPase ni les canaux de fuite \ce{K+}.
\end{itemize}

\courssub{Démarche et corrigé commenté}
\begin{exoresolu}[Corrigé détaillé de l'enquête électrophysiologique]
\textbf{1. Analyse du Document 1 (Potentiel de repos et Potentiel d'action)}

\textbf{1.a Description des tracés.}
\begin{itemize}
    \item \textbf{Potentiel de repos (Doc 1a) :} Tracé plat, stable à \textbf{-72 mV} (valeur moyenne entre -70 et -75 mV pour un axone myélinisé de mammifère à 37°C). Aucune fluctuation spontanée n'est visible, signe d'une membrane stable au repos.
    \item \textbf{Potentiel d'action (Doc 1b) :} Réponse « tout ou rien » à une stimulation brève supra-seuil. Phases successives :
    \begin{enumerate}
        \item \textbf{Dépolarisation rapide :} Passage de -72 mV à \textbf{+35 mV} (overshoot) en \textasciitilde 0,5 ms. Pente ascendante très raide.
        \item \textbf{Repolarisation :} Retour rapide vers les valeurs négatives, franchissant le potentiel de repos.
        \item \textbf{Hyperpolarisation (potentiel de suite) :} Dépassement négatif jusqu'à \textbf{-90 mV} environ, suivi d'un retour exponentiel lent au potentiel de repos (\textasciitilde 2-3 ms).
    \end{enumerate}
\end{itemize}

\textbf{1.b Calcul des potentiels d'équilibre de Nernst (37°C).}
Formule : $E_{ion} = \frac{60}{z} \log_{10} \frac{[ion]_{out}}{[ion]_{in}}$.
\begin{align*}
E_{\ce{Na+}} &= \frac{60}{+1} \log_{10} \frac{145}{15} = 60 \times \log_{10}(9,67) \approx 60 \times 0,985 = \mathbf{+59,1\; mV} \\
E_{\ce{K+}}  &= \frac{60}{+1} \log_{10} \frac{4}{140} = 60 \times \log_{10}(0,0286) \approx 60 \times (-1,544) = \mathbf{-92,6\; mV} \\
E_{\ce{Cl-}} &= \frac{60}{-1} \log_{10} \frac{120}{10} = -60 \times \log_{10}(12) \approx -60 \times 1,079 = \mathbf{-64,7\; mV}
\end{align*}
\textit{Remarque :} Les valeurs arrondies usuelles sont $E_{\ce{Na+}} \approx +60$ mV, $E_{\ce{K+}} \approx -90$ mV, $E_{\ce{Cl-}} \approx -65$ mV.

\textbf{1.c Interprétation ionique du potentiel de repos (-72 mV).}
Le potentiel de repos ($V_m = -72$ mV) est proche de $E_{\ce{K+}}$ (-92,6 mV) mais s'en écarte vers des valeurs moins négatives. Selon l'équation de Goldman-Hodgkin-Katz (GHK), $V_m$ est une moyenne pondérée des potentiels d'équilibre par les \textbf{perméabilités membranaires} ($P_{ion}$).
Au repos : $P_{\ce{K+}} \gg P_{\ce{Na+}}, P_{\ce{Cl-}}$ (canaux de fuite \ce{K+} ouverts, canaux \ce{Na+} fermés, pompe \ce{Na+/K+} ATPase active : 3 \ce{Na+}$_{out}$ / 2 \ce{K+}$_{in}$).
$V_m$ est « tiré » vers $E_{\ce{K+}}$ par la forte perméabilité au \ce{K+}, mais la faible perméabilité résiduelle au \ce{Na+} ($P_{\ce{Na+}} \approx 0,01 - 0,05 P_{\ce{K+}}$) et l'activité électrogène de la pompe (contribution \textasciitilde -3 à -5 mV) le maintiennent à -72 mV, moins négatif que $E_{\ce{K+}}$.

\textbf{1.d Mécanismes ioniques du potentiel d'action (Hodgkin-Huxley).}
\begin{itemize}
    \item \textbf{Dépolarisation :} Ouverture brutale des \textbf{canaux \ce{Na+} voltage-dépendants} (activation $m^3$). $P_{\ce{Na+}}$ augmente \textasciitilde 1000 fois. $V_m$ tend vers $E_{\ce{Na+}}$ (+60 mV). Influx massif de \ce{Na+} ($I_{\ce{Na+}}$ entrant).
    \item \textbf{Repolarisation :} Double mécanisme : (1) \textbf{Inactivation} des canaux \ce{Na+} (fermeture balle-et-chaîne, variable $h \to 0$) stoppant l'influx \ce{Na+} ; (2) Ouverture retardée des \textbf{canaux \ce{K+} voltage-dépendants} (activation $n^4$). $P_{\ce{K+}}$ augmente fort. Efflux massif de \ce{K+} ($I_{\ce{K+}}$ sortant). $V_m$ redescend vers $E_{\ce{K+}}$.
    \item \textbf{Hyperpolarisation :} Les canaux \ce{K+} restent ouverts un certain temps après le retour de $V_m$ au repos (inactivation lente). $P_{\ce{K+}}$ reste $> P_{\ce{K+}}^{repos}$. $V_m$ se rapproche temporairement de $E_{\ce{K+}}$ (-90 mV). La pompe \ce{Na+/K+} restaure ensuite les gradients.
\end{itemize}

\textbf{2. Analyse du Document 2 (Courbes de recrutement : Seuil et Amplitude)}

\textbf{2.a Définition du seuil.}
\begin{itemize}
    \item \textbf{Fibre (axone unique) :} Intensité minimale de courant ($I_{seuil}$) nécessaire pour amener le potentiel membranaire au \textbf{seuil de dépolarisation} (\textasciitilde -55 mV), déclenchant l'ouverture regenerative des canaux \ce{Na+}. C'est une propriété « tout ou rien ».
    \item \textbf{Nerf (faisceau de fibres) :} Le seuil du nerf correspond au seuil de la fibre la plus excitable (plus gros diamètre, myéline épaisse). L'amplitude de la réponse composée (somme des PA synchrones) croît avec l'intensité du stimulus car on \textbf{recrute} progressivement des fibres de seuils plus élevés (fibres plus fines).
\end{itemize}

\textbf{2.b Comparaison Sujet Sain vs Patient.}
\begin{center}
\begin{tabular}{|l|c|c|}\hline
\textbf{Paramètre} & \textbf{Sujet Sain} & \textbf{Patient M. Mbarga} \\ \hline
Seuil d'intensité ($I_{seuil}$) & \SI{0,3}{mA} & \SI{1,2}{mA} (\textbf{x4}) \\ \hline
Amplitude max ($A_{max}$) & \SI{10}{mV} & \SI{3}{mV} (\textbf{/3,3}) \\ \hline
\end{tabular}
\end{center}
Le patient présente une \textbf{élévation nette du seuil} (hypoexcitabilité) et une \textbf{réduction majeure de l'amplitude maximale}.

\textbf{2.c Explication physiopathologique.}
\begin{itemize}
    \item \textbf{Seuil élevé (x4) :} Suggère une \textbf{atteinte démyélinisante}. La perte de myéline expose la membrane axonnaire (nœuds de Ranvier élargis, paranodes démyélinisés) : fuite des courants locaux (diminution de la résistance membranaire $R_m$), augmentation de la capacitance $C_m$. Il faut un courant injecté plus fort ($I = C_m \frac{dV}{dt} + \frac{V}{R_m}$) pour charger la membrane jusqu'au seuil de dépolarisation. La chronaxie est augmentée.
    \item \textbf{Amplitude réduite (/3,3) :} Suggère une \textbf{perte axonale (dégénérescence axonale secondaire)}. Moins de fibres fonctionnelles sont disponibles pour être recrutées. L'amplitude de la réponse composée est proportionnelle au nombre d'axones myélinisés intacts. Dans une neuropathie diabétique/alcoolique mixte, on observe souvent une \textbf{neuropathie axonale prédominante avec démyélinisation secondaire} (ou l'inverse selon l'ancienneté). Ici, la conjonction des deux signes oriente vers une atteinte mixte sévère.
\end{itemize}

\textbf{3. Analyse du Document 3 (Vitesse de conduction)}

\textbf{3.a Calcul des vitesses.}
Distance inter-électrodes $d = \SI{8,0}{cm} = \SI{0,080}{m}$.
\begin{itemize}
    \item \textbf{Sujet sain :} $\Delta t = \SI{4,0}{ms} = 4,0 \times 10^{-3}$ s.
    $v_{sain} = \frac{0,080}{4,0 \times 10^{-3}} = \mathbf{20,0\; m/s}$.
    \item \textbf{Patient :} $\Delta t = \SI{16,0}{ms} = 16,0 \times 10^{-3}$ s.
    $v_{patient} = \frac{0,080}{16,0 \times 10^{-3}} = \mathbf{5,0\; m/s}$.
\end{itemize}

\textbf{3.b Comparaison aux références.}
\begin{itemize}
    \item $v_{sain} = 20$ m/s : Correspond à des fibres \textbf{myélinisées de petit calibre} (fibres A$\delta$ ou A$\beta$ fines, ou fibres motrices gamma). Le nerf sciatique poplité externe contient des fibres motrices (A$\alpha$ \textasciitilde 80-100 m/s) et sensitives (A$\beta$ \textasciitilde 50-70 m/s). Une vitesse de 20 m/s sur un trajet court (8 cm) avec électrodes de surface peut sous-estimer la vitesse réelle (latence distale, temps de couplage neuromusculaire non soustrait ici). Cependant, l'ordre de grandeur confirme une conduction \textbf{myélinisée (saltatoire)}.
    \item $v_{patient} = 5$ m/s : Vitesse \textbf{4 fois plus lente}. Elle se situe dans la gamme des fibres \textbf{amyéliniques (fibres C : 0,5-2 m/s) ou fibres myélinisées sévèrement démyélinisées}. Une conduction à 5 m/s sur un nerf normalement myélinisé est le signe pathognomonique d'une \textbf{démyélinisation segmentaire/diffuse} avec perte de la conduction saltatoire. Le courant doit se propager de façon continue (mode passif) le long de zones dénudées, ce qui est lent.
\end{itemize}

\textbf{3.c Nature de l'atteinte.}
La chute drastique de la vitesse (20 $\to$ 5 m/s) est l'argument majeur pour une \textbf{neuropathie démyélinisante} (ou composante démyélinisante majeure). La myéline assure l'isolation électrique (haute $R_m$, basse $C_m$) permettant la conduction saltatoire (saut du nœud au nœud). Sa destruction force le courant à traverser la membrane sur toute la surface démyélinisée (conduction continue), augmentant la constante de temps $\tau = R_m C_m$ et diminuant la constante de longueur $\lambda = \sqrt{R_m / R_i}$, d'où une vitesse $v \propto \lambda / \tau$ effondrée.

\textbf{4. Analyse du Document 4 (Effet de la Lidocaïne)}

\textbf{4.a Description de l'effet.}
La lidocaïne \textbf{abolitionne le potentiel d'action} (plus de dépolarisation regenerative, plus d'overshoot). Seule persiste une \textbf{réponse locale dépolarisante passive} (amplitude faible, \textasciitilde -40 mV, pas de tout-ou-rien), qui décroît avec la distance. Le \textbf{potentiel de repos reste inchangé à -72 mV}.

\textbf{4.b Mécanisme moléculaire (Blocage \emph{use-dependent}).}
La lidocaïne (base faible, pKa \textasciitilde 7,9) traverse la membrane sous forme non ionisée (base libre), puis se protonne dans le cytoplasme (forme cationique). Elle se lie à son \textbf{site récepteur intracellulaire} sur le canal \ce{Na+} voltage-dépendant (domaine S6 des répétitions I, III, IV), \textbf{uniquement lorsque le canal est dans l'état OUVERT ou INACTIVÉ} (dépendant de l'état/usage).
\begin{itemize}
    \item Elle empêche le passage des ions \ce{Na+} (blocage du pore).
    \item Elle stabilise l'état inactivé, décalant la courbe d'inactivation vers des potentiels plus négatifs.
    \item \textbf{Conséquence :} Impossible d'ouvrir suffisamment de canaux \ce{Na+} pour atteindre le seuil critique ($g_{\ce{Na+}} > g_{\ce{K+}}$). La dépolarisation regenerative (boucle de rétroaction positive) est bloquée.
\end{itemize}

\textbf{4.c Pourquoi le potentiel de repos est intact ?}
Le potentiel de repos dépend des \textbf{canaux de fuite \ce{K+} (Kir, TREK)} et de la \textbf{pompe \ce{Na+/K+} ATPase}. La lidocaïne n'affecte ni les canaux de fuite \ce{K+}, ni la pompe ATPase, ni les canaux \ce{Cl-}. Elle est spécifique des canaux \ce{Na+} voltage-dépendants (Nav1.1 à Nav1.9). Au repos, ces canaux sont \textbf{fermés} (état de repos), et la lidocaïne a une très faible affinité pour l'état fermé. Donc pas de blocage au repos.

\textbf{5. Synthèse clinique : Compte-rendu d'ENMG pour M. Mbarga}

\textit{Objet : Compte-rendu d'électroneuromyographie des membres inférieurs - Neuropathie périphérique mixte démyélinisante et axonale sévère.}

Monsieur le Médecin,

L'examen électrophysiologique du nerf sciatique poplité externe de M. Mbarga met en évidence une \textbf{atteinte sévère des fibres nerveuses périphériques de type mixte (sensitivo-motrice)}.

L'analyse des vitesses de conduction motrice (VCM) montre un \textbf{ralentissement marqué} (\SI{5}{m/s} vs \SI{20}{m/s} côté controlatéral/sain), témoignant d'une \textbf{conduction non saltatoire} compatible avec une \textbf{démyélinisation segmentaire diffuse} (allongement des nœuds de Ranvier, perte de l'isolation myélique). Parallèlement, l'étude du recrutement (Document 2) révèle une \textbf{élévation du seuil d'excitabilité} (x4) cohérente avec la démyélinisation, et une \textbf{réduction de l'amplitude maximale de la réponse composée} (\SI{3}{mV} vs \SI{10}{mV}), signe d'une \textbf{perte axonale significative} (dégénérescence de Waller secondaire ou axonopathie primitive).

Le mécanisme biophysique est double : la démyélinisation augmente la capacitance membranaire et diminue la résistance, nécessitant un courant de charge plus important (seuil élevé) et ralentissant la propagation (vitesse effondrée) ; la perte axonale réduit le nombre d'unités motrices recrutables (amplitude basse).

L'essai pharmacologique (lidocaïne, Doc 4) confirme l'intégrité de la machinerie ionique de base (pompe \ce{Na+/K+}, canaux de fuite \ce{K+}) mais démontre la vulnérabilité des canaux \ce{Na+} voltage-dépendants, cibles de l'anesthésie locale.

\textbf{Conclusion étiologique :} Le tableau clinique (diabète déséquilibré, alcoolisme chronique) et l'ENMG orientent vers une \textbf{polyneuropathie diabétique et alcoolique mixte, à composante démyélinisante et axonale sévère}.

\textbf{Recommandations pour préserver le fonctionnement neuronal :}
\begin{enumerate}
    \item \textbf{Équilibrage glycémique strict} (HbA1c < 7\%) pour limiter la glycation des protéines structurales (myéline, cytosquelette) et le stress oxydatif mitochondrial (voie des polyols, AGEs).
    \item \textbf{Sevrage alcoolique accompagné} (suppression de l'éthanol neurotoxique et correction de la carence en \textbf{vitamine B1 (thiamine)} et \textbf{B12}, cofacteurs essentiels du métabolisme énergétique axonal et de la synthèse de myéline).
    \item \textbf{Prévention des traumatismes} (chaussures adaptées, inspection plantaire quotidienne) compte tenu de l'insensibilité (risque d'ulcères plantaires perforants).
    \item \textbf{Kinésithérapie neuroproprioceptive} pour entretenir la plasticité des circuits moteurs résiduels.
\end{enumerate}
La préservation de l'intégrité structurale (myéline, axone, mitochondries) et fonctionnelle (canaux ioniques, pompes, transport axonal) des neurones est la condition \emph{sine qua non} du maintien de l'autonomie motrice et sensitive. Tout retard thérapeutique aggrave irréversiblement le pronostic fonctionnel.

Respectueusement, \textit{Le Neurophysiologiste.}
\tcblower
\textbf{Commentaires pédagogiques (pour l'enseignant) :}
\begin{itemize}
    \item \textbf{Doc 1 :} Vérifier que l'élève utilise la formule de Nernst avec $z$ signé. Insister sur le fait que $V_m \neq E_K$ à cause de $P_{Na} > 0$ et de la pompe électrogène. Pour le PA, exiger le vocabulaire : « ouverture canaux Na+ voltage-dépendants $\to$ influx Na+ $\to$ dépolarisation », « inactivation Na+ + ouverture K+ retardée $\to$ efflux K+ $\to$ repolarisation », « inactivation lente K+ $\to$ hyperpolarisation ».
    \item \textbf{Doc 2 :} Distinguer clairement seuil de la fibre (biophysique) et seuil du nerf (recrutement). L'amplitude max reflète le nombre d'axones. L'élève doit citer « démyélinisation $\to$ seuil $\uparrow$ » et « perte axonale $\to$ amplitude $\downarrow$ ».
    \item \textbf{Doc 3 :} Calcul simple mais attention aux unités (cm $\to$ m, ms $\to$ s). L'interprétation de 5 m/s est cruciale : c'est la preuve de la conduction continue (démyélinisation). Relier à la période réfractaire : la démyélinisation allonge la période réfractaire absolue (récupération canaux Na+ plus lente) et limite la fréquence maximale de décharge.
    \item \textbf{Doc 4 :} Mécanisme « use-dependent » (dépendant de l'usage/état) souvent mal compris. Préciser : forme basique traverse membrane $\to$ forme cationique intracellulaire $\to$ blocage pore côté cytoplasmique $\to$ spécifique état Ouvert/Inactivé. Repos = état Fermé $\to$ pas de blocage.
    \item \textbf{Synthèse :} Évaluer la capacité à écrire un texte médical structuré, utilisant le vocabulaire précis (démyélinisation, dégénérescence axonale, conduction saltatoire, recrutement) et liant biophysique et prévention.
\end{itemize}
\end{exoresolu}

\courssub{Bilan de l'activité}
Cette activité d'intégration a permis de mettre en évidence, à travers une démarche d'investigation clinique réaliste (contexte camerounais : diabète, alcool artisanal, hôpital régional), l'articulation fondamentale entre :
\begin{enumerate}
    \item \textbf{La biophysique de la membrane} (potentiels de repos et d'action, équations de Nernst et Goldman, canaux ioniques voltage-dépendants) ;
    \item \textbf{L'électrophysiologie clinique} (seuil, amplitude, vitesse de conduction, période réfractaire) comme fenêtre fonctionnelle sur l'état structural du nerf ;
    \item \textbf{La physiopathologie} (démyélinisation vs axonalopathie) distinguée par des critères quantitatifs (vitesse, seuil, amplitude) ;
    \item \textbf{La pharmacologie ciblée} (lidocaïne) validant le rôle clé des canaux \ce{Na+} dans la genèse du message nerveux.
\end{enumerate}
Elle ancre la \textbf{sensibilisation} demandée par le programme : le neurone est une machine moléculaire de haute précision (pompes, canaux, myéline, mitochondria, transport axonal) dont le dysfonctionnement — métabolique (diabète), toxique (alcool, carences), traumatique — entraîne une perte irréversible de l'information nerveuse. La prévention (éducation sanitaire, dépistage précoce, prise en charge étiologique) est la seule stratégie pour préserver ce capital neuronal non renouvelable. Cette compétence d'analyse de documents expérimentaux et cliniques est directement évaluée au Baccalauréat (épreuves écrites et pratiques).

\newpage

\coursec{Méthodes et exercices résolus}

\courssub{Méthodes et exercices résolus}

\begin{methode}[Mise en évidence et interprétation ionique du potentiel de repos]
\textbf{Objectif :} Réaliser le schéma du montage expérimental (technique de la microélectrode) et expliquer l'origine ionique du potentiel de repos (environ \SI{-70}{mV} à \SI{-90}{mV}).

\textbf{Étapes de la démarche :}
\begin{enumerate}
    \item \textbf{Monter l'expérience :} Représenter une fibre nerveuse (axone) baignée dans un liquide extracellulaire. Placer une \textbf{microélectrode intracellulaire} (pointe fine < \SI{1}{\micro\meter}, remplie de solution conductrice \ce{KCl} \SI{3}{M}) à l'intérieur de l'axoplasme, et une \textbf{électrode de référence extracellulaire} dans le bain. Relier les deux à un \textbf{voltmètre à haute impédance} (ou oscilloscope) capable de mesurer des différences de potentiel de l'ordre du millivolt sans perturber la membrane.
    \item \textbf{Observer le tracé :} À la perforation de la membrane, le voltmètre indique brutalement une tension négative (ex. \SI{-70}{mV}) par rapport au milieu extracellulaire pris comme référence (0 mV). C'est le \cle{potentiel de repos} ($V_{repos}$).
    \item \textbf{Interpréter ioniquement (raisonnement causal) :}
    \begin{itemize}
        \item \textbf{Gradient de concentration :} La pompe \ce{Na+/K+}-ATPase (3 \ce{Na+}$_{sortent}$ / 2 \ce{K+}$_{entrent}$ par ATP hydrolysé) maintient $[\ce{Na+}]_{in} \ll [\ce{Na+}]_{out}$ et $[\ce{K+}]_{in} \gg [\ce{K+}]_{out}$.
        \item \textbf{Perméabilité différentielle au repos :} La membrane est \textbf{beaucoup plus perméable aux \ce{K+} qu'aux \ce{Na+}} (canaux \ce{K+} « fuites » ouverts, canaux \ce{Na+} voltage-dépendants fermés).
        \item \textbf{Force motrice :} Les \ce{K+} sortent de la cellule suivant leur gradient de concentration (force chimique), emportant des charges positives $\rightarrow$ l'intérieur devient négatif. Cette négativité attire les \ce{K+} vers l'intérieur (force électrique).
        \item \textbf{Équilibre électrochimique :} Le potentiel de repos se stabilise proche du \cle{potentiel d'équilibre du potassium} ($E_K$), donné par l'équation de Nernst : $E_K = \frac{RT}{zF} \ln \frac{[\ce{K+}]_{out}}{[\ce{K+}]_{in}} \approx \SI{-90}{mV}$ (à \SI{37}{\celsius}). L'écart avec $V_{repos}$ (ex. \SI{-70}{mV}) s'explique par la faible perméabilité résiduelle aux \ce{Na+} (entrée de \ce{Na+} qui dépolarise légèrement).
    \end{itemize}
\end{enumerate}

\textbf{Formules à retenir :}
\begin{itemize}
    \item Équation de Nernst (ion $X$ de charge $z$) : $E_X = \frac{0.061}{z} \log_{10} \frac{[X]_{out}}{[X]_{in}}$ (en V, à \SI{37}{\celsius}).
    \item Équation de Goldman-Hodgkin-Katz (potentiel de membrane $V_m$) : $V_m = \frac{RT}{F} \ln \frac{P_K[\ce{K+}]_{out} + P_{Na}[\ce{Na+}]_{out} + P_{Cl}[\ce{Cl-}]_{in}}{P_K[\ce{K+}]_{in} + P_{Na}[\ce{Na+}]_{in} + P_{Cl}[\ce{Cl-}]_{out}}$.
\end{itemize}

\textbf{Réflexes Bac :} Préciser « intracellulaire négatif par rapport à l'extracellulaire ». Citer la pompe \ce{Na+/K+}-ATPase comme maintien des gradients \textbf{à long terme}, pas comme générateur direct du voltage instantané. Distinguer \textbf{gradient de concentration} (cause initiale) et \textbf{perméabilité} (cause proximale du signe du potentiel).
\end{methode}

\begin{exoresolu}[Analyse d'un montage de mesure du potentiel de repos]
\textbf{Énoncé :} On souhaite mettre en évidence le potentiel de repos d'un axone de grenouille ($d \approx \SI{500}{\micro\meter}$). Le milieu extracellulaire est une solution de Ringer. L'axoplasme contient \SI{120}{mmol/L} \ce{K+} et \SI{15}{mmol/L} \ce{Na+}. Le milieu extracellulaire contient \SI{2.5}{mmol/L} \ce{K+} et \SI{115}{mmol/L} \ce{Na+}.
\begin{enumerate}
    \item Schématiser le montage expérimental complet en légendant chaque élément.
    \item Calculer le potentiel d'équilibre du potassium ($E_K$) et du sodium ($E_{Na}$) à \SI{20}{\celsius} ($RT/F \approx \SI{25.3}{mV}$).
    \item Le potentiel de repos mesuré est de \SI{-88}{mV}. Comparer cette valeur à $E_K$ et $E_{Na}$. En déduire l'ion principalement responsable du potentiel de repos et justifier par la perméabilité membranaire.
\end{enumerate}

\tcblower
\textbf{Solution détaillée :}

\textbf{1. Schéma du montage (description textuelle pour le code LaTeX/TikZ) :}
\begin{popfigure}
\begin{tikzpicture}[scale=0.8, every node/.style={font=\small}]
% Bain extracellulaire
\draw[fill=popblue!10, draw=popblue, thick] (0,0) rectangle (8,5);
\node[align=center, text width=7cm] at (4, 4.3) {\textbf{Milieu extracellulaire (Ringer)} \\ $[\ce{K+}] = 2,5$ mmol/L ; $[\ce{Na+}] = 115$ mmol/L};

% Axone
\draw[fill=poporange!20, draw=poporange, thick] (2,1.5) rectangle (6,3.5);
\node[align=center] at (4, 2.5) {\textbf{Axone} \\ $[\ce{K+}] = 120$ mmol/L ; $[\ce{Na+}] = 15$ mmol/L};

% Microélectrode intracellulaire
\draw[thick, -{Stealth[length=3mm]}] (4, 3.5) -- (4, 5.5) node[above, align=center] {Microélectrode intracellulaire \\ (pointe < 1 µm, \ce{KCl} 3M)};
\draw[fill=white, draw=black] (3.8, 3.5) rectangle (4.2, 3.6); % pointe dans membrane

% Électrode de référence
\draw[thick, -{Stealth[length=3mm]}] (7, 2.5) -- (8.5, 2.5) node[right, align=center] {Électrode de référence \\ extracellulaire (Ag/AgCl)};

% Voltmetre / Oscilloscope
\draw[fill=popmuted!20, draw=popdark, thick, rounded corners] (4, 6) rectangle (6, 7.5);
\node at (5, 6.75) {\textbf{Voltmètre / Oscilloscope}};
\node at (5, 6.25) {Impédance d'entrée $> 10^{12}\ \Omega$};
\draw[thick] (4, 6) -- (4, 5.5); % connexion
\draw[thick] (6, 6.75) -- (7, 6.75) -- (7, 2.5); % connexion ref

% Affichage
\node[align=center, draw, fill=popcream, rounded corners] at (5, 8.2) {Affichage : \textbf{-88 mV} \\ (Intracellulaire négatif)};
\end{tikzpicture}
\legende{Montage classique à microélectrode pour l'enregistrement du potentiel de repos. L'électrode intracellulaire pénètre l'axoplasme ; l'électrode de référence est dans le bain. Le voltmètre à haute impédance mesure $V_{in} - V_{out}$.}
\end{popfigure}

\textbf{2. Calcul des potentiels d'équilibre (Équation de Nernst à \SI{20}{\celsius}) :}
Formule : $E_X = \frac{0.0253}{z} \ln \frac{[X]_{out}}{[X]_{in}}$ (en Volt) ou $E_X = \frac{58}{z} \log_{10} \frac{[X]_{out}}{[X]_{in}}$ (en mV, à \SI{20}{\celsius} approx). Utilisons la forme $\ln$ avec $RT/F = 25.3$ mV.
\begin{itemize}
    \item Pour \ce{K+} ($z=+1$) :
    \[ E_K = 25.3 \times \ln \left( \frac{2.5}{120} \right) = 25.3 \times \ln(0.02083) = 25.3 \times (-3.87) \approx \mathbf{-98.0\ mV} \]
    \item Pour \ce{Na+} ($z=+1$) :
    \[ E_{Na} = 25.3 \times \ln \left( \frac{115}{15} \right) = 25.3 \times \ln(7.67) = 25.3 \times 2.04 \approx \mathbf{+51.6\ mV} \]
\end{itemize}

\textbf{3. Interprétation :}
\begin{itemize}
    \item $V_{repos} = \SI{-88}{mV}$.
    \item $E_K = \SI{-98}{mV}$ ; $E_{Na} = \SI{+52}{mV}$.
    \item $V_{repos}$ est \textbf{très proche de $E_K$} (différence de \SI{10}{mV}) et \textbf{très loin de $E_{Na}$} (différence de \SI{140}{mV}).
    \item \textbf{Conclusion :} Au repos, la membrane est \textbf{essentiellement perméable aux ions \ce{K+}} (canaux « fuites » \ce{K+} ouverts). Les \ce{K+} sortent jusqu'à ce que la force électrique s'oppose exactement à la force chimique (équilibre électrochimique proche de $E_K$). La faible perméabilité aux \ce{Na+} ($P_{Na} \ll P_K$) explique que $V_{repos}$ soit légèrement moins négatif que $E_K$ (entrée de \ce{Na+} dépolarisante).
\end{itemize}

\textbf{Phrase de conclusion :} Le potentiel de repos de \SI{-88}{mV} résulte de la prédominance de la perméabilité membranaire aux ions potassium ($P_K \gg P_{Na}$), ramenant le potentiel membranaire vers le potentiel d'équilibre du potassium ($E_K \approx \SI{-98}{mV}$), maintenu par l'activité continue de la pompe \ce{Na+/K+}-ATPase.
\end{exoresolu}

\begin{methode}[Enregistrement et description du potentiel d'action]
\textbf{Objectif :} Schématiser le montage (identique au repos + stimulation) et décrire rigoureusement les phases du potentiel d'action (PA) sur l'oscillogramme.

\textbf{Étapes de la démarche :}
\begin{enumerate}
    \item \textbf{Monter l'expérience :} Reprendre le montage « potentiel de repos » (microélectrode intracellulaire + référence + oscilloscope). Ajouter un \textbf{stimulateur électrique} (générateur de courant) avec deux électrodes de stimulation placées en amont de la zone d'enregistrement (ou une électrode de stimulation traversante si fibre unique). Le stimulateur délivre une impulsion de courant carrée (dépolarisante = injection de charges positives à l'intérieur).
    \item \textbf{Déclencher et enregistrer :} Appliquer un stimulus \textbf{sous-seuil} (pas de PA), puis un stimulus \textbf{seuil} (premier PA), puis \textbf{supra-seuil} (PA de même amplitude). Observer l'invariance de l'amplitude (loi du tout ou rien).
    \item \textbf{Décrire la courbe (4 phases caractéristiques) :}
    \begin{enumerate}
        \item \textbf{Phase de dépolarisation (montée) :} Ouverture brutale des canaux \ce{Na+} voltage-dépendants $\rightarrow$ entrée massive de \ce{Na+} (gradient électrochimique favorable). $V_m$ passe de \SI{-70}{mV} vers \textbf{+30 mV} (dépassement du 0). \textit{Pente maximale = vitesse de dépolarisation (critère de conduction).}
        \item \textbf{Phase de repolarisation (descente) :} \textbf{Fermeture (inactivation) des canaux \ce{Na+}} + \textbf{Ouverture retardée des canaux \ce{K+} voltage-dépendants} $\rightarrow$ sortie massive de \ce{K+}. $V_m$ redescend vers valeurs négatives.
        \item \textbf{Phase d'hyperpolarisation (potentiel « après-potentiel » négatif) :} Les canaux \ce{K+} restent ouverts un peu après le retour au potentiel de repos (inactivation lente) $\rightarrow$ sortie d'excès de \ce{K+} $\rightarrow$ $V_m$ devient plus négatif que $V_{repos}$ (ex. \SI{-85}{mV}).
        \item \textbf{Retour au repos :} Fermeture des canaux \ce{K+} voltage-dépendants + activité de la pompe \ce{Na+/K+}-ATPase (restauration gradients à long terme).
    \end{enumerate}
    \item \textbf{Notions clés à annoter sur le schéma :} Amplitude ($\approx \SI{100}{mV}$), Seuil (\SI{-55}{mV} env.), Durée (\SI{1}{ms} à \SI{2}{ms} pour fibre myélinisée), Période réfractaire absolue (pendant PA) et relative (pendant hyperpolarisation).
\end{enumerate}

\textbf{Réflexes Bac :} Ne pas confondre « seuil d'excitation » (niveau de tension \SI{-55}{mV}) et « intensité seuil » (courant minimal à injecter). La loi du tout ou rien : \textbf{amplitude et forme du PA sont indépendantes de l'intensité du stimulus} (pour $I \ge I_{seuil}$). Citer les canaux \textbf{voltage-dépendants} (pas « dépendants du ligand » ici).
\end{methode}

\begin{exoresolu}[Analyse d'un oscillogramme de potentiel d'action]
\textbf{Énoncé :} L'oscillogramme ci-dessous (Figure 1) montre l'enregistrement intracellulaire du potentiel de membrane d'un axone myélinisé de mammifère stimulé par une impulsion de courant carrée de \SI{0.5}{ms}.

\begin{popfigure}
\begin{tikzpicture}[scale=0.7]
\begin{axis}[
    width=\linewidth, height=8cm,
    axis lines=middle,
    xlabel={Temps (ms)}, ylabel={Potentiel de membrane (mV)},
    xmin=0, xmax=10, ymin=-90, ymax=50,
    xtick={0,1,2,3,4,5,6,7,8,9,10}, ytick={-70,-55,0,30},
    yticklabels={-70,-55,0,+30},
    tick label style={font=\small},
    xlabel style={anchor=north}, ylabel style={anchor=south},
    domain=0:10, samples=500,
    clip=false
]
% Potentiel de repos
\addplot[popmuted, thick, domain=0:1] {-70};
% Action Potential - simplified shape using coordinates
\addplot[popcurve, very thick, smooth, tension=0.7] coordinates {
    (0, -70) (1, -70) 
    (1.2, -60) (1.4, -20) (1.5, 10) (1.55, 30) % Depol
    (1.6, 25) (1.7, 0) (1.9, -50) (2.1, -75) (2.3, -85) % Repol + Hyperpol
    (2.8, -80) (3.5, -75) (5, -72) (10, -70) % Return
};
% Threshold line
\addplot[dashed, poporange, thick] coordinates {(0, -55) (10, -55)} node[pos=0.05, left, poporange] {Seuil (\SI{-55}{mV})};
% Rest line
\addplot[dashed, popblue, thick] coordinates {(0, -70) (10, -70)} node[pos=0.05, left, popblue] {Repos (\SI{-70}{mV})};
% Stimulus marker
\draw[popgreen, thick, -{Stealth[length=2mm]}] (1, -95) -- (1.5, -95) node[midway, below] {Stimulus \SI{0.5}{ms}};
\end{axis}
\end{tikzpicture}
\legende{Oscillogramme d'un potentiel d'action enregistré en intracellulaire.}
\end{popfigure}

\begin{enumerate}
    \item Identifier sur la courbe les 4 phases du potentiel d'action et donner leurs valeurs caractéristiques (potentiels, durée).
    \item Expliquer les mécanismes ioniques (mouvements d'ions, canaux) responsables de la phase de dépolarisation et de la phase de repolarisation.
    \item L'expérience est répétée avec un stimulus d'intensité double. Prédire la modification de la courbe et justifier (Loi du tout ou rien).
    \item Calculer la vitesse maximale de dépolarisation ($dV/dt_{max}$) si le potentiel passe de \SI{-55}{mV} à \SI{+30}{mV} en \SI{0.2}{ms}.
\end{enumerate}

\tcblower
\textbf{Solution détaillée :}

\textbf{1. Identification des phases (lecture graphique) :}
\begin{itemize}
    \item \textbf{Phase 1 : Dépolarisation (montée).} Début : \SI{-70}{mV} (repos). Seuil franchi vers \SI{-55}{mV}. Pic (overshoot) : \textbf{+30 mV}. Durée montée : \textbf{$\approx \SI{0.5}{ms}$} (de \SI{-55}{mV} à \SI{+30}{mV}).
    \item \textbf{Phase 2 : Repolarisation (descente).} De \SI{+30}{mV} jusqu'au franchissement du potentiel de repos (\SI{-70}{mV}). Durée : \textbf{$\approx \SI{0.6}{ms}$}.
    \item \textbf{Phase 3 : Hyperpolarisation (après-potentiel négatif).} Minimum : \textbf{\SI{-85}{mV}}. Durée : \textbf{$\approx \SI{1.5}{ms}$} (retour au repos).
    \item \textbf{Phase 4 : Retour au repos.} Stabilisation à \textbf{\SI{-70}{mV}}.
    \item \textbf{Durée totale du PA (à mi-hauteur) :} $\approx \SI{1}{ms}$.
\end{itemize}

\textbf{2. Mécanismes ioniques :}
\begin{itemize}
    \item \textbf{Dépolarisation :} L'atteinte du seuil (\SI{-55}{mV}) provoque l'\textbf{ouverture brutale des canaux \ce{Na+} voltage-dépendants (canaux \ce{Na_v})}. La force électrochimique (gradient de concentration \textbf{ET} gradient électrique) pousse les \ce{Na+} vers l'intérieur. L'entrée massive de charges positives ($I_{Na}$ entrant) dépolarise la membrane. C'est un \textbf{emballement positif} (boucle de rétroaction positive : dépolarisation $\rightarrow$ ouverture \ce{Na+} $\rightarrow$ entrée \ce{Na+} $\rightarrow$ dépolarisation).
    \item \textbf{Repolarisation :} Deux événements simultanés :
    \begin{enumerate}
        \item \textbf{Inactivation des canaux \ce{Na+}} (boucle « ball-and-chain » se referme, \textbf{réfractaire absolu}) : arrêt de l'entrée de \ce{Na+}.
        \item \textbf{Ouverture retardée des canaux \ce{K+} voltage-dépendants (canaux \ce{K_v})} : sortie massive de \ce{K+} ($I_K$ sortant) suivant son gradient électrochimique. Cette sortie de charges positives ramène $V_m$ vers $E_K$.
    \end{enumerate}
\end{itemize}

\textbf{3. Prédiction stimulus double intensité (Loi du tout ou rien) :}
\begin{itemize}
    \item \textbf{Prédiction :} La courbe du potentiel d'action sera \textbf{strictement identique} (même amplitude, même forme, même durée, même seuil).
    \item \textbf{Justification :} Le potentiel d'action est une réponse \textbf{réactive} de la membrane (ouverture canaux voltage-dépendants). Une fois le seuil de tension atteint, l'emballement \ce{Na+} est auto-déclenché et saturé (tous les canaux \ce{Na+} disponibles s'ouvrent). L'intensité du stimulus ne module que la \textbf{latence} (délai avant le PA) et la \textbf{fréquence} des PA si le stimulus est maintenu, jamais l'amplitude d'un PA unique.
\end{itemize}

\textbf{4. Calcul de la vitesse maximale de dépolarisation :}
\[ \left( \frac{dV}{dt} \right)_{max} \approx \frac{\Delta V}{\Delta t} = \frac{V_{pic} - V_{seuil}}{\Delta t_{montée}} = \frac{30 - (-55)}{0.2} = \frac{85}{0.2} = \mathbf{425\ V/s} \quad (\text{ou } \mathbf{425\ mV/ms}) \]
\textit{Note :} Cette valeur élevée (\SI{400}{V/s} typique fibre myélinisée) témoigne de la densité élevée en canaux \ce{Na+} aux nœuds de Ranvier.

\textbf{Phrase de conclusion :} Le potentiel d'action est une onde de dépolarisation-répolarisation stéréotypée (\SI{100}{mV} d'amplitude, \SI{1}{ms} de durée) générée par l'ouverture séquentielle et voltage-dépendante des canaux \ce{Na+} (dépolarisation) puis \ce{K+} (repolarisation), obéissant à la loi du tout ou rien.
\end{exoresolu}

\begin{methode}[Interprétation des courbes de perméabilité membranaire ($P_{Na}$, $P_K$) lors d'un PA]
\textbf{Objectif :} Relier les variations de perméabilité ($P_{Na}$, $P_K$) aux phases du potentiel d'action et aux courants ioniques.

\textbf{Étapes de la démarche :}
\begin{enumerate}
    \item \textbf{Identifier les axes :} Abscisse = Temps (ms). Ordonnée = Perméabilité relative ($P_{Na}$, $P_K$) ou Conductance ($g_{Na}$, $g_K$). Souvent normalisées au max ($P/P_{max}$).
    \item \textbf{Analyser la courbe $P_{Na}$ (ou $g_{Na}$) :}
    \begin{itemize}
        \item \textbf{Au repos :} $P_{Na} \approx 0$ (canaux fermés).
        \item \textbf{Au seuil :} Ouverture brutale $\rightarrow$ pic de $P_{Na}$ \textbf{très précoce} (coïncide avec le début de la dépolarisation).
        \item \textbf{Pendant la dépolarisation :} $P_{Na}$ reste maximal.
        \item \textbf{Au pic du PA / début repolarisation :} \textbf{Chute brutale de $P_{Na}$} (inactivation des canaux \ce{Na+}). $P_{Na}$ revient à 0 avant la fin de la repolarisation.
    \end{itemize}
    \item \textbf{Analyser la courbe $P_K$ (ou $g_K$) :}
    \begin{itemize}
        \item \textbf{Au repos :} $P_K = P_{K,repos}$ (canaux fuites, non voltage-dépendants).
        \item \textbf{Au seuil :} Pas de changement immédiat.
        \item \textbf{Pendant la dépolarisation / pic :} \textbf{Augmentation progressive et retardée} de $P_K$ (activation lente des canaux \ce{K_v}).
        \item \textbf{Pendant la repolarisation / hyperpolarisation :} $P_K$ \textbf{atteint un maximum} (souvent $> P_{Na,max}$) puis \textbf{décroît lentement} (inactivation lente / déactivation).
        \item \textbf{Retour au repos :} $P_K$ revient à sa valeur basale ($P_{K,repos}$).
    \end{itemize}
    \item \textbf{Synthétiser (Tableau de correspondance) :}
    \begin{center}
    \begin{tabularx}{\linewidth}{|l|>{\centering\arraybackslash}X|>{\centering\arraybackslash}X|}\hline
    \rowcolor{popblueL} \textbf{Phase PA} & \textbf{$P_{Na}$ (Canaux \ce{Na_v})} & \textbf{$P_K$ (Canaux \ce{K_v} + fuites)} \\ \hline
    Repos & $\approx 0$ & Basale ($P_{K,leak}$) \\ \hline
    Dépolarisation & \textbf{Pic maximal brutal} (Ouverture $m^3$) & Croissance lente (Ouverture $n^4$) \\ \hline
    Pic / Repolarisation & \textbf{Chute à 0} (Inactivation $h \to 0$) & \textbf{Pic maximal} \\ \hline
    Hyperpolarisation & 0 & Décroissance lente (Déactivation) \\ \hline
    Retour repos & 0 & Basale \\ \hline
    \end{tabularx}
    \end{center}
\end{enumerate}

\textbf{Formules clés (Modèle Hodgkin-Huxley simplifié) :}
$I_{Na} = g_{Na} (V_m - E_{Na})$ avec $g_{Na} = \bar{g}_{Na} m^3 h$.
$I_K = g_K (V_m - E_K)$ avec $g_K = \bar{g}_K n^4$.
$m$ (activation \ce{Na+}) : rapide, suit $V_m$. $h$ (inactivation \ce{Na+}) : lent, suit $V_m$ inversément. $n$ (activation \ce{K+}) : lent, suit $V_m$.

\textbf{Piège fréquent :} Confondre « Perméabilité au repos » (fuites \ce{K+}) et « Perméabilité voltage-dépendante ». Ne pas oublier que $P_K$ au repos n'est pas nul !
\end{methode}

\begin{exoresolu}[Lien perméabilité - potentiel d'action (Document Hodgkin-Huxley)]
\textbf{Énoncé :} Le document ci-dessous (Figure 2) représente l'évolution dans le temps de la perméabilité membranaire au sodium ($P_{Na}$) et au potassium ($P_K$) lors d'un potentiel d'action enregistré sur un axone géant de calmar (température \SI{6.3}{\celsius}).

\begin{popfigure}
\begin{tikzpicture}[scale=0.7]
\begin{axis}[
    width=\linewidth, height=7cm,
    axis lines=middle,
    xlabel={Temps (ms)}, ylabel={Perméabilité relative ($P/P_{max}$)},
    xmin=0, xmax=12, ymin=-0.1, ymax=1.1,
    xtick={0,1,2,3,4,5,6,7,8,9,10}, ytick={0,0.5,1},
    tick label style={font=\small},
    legend pos=north east, legend style={font=\small, fill=white, draw=popline},
    domain=0:12, samples=500
]
% P_Na : fast rise, fast fall
\addplot[poporange, very thick, smooth, tension=0.5, name path=PNa] coordinates {
    (0,0) (0.5,0) (0.6,0.1) (0.7,0.5) (0.8,0.9) (0.9,1.0) (1.0,1.0) (1.1,0.8) (1.2,0.4) (1.3,0.1) (1.5,0) (12,0)
};
% P_K : slow rise, slow fall
\addplot[popblue, very thick, smooth, tension=0.5, name path=PK] coordinates {
    (0,0.05) (0.5,0.05) (0.7,0.1) (1.0,0.3) (1.5,0.7) (2.0,1.0) (3.0,1.0) (4.0,0.8) (6.0,0.4) (8.0,0.15) (12,0.05)
};
% Phase markers
\draw[dashed, popmuted] (0.5, 0) -- (0.5, 1.1) node[above, font=\tiny, rotate=90, anchor=south] {Seuil};
\draw[dashed, popmuted] (1.0, 0) -- (1.0, 1.1) node[above, font=\tiny, rotate=90, anchor=south] {Pic PA};
\draw[dashed, popmuted] (2.0, 0) -- (2.0, 1.1) node[above, font=\tiny, rotate=90, anchor=south] {Fin Repolarisation};
% Labels
\node[poporange, font=\small] at (1.5, 1.1) {$P_{Na}$};
\node[popblue, font=\small] at (3.5, 1.1) {$P_K$};
\end{axis}
\end{tikzpicture}
\legende{Variations de la perméabilité membranaire relative au Na$^+$ (orange) et au K$^+$ (bleu) au cours d'un potentiel d'action (d'après Hodgkin \& Huxley, 1952).}
\end{popfigure}

\begin{enumerate}
    \item Relever les valeurs de $P_{Na}$ et $P_K$ au repos ($t=0$). Commenter.
    \item Déterminer l'instant (en ms) où $P_{Na}$ est maximale. Comparer cet instant au pic du potentiel d'action (supposé à $t=1.0$ ms). Expliquer pourquoi $P_{Na}$ chute \textbf{avant} la fin de la repolarisation.
    \item Déterminer l'instant où $P_K$ est maximale. Comparer à l'instant du pic de $P_{Na}$. Justifier ce décalage par la cinétique des canaux.
    \item À $t = 3$ ms, $P_K$ est encore élevée ($> 0.5$) alors que $V_m$ est revenu au repos (supposé). Expliquer l'hyperpolarisation (après-potentiel négatif) par ce décalage.
    \item \textbf{Calcul de courant :} Sachant qu'au pic du PA ($t=1.0$ ms), $P_{Na} = P_{Na,max}$ et $P_K = 0.3 P_{K,max}$. Si $\bar{g}_{Na} = \SI{120}{mS/cm^2}$, $\bar{g}_K = \SI{36}{mS/cm^2}$, $E_{Na} = \SI{+50}{mV}$, $E_K = \SI{-77}{mV}$, $V_m = \SI{+30}{mV}$. Calculer le courant sodium $I_{Na}$ et le courant potassium $I_K$ (densité de courant en $\mu A/cm^2$) à cet instant. Conclure sur le sens du courant membranaire total.
\end{enumerate}

\tcblower
\textbf{Solution détaillée :}

\textbf{1. Perméabilités au repos ($t=0$) :}
Lecture graphique : $P_{Na}(0) \approx \mathbf{0}$. $P_K(0) \approx \mathbf{0.05}$ (5\% de $P_{K,max}$).
\textbf{Commentaire :} Au repos, les canaux \ce{Na+} voltage-dépendants sont \textbf{fermés} ($P_{Na} \approx 0$). La perméabilité au \ce{K+} n'est pas nulle car elle est assurée par les \textbf{canaux \ce{K+} « fuites » (leak)} non voltage-dépendants, responsables du potentiel de repos proche de $E_K$.

\textbf{2. Pic de $P_{Na}$ et inactivation :}
Lecture : $P_{Na}$ atteint son maximum (1.0) vers \textbf{$t = 0.9 - 1.0$ ms}. Cela coïncide avec le \textbf{pic du potentiel d'action} ($V_m \approx \SI{+30}{mV}$).
\textbf{Explication de la chute précoce :} $P_{Na}$ chute brutalement dès $t \approx 1.0$ ms (inactivation variable $h \to 0$). Les canaux \ce{Na+} entrent en \textbf{état inactivé (réfractaire absolu)}. Ils ne peuvent plus s'ouvrir tant que la membrane n'est pas repolarisée. Cela \textbf{stoppe l'entrée de \ce{Na+}} et permet à la repolarisation (sortie \ce{K+}) de s'opérer sans opposition. Si $P_{Na}$ restait haute, la membrane resterait dépolarisée (blocage).

\textbf{3. Pic de $P_K$ et décalage :}
Lecture : $P_K$ atteint son maximum (1.0) vers \textbf{$t = 2.0 - 3.0$ ms}.
\textbf{Comparaison :} Le pic de $P_K$ est \textbf{retardé d'environ 1 à 2 ms} par rapport au pic de $P_{Na}$.
\textbf{Justification cinétique :} L'activation des canaux \ce{K+} (variable $n^4$) est \textbf{plus lente} (constante de temps $\tau_n > \tau_m$) que l'activation des canaux \ce{Na+} (variable $m^3$). Ce décalage est \textbf{essentiel} : il assure que la sortie de \ce{K+} (repolarisation) survient \textbf{après} l'entrée de \ce{Na+} (dépolarisation), permettant l'existence d'un pic d'amplitude suffisante.

\textbf{4. Hyperpolarisation ($t=3$ ms) :}
À $t=3$ ms, $V_m$ a traversé le potentiel de repos (\SI{-70}{mV}) mais $P_K$ est encore $\approx 0.8-1.0$ (bien supérieure à la perméabilité au repos $0.05$).
La force motrice pour le \ce{K+} est $(V_m - E_K)$. Comme $V_m$ (ex. \SI{-70}{mV}) $> E_K$ (\SI{-77}{mV}), la force motrice est \textbf{positive (sortante)}.
Comme $P_K$ est encore très grande, le courant sortant de \ce{K+} ($I_K = g_K(V_m-E_K)$) reste \textbf{important}, tirant $V_m$ \textbf{au-delà de $E_K$ (vers \SI{-85}{mV})} : c'est l'hyperpolarisation. Ce n'est que lorsque $P_K$ redescendra vers sa valeur basale que $V_m$ reviendra au repos.

\textbf{5. Calcul des courants au pic ($t=1.0$ ms, $V_m = \SI{+30}{mV}$) :}
Données : $\bar{g}_{Na} = 120$ mS/cm², $\bar{g}_K = 36$ mS/cm². $P_{Na}=1 \Rightarrow g_{Na} = 120$ mS/cm². $P_K=0.3 \Rightarrow g_K = 0.3 \times 36 = 10.8$ mS/cm².
Formule : $I = g (V_m - E)$ (en $\mu A/cm^2$ si $g$ en mS/cm² et $V$ en mV).
\begin{itemize}
    \item $I_{Na} = 120 \times (30 - 50) = 120 \times (-20) = \mathbf{-2400\ \mu A/cm^2}$.
    \item $I_K = 10.8 \times (30 - (-77)) = 10.8 \times 107 = \mathbf{+1155.6\ \mu A/cm^2}$.
\end{itemize}
\textbf{Interprétation :} $I_{Na}$ est \textbf{négatif} (courant \textbf{entrant}, dépolarisant, charges + entrent). $I_K$ est \textbf{positif} (courant \textbf{sortant}, repolarisant, charges + sortent).
Courant total $I_m = I_{Na} + I_K = -2400 + 1155.6 = \mathbf{-1244.4\ \mu A/cm^2}$.
Le courant total est \textbf{entrant (négatif)}. Cela signifie qu'au pic du PA, la membrane se charge encore légèrement en charges positives (capacité membranaire), mais l'inactivation imminente de $g_{Na}$ va inverser le signe du courant total pour permettre la repolarisation.

\textbf{Phrase de conclusion :} La dynamique temporelle décalée de $P_{Na}$ (rapide, transitoire) et $P_K$ (lente, soutenue) orchestrée par la voltage-dépendance de leurs canaux respectifs, est la cause biophysique directe de la forme stéréotypée du potentiel d'action.
\end{exoresolu}

\begin{methode}[Codage du message nerveux au niveau du récepteur sensoriel]
\textbf{Objectif :} Interpréter les tracés (potentiel de récepteur $\rightarrow$ potentiel d'action) et expliquer le codage de l'intensité du stimulus (fréquence des PA).

\textbf{Étapes de la démarche :}
\begin{enumerate}
    \item \textbf{Identifier les deux zones d'enregistrement :}
    \begin{itemize}
        \item \textbf{Zone réceptrice (dendrites/terminaisons) :} Enregistrement du \textbf{Potentiel de Récepteur (PR)} ou Potentiel Générateur. C'est un \textbf{potentiel gradué} (amplitude $\propto$ intensité stimulus), \textbf{local}, \textbf{non propagé activement}, \textbf{sans seuil} (pas de tout-ou-rien), \textbf{sans période réfractaire}. Polarité : dépolarisation (généralement).
        \item \textbf{Zone de déclenchement (cône d'émergence / premier nœud de Ranvier) :} Enregistrement des \textbf{Potentiels d'Action (PA)}. \textbf{Tout-ou-rien}, \textbf{propagés}, \textbf{seuil}, \textbf{période réfractaire}.
    \end{itemize}
    \item \textbf{Analyser la relation Stimulus $\rightarrow$ PR :}
    \begin{itemize}
        \item Stimulus faible $\rightarrow$ PR faible (amplitude faible).
        \item Stimulus fort $\rightarrow$ PR fort (amplitude grande, saturation possible).
        \item \textbf{Mécanisme :} Canaux \textbf{méchanosensibles / chimiosensibles / thermosensibles} (canaux dépendants du stimulus, PAS voltage-dépendants). Ouverture $\rightarrow$ entrée cations (\ce{Na+}, \ce{Ca^{2+}}) $\rightarrow$ dépolarisation locale (courant récepteur).
    \end{itemize}
    \item \textbf{Analyser la relation PR $\rightarrow$ PA (Codage) :}
    \begin{itemize}
        \item Le PR se propage \textbf{électrotoniquement} (atténuation exponentielle) jusqu'au cône d'émergence (zone à forte densité canaux \ce{Na_v}, seuil bas).
        \item Si $PR_{cône} \ge \text{Seuil}$ $\rightarrow$ Déclenchement d'un PA.
        \item \textbf{Codage de l'intensité :} \textbf{Fréquence des PA} ($f$) codée par l'amplitude du PR. $f \propto$ Amplitude PR $\propto$ Intensité Stimulus (relation souvent logarithmique ou sigmoïde : loi de Weber-Fechner / Stevens).
        \item \textbf{Codage de la durée :} Durée de la salve de PA = Durée du stimulus (si stimulus maintenu).
        \item \textbf{Adaptation :} Diminution de la fréquence des PA malgré stimulus constant (inactivation canaux récepteurs, accommodation canaux \ce{Na+} au cône). Récepteurs \textbf{toniques} (peu d'adaptation, ex. propriocepteurs) vs \textbf{phasiques} (forte adaptation, ex. corpuscules de Pacini).
    \end{itemize}
\end{enumerate}

\textbf{Schéma type à savoir dessiner :} Tracé double : Haut = PR (gradué, lent). Bas = PA (tout-ou-rien, fréquent si PR fort). Flèche « Propagation électrotonique » entre les deux sites.

\textbf{Réflexes Bac :} Ne jamais dire « le potentiel de récepteur se propage ». Dire : « Le potentiel de récepteur se propage \textbf{électrotoniquement} (passivement, avec atténuation) jusqu'à la zone de déclenchement ». Distinguer \textbf{seuil de tension} (propriété membrane, \SI{-55}{mV}) et \textbf{seuil d'intensité du stimulus} (intensité minimale pour que PR atteigne le seuil de tension).
\end{methode}

\begin{exoresolu}[Du stimulus à la fréquence : analyse d'un tracé récepteur]
\textbf{Énoncé :} On enregistre l'activité électrique d'un neurone sensoriel (récepteur au toucher type corpuscule de Merkel - récepteur tonique) soumis à des pressions mécaniques croissantes. La Figure 3 montre deux enregistrements simultanés : en haut, le potentiel de récepteur (PR) au niveau de la terminaison sensorielle ; en bas, les potentiels d'action (PA) au niveau du premier nœud de Ranvier.

\begin{popfigure}
\begin{tikzpicture}[scale=0.75]
% Two stacked axes
\begin{axis}[
    width=\linewidth, height=5cm,
    axis lines=left,
    xlabel={Temps (ms)}, ylabel={PR (mV)},
    xmin=0, xmax=50, ymin=-80, ymax=20,
    xtick={0,10,20,30,40,50}, ytick={-70,-55,-40,-20,0,10},
    tick label style={font=\small},
    domain=0:50, samples=500,
    clip=false,
    y axis line style={-},
    x axis line style={-},
    axis x line*=bottom,
    axis y line*=left,
    name=topAxis
]
% Stimulus bar
\draw[popgreen, thick] (5, 15) -- (45, 15) node[midway, above] {Pression mécanique maintenue};
\draw[popgreen, thick, -{Stealth[length=2mm]}] (5, 13) -- (5, 15);
\draw[popgreen, thick, -{Stealth[length=2mm]}] (45, 13) -- (45, 15);

% Stim 1 (Weak) PR: Plateau -40mV
\addplot[popblue, very thick] coordinates {(0,-70) (5,-70) (6,-40) (45,-40) (46,-70) (50,-70)};
% Stim 2 (Medium) PR: Plateau -20mV (shifted up for visibility? No, same axis, different color/style)
% We'll plot only one representative trace for clarity, or use the data textually.
% Since the exercise gives textual data, we just show a schematic PR trace.
\node[popblue, anchor=south west] at (5, -35) {Stimulus 2 : Plateau \SI{-20}{mV}};
\node[poporange, anchor=south west] at (5, -15) {Stimulus 3 : Plateau \SI{+10}{mV}};
\end{axis}

\begin{axis}[
    width=\linewidth, height=4cm,
    at={(topAxis.below south west)}, anchor=above north west,
    xshift=0pt, yshift=-1.5cm,
    axis lines=left,
    xlabel={Temps (ms)}, ylabel={PA (mV)},
    xmin=0, xmax=50, ymin=-80, ymax=50,
    xtick={0,10,20,30,40,50}, ytick={-70,-55,0,30},
    yticklabels={-70,-55,0,+30},
    tick label style={font=\small},
    domain=0:50, samples=500,
    clip=false,
    y axis line style={-},
    x axis line style={-},
]
% Stim 1: 1 AP at start
\addplot[popblue, very thick] coordinates {
    (0,-70) (5,-70) 
    (6,-70) (6.1,30) (6.3,-70) % 1 AP
    (6.3,-70) (50,-70)
};
% Stim 2: Train at ~50Hz (T=20ms) starting after latency
% Latency ~5ms. APs at 11, 31, 51... but max 50ms. So 11, 31.
\addplot[poporange, very thick] coordinates {
    (0,-70) (10,-70)
    (11,-70) (11.1,30) (11.3,-70)
    (31,-70) (31.1,30) (31.3,-70)
    (31.3,-70) (50,-70)
};
% Stim 3: Train at ~150Hz (T=6.7ms)
\addplot[popgreen, very thick] coordinates {
    (0,-70) (15,-70)
    (16,-70) (16.1,30) (16.3,-70)
    (23,-70) (23.1,30) (23.3,-70)
    (30,-70) (30.1,30) (30.3,-70)
    (37,-70) (37.1,30) (37.3,-70)
    (44,-70) (44.1,30) (44.3,-70)
    (44.3,-70) (50,-70)
};
% Labels for traces
\node[popblue, font=\small, anchor=west] at (0, -70) {Stim 1};
\node[poporange, font=\small, anchor=west] at (0, -75) {Stim 2 (50 Hz)};
\node[popgreen, font=\small, anchor=west] at (0, -80) {Stim 3 (150 Hz)};
\end{axis}
\end{tikzpicture}
\legende{Enregistrement simultané schématique : Potentiel de Récepteur (haut, gradué) et Potentiels d'Action (bas, tout-ou-rien) pour trois intensités de stimulus.}
\end{popfigure}

\textit{Données chiffrées extraites de la figure (simulées) :}
\begin{itemize}
    \item \textbf{Stimulus 1 (Faible) :} PR plateau à \SI{-40}{mV} (dépolarisation de \SI{30}{mV} par rapport au repos \SI{-70}{mV}). $\rightarrow$ \textbf{1 seul PA} au début du stimulus, puis silence.
    \item \textbf{Stimulus 2 (Moyen) :} PR plateau à \SI{-20}{mV} (dépolarisation \SI{50}{mV}). $\rightarrow$ \textbf{Salve de PA} : Fréquence initiale \SI{200}{Hz}, se stabilisant à \SI{50}{Hz} (adaptation lente).
    \item \textbf{Stimulus 3 (Fort) :} PR plateau à \SI{+10}{mV} (dépolarisation \SI{80}{mV}, saturation). $\rightarrow$ \textbf{Salve de PA} : Fréquence initiale \SI{500}{Hz}, se stabilisant à \SI{150}{Hz}.
\end{itemize}

\begin{enumerate}
    \item Qualifier la nature du Potentiel de Récepteur (gradué/propagé, seuil, tout-ou-rien). Justifier par la courbe « Stimulus 1 ».
    \item Expliquer pourquoi le Stimulus 1 ne déclenche qu'un seul PA alors que le stimulus est maintenu.
    \item Calculer la période des PA pour le Stimulus 2 en régime stable (\SI{50}{Hz}). Exprimer en ms.
    \item Décrire l'évolution de la fréquence des PA en fonction de l'intensité du stimulus. Donner la relation mathématique type (loi de Weber-Fechner ou Stevens).
    \item Expliquer le mécanisme biophysique de l'adaptation de la fréquence observée pour le Stimulus 2 (baisse de 200 à 50 Hz).
\end{enumerate}

\tcblower
\textbf{Solution détaillée :}

\textbf{1. Nature du Potentiel de Récepteur (PR) :}
\begin{itemize}
    \item \textbf{Gradué :} Son amplitude varie avec l'intensité du stimulus (\SI{-40}{mV} pour stimulus faible, \SI{+10}{mV} pour stimulus fort). Pas de « tout-ou-rien ».
    \item \textbf{Local / Non propagé activement :} Il ne se propage pas sous forme d'onde régénérative. Il se propage \textbf{électrotoniquement} (passivement, avec décroissance exponentielle de l'amplitude) le long de la fibre jusqu'au cône d'émergence.
    \item \textbf{Sans seuil :} Il apparaît dès le stimulus le plus faible (pas de « seuil d'intensité » pour le PR lui-même, seulement pour le PA).
    \item \textbf{Sans période réfractaire :} Il suit fidèlement l'enveloppe du stimulus (plateau maintenu tant que le stimulus est maintenu).
\end{itemize}

\textbf{2. Un seul PA pour Stimulus 1 (Faible) :}
Le PR atteint un plateau à \SI{-40}{mV}. La dépolarisation électrotonique arrivant au cône d'émergence (seuil \SI{-55}{mV}) est suffisante pour \textbf{dépasser le seuil une fois} au tout début (pic initial du PR souvent plus haut que le plateau). Ensuite, le niveau de plateau (\SI{-40}{mV} au niveau récepteur) s'atténue électrotoniquement pour devenir \textbf{inférieur au seuil} (\SI{-55}{mV}) au niveau du cône. Aucun nouveau PA n'est déclenché. C'est le \textbf{seuil d'intensité du stimulus} : l'intensité 1 est proche du seuil minimal pour déclencher *un* PA.

\textbf{3. Calcul de la période (Stimulus 2, régime stable) :}
Fréquence $f = \SI{50}{Hz}$.
Période $T = \frac{1}{f} = \frac{1}{50} = \mathbf{0.02\ s = 20\ ms}$.
\textit{Vérification cohérence :} Période réfractaire absolue $\approx \SI{1}{ms}$, relative $\approx \SI{3-5}{ms}$. $T = \SI{20}{ms} \gg$ période réfractaire $\rightarrow$ physiologiquement possible.

\textbf{4. Relation Fréquence - Intensité Stimulus :}
\begin{itemize}
    \item Stimulus Faible $\rightarrow$ 0 Hz (régime stable) / 1 PA (transitoire).
    \item Stimulus Moyen $\rightarrow$ \SI{50}{Hz} (stable).
    \item Stimulus Fort $\rightarrow$ \SI{150}{Hz} (stable).
    \item \textbf{Conclusion :} La \textbf{fréquence des PA augmente avec l'intensité du stimulus}. C'est le \textbf{codage fréquentiel} (loi du tout-ou-rien pour l'amplitude du PA, codage de l'intensité par la fréquence).
    \item \textbf{Relation mathématique :} Souvent décrite par la \textbf{loi de Stevens (loi de puissance)} : $f = k (I - I_0)^n$ ou \textbf{loi de Weber-Fechner (logarithmique)} : $f = k \log(I/I_0)$. Au Bac, on attend : « La fréquence est une fonction croissante de l'intensité du stimulus, souvent logarithmique ou sigmoïde ».
\end{itemize}

\textbf{5. Mécanisme de l'adaptation (baisse 200 $\to$ 50 Hz) :}
Deux mécanismes concomitants :
\begin{enumerate}
    \item \textbf{Adaptation du récepteur (niveau terminaison) :} Les canaux mécanosensibles (ex. canaux \ce{PIEZO2}) s'\textbf{inactivent} partiellement malgré le maintien de la pression. L'amplitude du PR \textbf{diminue} au cours du temps (le plateau \SI{-20}{mV} est plus bas que le pic initial). Moins de dépolarisation arrivant au cône $\rightarrow$ fréquence plus basse.
    \item \textbf{Accommodation du générateur de PA (niveau cône) :} La dépolarisation maintenue (même si PR stable) provoque une \textbf{inactivation partielle des canaux \ce{Na+} voltage-dépendants} au cône (variable $h$ diminue). Le \textbf{seuil effectif augmente} (il faut plus de dépolarisation pour ouvrir les canaux \ce{Na+} restants disponibles). Cela allonge l'intervalle inter-spikes $\rightarrow$ baisse de fréquence.
\end{enumerate}
\textit{Note :} Pour un récepteur \textbf{tonique} (Merkel), l'adaptation est \textbf{lente} (fréquence se stabilise $> 0$). Pour un récepteur \textbf{phasique} (Pacini), l'adaptation est \textbf{rapide} (fréquence tombe à 0 en quelques ms, seul le début et la fin du stimulus codés).

\textbf{Phrase de conclusion :} Le récepteur sensoriel transcode l'intensité physique du stimulus (pression) en fréquence de potentiels d'action via un potentiel de récepteur gradué ; l'adaptation fréquentielle résulte de l'inactivation des canaux sensoriels et de l'accommodation des canaux voltage-dépendants du déclencheur.
\end{exoresolu}

\begin{methode}[Seuil d'excitation et relation Intensité-Durée (Force-Durée)]
\textbf{Objectif :} Interpréter la courbe de réponse d'un nerf/fibre à des stimuli d'intensité croissante (détermination du seuil) et la courbe Intensité-Durée (chronaxie, rhéobase).

\textbf{Étapes de la démarche :}
\begin{enumerate}
    \item \textbf{Courbe Réponse vs Intensité Stimulus (à durée fixe) :}
    \begin{itemize}
        \item Abscisse : Intensité du stimulus $I$ (\SI{}{\micro A} ou \SI{}{mA}). Ordonnée : Réponse (Amplitude PA ou Nombre de fibres excitées / \% fibres).
        \item \textbf{Fibre unique :} Tout-ou-rien. Seuil net $I_{seuil}$. En dessous : 0 PA. Au-dessus : 1 PA (amplitude constante).
        \item \textbf{Nerf (faisceau de fibres) :} Courbe \textbf{sigmoïde (S)}. Recrutement progressif des fibres (diamètres variés $\rightarrow$ seuils variés). Fibres grosses (myélinisées, faible seuil) recrutées en premier. Fibres fines (amylinisées, fort seuil) recrutées en dernier.
        \item \textbf{Seuil du nerf ($I_{seuil}$) :} Intensité déclenchant la réponse dans 50\% des fibres (ou première réponse detectable).
        \item \textbf{Supramaximal :} Intensité $\ge 1.5 \times I_{max}$ (toutes fibres recrutées). Utilisé en clinique (électroneuromyographie - ENMG).
    \end{itemize}
    \item \textbf{Courbe Intensité-Durée (Force-Durée) :} Relation entre l'intensité minimale $I$ pour déclencher un PA et la durée $d$ de l'impulsion de courant (carrée).
    \begin{itemize}
        \item \textbf{Équation de Lapicque :} $I = I_{rh} \left( 1 + \frac{t_{ch}}{d} \right)$.
        \item \textbf{Rhéobase ($I_{rh}$) :} Intensité minimale pour exciter avec une durée \textbf{infinie} (asymptote horizontale). Reflète le \textbf{seuil de tension} de la membrane ($V_{seuil}$) et la résistance d'entrée. Unité : \SI{}{\micro A} ou \SI{}{mA}.
        \item \textbf{Chronaxie ($t_{ch}$) :} Durée \textbf{minimale} pour exciter avec une intensité **$2 \times I_{rh}$\textbf{. Reflète la }constante de temps membranaire** ($\tau_m = R_m C_m$) et la cinétique d'ouverture des canaux \ce{Na+}. Unité : \SI{}{\micro s} ou \SI{}{ms}.
        \item \textbf{Interprétation physiologique :}
            \begin{itemize}
                \item Courte chronaxie (\SI{< 100}{\micro s}) = Fibres \textbf{myélinisées} (grosses, $\tau_m$ faible, canaux \ce{Na+} denses aux nœuds, réponse rapide). Ex: Fibres motrices $\alpha$, A$\alpha$ (\SI{50-100}{\micro s}).
                \item Longue chronaxie (\SI{> 1}{ms}) = Fibres \textbf{amylinisées} (fines, $\tau_m$ forte, canaux \ce{Na+} diffus, réponse lente). Ex: Fibres C (douleur diffuse) \SI{0.4-1}{ms}.
            \end{itemize}
    \end{itemize}
    \item \textbf{Calculs types :} Donner $I_{rh}$ et $t_{ch}$, calculer $I$ pour une $d$ donnée. Ou lire $I_{rh}$ et $t_{ch}$ sur le graphique (asymptote et point $I=2I_{rh}$).
\end{enumerate}

\textbf{Formules à retenir :}
$I = I_{rh} \left( 1 + \frac{t_{ch}}{d} \right)$  (Lapicque, approximation pour impulsions carrées).
$Q_{min} = I \times d$ (Charge minimale). Minimum de charge pour $d = t_{ch}$ : $Q_{min} = 2 I_{rh} t_{ch}$.

\textbf{Pièges :} Unités ($d$ en ms ou $\mu$s ? $I$ en mA ou $\mu$A ?). Confondre Rhéobase (intensité) et Chronaxie (durée). Oublier que la courbe $I(d)$ est une \textbf{hyperbole} (branche convexe).
\end{methode}

\begin{exoresolu}[Détermination de la rhéobase et de la chronaxie d'un nerf moteur]
\textbf{Énoncé :} On étudie l'excitabilité du nerf sciatique d'une grenouille (fibres motrices myélinisées A$\alpha$) à \SI{20}{\celsius}. On utilise des impulsions de courant carrées de durée variable $d$. On détermine l'intensité seuil $I_{seuil}$ (en \SI{}{\micro A}) pour chaque durée. Résultats :

\begin{center}
\begin{tabularx}{\linewidth}{|c|>{\centering\arraybackslash}X|>{\centering\arraybackslash}X|>{\centering\arraybackslash}X|>{\centering\arraybackslash}X|>{\centering\arraybackslash}X|>{\centering\arraybackslash}X|}\hline
\rowcolor{popblueL}
$d$ (\SI{}{\micro s}) & 50 & 100 & 200 & 500 & 1000 & 2000 \\ \hline
$I_{seuil}$ (\SI{}{\micro A}) & 120 & 70 & 45 & 30 & 24 & 22 \\ \hline
\end{tabularx}
\end{center}

\begin{enumerate}
    \item Tracer la courbe $I_{seuil} = f(d)$ (qualitatif : hyperbole décroissante).
    \item Déterminer graphiquement (ou par calcul/régression) la \textbf{rhéobase} ($I_{rh}$) et la \textbf{chronaxie} ($t_{ch}$).
    \item Vérifier la cohérence avec la formule de Lapicque pour $d = \SI{200}{\micro s}$.
    \item Calculer la charge minimale $Q_{min}$ nécessaire pour l'excitation et la durée $d_{opt}$ qui la permet.
    \item Comparer la chronaxie trouvée à celle d'une fibre sensitive amylinisée (type C, chronaxie $\approx \SI{400}{\micro s}$). Expliquer la différence par la structure de la fibre (myélinisation, diamètre, densité canaux).
\end{enumerate}

\tcblower
\textbf{Solution détaillée :}

\textbf{1. Tracé de la courbe :}
Points : (50, 120), (100, 70), (200, 45), (500, 30), (1000, 24), (2000, 22).
La courbe est une \textbf{branche d'hyperbole convexe}, décroissante, avec asymptote horizontale vers \SI{20-22}{\micro A}.

\textbf{2. Détermination de $I_{rh}$ et $t_{ch}$ :}
\begin{itemize}
    \item \textbf{Rhéobase $I_{rh}$ :} Asymptote pour $d \to \infty$. Lecture sur le tableau : pour $d=2000\ \mu s$, $I=22\ \mu A$ ; pour $d=1000\ \mu s$, $I=24\ \mu A$. La valeur limite est **$I_{rh} \approx \mathbf{22\ \mu A}$**.
    \item \textbf{Chronaxie $t_{ch}$ :} Durée pour laquelle $I_{seuil} = 2 \times I_{rh} = \mathbf{44\ \mu A}$.
    Lecture tableau : $I=45\ \mu A$ pour $d=200\ \mu s$ ; $I=30\ \mu A$ pour $d=500\ \mu s$.
    $44\ \mu A$ est très proche de la valeur à $d=200\ \mu s$.
    \textbf{Résultat retenu : $I_{rh} \approx \mathbf{22\ \mu A}$ ; $t_{ch} \approx \mathbf{200\ \mu s}$.} (Valeurs typiques fibre A$\alpha$ grenouille à \SI{20}{\celsius} : Rhéobase ~20-30 $\mu$A, Chronaxie ~100-200 $\mu$s).
\end{itemize}

\textbf{3. Vérification Lapicque pour $d = \SI{200}{\micro s}$ :}
Avec $I_{rh}=22\ \mu A$, $t_{ch}=200\ \mu s$.
$I_{calc} = 22 \times (1 + 200/200) = 22 \times 2 = \mathbf{44\ \mu A}$.
Valeur expérimentale : $\mathbf{45\ \mu A}$.
\textbf{Accord excellent} (erreur < 3\%). Le modèle de Lapicque décrit bien l'excitabilité de cette fibre myélinisée.

\textbf{4. Charge minimale $Q_{min}$ et $d_{opt}$ :}
La charge $Q = I \times d = I_{rh} (d + t_{ch})$.
C'est une fonction \textbf{linéaire croissante} de $d$ ($Q = I_{rh}d + I_{rh}t_{ch}$).
$\rightarrow$ **Pas de minimum absolu pour $Q$ en fonction de $d$** (plus $d$ est court, plus $Q$ est faible, mais $I$ devient énorme, limité par la capacité du stimulateur et l'électrolyse).
\textit{Note pédagogique :} On retient souvent que l'excitation est la plus « économique » en énergie ($W = I^2 R d$) pour $d \approx t_{ch}$.
Calculons $Q$ pour $d = t_{ch} = 200\ \mu s$ : $I = 44\ \mu A$. $Q = 44 \times 10^{-6} \times 200 \times 10^{-6} = \mathbf{8.8 \times 10^{-9}\ C = 8.8\ nC}$.
Pour $d = 1000\ \mu s$ : $I = 22(1+0.2) = 26.4\ \mu A$. $Q = 26.4 \times 1000 = 26.4\ nC$.
$\rightarrow$ **La charge injectée est minimale pour les courtes durées (proches de $t_{ch}$).**

\textbf{5. Comparaison fibre A$\alpha$ (myélinisée) vs Fibre C (amylinisée) :}
\begin{itemize}
    \item \textbf{Fibre A$\alpha$ (motrice) :} $t_{ch} \approx \mathbf{200\ \mu s}$ (courte).
    \item \textbf{Fibre C (douleur diffuse) :} $t_{ch} \approx \mathbf{400\ \mu s \text{ à } 1\ ms}$ (longue).
    \item \textbf{Explication :}
    \begin{enumerate}
        \item \textbf{Constante de temps membranaire $\tau_m = R_m C_m$ :} Fibre C : diamètre petit ($\sim \SI{1}{\micro m}$) $\rightarrow$ $R_m$ (résistance membranaire spécifique) $\times$ surface $\rightarrow$ $R_{in}$ (résistance d'entrée) \textbf{élevée}. $C_m$ proportionnelle à surface. $\tau_m$ \textbf{grande}. Fibre A$\alpha$ : diamètre grand ($\sim \SI{20}{\micro m}$) + \textbf{myéline} (capacité nœudale très faible, $C_m$ nœud $\ll$ $C_m$ fibre nue). $\tau_m$ nœudale \textbf{très faible}.
        \item \textbf{Cinétique canaux \ce{Na+} :} Densité énorme aux nœuds de Ranvier (fibre A$\alpha$) $\rightarrow$ courant \ce{Na+} maximal énorme $\rightarrow$ dépolarisation ultra-rapide $\rightarrow$ seuil atteint vite $\rightarrow$ chronaxie courte. Fibre C : canaux \ce{Na+} diffus, densité faible $\rightarrow$ courant \ce{Na+} lent $\rightarrow$ chronaxie longue.
    \end{enumerate}
    \textbf{Application clinique (ENMG) :} On utilise des impulsions courtes (\SI{100}{\micro s}) pour stimuler sélectivement les fibres motrices (A$\alpha$, chronaxie courte) sans activer les fibres sensitives douloureuses (C, chronaxie longue, seuil plus haut pour courtes durées).
\end{itemize}

\textbf{Phrase de conclusion :} La chronaxie (\SI{200}{\mu s}) reflète la rapidité de la réponse membranaire : les fibres myélinisées de gros diamètre, à faible constante de temps et forte densité de canaux \ce{Na+} aux nœuds, s'excitent par des impulsions brèves et intenses, contrairement aux fibres amylinisées fines qui nécessitent des impulsions plus longues (chronaxie plus grande).
\end{exoresolu}

\begin{methode}[Période réfractaire et vitesse de propagation]
\textbf{Objectif :} Interpréter le montage à double stimulation (deux cathodes) pour mettre en évidence la période réfractaire (absolue et relative) et calculer la vitesse de conduction.

\textbf{Étapes de la démarche :}
\begin{enumerate}
    \item \textbf{Montage « Double Stimulation » (ou double enregistrement) :}
    \begin{itemize}
        \item \textbf{Méthode 1 (Période réfractaire - 1 site enregistrement, 2 stimuli) :} Une électrode de stimulation $S_1$ (conditionnant) et $S_2$ (test) au même endroit (ou très proche). Un enregistrement en aval. On varie l'intervalle $\Delta t$ entre $S_1$ et $S_2$.
        \item \textbf{Méthode 2 (Vitesse - 1 stimulus, 2 sites enregistrement) :} Un stimulus $S$. Deux électrodes d'enregistrement $E_1$ (proximale) et $E_2$ (distale) espacées d'une distance connue $D$. On mesure le délai $\Delta t$ entre les deux PA enregistrés.
    \end{itemize}
    \item \textbf{Analyse Période Réfractaire (Méthode 1) :}
    \begin{itemize}
        \item Tracer : Abscisse = Intervalle $\Delta t$ (ms). Ordonnée = Réponse à $S_2$ (Oui/Non ou Amplitude/Latence).
        \item \textbf{Période Réfractaire Absolue (PRA) :} $\Delta t$ pendant lequel \textbf{aucun PA} n'est évoqué par $S_2$, quelle que soit son intensité (même supramaximale). Correspond à l'**inactivation des canaux \ce{Na+} ($h \approx 0$)**. Durée $\approx$ durée du PA (\SI{0.5-2}{ms}).
        \item \textbf{Période Réfractaire Relative (PRR) :} $\Delta t$ suivant la PRA. Un PA \textbf{peut} être évoqué par $S_2$ **mais seulement si son intensité est $>$ Seuil normal\textbf{ (seuil élevé). Correspond à l'}hyperpolarisation (canaux \ce{K+} ouverts)\textbf{ + }récupération incomplète canaux \ce{Na+} ($h < 1$)**. Durée $\approx$ durée de l'hyperpolarisation (\SI{5-20}{ms}).
        \item \textbf{Période Supra-normale :} Parfois, excitabilité légèrement *augmentée* (seuil abaissé) juste après PRR. Rarement demandé au Bac.
    \end{itemize}
    \item \textbf{Calcul Vitesse de Propagation ($v$) :}
    \[ v = \frac{D}{\Delta t} \]
    $D$ = distance entre $E_1$ et $E_2$ (en m ou mm). $\Delta t$ = différence de latence (temps entre stimulus et PA) entre $E_2$ et $E_1$ (en s ou ms).
    \textbf{Unités :} $v$ en \SI{}{m/s}.
    \item \textbf{Facteurs influençant $v$ (à citer) :}
    \begin{enumerate}
        \item \textbf{Myélinisation :} Fibres myélinisées (conduction saltatoire, $v \propto D$) $\gg$ Fibres amylinisées (conduction continue, $v \propto \sqrt{D}$). \textit{Facteur principal.}
        \item \textbf{Diamètre de la fibre (axone) :} $v$ croît avec le diamètre.
        \item \textbf{Température :} $v$ croît avec $T$ (cinétique canaux plus rapide, $Q_{10} \approx 1,5-2$). Hypothermie $\rightarrow$ ralentissement.
        \item \textbf{État métabolique / Pathologie :} Démyélinisation (SEP, Guillain-Barré) $\rightarrow$ chute drastique de $v$, bloc de conduction, dispersion des latences.
    \end{enumerate}
\end{enumerate}

\textbf{Réflexes Bac :} Unités ! $D$ en m, $\Delta t$ en s $\rightarrow$ $v$ en m/s. Si $D$ en mm, $\Delta t$ en ms $\rightarrow$ $v$ en m/s (car $1\ mm/ms = 1\ m/s$). Ne pas confondre \textbf{PRA} (inactivation \ce{Na+}, infranchissable) et \textbf{PRR} (hyperpolarisation \ce{K+}, franchissable si stimulus plus fort). Sens de propagation : \textbf{orthodrome} (corps cellulaire $\rightarrow$ terminaison) physiologique ; \textbf{antidrome} (inverse) possible en stimulation électrique artificielle.
\end{methode}

\begin{exoresolu}[Mesure de la vitesse de conduction et analyse de la période réfractaire]
\textbf{Énoncé :} On réalise deux expériences sur un nerf sciatique de grenouille (fibres A$\alpha$, diamètre \SI{15}{\micro m}, myélinisé) à \SI{20}{\celsius}.

\textbf{Expérience A (Vitesse) :} Stimulation unique en $S$. Enregistrement en $E_1$ (distance $S-E_1 = \SI{30}{mm}$) et $E_2$ (distance $S-E_2 = \SI{80}{mm}$).
Latence en $E_1$ : $L_1 = \SI{1,2}{ms}$. Latence en $E_2$ : $L_2 = \SI{3,2}{ms}$.

\textbf{Expérience B (Période réfractaire - Montage double stimulation) :} Deux stimuli $S_1$ (conditionnant, supramaximal) et $S_2$ (test, intensité variable) appliqués au même point. Enregistrement du PA en $E$ (loin). On varie l'intervalle $\Delta t = t_{S_2} - t_{S_1}$.
Résultats pour $S_2$ à intensité \textbf{seuil} ($I_{seuil}$) :
\begin{itemize}
    \item $\Delta t = \SI{0,5}{ms}$ : Pas de PA évoqué par $S_2$.
    \item $\Delta t = \SI{1,0}{ms}$ : Pas de PA évoqué par $S_2$.
    \item $\Delta t = \SI{1,5}{ms}$ : Pas de PA évoqué par $S_2$.
    \item $\Delta t = \SI{2,0}{ms}$ : PA évoqué par $S_2$ (amplitude normale).
    \item $\Delta t = \SI{3,0}{ms}$ : PA évoqué par $S_2$ (amplitude normale).
\end{itemize}
Résultats pour $S_2$ à intensité \textbf{$2 \times I_{seuil}$} (supramaximal) :
\begin{itemize}
    \item $\Delta t = \SI{1,0}{ms}$ : Pas de PA.
    \item $\Delta t = \SI{1,5}{ms}$ : PA évoqué (amplitude normale).
    \item $\Delta t = \SI{2,0}{ms}$ : PA évoqué.
\end{itemize}

\begin{enumerate}
    \item Calculer la vitesse de conduction $v$ du nerf (Expérience A). Exprimer en \SI{}{m/s}.
    \item Déterminer la durée de la \textbf{Période Réfractaire Absolue (PRA)} et de la \textbf{Période Réfractaire Relative (PRR)} à partir de l'Expérience B. Justifier par les mécanismes ioniques.
    \item Prédire l'effet d'une baisse de température à \SI{10}{\celsius} sur la vitesse $v$ et sur la durée de la PRA. Justifier.
    \item Dans le contexte d'une sclérose en plaques (démélinisation), décrire les modifications attendues sur la vitesse $v$ et sur la fidélité de la conduction (répétition des stimuli à haute fréquence).
\end{enumerate}

\tcblower
\textbf{Solution détaillée :}

\textbf{1. Calcul de la vitesse de conduction ($v$) :}
Distance entre les deux points d'enregistrement : $D = d_{S-E_2} - d_{S-E_1} = 80 - 30 = \mathbf{\SI{50}{mm}}$.
Différence de latence (temps de parcours entre $E_1$ et $E_2$) : $\Delta t = L_2 - L_1 = 3,2 - 1,2 = \mathbf{\SI{2,0}{ms}}$.
\[ v = \frac{D}{\Delta t} = \frac{\SI{50}{mm}}{\SI{2,0}{ms}} = \mathbf{\SI{25}{m/s}} \]
\textit{Vérification :} Cohérent pour une fibre A$\alpha$ de grenouille à \SI{20}{\celsius} (mammifère \SI{37}{\celsius} $\approx$ \SI{80-100}{m/s}). Facteur température $Q_{10} \approx 1,7$ : $25 \times 1,7^{1,7} \approx 25 \times 2,5 \approx 60$ m/s (ordre de grandeur correct).

\textbf{2. Détermination PRA et PRR :}
\begin{itemize}
    \item \textbf{PRA :} Intervalle pendant lequel \textbf{même un stimulus supramaximal ($2 I_{seuil}$) ne déclenche pas de PA}.
    Données $2 I_{seuil}$ : $\Delta t = 1,0$ ms $\rightarrow$ Pas de PA. $\Delta t = 1,5$ ms $\rightarrow$ PA évoqué.
    $\rightarrow$ \textbf{PRA $\approx \mathbf{1,2 - 1,5\ ms}$} (fin de la dépolarisation / début repolarisation, canaux \ce{Na+} inactivés).
    \item \textbf{PRR :} Intervalle suivant la PRA où un stimulus \textbf{supramaximal} déclenche un PA, mais un stimulus \textbf{seuil} ne le déclenche \textbf{pas} (seuil élevé).
    Données $I_{seuil}$ : $\Delta t = 1,5$ ms $\rightarrow$ Pas de PA. $\Delta t = 2,0$ ms $\rightarrow$ PA évoqué.
    Données $2 I_{seuil}$ : $\Delta t = 1,5$ ms $\rightarrow$ PA évoqué.
    $\rightarrow$ \textbf{PRR $\approx$ de \SI{1,5}{ms} à \SI{2,0}{ms}} (durée $\approx \mathbf{0,5\ ms}$).
    \item \textbf{Mécanismes ioniques :}
    \begin{itemize}
        \item \textbf{PRA (0 à 1,5 ms) :} Canaux \ce{Na+} voltage-dépendants en \textbf{état inactivé} (particule $h$ bloquée). Impossible à ouvrir par n'importe quel stimulus. Correspond à la phase de dépolarisation + début repolarisation.
        \item \textbf{PRR (1,5 à 2,0 ms) :} Canaux \ce{Na+} en \textbf{récupération de l'inactivation} ($h$ augmente vers 1) MAIS canaux \ce{K+} voltage-dépendants \textbf{encore ouverts} $\rightarrow$ \textbf{Hyperpolarisation} ($V_m < V_{repos}$). Il faut un stimulus plus fort (dépolarisation plus grande) pour amener $V_m$ au seuil. À $\Delta t = 2,0$ ms, hyperpolarisation finie, canaux \ce{Na+} récupérés $\rightarrow$ seuil normal rétabli.
    \end{itemize}
\end{itemize}

\textbf{3. Effet de la température (\SI{10}{\celsius} au lieu de \SI{20}{\celsius}) :}
\begin{itemize}
    \item \textbf{Vitesse $v$ :} \textbf{Diminue}. Les réactions enzymatiques (pompes) et transitions conformationnelles des canaux (ouverture/fermeture) sont ralenties ($Q_{10} \approx 2-3$ pour la cinétique des canaux). La dépolarisation est plus lente $\rightarrow$ le PA met plus de temps à se propager d'un nœud à l'autre. $v$ divisé par $\approx 2$ à \num{3}.
    \item \textbf{Durée PRA :} \textbf{Augmente}. L'inactivation des canaux \ce{Na+} (processus dépendant de la température) met plus de temps à se rétablir. La repolarisation (sortie \ce{K+}) est plus lente. Le PA s'élargit. PRA $\approx$ durée du PA $\rightarrow$ PRA plus longue.
\end{itemize}

\textbf{4. Sclérose en plaques (Démélinisation) :}
\begin{itemize}
    \item \textbf{Vitesse $v$ :} \textbf{Chute drastique}. Perte de l'isolant myélique $\rightarrow$ capacité membranaire $C_m$ augmente enormément aux zones démyélinisées, résistance $R_m$ chute. Constante de temps $\tau_m$ augmente, constante de longueur $\lambda$ diminue. La conduction saltatoire est rompue $\rightarrow$ conduction continue lente ou \textbf{bloc de conduction} ($v \to 0$).
    \item \textbf{Fidélité à haute fréquence :} \textbf{Défaillance}. Lors de stimulations répétées (ex. \SI{50}{Hz}), la zone démyélinisée ne peut pas suivre.
    \begin{itemize}
        \item Accumulation de \ce{K+} extracellulaire dans l'espace périneural restreint (pas de myéline pour tamponner) $\rightarrow$ dépolarisation permanente $\rightarrow$ inactivation \ce{Na+} $\rightarrow$ bloc.
        \item Période réfractaire prolongée (élargissement PA) $\rightarrow$ impossibilité de suivre des fréquences élevées.
        \item \textbf{Dispersion des latences} : fibres touchées différemment $\rightarrow$ potentiel d'action composé (neurogramme) évasé, polyphasique.
    \end{itemize}
    \textit{Signe clinique :} Fatigabilité musculaire à l'effort, troubles visuels (névrite optique) déclenchés par la chaleur (bain chaud, exercice - phénomène d'Uhthoff).
\end{itemize}

\textbf{Phrase de conclusion :} La vitesse de conduction (\SI{25}{m/s} ici) dépendcritiquement de la myélinisation et du diamètre ; la période réfractaire absolue (\SI{1,5}{ms}) garantit le sens unique de propagation et limite la fréquence maximale (\SI{660}{Hz} théorique ici) ; la démyélinisation pathologique détruit la conduction saltatoire, ralentissant ou bloquant l'influx nerveux.
\end{exoresolu}

\begin{methode}[Sens de propagation et mode de propagation (Saltatoire vs Continue)]
\textbf{Objectif :} Expliquer le sens unique de propagation (orthodrome) et contraster conduction saltatoire (myélinisée) et continue (amylinisée).

\textbf{Étapes de la démarche :}
\begin{enumerate}
    \item \textbf{Sens de propagation (Orthodrome) :}
    \begin{itemize}
        \item \textbf{Physiologique :} Corps cellulaire / Zone de déclenchement (cône d'émergence) $\rightarrow$ Axone $\rightarrow$ Terminaisons synaptiques.
        \item \textbf{Base biophysique :} La \textbf{période réfractaire absolue (PRA)} empêche le retour en arrière. Juste après le passage du PA, la zone proximale est en PRA (canaux \ce{Na+} inactivés) : elle ne peut pas être redépolarisée par le courant local circulant vers l'amont. La zone distale, au repos, est excitable. Le courant local ($I_{local}$) ne peut donc circuler que vers l'aval (dépolarisation de la zone au repos).
        \item \textbf{Stimulation électrique artificielle :} Peut déclencher un PA n'importe où. La propagation est alors \textbf{bidirectionnelle} (orthodrome + antidrome) depuis le point de stimulation.
    \end{itemize}
    \item \textbf{Mode de propagation :}
    \begin{center}
    \begin{tabularx}{\linewidth}{|l|X|X|}\hline
    \rowcolor{popblueL} \textbf{Paramètre} & \textbf{Conduction Continue (Fibre amylinisée, ex. fibre C)} & \textbf{Conduction Saltatoire (Fibre myélinisée, ex. fibre A$\alpha$)} \\ \hline
    \textbf{Myéline} & Absente & Présente (feuillets de Schwann / oligodendrocytes) \\ \hline
    \textbf{Canaux \ce{Na+}} & Répartis uniformément sur toute la membrane & Concentrés aux \textbf{nœuds de Ranvier} (densité $\sim 1000-2000/\mu m^2$) \\ \hline
    \textbf{Dépolarisation} & Continue, point par point ($\Delta x \sim \lambda$) & Discontinue : \textbf{saut} de nœud en nœud \\ \hline
    \textbf{Courant local} & Faible, circule sur courte distance & Fort, circule \textbf{sous la myéline} (faible $C_m$, forte $R_m$) d'un nœud au suivant \\ \hline
    \textbf{Vitesse $v$} & Faible : $v \propto \sqrt{d}$ ($d$ = diamètre) & Élevée : $v \propto d$ (linéaire avec le diamètre) \\ \hline
    \textbf{Exemples / Valeurs} & Fibres C (douleur diffuse) : \SI{0,5-2}{m/s} & Fibres A$\alpha$ (motrices) : \SI{80-120}{m/s} (homme, \SI{37}{\celsius}) \\ \hline
    \textbf{Coût métabolique} & Élevé (beaucoup d'ions \ce{Na+}/\ce{K+} à repomper par $cm^2$) & Faible (échanges ioniques limités aux nœuds) \\ \hline
    \end{tabularx}
    \end{center}
    \item \textbf{Schéma à savoir tracer :} Axone myélinisé avec nœuds de Ranvier. Flèches « Courant local » sous la myéline (flèches larges) et « Courant ionique \ce{Na+}/\ce{K+} » aux nœuds (flèches courtes traversant la membrane). Indiquer sens orthodrome.
\end{enumerate}

\textbf{Formules clés :}
\begin{itemize}
    \item Vitesse continue : $v \approx \frac{k}{\sqrt{R_i C_m}} \sqrt{d}$ ($R_i$ résistivité axoplasmique).
    \item Vitesse saltatoire : $v \approx \frac{L_{internoeud}}{\tau_{noeud}} \propto d$ ($L_{internoeud} \propto d$).
    \item Constante de longueur (espace) : $\lambda = \sqrt{\frac{R_m}{R_i}} \cdot \sqrt{\frac{d}{4}}$. Myéline $\uparrow R_m \Rightarrow \uparrow \lambda \Rightarrow$ saut plus long.
\end{itemize}

\textbf{Réflexes Bac :} Ne pas dire « le courant saute ». Dire : « La \textbf{dépolarisation active} (ouverture canaux \ce{Na+}) \textbf{ne se produit qu'aux nœuds de Ranvier} ; le \textbf{courant local passif} circule sous la myéline pour dépolariser le nœud suivant au seuil. » La myéline agit comme un \textbf{isolant capacitif} (diminue $C_m$ 100x) et \textbf{résistif} (augmente $R_m$ 1000x).
\end{methode}

\begin{exoresolu}[Analyse comparative : conduction continue vs saltatoire]
\textbf{Énoncé :} On compare la conduction de l'influx nerveux dans une fibre sensitive amylinisée de type C (diamètre \SI{1}{\micro m}, pas de myéline) et une fibre motrice myélinisée de type A$\alpha$ (diamètre \SI{20}{\micro m}, gaine de myéline épaisse, longueur internodale \SI{2}{mm}) chez l'homme à \SI{37}{\celsius}.

\begin{enumerate}
    \item Pour chaque type de fibre, décrire le mode de propagation (continue ou saltatoire) et expliquer le rôle de la myéline et la localisation des canaux \ce{Na+} voltage-dépendants.
    \item La vitesse de conduction mesurée est de \SI{1}{m/s} pour la fibre C et \SI{100}{m/s} pour la fibre A$\alpha$. Vérifier les lois d'échelle : $v \propto \sqrt{d}$ pour la fibre C et $v \propto d$ pour la fibre A$\alpha$ en calculant les coefficients de proportionnalité.
    \item Calculer le temps nécessaire pour qu'un influx parcoure \SI{1}{m} dans chaque fibre. Discuter de l'intérêt physiologique de la myélinisation pour le contrôle moteur.
    \item Une pathologie démyélinisante (sclérose en plaques) atteint la fibre A$\alpha$. La gaine de myéline est détruite sur un segment de \SI{5}{mm} (plaque), exposant l'axone nu (diamètre \SI{20}{\micro m}, canaux \ce{Na+} diffusés). Prédire l'effet sur la conduction à cet endroit (vitesse, sécurité de la conduction, forme du PA). Justifier par les constantes de temps et de longueur.
\end{enumerate}

\tcblower
\textbf{Solution détaillée :}

\textbf{1. Mode de propagation et rôle de la myéline :}
\begin{itemize}
    \item \textbf{Fibre C (amylinisée) : Conduction CONTINUE.}
    \begin{itemize}
        \item Canaux \ce{Na+} voltage-dépendants : \textbf{Répartis uniformément} sur toute la surface de l'axolemme.
        \item Mécanisme : Dépolarisation active en chaque point. Le courant local ($I_{local}$) circule dans l'axoplasme vers la zone adjacente au repos, la dépolarise au seuil $\rightarrow$ ouverture canaux \ce{Na+} $\rightarrow$ nouveau PA. Propagation « pas à pas » (échelle de la constante de longueur $\lambda \approx \SI{1-2}{mm}$).
        \item Pas de myéline $\rightarrow$ Capacité membranaire $C_m$ élevée (\SI{1}{\micro F/cm^2}), Résistance $R_m$ faible $\rightarrow$ Courant local faible, $\lambda$ courte.
    \end{itemize}
    \item \textbf{Fibre A$\alpha$ (myélinisée) : Conduction SALTOIRE.}
    \begin{itemize}
        \item Canaux \ce{Na+} voltage-dépendants : \textbf{Concentrés quasi-exclusivement aux nœuds de Ranvier} (largeur \SI{1}{\micro m}, espacés de \SI{2}{mm}).
        \item Myéline : \textbf{Isolant électrique} (lipides). $C_m$ nœudale $\approx \SI{1}{\micro F/cm^2}$ mais \textbf{sous myéline $C_m \approx \SI{0,01}{\micro F/cm^2}$} (100x moins). $R_m$ \textbf{sous myéline $\approx \SI{1000}{k\Omega.cm^2}$} (1000x plus).
        \item Mécanisme : PA \textbf{uniquement aux nœuds}. Courant local \textbf{passif, rapide, sans perte} sous la myéline (grande $\lambda \approx \SI{5-10}{mm} > L_{internoeud}$) $\rightarrow$ dépolarise le nœud suivant \textbf{au-delà du seuil} $\rightarrow$ déclenchement PA au nœud suivant. « Saut » de nœud en nœud.
    \end{itemize}
\end{itemize}

\textbf{2. Vérification des lois d'échelle :}
\begin{itemize}
    \item \textbf{Fibre C :} $v = k_C \sqrt{d}$. $d = \SI{1}{\mu m}$. $v = \SI{1}{m/s}$.
    $k_C = \frac{v}{\sqrt{d}} = \frac{1}{\sqrt{1}} = \mathbf{1\ m/s/\mu m^{1/2}}$.
    \item \textbf{Fibre A$\alpha$ :} $v = k_A d$. $d = \SI{20}{\mu m}$. $v = \SI{100}{m/s}$.
    $k_A = \frac{v}{d} = \frac{100}{20} = \mathbf{5\ m/s/\mu m}$.
    \item \textbf{Comparaison :} Pour $d = \SI{20}{\mu m}$, fibre C hypothétique : $v = 1 \times \sqrt{20} \approx \SI{4,5}{m/s}$. Fibre A$\alpha$ réelle : \SI{100}{m/s}. \textbf{Gain de facteur 22} grâce à la myélinisation (saltatoire + gros diamètre).
\end{itemize}

\textbf{3. Temps de parcours sur \SI{1}{m} :}
\begin{itemize}
    \item Fibre C : $t_C = \frac{L}{v} = \frac{1}{1} = \mathbf{\SI{1}{s}}$.
    \item Fibre A$\alpha$ : $t_A = \frac{1}{100} = \mathbf{\SI{10}{ms}}$.
    \item \textbf{Intérêt physiologique :} Pour un motoneurone innervant un muscle du pied (longueur nerf $\approx \SI{1}{m}$), le commandement moteur arrive en \SI{10}{ms} (myélinisé) au lieu de \SI{1}{s} (amylinisé). \textbf{Gain de \SI{990}{ms}}. Essentiel pour la \textbf{réactivité}, la \textbf{coordination fine}, l'\textbf{équilibre} (réflexes posturaux < \SI{50}{ms}) et la \textbf{fuite face au danger}. La myélinisation est une adaptation majeure des vertébrés pour la vitesse et l'économie d'énergie (pompe \ce{Na+/K+} 100x moins sollicitée par $cm$ d'axone).
\end{itemize}

\textbf{4. Effet d'une plaque démyélinisante (\SI{5}{mm}) sur fibre A$\alpha$ :}
\begin{itemize}
    \item \textbf{Structure locale :} Axone nu, diamètre \SI{20}{\mu m}, canaux \ce{Na+} \textbf{diffus} (densité faible, type fibre amylinisée), $C_m$ et $R_m$ revenus aux valeurs « nues ».
    \item \textbf{Vitesse dans la plaque :} Conduction \textbf{continue} imposée. $v_{plaque} \approx k_C \sqrt{d} = 1 \times \sqrt{20} \approx \mathbf{\SI{4,5}{m/s}}$ (au lieu de \SI{100}{m/s}). \textbf{Ralentissement 20x}.
    \item \textbf{Sécurité de la conduction (Facteur de sécurité) :}
    \begin{itemize}
        \item Normal (saltatoire) : Courant local énorme (grosse $\lambda$) $\rightarrow$ nœud suivant dépolarisé \textbf{largement au-dessus du seuil} (facteur de sécurité $\gg 1$).
        \item Plaque : Courant local faible (petite $\lambda$ nu) $\rightarrow$ dépolarisation du segment suivant \textbf{juste au seuil ou en dessous}.
        \item \textbf{Risque :} \textbf{Bloc de conduction} (échec de propagation) surtout à haute fréquence (accumulation \ce{K+}, inactivation \ce{Na+}) ou si la plaque est longue ($> \lambda$).
    \end{itemize}
    \item \textbf{Forme du PA dans la plaque :} \textbf{Élargi} (durée $\uparrow$). Dépolarisation plus lente (densité canaux \ce{Na+} faible), repolarisation plus lente (canaux \ce{K+} diffus). Amplitude potentiellement réduite (courant \ce{Na+} maximal plus faible).
    \item \textbf{Conséquence clinique :} \textbf{Dispersion des vitesses} (conduction saltatoire rapide $\to$ lent dans plaque $\to$ rapide). Le potentiel d'action composé (enregistrement nerveux) devient \textbf{polyphasique, évasé, de latence prolongée}. Symptômes : troubles de la conduction (paresthésies, faiblesse, ataxie) aggravés par la chaleur (phénomène d'Uhthoff : chaleur $\uparrow$ $\rightarrow$ cinétique canaux $\uparrow$ $\rightarrow$ inactivation $\uparrow$ $\rightarrow$ bloc).
\end{itemize}

\textbf{Phrase de conclusion :} La myélinisation permet une conduction saltatoire rapide (\SI{100}{m/s}), économe en énergie et sûre ; sa destruction impose une conduction continue lente (\SI{4,5}{m/s}), coûteuse et à haut risque de bloc, expliquant les déficits neurologiques de la sclérose en plaques.
\end{exoresolu}

\begin{attention}[Les 5 erreurs qui coûtent des points au Bac]
\begin{enumerate}
    \item \textbf{Confondre « Seuil d'intensité du stimulus » et « Seuil de tension (de dépolarisation) ».}
    \textit{Erreur :} « Le seuil du neurone est de 10 µA. »
    \textit{Correct :} « L'\textbf{intensité seuil} du stimulus est de \SI{10}{\micro A} (pour une durée donnée) ; le \textbf{seuil de tension} de la membrane est d'environ \SI{-55}{mV}. » Le seuil de tension est une propriété biophysique de la membrane (canaux \ce{Na+}) ; l'intensité seuil dépend de l'électrode, de la durée, de la distance, de la chronaxie.
    
    \item \textbf{Dire que « le potentiel de récepteur se propage » ou « est un potentiel d'action ».}
    \textit{Erreur :} « Le potentiel de récepteur se propage le long de l'axone. »
    \textit{Correct :} « Le potentiel de récepteur est un \textbf{potentiel gradué, local, non propagé activement}. Il se propage \textbf{électrotoniquement} (passivement, avec atténuation) jusqu'à la zone de déclenchement (cône d'émergence) où il peut déclencher un \textbf{potentiel d'action} (tout-ou-rien, propagé activement). »
    
    \item \textbf{Inverser les mécanismes de la PRA et de la PRR (ou les confondre).}
    \textit{Erreur :} « Pendant la PRA, les canaux K+ sont ouverts donc on ne peut pas exciter. » / « La PRR est due à l'inactivation des canaux Na+. »
    \textit{Correct :} \textbf{PRA} = \textbf{Inactivation des canaux \ce{Na+}} ($h=0$, infranchissable). \textbf{PRR} = \textbf{Hyperpolarisation (canaux \ce{K+} ouverts) + récupération incomplète canaux \ce{Na+}} ($h < 1$, franchissable si stimulus $>$ seuil). La PRA précède la PRR.
    
    \item \textbf{Oublier les unités ou faire des erreurs de conversion dans le calcul de la vitesse ($v = D/\Delta t$).}
    \textit{Erreur :} $D = \SI{50}{mm}, \Delta t = \SI{2}{ms} \Rightarrow v = 25\ mm/ms$ (sans convertir en m/s) ou $v = 0,025\ m/s$.
    \textit{Correct :} \textbf{Travailler en unités cohérentes SI ou mm/ms.}
    $v = \frac{\SI{50}{mm}}{\SI{2}{ms}} = \SI{25}{mm/ms} = \mathbf{\SI{25}{m/s}}$ (car $1\ mm/ms = 1\ m/s$).
    Ou : $D = \SI{0,05}{m}, \Delta t = \SI{0,002}{s} \Rightarrow v = \frac{0,05}{0,002} = \SI{25}{m/s}$.
    \textbf{Vérifier l'ordre de grandeur :} Fibre myélinisée = dizaines de m/s. Fibre amylinisée = unités ou dixièmes de m/s.
    
    \item \textbf{Affirmer que « l'amplitude du potentiel d'action augmente avec l'intensité du stimulus » (violation de la loi du tout-ou-rien).}
    \textit{Erreur :} « Un stimulus plus fort donne un potentiel d'action plus grand. »
    \textit{Correct :} « L'amplitude du potentiel d'action est \textbf{constante} (loi du tout-ou-rien) pour tout stimulus d'intensité $\ge$ seuil. Seule la \textbf{fréquence} des potentiels d'action code l'intensité du stimulus (codage fréquentiel). » Préciser : « Amplitude invariante, Fréquence variable. »
\end{enumerate}
\end{attention}

\newpage

\coursec{Exercices}

\courssub{Je vérifie mes connaissances}

\begin{exercice}[QCM : Mécanismes ioniques du potentiel de repos et d'action]
\textbf{1.} Le potentiel de repos d'un neurone typique est d'environ $-70\,\text{mV}$. Cette valeur est principalement due à :
\begin{enumerate}[label=\Alph*.]
\item une forte concentration intracellulaire en $\text{Na}^+$ et une perméabilité élevée au $\text{Na}^+$ au repos.
\item une forte concentration extracellulaire en $\text{K}^+$ et une perméabilité nulle au $\text{K}^+$ au repos.
\item \textbf{une forte concentration intracellulaire en $\text{K}^+$, une forte concentration extracellulaire en $\text{Na}^+$ et une perméabilité membranaire au repos beaucoup plus grande pour $\text{K}^+$ que pour $\text{Na}^+$.}
\item l'absence de pompe $\text{Na}^+/\text{K}^+$-ATPase.
\end{enumerate}

\textbf{2.} Lors de la phase de dépolarisation du potentiel d'action, l'événement majeur est :
\begin{enumerate}[label=\Alph*.]
\item l'ouverture massive des canaux $\text{K}^+$ voltage-dépendants.
\item la fermeture des canaux $\text{Na}^+$ voltage-dépendants.
\item \textbf{l'ouverture massive et transitoire des canaux $\text{Na}^+$ voltage-dépendants.}
\item l'activation de la pompe $\text{Na}^+/\text{K}^+$-ATPase.
\end{enumerate}

\textbf{3.} La période réfractaire absolue correspond au moment où :
\begin{enumerate}[label=\Alph*.]
\item un second potentiel d'action peut être déclenché par un stimulus plus intense que le seuil.
\item \textbf{aucun second potentiel d'action ne peut être déclenché, quel que soit l'intensité du stimulus.}
\item la membrane est hyperpolarisée par rapport au potentiel de repos.
\item les canaux $\text{Na}^+$ sont tous fermés et activables.
\end{enumerate}

\textbf{4.} La conduction saltatoire, observée sur les fibres myélinisées, se caractérise par :
\begin{enumerate}[label=\Alph*.]
\item une propagation continue le long de toute la membrane.
\item \textbf{une propagation par sauts d'un nœud de Ranvier au suivant, augmentant la vitesse.}
\item une vitesse de propagation plus lente que sur les fibres amyélinisées de même diamètre.
\item l'absence de dépolarisation au niveau des nœuds de Ranvier.
\end{enumerate}
\end{exercice}

\begin{exercice}[Vrai ou Faux ? Justifiez chaque réponse par une phrase précise.]
\begin{enumerate}
\item \textbf{Au potentiel de repos, les forces électrochimiques s'exerçant sur $\text{K}^+$ et $\text{Na}^+$ sont nulles.}
\par \textbf{Faux.} Au potentiel de repos ($V_r \approx -70\,\text{mV}$), la force électrochimique sur $\text{K}^+$ est quasi nulle ($V_r \approx E_{\text{K}}$), mais elle est forte et orientée vers l'intérieur pour $\text{Na}^+$ ($V_r \ll E_{\text{Na}}$), ce qui explique l'entrée permanente de $\text{Na}^+$ compensée par la pompe.
\item \textbf{Le seuil d'excitation d'une fibre nerveuse correspond à la valeur de potentiel membranaire à laquelle le courant entrant $\text{Na}^+$ devient supérieur au courant sortant $\text{K}^+$.}
\par \textbf{Vrai.} C'est le point de bascule où la rétroaction positive (ouverture des canaux $\text{Na}^+$ voltage-dépendants) l'emporte sur les courants sortants ($\text{K}^+$, fuite), entraînant l'emballement dépolarisant.
\item \textbf{La vitesse de propagation du potentiel d'action est proportionnelle au diamètre de la fibre pour les fibres amyélinisées, et proportionnelle au diamètre pour les fibres myélinisées.}
\par \textbf{Faux.} Pour les fibres \emph{amyélinisées}, la vitesse est proportionnelle à la \emph{racine carrée} du diamètre ($v \propto \sqrt{d}$). Pour les fibres \emph{myélinisées}, elle est approximativement \emph{proportionnelle au diamètre} ($v \propto d$).
\item \textbf{Un potentiel de récepteur est un potentiel gradué, non propagé, dont l'amplitude est proportionnelle à l'intensité du stimulus (dans une certaine plage).}
\par \textbf{Vrai.} C'est un potentiel local (électrotonique), sans seuil, qui code l'intensité du stimulus par son amplitude (codage analogique) avant d'être converti en fréquence de potentiels d'action (codage numérique) au niveau de la zone de déclenchement.
\item \textbf{La rheobase est la durée minimale de stimulation nécessaire pour déclencher un potentiel d'action avec un courant de seuil minimal.}
\par \textbf{Faux.} La \textbf{rheobase} ($I_{\infty}$) est l'\textbf{intensité minimale} de courant (pour une durée infinie) capable de déclencher un potentiel d'action. Le \textbf{temps utile} (chronaxie, $t_u$) est la durée minimale nécessaire avec un courant d'intensité $2 \times \text{rheobase}$.
\end{enumerate}
\end{exercice}

\begin{exercice}[Définitions et vocabulaire précis]
Donnez une définition rigoureuse (une à deux phrases) pour chacun des termes suivants, en utilisant le vocabulaire scientifique attendu au Baccalauréat :
\begin{enumerate}
\item \textbf{Seuil d'excitation (d'une fibre nerveuse) :} Valeur critique du potentiel membranaire (environ $-55\,\text{mV}$) à laquelle l'ouverture des canaux $\text{Na}^+$ voltage-dépendants devient autocatalytique (courant entrant $\text{Na}^+$ $>$ courant sortant $\text{K}^+$), déclenchant inévitablement un potentiel d'action.
\item \textbf{Période réfractaire absolue :} Intervalle de temps immédiatement après le pic du potentiel d'action pendant lequel les canaux $\text{Na}^+$ sont inactivés (boule de verrouillage fermée) ; aucun second potentiel d'action ne peut être initié, quelle que soit l'intensité du stimulus.
\item \textbf{Conduction saltatoire :} Mode de propagation du potentiel d'action le long des fibres myélinisées où la dépolarisation « saute » d'un nœud de Ranvier à l'autre, la gaine de myéline isolant l'axone et concentrant les canaux $\text{Na}^+$ aux nœuds, ce qui accélère la conduction et économise l'énergie.
\item \textbf{Potentiel de récepteur (ou potentiel générateur) :} Potentiel gradué, local et non propagé, généré au niveau d'un récepteur sensoriel par transduction de l'énergie du stimulus (ouverture de canaux mécanosensibles, chimiques, etc.) ; son amplitude est proportionnelle à l'intensité du stimulus.
\item \textbf{Rheobase et Chronaxie :} La \textbf{rheobase} ($I_{\infty}$) est l'intensité seuil minimale pour une stimulation de durée infinie. La \textbf{chronaxie} ($t_u$) est la durée de stimulation nécessaire pour déclencher une réponse avec une intensité égale à $2 \times \text{rheobase}$ ; elle caractérise l'excitabilité pour les stimuli brefs.
\item \textbf{Codage de l'intensité du stimulus dans le message nerveux (loi de la fréquence) :} L'intensité du stimulus est codée par la \textbf{fréquence d'émission des potentiels d'action} (nombre de PA par seconde) le long de la fibre afferente ; plus le stimulus est intense, plus la fréquence instantanée est élevée (jusqu'à la fréquence maximale limitée par la période réfractaire).
\end{enumerate}
\end{exercice}

\begin{exercice}[Calculs directs : Vitesse de propagation et équation de Nernst]
\textbf{Données :} $R = 8,31\,\text{J}\cdot\text{mol}^{-1}\cdot\text{K}^{-1}$ ; $F = 96\,500\,\text{C}\cdot\text{mol}^{-1}$ ; $T = 293\,\text{K}$ ($20^\circ\text{C}$) ; $\frac{RT}{F} \approx 25,3\,\text{mV}$ ; $\log_{10} \approx 2,303 \times \ln$.
\begin{enumerate}
\item On enregistre un potentiel d'action au niveau de deux électrodes placées sur un axone de calmar géant (amyélinisé) distantes de $d = 10\,\text{cm}$. Le décalage temporel entre les deux pics est $\Delta t = 2,5\,\text{ms}$. Calculer la vitesse de propagation $v$ en $\text{m}\cdot\text{s}^{-1}$.
\par \textbf{Solution :} $v = \frac{d}{\Delta t} = \frac{0,10\,\text{m}}{2,5 \times 10^{-3}\,\text{s}} = 40\,\text{m}\cdot\text{s}^{-1}$.
\item Sur une fibre myélinisée de même diamètre, la vitesse mesurée est de $50\,\text{m}\cdot\text{s}^{-1}$. Comparer et expliquer brièvement la différence.
\par \textbf{Comparaison :} La fibre myélinisée conduit plus vite ($50\,\text{m/s}$ vs $40\,\text{m/s}$) malgré un diamètre identique. \textbf{Explication :} La gaine de myéline augmente la résistance membranaire ($R_m \uparrow$) et diminue la capacité membranaire ($C_m \downarrow$), ce qui augmente la constante de longueur ($\lambda \uparrow$) et diminue la constante de temps ($\tau \downarrow$), permettant au courant local de se propager plus loin et plus vite (conduction saltatoire).
\item Calculer le potentiel d'équilibre du potassium ($E_K$) sachant que $[\text{K}^+]_i = 150\,\text{mmol}\cdot\text{L}^{-1}$ et $[\text{K}^+]_o = 5\,\text{mmol}\cdot\text{L}^{-1}$. Utiliser la formule de Nernst simplifiée à $37^\circ\text{C}$ : $E_X = 61,5 \log_{10}\frac{[X]_o}{[X]_i}$ (en mV). Commenter par rapport au potentiel de repos ($V_r \approx -70\,\text{mV}$).
\par \textbf{Calcul :} $E_{\text{K}} = 61,5 \times \log_{10}\left(\frac{5}{150}\right) = 61,5 \times \log_{10}\left(\frac{1}{30}\right) = 61,5 \times (-1,477) \approx -90,8\,\text{mV}$.
\par \textbf{Commentaire :} $V_r (-70\,\text{mV})$ est moins négatif que $E_{\text{K}} (-91\,\text{mV})$. Cela prouve que la membrane au repos n'est pas \emph{exclusivement} perméable au $\text{K}^+$ ; une perméabilité résiduelle au $\text{Na}^+$ ($P_{\text{Na}} > 0$) « tire » le potentiel vers $E_{\text{Na}}$ (positif), le dépolarisant par rapport à $E_{\text{K}}$.
\end{enumerate}
\end{exercice}

\courssub{Je m'entraîne}

\begin{aretenir}[Rappels indispensables pour les exercices]
\begin{itemize}
\item \textbf{Équation de Nernst (à $20^\circ\text{C}$) :} $E_X = 58 \log_{10}\frac{[X]_o}{[X]_i}$ (mV) pour un cation. Pour un anion ($\text{Cl}^-$), $E_{\text{Cl}} = -58 \log_{10}\frac{[\text{Cl}^-]_o}{[\text{Cl}^-]_i} = 58 \log_{10}\frac{[\text{Cl}^-]_i}{[\text{Cl}^-]_o}$.
\item \textbf{Équation de Goldman-Hodgkin-Katz (GHK) simplifiée :} $V_m \approx 58 \log_{10} \frac{P_{\text{K}}[\text{K}^+]_o + P_{\text{Na}}[\text{Na}^+]_o + P_{\text{Cl}}[\text{Cl}^-]_i}{P_{\text{K}}[\text{K}^+]_i + P_{\text{Na}}[\text{Na}^+]_i + P_{\text{Cl}}[\text{Cl}^-]_o}$.
\item \textbf{Relation Intensité-Durée (Lapicque) :} $I = I_{\infty} \left(1 + \frac{t_u}{d}\right)$ où $I_{\infty}$ = rheobase, $t_u$ = chronaxie.
\item \textbf{Vitesse de propagation :} $v = \frac{\text{distance entre électrodes}}{\text{décalage temporel}}$.
\end{itemize}
\end{aretenir}

\begin{exercice}[Interprétation de concentrations ioniques et potentiel de repos]
On réalise l'expérience classique sur l'axone géant de calmar. On mesure les concentrations ioniques internes et externes (tableau ci-dessous) et le potentiel de repos $V_r = -70\,\text{mV}$.

\begin{center}
\begin{tabular}{|c|c|c|}\hline
Ion & Concentration interne (mmol/L) & Concentration externe (mmol/L) \\ \hline
$\text{Na}^+$ & 50 & 440 \\ \hline
$\text{K}^+$ & 400 & 20 \\ \hline
$\text{Cl}^-$ & 40 & 560 \\ \hline
\end{tabular}
\end{center}

\begin{enumerate}
\item Calculer les potentiels d'équilibre $E_{\text{Na}}$, $E_{\text{K}}$ et $E_{\text{Cl}}$ à $20^\circ\text{C}$ (utiliser $58 \log_{10}\frac{[X]_o}{[X]_i}$).
\par \textbf{Solution :}
\begin{align*}
E_{\text{Na}} &= 58 \log_{10}\left(\frac{440}{50}\right) = 58 \log_{10}(8,8) \approx 58 \times 0,944 = \mathbf{+54,8\,mV} \\
E_{\text{K}} &= 58 \log_{10}\left(\frac{20}{400}\right) = 58 \log_{10}(0,05) \approx 58 \times (-1,301) = \mathbf{-75,5\,mV} \\
E_{\text{Cl}} &= 58 \log_{10}\left(\frac{40}{560}\right) = 58 \log_{10}(0,0714) \approx 58 \times (-1,146) = \mathbf{-66,5\,mV}
\end{align*}
\item Comparer $V_r$ à $E_{\text{K}}$, $E_{\text{Na}}$ et $E_{\text{Cl}}$. Que pouvez-vous en déduire sur les perméabilités relatives ($P_{\text{K}}, P_{\text{Na}}, P_{\text{Cl}}$) de la membrane au repos ? Justifiez par l'équation de Goldman-Hodgkin-Katz (qualitativement).
\par \textbf{Comparaison :} $V_r (-70\,\text{mV})$ est très proche de $E_{\text{K}} (-75,5\,\text{mV})$ et $E_{\text{Cl}} (-66,5\,\text{mV})$, mais très loin de $E_{\text{Na}} (+54,8\,\text{mV})$.
\par \textbf{Déduction :} $P_{\text{K}} \gg P_{\text{Na}}$ (la membrane est bien plus perméable au $\text{K}^+$ qu'au $\text{Na}^+$). $P_{\text{Cl}}$ est significative (proche de $P_{\text{K}}$) car $V_r$ est aussi proche de $E_{\text{Cl}}$.
\par \textbf{Justification GHK :} Dans l'équation $V_m \approx 58 \log \frac{P_{\text{K}}[\text{K}^+]_o + P_{\text{Na}}[\text{Na}^+]_o + P_{\text{Cl}}[\text{Cl}^-]_i}{P_{\text{K}}[\text{K}^+]_i + P_{\text{Na}}[\text{Na}^+]_i + P_{\text{Cl}}[\text{Cl}^-]_o}$, si $P_{\text{K}}$ domine le numérateur et le dénominateur, le rapport tend vers $[\text{K}^+]_o/[\text{K}^+]_i$, donnant $V_m \approx E_{\text{K}}$. La légère dépolarisation par rapport à $E_{\text{K}}$ est due au terme $P_{\text{Na}}[\text{Na}^+]_o$ au numérateur.
\item Prédire l'effet sur $V_r$ d'une augmentation brutale de $[\text{K}^+]_o$ à $100\,\text{mmol}\cdot\text{L}^{-1}$. Calculer le nouveau $E_{\text{K}}$ et expliquer le risque physiologique (ex : hyperkaliémie).
\par \textbf{Nouveau $E_{\text{K}}$ :} $E_{\text{K}}' = 58 \log_{10}\left(\frac{100}{400}\right) = 58 \log_{10}(0,25) \approx 58 \times (-0,602) = \mathbf{-34,9\,mV}$.
\par \textbf{Effet sur $V_r$ :} $V_r$ se dépolarise fortement (se rapproche de $-35\,\text{mV}$).
\par \textbf{Risque physiologique :} Cette dépolarisation massive inactivera les canaux $\text{Na}^+$ voltage-dépendants (état inactivé stable), rendant les cellules \emph{inexcitables} (bloc de la conduction nerveuse, paralysie musculaire, arythmies cardiaques fatales). C'est le danger de l'hyperkaliémie sévère.
\end{enumerate}
\end{exercice}

\begin{exercice}[Analyse des courbes de conductances $g_{\text{Na}}$ et $g_{\text{K}}$ pendant un potentiel d'action]
On vous donne le tracé schématique ci-dessous représentant la variation du potentiel membranaire $V_m$ (en noir), de la conductance sodique $g_{\text{Na}}$ (en rouge) et de la conductance potassique $g_{\text{K}}$ (en bleu) en fonction du temps lors d'un potentiel d'action.

\begin{popfigure}
\begin{tikzpicture}[scale=0.8, font=\small]
% Axes
\draw[->] (0,0) -- (12,0) node[right] {$t\,(\text{ms})$};
\draw[->] (0,-1) -- (0,5) node[above] {$V_m\,(\text{mV})$};
\draw[->] (0,-3) -- (0,1) node[above] {$g\,(\text{mS/cm}^2)$};
% Time markers
\foreach \x/\label in {1/0, 2/1, 3/2, 4/3, 5/4, 6/5} \draw (\x, -0.1) -- (\x, 0.1) node[below] {\label};
\foreach \x/\label in {1/0, 2/1, 3/2, 4/3, 5/4, 6/5} \draw (\x, -3.1) -- (\x, -2.9) node[below] {\label};
% Vm curve (black)
\draw[popdark, thick, domain=1:6, samples=200] plot (\x, { 
    (\x < 1.5) ? -0.7 : 
    (\x < 2.0) ? -0.7 + 5.5*(\x-1.5)/0.5 : 
    (\x < 2.5) ? 4.8 - 3.0*(\x-2.0)/0.5 : 
    (\x < 3.5) ? 1.8 - 2.5*(\x-2.5)/1.0 : 
    (\x < 5.0) ? -0.7 - 0.3*(\x-3.5)/1.5 : -0.7 
});
% gNa curve (red)
\draw[poppink, thick, domain=1:6, samples=200] plot (\x, { 
    (\x < 1.8) ? 0 : 
    (\x < 2.1) ? 3.5*(\x-1.8)/0.3 : 
    (\x < 2.6) ? 3.5 - 3.5*(\x-2.1)/0.5 : 0 
});
% gK curve (blue)
\draw[popblue, thick, domain=1:6, samples=200] plot (\x, { 
    (\x < 2.0) ? 0 : 
    (\x < 3.5) ? 2.5*(\x-2.0)/1.5 : 
    (\x < 5.5) ? 2.5 - 2.5*(\x-3.5)/2.0 : 0 
});
% Labels
\node[popdark] at (0.5, 4.5) {$V_m$};
\node[poppink] at (0.5, 0.5) {$g_{\text{Na}}$};
\node[popblue] at (0.5, -2.5) {$g_{\text{K}}$};
\draw[dashed, thin] (2, -3) -- (2, 5);
\draw[dashed, thin] (2.5, -3) -- (2.5, 5);
\draw[dashed, thin] (3.5, -3) -- (3.5, 5);
\end{tikzpicture}
\legende{Variations schématiques du potentiel membranaire et des conductances $g_{\text{Na}}$ et $g_{\text{K}}$ au cours d'un potentiel d'action.}
\end{popfigure}

\begin{enumerate}
\item Identifier sur la figure les phases : dépolarisation, repolarisation, hyperpolarisation (post-potentiel).
\par \textbf{Réponse :} 
\begin{itemize}
\item \textbf{Dépolarisation :} $t \approx 1,5$ à $2,0\,\text{ms}$ (montée brutale de $V_m$).
\item \textbf{Repolarisation :} $t \approx 2,0$ à $3,5\,\text{ms}$ (redescente de $V_m$ vers $V_r$).
\item \textbf{Hyperpolarisation (post-potentiel) :} $t \approx 3,5$ à $5,0\,\text{ms}$ ($V_m$ plus négatif que $V_r$).
\end{itemize}
\item Relier l'évolution de $g_{\text{Na}}$ à la phase de dépolarisation (mécanisme à rétroaction positive). Pourquoi $g_{\text{Na}}$ chute-t-elle brutalement après le pic ?
\par \textbf{Lien :} L'ouverture des canaux $\text{Na}^+$ voltage-dépendants ($g_{\text{Na}} \uparrow$) laisse entrer $\text{Na}^+$ ($V_m \uparrow$), ce qui ouvre encore plus de canaux $\text{Na}^+$ : c'est la \textbf{rétroaction positive} (boucle de Hodgkin-Huxley).
\par \textbf{Chute de $g_{\text{Na}}$ :} Deux mécanismes simultanés : (1) \textbf{Inactivation} (boule de verrouillage intracellulaire qui bouche le pore, dépendante du temps et du voltage) ; (2) \textbf{Fermeture} (désactivation) des portes d'activation lors de la repolarisation. L'inactivation est la cause majeure de l'arrêt brutal de l'entrée de $\text{Na}^+$.
\item Relier l'évolution de $g_{\text{K}}$ à la phase de repolarisation et d'hyperpolarisation. Pourquoi $g_{\text{K}}$ met-elle plus de temps à revenir à zéro ?
\par \textbf{Lien :} L'ouverture retardée des canaux $\text{K}^+$ voltage-dépendants ($g_{\text{K}} \uparrow$ vers $t=2\,\text{ms}$) augmente le courant sortant de $\text{K}^+$, repoussant $V_m$ vers $E_{\text{K}}$ (repolarisation). Comme $g_{\text{K}}$ reste élevée après le retour à $V_r$, $V_m$ dépasse $V_r$ vers $E_{\text{K}}$ (hyperpolarisation).
\par \textbf{Lenteur :} Les canaux $\text{K}^+$ (type retardé) n'ont \emph{pas de mécanisme d'inactivation rapide} (pas de boule de verrouillage) ; leur fermeture (désactivation) est lente et suit le retour de $V_m$ au repos.
\item Expliquer l'origine de la période réfractaire absolue et relative à partir de ces courbes.
\par \textbf{Absolue ($t \approx 1,5$ à $3,5\,\text{ms}$) :} Correspond à la période où $g_{\text{Na}}$ est non-nulle (canaux ouverts) puis \emph{inactivés} (canaux fermés mais non activables). Impossible d'ouvrir de nouveaux canaux $\text{Na}^+$.
\par \textbf{Relative ($t \approx 3,5$ à $5,0\,\text{ms}$) :} Correspond à l'hyperpolarisation ($g_{\text{K}}$ encore élevée). Les canaux $\text{Na}^+$ sont revenus à l'état \emph{fermés activables} (récupération de l'inactivation), mais $V_m$ est éloigné du seuil (proche de $E_{\text{K}}$). Un stimulus \emph{plus fort} que le seuil habituel peut déclencher un PA.
\end{enumerate}
\end{exercice}

\begin{exercice}[Courbe Intensité-Durée : Rheobase et Chronaxie]
On stimule une fibre nerveuse isolée par des impulsions carrées de durée variable $d$ et on mesure l'intensité minimale $I$ (en $\mu\text{A}$) nécessaire pour déclencher un potentiel d'action. Résultats :

\begin{center}
\begin{tabular}{|c|c|c|c|c|c|c|}\hline
Durée $d\,(\text{ms})$ & 0,1 & 0,2 & 0,5 & 1,0 & 2,0 & 5,0 \\ \hline
Intensité seuil $I\,(\mu\text{A})$ & 50 & 30 & 15 & 10 & 7,5 & 6 \\ \hline
\end{tabular}
\end{center}

\begin{enumerate}
\item Tracer la courbe $I = f(d)$ en axes linéaires (papier millimétré). Indiquer les unités.
\par \textbf{Consigne de tracé :} Abscisses : $d\,(\text{ms})$ de 0 à 6. Ordonnées : $I\,(\mu\text{A})$ de 0 à 55. Points : $(0,1; 50)$, $(0,2; 30)$, $(0,5; 15)$, $(1; 10)$, $(2; 7,5)$, $(5; 6)$. Courbe hyperbolique décroissante, asymptote horizontale vers $6\,\mu\text{A}$.
\item Déterminer graphiquement la \textbf{rheobase} ($I_{\infty}$) et le \textbf{temps utile} (chronaxie, $t_u$). Rappel : $t_u$ est la durée nécessaire pour déclencher une réponse avec $I = 2 \times \text{rheobase}$.
\par \textbf{Rheobase :} Asymptote pour $d \to \infty$. Lecture graphique : $I_{\infty} \approx \mathbf{6\,\mu\text{A}}$.
\par \textbf{Chronaxie :} $I = 2 \times 6 = 12\,\mu\text{A}$. Sur la courbe, $I=12\,\mu\text{A}$ correspond à $d \approx \mathbf{0,75\,\text{ms}}$ (entre $0,5$ et $1,0\,\text{ms}$).
\item Écrire la relation mathématique reliant $I$, $d$, $I_{\infty}$ et $t_u$ (équation de Lapicque). Vérifier si les points expérimentaux la valident.
\par \textbf{Équation :} $I = I_{\infty} \left(1 + \frac{t_u}{d}\right)$.
\par \textbf{Vérification (exemple pour $d=1\,\text{ms}$) :} $I_{\text{théo}} = 6 \left(1 + \frac{0,75}{1}\right) = 6 \times 1,75 = 10,5\,\mu\text{A}$. Expérimental : $10\,\mu\text{A}$. Bon accord. Pour $d=0,5\,\text{ms}$ : $I_{\text{théo}} = 6(1+1,5) = 15\,\mu\text{A}$. Expérimental : $15\,\mu\text{A}$. La loi de Lapicque décrit bien les données.
\item On compare cette fibre à une autre fibre dont la chronaxie est deux fois plus grande ($t_u' = 1,5\,\text{ms}$). Laquelle est la plus excitable pour des stimuli brefs ? Pour des stimuli longs ? Justifier.
\par \textbf{Stimuli longs ($d \gg t_u$) :} $I \approx I_{\infty}$. Les deux fibres ont même rheobase ($\approx 6\,\mu\text{A}$) $\to$ \textbf{même excitabilité}.
\par \textbf{Stimuli brefs ($d \ll t_u$) :} $I \approx I_{\infty} \frac{t_u}{d}$. L'intensité seuil est proportionnelle à $t_u$. La fibre avec $t_u' = 2 t_u$ nécessite une intensité \textbf{deux fois plus grande} $\to$ \textbf{moins excitable} pour les stimuli brefs.
\end{enumerate}
\end{exercice}

\begin{exercice}[Réponse d'une fibre isolée vs réponse d'un nerf entier : Tout-ou-rien et Recrutement]
On stimule électriquement : (A) une fibre nerveuse unique isolée ; (B) un nerf mixte (contenant des centaines de fibres de diamètres variés). On enregistre la réponse (amplitude du potentiel d'action composé pour le nerf) en fonction de l'intensité du stimulus $I$.

\begin{popfigure}
\begin{tikzpicture}[scale=0.7, font=\small]
% Graph A: Single Fiber
\begin{scope}[xshift=0cm]
\draw[->] (0,0) -- (5,0) node[right] {$I\,(\text{mA})$};
\draw[->] (0,0) -- (0,4) node[above] {Amplitude};
\draw[popdark, very thick] (0,0) -- (1.5,0) node[midway, below] {0};
\draw[popdark, very thick] (1.5,0) -- (1.5,3);
\draw[popdark, very thick] (1.5,3) -- (5,3);
\draw[dashed] (1.5,0) -- (1.5,-0.3) node[below] {$I_s$};
\node at (2.5, 3.5) {\textbf{(A) Fibre unique}};
\end{scope}
% Graph B: Whole Nerve
\begin{scope}[xshift=7cm]
\draw[->] (0,0) -- (5,0) -- (5,0) node[right] {$I\,(\text{mA})$};
\draw[->] (0,0) -- (0,4) node[above] {Amplitude};
\draw[popblue, very thick] (0,0) .. controls (1,0.5) and (2,2) .. (3,3.5) .. controls (4,3.8) .. (5,3.9);
\draw[dashed] (0.5,0) -- (0.5,-0.3) node[below] {$I_{s1}$};
\draw[dashed] (2.5,0) -- (2.5,-0.3) node[below] {$I_{s2}$};
\node at (2.5, 3.5) {\textbf{(B) Nerf entier}};
\end{scope}
\end{tikzpicture}
\legende{Relation stimulus-réponse pour une fibre unique (A) et un nerf entier (B).}
\end{popfigure}

\begin{enumerate}
\item Décrire et expliquer la forme de la courbe (A). Quel principe fondamental illustre-t-elle ?
\par \textbf{Description :} Amplitude nulle pour $I < I_s$ ; amplitude maximale et constante pour $I \ge I_s$.
\par \textbf{Explication :} Le potentiel d'action est un phénomène \textbf{tout-ou-rien} : s'il est déclenché, son amplitude est indépendante de l'intensité du stimulus (dépend seulement des gradients ioniques et propriétés membranaires).
\par \textbf{Principe :} \textbf{Loi du tout-ou-rien} (à l'échelle de la fibre unique).
\item Décrire et expliquer la forme de la courbe (B). Pourquoi observe-t-on une augmentation progressive de l'amplitude ? Définir le terme \textbf{recrutement}.
\par \textbf{Description :} Courbe sigmoïde croissante. Seuil $I_{s1}$ pour les premières fibres, amplitude maximale (plateau) pour $I$ fort (toutes fibres excitées).
\par \textbf{Explication :} Un nerf contient des fibres de diamètres et seuils différents. L'augmentation de $I$ excite progressivement plus de fibres.
\par \textbf{Définition :} Le \textbf{recrutement} est l'augmentation du nombre de fibres excitables (et donc de l'amplitude du potentiel d'action composé) lorsque l'intensité du stimulus augmente.
\item Sur la courbe (B), les seuils $I_{s1}$ et $I_{s2}$ correspondent à l'excitation de quelle population de fibres (gros ou petit diamètre, myélinisées ou non) ? Justifier par la relation diamètre/seuil.
\par \textbf{Réponse :} $I_{s1}$ (seuil le plus bas) $\to$ \textbf{grosses fibres myélinisées} (type A$\alpha$, A$\beta$). $I_{s2}$ (seuil plus haut) $\to$ \textbf{fibres de plus petit diamètre} (myélinisées fines A$\gamma$, A$\delta$ ou amyélinisées C).
\par \textbf{Justification :} La résistance axoplasmique $R_a \propto 1/d^2$. Les grosses fibres ont une résistance interne plus faible, le courant local se propage mieux, la dépolarisation au nœud adjacent est plus forte $\to$ \textbf{seuil d'excitation plus bas} (plus excitables).
\item Si l'on bloque les canaux $\text{Na}^+$ voltage-dépendants par la tétrodotoxine (TTX) sur le nerf entier, que devient la courbe (B) ?
\par \textbf{Réponse :} La courbe (B) devient une \textbf{ligne plate à amplitude nulle} (pas de potentiel d'action composé possible, aucune fibre ne peut générer de PA).
\end{enumerate}
\end{exercice}

\courssub{Je me prépare au Bac}

\begin{exercice}[Sujet type Bac : Mise en évidence et interprétation ionique du potentiel de repos (Axone de calmar)]
\textbf{Contexte :} L'axone géant de calmar (\emph{Dosidicus gigas}) est un modèle historique en neurophysiologie (diamètre $\approx 500\,\mu\text{m}$). On utilise un montage à deux électrodes intracellulaires (microélectrodes en verre remplies de $\text{KCl}\,3\,\text{M}$) reliées à un oscilloscope à double trace ou un voltmètre différentiel.

\begin{popfigure}
\begin{tikzpicture}[scale=0.7, font=\small]
% Axon
\draw[fill=popcream] (0,1) rectangle (10,3);
\draw[popdark, thick] (0,2) -- (10,2); % axis
% Electrodes
\draw[popblue, thick] (2,2) -- (2,5) node[above] {$\text{E}_1$ (Référence)};
\draw[poppink, thick] (7,2) -- (7,5) node[above] {$\text{E}_2$ (Mesure)};
% Voltmeter
\draw[popdark, thick] (5,5) rectangle (7,6.5);
\node at (6, 5.75) {Voltmètre différentiel / Oscilloscope};
\draw[popblue] (2,5) -- (5,5);
\draw[poppink] (7,5) -- (7,6.5);
% Ground
\draw[popmuted, thick] (5,1) -- (5,0) node[below] {Bain extracellulaire (Référence)};
% Ion concentrations box
\node[draw, fill=popblueL, rounded corners, align=left] at (12, 2) {
\textbf{Tableau des concentrations (mmol/L)} \\
$[\text{Na}^+]_i = 50$ ; $[\text{Na}^+]_o = 440$ \\
$[\text{K}^+]_i = 400$ ; $[\text{K}^+]_o = 20$ \\
$[\text{Cl}^-]_i = 40$ ; $[\text{Cl}^-]_o = 560$ \\
Température $20^\circ\text{C}$
};
\end{tikzpicture}
\legende{Montage expérimental pour l'enregistrement du potentiel de repos sur axone géant de calmar.}
\end{popfigure}

\begin{enumerate}
\item \textbf{Montage et observation :} Schématiser le circuit électrique équivalent de ce montage (membrane = capacitance $C_m$ en parallèle avec résistances ioniques $R_{\text{Na}}, R_{\text{K}}, R_{\text{Cl}}$ ; électrodes = résistances $R_e$ ; bain = résistance $R_s$). Quelle est la polarité du signal affiché si $\text{E}_2$ est à l'intérieur et $\text{E}_1$ dans le bain ? Quelle valeur tensionnelle observe-t-on au repos ?
\par \textbf{Circuit équivalent :} 
\begin{center}
\begin{tikzpicture}[scale=0.8, font=\small,
  resbox/.style={draw, rectangle, minimum width=1.4cm, minimum height=0.6cm}]
\draw (0,0) -- (0.6,0) node[resbox, right=0pt] (Re1) {$R_{e1}$};
\node[above] at (2,0.4) {E1 (Bain)};
\draw (Re1.east) -- (2.6,0) node[resbox, right=0pt] (Rs) {$R_s$};
\draw (Rs.east) -- (4.2,0) -- (4.2,-0.6);
\node[resbox] (Rk) at (4.2,-1.1) {$R_{\text{K}}$};
\node[resbox] (Rna) at (4.2,-1.7) {$R_{\text{Na}}$};
\node[resbox] (Rcl) at (4.2,-2.3) {$R_{\text{Cl}}$};
\node[resbox] (Cm) at (5.4,-1.7) {$C_m$};
\draw (4.2,-0.6) -- (5.4,-0.6) -- (5.4,-1.4);
\draw (4.2,-2.6) -- (5.4,-2.6) -- (5.4,-2.0);
\draw (Rk.south) -- (4.2,-2.6);
\draw (4.2,-2.6) -- (2.6,-2.6) node[resbox, left=0pt] (Re2) {$R_{e2}$};
\node[below] at (2,-3.0) {E2 (Intra)};
\draw (Re2.west) -- (0,-2.6);
\draw (0,0) -- (0,-2.6);
\node at (3, 1) {Voltmètre différentiel};
\end{tikzpicture}
\end{center}
\par \textbf{Polarité :} $\text{E}_2$ (intracellulaire) est négative par rapport à $\text{E}_1$ (extracellulaire/bain). Le voltmètre différentiel ($V = V_{\text{E2}} - V_{\text{E1}}$) affiche une tension \textbf{négative}.
\par \textbf{Valeur observée :} $V_r \approx \mathbf{-70\,\text{mV}}$ (intérieur négatif).
\item \textbf{Interprétation ionique :} À l'aide des concentrations données et de l'équation de Goldman-Hodgkin-Katz (forme simplifiée : $V_m \approx 58 \log_{10} \frac{P_{\text{K}}[\text{K}^+]_o + P_{\text{Na}}[\text{Na}^+]_o + P_{\text{Cl}}[\text{Cl}^-]_i}{P_{\text{K}}[\text{K}^+]_i + P_{\text{Na}}[\text{Na}^+]_i + P_{\text{Cl}}[\text{Cl}^-]_o}$), montrer que le potentiel de repos mesuré ($-70\,\text{mV}$) implique $P_{\text{K}} \gg P_{\text{Na}}$ au repos. Calculer $E_{\text{K}}$ et $E_{\text{Na}}$ et comparer à $V_r$.
\par \textbf{Calculs Nernst ($20^\circ\text{C}$) :}
$E_{\text{K}} = 58 \log_{10}(20/400) = 58 \log_{10}(0,05) \approx \mathbf{-75,5\,mV}$.
$E_{\text{Na}} = 58 \log_{10}(440/50) = 58 \log_{10}(8,8) \approx \mathbf{+54,8\,mV}$.
$E_{\text{Cl}} = 58 \log_{10}(40/560) \approx \mathbf{-66,5\,mV}$.
\par \textbf{Comparaison :} $V_r (-70\,\text{mV})$ est très proche de $E_{\text{K}}$ et $E_{\text{Cl}}$, très loin de $E_{\text{Na}}$.
\par \textbf{Démonstration GHK :} Si $P_{\text{Na}}$ était comparable à $P_{\text{K}}$, le terme $P_{\text{Na}}[\text{Na}^+]_o$ (numérateur) serait énorme ($P_{\text{Na}} \times 440$) et tirerait $V_m$ vers des valeurs positives. Pour que $V_m \approx -70\,\text{mV}$, il faut que le numérateur soit dominé par $P_{\text{K}}[\text{K}^+]_o$ et le dénominateur par $P_{\text{K}}[\text{K}^+]_i$. Donc \textbf{$P_{\text{K}} \gg P_{\text{Na}}$} (typiquement $P_{\text{K}} / P_{\text{Na}} \approx 20 \text{ à } 100$ au repos).
\item \textbf{Rôle de la pompe :} Expliquer pourquoi l'inhibition de la $\text{Na}^+/\text{K}^+$-ATPase (par l'ouabaïne ou l'absence d'ATP) entraîne une disparition progressive du potentiel de repos, alors que la pompe est électrogène (3 $\text{Na}^+$ sortent, 2 $\text{K}^+$ entrent).
\par \textbf{Explication :} La pompe est \emph{indirectement} responsable du potentiel de repos. Elle maintient les \textbf{gradients de concentration} ($[\text{K}^+]_i$ haut, $[\text{Na}^+]_i$ bas) qui sont la \textbf{source d'énergie} (force électrochimique) de la diffusion passive des ions à travers les canaux. Sans pompe, les fuites passives ($\text{K}^+$ sort, $\text{Na}^+$ entre) effacent les gradients $\to$ $E_{\text{K}}$ et $E_{\text{Na}}$ tendent vers 0 $\to$ $V_r$ tend vers 0. La contribution électrogène directe (3 $\text{Na}^+$ out / 2 $\text{K}^+$ in $\approx -2$ à $-5\,\text{mV}$) est mineure par rapport à l'effet des gradients.
\item \textbf{Application clinique :} L'hyperkaliémie (augmentation de $[\text{K}^+]_o$) dépolarise les membranes cardiaques et nerveuses. En utilisant l'équation de Nernst pour $\text{K}^+$, expliquer le danger d'arythmies cardiaques et de paralysie musculaire.
\par \textbf{Mécanisme :} $\uparrow [\text{K}^+]_o \implies \uparrow \frac{[\text{K}^+]_o}{[\text{K}^+]_i} \implies E_{\text{K}}$ \textbf{moins négatif} (ex: $-90 \to -60\,\text{mV}$). $V_r$ suit $E_{\text{K}}$ (car $P_{\text{K}}$ dominant) $\implies$ \textbf{dépolarisation} de la membrane au repos.
\par \textbf{Conséquence :} Cette dépolarisation rapproche $V_r$ du seuil d'inactivation des canaux $\text{Na}^+$ voltage-dépendants ($h_\infty$ chute). Les canaux $\text{Na}^+$ passent à l'état \emph{inactivé} (non activables). $\to$ \textbf{Inexcitabilité} : bloc de la conduction dans les nerfs (paralysie), bloc de la propagation dans le myocarde (arythmies, fibrillation ventriculaire, arrêt cardiaque).
\end{enumerate}
\end{exercice}

\begin{exercice}[Sujet type Bac : Codage de l'intensité du stimulus au niveau d'un récepteur sensoriel (Corpuscule de Pacini)]
\textbf{Contexte :} On étudie un corpuscule de Pacini (mécanorécepteur de pression profonde/vibration) isolé. On applique des pressions croissantes ($P_1 < P_2 < P_3$) sur sa capsule conjonctive et on enregistre le potentiel de récepteur ($V_{\text{rec}}$) et les potentiels d'action (PA) émis par la fibre afferente attachée.

\begin{popfigure}
\begin{tikzpicture}[scale=0.65, font=\small]
% Pressure stimulus
\draw[->] (0,0) -- (14,0) node[right] {$t\,(\text{ms})$};
\draw[->] (0,-3) -- (0,-1) node[left] {Pression};
\draw[poporange, thick] (1,-2.5) -- (3,-2.5) -- (3,-1.5) -- (5,-1.5) -- (5,-2.5) -- (7,-2.5); % P1
\draw[poporange, thick] (8,-2.5) -- (10,-2.5) -- (10,-1.0) -- (12,-1.0) -- (12,-2.5) -- (14,-2.5); % P2
% Vrec P1
\draw[->] (0,-5) -- (0,-3.5) node[left] {$V_{\text{rec}}\,(\text{mV})$};
\draw[popdark, thick, domain=1:7, samples=100] plot (\x, {-4.5 + 1.5*exp(-(\x-3)^2/1)});
\node at (0.5, -3.8) {$P_1$};
% AP P1
\draw[->] (0,-7) -- (0,-5.5) node[left] {PA};
\foreach \t in {2.5, 3.2, 4.0} \draw[poppink, very thick] (\t, -6.5) -- (\t, -5.8);
% Vrec P2
\draw[->] (7,-5) -- (7,-3.5) node[left] {$V_{\text{rec}}\,(\text{mV})$};
\draw[popdark, thick, domain=8:14, samples=100] plot (\x, {-4.5 + 3.0*exp(-(\x-10)^2/1)});
\node at (7.5, -3.8) {$P_2$};
% AP P2
\draw[->] (7,-7) -- (7,-5.5) node[left] {PA};
\foreach \t in {9.0, 9.5, 10.0, 10.5, 11.0, 11.5} \draw[poppink, very thick] (\t, -6.5) -- (\t, -5.8);
% Labels
\node at (-1.5, -2.5) {Stimulus};
\end{tikzpicture}
\legende{Enregistrements intracellulaires (potentiel de récepteur) et extracellulaires (potentiels d'action) sur une fibre afferente de corpuscule de Pacini soumis à deux pressions $P_1$ et $P_2$ ($P_2 > P_1$).}
\end{popfigure}

\begin{enumerate}
\item \textbf{Nature du potentiel de récepteur :} Décrire les caractéristiques du potentiel de récepteur (amplitude, durée, propagation) en le comparant au potentiel d'action. Est-ce un potentiel « tout-ou-rien » ou gradué ?
\par \textbf{Caractéristiques :} 
\begin{itemize}
\item \textbf{Amplitude variable (graduée)} : proportionnelle à l'intensité du stimulus ($P_2 > P_1 \implies V_{\text{rec}}(P_2) > V_{\text{rec}}(P_1)$).
\item \textbf{Non propagé} : potentiel local (électrotonique), se dégrade avec la distance (constante de longueur $\lambda$), ne se régénère pas.
\item \textbf{Pas de seuil} : pas de tout-ou-rien.
\item \textbf{Durée} : suit l'enveloppe temporelle du stimulus (ici adaptation rapide).
\end{itemize}
\par \textbf{Conclusion :} C'est un potentiel \textbf{gradué} (analogique), contrairement au PA qui est \textbf{tout-ou-rien} (numérique, amplitude fixe, propagé sans décrément).
\item \textbf{Transduction et codage :} Expliquer le mécanisme de transduction (canaux mécanosensibles $\to$ entrée $\text{Na}^+/\text{Ca}^{2+}$ $\to$ dépolarisation). Comment l'intensité de la pression ($P_1$ vs $P_2$) est-elle codée dans le message nerveux émis par la fibre ? (Préciser : fréquence instantanée, nombre de PA, latence).
\par \textbf{Transduction :} Pression $\to$ déformation capsule $\to$ étirement membrane $\to$ ouverture canaux \textbf{mécanosensibles} (non voltage-dépendants, perméables $\text{Na}^+$, $\text{Ca}^{2+}$) $\to$ entrée cations $\to$ \textbf{dépolarisation} ($V_{\text{rec}}$).
\par \textbf{Codage (loi de la fréquence) :} 
\begin{itemize}
\item \textbf{Fréquence instantanée} plus élevée pour $P_2$ (intervalles inter-spikes plus courts).
\item \textbf{Nombre total de PA} plus grand pour $P_2$ sur la durée du stimulus.
\item \textbf{Latence} (délai 1er PA) plus courte pour $P_2$ (dépolarisation plus rapide jusqu'au seuil).
\end{itemize}
\item \textbf{Seuil et adaptation :} On observe que pour une pression maintenue constante, la fréquence des PA diminue rapidement puis s'arrête (adaptation rapide). Relier cette propriété à la structure du corpuscule (lamelles conjonctives) et au type d'information transmise (début/fin de stimulus, vibration).
\par \textbf{Structure :} Le corpuscule de Pacini est entouré de \textbf{lamelles conjonctives} (comme un oignon) remplies de fluide visqueux.
\par \textbf{Mécanisme adaptation :} Pression maintenue $\to$ le fluide se redistribue lentement $\to$ la tension mécanique sur la membrane du neurone \textbf{disparaît} $\to$ fermeture des canaux mécanosensibles $\to$ $V_{\text{rec}}$ redescend $\to$ arrêt des PA.
\par \textbf{Information codée :} Détecteur de \textbf{changements} (transitoires) : début/fin de contact, \textbf{vibrations} (haute fréquence, 200-300 Hz). Ne code pas la pression statique maintenue.
\item \textbf{Pathologie :} Une neuropathie périphérique (ex : diabète non équilibré au Cameroun) détruit les grosses fibres myélinisées (type A$\beta$, Pacini). Quels signes cliniques le patient présentera-t-il ? (Perte de quelle sensibilité ?)
\par \textbf{Signes cliniques :} Perte de la \textbf{sensibilité vibratoire} (diapason 128 Hz aux malléoles/genoux), perte de la \textbf{sensibilité de position} (proprioception) $\to$ \textbf{ataxie sensitive} (marche en talonnant, signe de Romberg positif), perte du toucher fin/discriminatif. Douleur et température (fibres C, A$\delta$) souvent conservées au début.
\end{enumerate}
\end{exercice}

\begin{exercice}[Sujet type Bac : Pharmacologie des canaux ioniques, Vitesse de propagation et Conduction saltatoire]
\textbf{Partie A : Effet de la Tétrodotoxine (TTX) et de la Lidoçaïne}
On enregistre le potentiel d'action sur un axone myélinisé de mammifère (noeud de Ranvier) avant et après application de deux substances :
\begin{itemize}
\item \textbf{Condition 1 (Témoin) :} PA normal (amplitude $\approx 100\,\text{mV}$, durée $\approx 1\,\text{ms}$).
\item \textbf{Condition 2 (TTX $10^{-8}\,\text{M}$ extracellulaire) :} Aucun PA ne peut être déclenché, même pour un stimulus très fort. Le potentiel de repos est inchangé.
\item \textbf{Condition 3 (Lidoçaïne $1\,\text{mM}$ intracellulaire) :} Aucun PA ne peut être déclenché. Le potentiel de repos est inchangé.
\end{itemize}

\begin{enumerate}
\item Identifier la cible moléculaire précise de la TTX et de la Lidoçaïne. Sont-elles compétitives ? Sont-elles réversibles ?
\par \textbf{TTX :} Cible = \textbf{site de liaison extracellulaire sur le pore} du canal $\text{Na}^+$ voltage-dépendant (bouche externe, domaine P entre S5-S6). Bloque physiquement le passage des ions. \textbf{Non compétitive} (ne compete pas avec $\text{Na}^+$ pour un site de liaison, bouche le pore). \textbf{Réversible} (lavage).
\par \textbf{Lidoçaïne :} Cible = \textbf{site de liaison intracellulaire} (récepteur des anesthésiques locaux, segments S6 des domaines III/IV). Bloque le canal de l'intérieur (blocage « à l'état ouvert/inactivé »). \textbf{Non compétitive}. \textbf{Réversible} (dissociation dépendante du voltage/état).
\item Expliquer pourquoi la TTX appliquée \emph{extracellulairement} bloque le PA, alors que la Lidoçaïne doit être \emph{intracellulaire} (ou diffuser à travers la membrane sous forme basique non ionisée) pour agir. Préciser le site de liaison (pore externe vs face interne du canal).
\par \textbf{TTX :} Molécule \textbf{chargée positivement} (guanidinium) $\to$ ne traverse pas la bicouche lipidique. Son site de liaison est sur la \textbf{face extracellulaire} du pore (bouche externe). Application extracellulaire = accès direct.
\par \textbf{Lidoçaïne :} Base faible (pKa $\approx 7,9$). Existe sous forme \textbf{non ionisée (base libre, B)} lipophile et forme \textbf{ionisée (BH$^+$)} hydrophile. Seule la forme B traverse la membrane. Une fois intracellulaire, elle s'équilibre en BH$^+$ qui se lie au \textbf{site intracellulaire} (face cytoplasmique du pore, accessible quand le canal est ouvert/inactivé). D'où la nécessité d'une application intracellulaire directe ou d'une diffusion trans-membranaire (forme B) pour agir.
\item \textbf{Intérêt médical :} La Lidoçaïne est un anesthésique local. Expliquer pourquoi elle bloque préférentiellement la conduction dans les fibres de petit diamètre (douleur, température, fibres C) avant les fibres de gros diamètre (toucher, motricité, fibres A$\alpha$), bien que toutes expriment des canaux $\text{Na}_{\text{v}}1.7/1.8$. (Indice : géométrie, surface membranaire, fréquence de stimulation).
\par \textbf{Explication (Blocage sélectif / Différentiel) :}
\begin{enumerate}
\item \textbf{Dépendance à l'utilisation (Use-dependence) :} La lidoçaïne se lie préférentiellement aux canaux \textbf{ouverts ou inactivés}. Les fibres nociceptives (C, A$\delta$) ont une \textbf{fréquence de décharge spontanée ou évoquée plus élevée} (douleur continue) $\to$ plus de temps passé en état ouvert/inactivé $\to$ blocage plus rapide et plus fort.
\item \textbf{Géométrie / Surface membranaire :} Les fibres de petit diamètre (C, amyélinisées) ont un \textbf{rapport surface/volume plus grand} et pas de myéline $\to$ plus de canaux $\text{Na}^+$ exposés par unité de longueur $\to$ capture plus efficace de la drogue diffusant depuis l'espace extracellulaire.
\item \textbf{Longueur du nœud / Sécurité factor :} Les fibres myélinisées (A$\alpha$) ont des nœuds espacés, un facteur de sécurité élevé et une conduction saltatoire rapide ; le blocage nécessite une occupation plus massive des canaux aux nœuds.
\end{enumerate}
\par \textbf{Résultat clinique :} Perte de la douleur (fibres C) et température (A$\delta$) \textbf{avant} la perte du toucher/pression (A$\beta$) et de la motricité (A$\alpha$).
\end{enumerate}

\textbf{Partie B : Vitesse de propagation et Myélinisation}
On mesure la vitesse de propagation $v$ sur deux types de fibres de diamètre identique ($d = 10\,\mu\text{m}$) à $20^\circ\text{C}$ :
\begin{itemize}
\item Fibre amyélinisée : $v_1 = 2\,\text{m}\cdot\text{s}^{-1}$.
\item Fibre myélinisée : $v_2 = 50\,\text{m}\cdot\text{s}^{-1}$.
\end{itemize}

\begin{enumerate}
\setcounter{enumi}{3}
\item Calculer le gain de vitesse apporté par la myéline. Rappeler les deux effets physiques majeurs de la gaine de myéline sur les paramètres électriques de la membrane (Capacité $C_m$ et Résistance $R_m$) qui expliquent ce gain.
\par \textbf{Gain :} $\frac{v_2}{v_1} = \frac{50}{2} = \mathbf{25}$ fois plus vite.
\par \textbf{Effets de la myéline (multicouches lipidiques) :}
\begin{enumerate}
\item \textbf{Diminution drastique de la capacité membranaire ($C_m \downarrow$) :} $C_m \propto \frac{1}{\text{épaisseur diélectrique}}$. La myéline augmente l'épaisseur $\times 100$ $\to$ $C_m$ divisée par $\sim 100$. $\implies$ Moins de charge nécessaire pour changer $V_m$ ($\Delta Q = C_m \Delta V$) $\to$ dépolarisation plus rapide.
\item \textbf{Augmentation massive de la résistance membranaire ($R_m \uparrow$) :} $R_m \propto \text{épaisseur}$. Fuites ioniques quasi nulles sous la myéline. $\implies$ Courant local se propage plus loin (constante de longueur $\lambda = \sqrt{\frac{R_m}{R_i}} \uparrow$) sans s'éteindre.
\end{enumerate}
\par \textbf{Conséquence combinée :} Constante de temps $\tau = R_m C_m$ diminue (charge plus rapide) et $\lambda$ augmente (portée plus longue) $\to$ propagation saltatoire ultra-rapide.
\item Schématiser le courant local (boucle de courant) au niveau d'un nœud de Ranvier actif (dépolarisé) et du nœud adjacent au repos. Expliquer pourquoi le potentiel d'action « saute » de nœud en nœud (conduction saltatoire) et pourquoi cela accélère la propagation.
\par \textbf{Schéma (description textuelle pour le dessin) :}
\begin{center}
\begin{tikzpicture}[scale=0.6, font=\small]
% Myelin sheath
\draw[fill=popblueL] (-3, 1) rectangle (-1, -1);
\draw[fill=popblueL] (1, 1) rectangle (3, -1);
\draw[fill=popblueL] (5, 1) rectangle (7, -1);
% Nodes
\draw[popdark, very thick] (-1, 1) -- (-1, -1) node[midway, right=2mm] {Nœud $n$ (Actif)};
\draw[popdark, very thick] (1, 1) -- (1, -1) node[midway, right=2mm] {Nœud $n+1$ (Repos)};
\draw[popdark, very thick] (5, 1) -- (5, -1);
% Axoplasm
\draw[poporange!50] (-3, 0) -- (7, 0);
% Current loops
\draw[->, poppink, thick] (-1, 0.5) .. controls (0, 1.5) .. (1, 0.5) node[midway, above] {$I_{\text{local}}$};
\draw[->, poppink, thick] (1, -0.5) .. controls (0, -1.5) .. (-1, -0.5);
\draw[->, popblue, thick] (-1, 0) -- (-2.5, 0) node[left] {$I_{\text{axoplasmique}}$};
\draw[->, popblue, thick] (1, 0) -- (2.5, 0);
% Charges
\node[poppink] at (-1, 0.8) {$+$}; \node[popblue] at (-1, -0.8) {$-$};
\node[popblue] at (1, 0.8) {$-$}; \node[poppink] at (1, -0.8) {$+$};
\node at (2, 2) {\textbf{Conduction saltatoire}};
\end{tikzpicture}
\end{center}
\par \textbf{Explication :} Au nœud $n$ (actif, dépolarisé $+30\,\text{mV}$), l'entrée massive de $\text{Na}^+$ crée un courant axial ($I_{\text{axo}}$) vers le nœud $n+1$ (au repos, $-70\,\text{mV}$). Ce courant traverse l'axoplasme (faible $R_i$ grâce au gros diamètre) et sort par la capacité/résistance du nœud $n+1$ (boucle de courant local $I_{\text{local}}$). Comme la myéline isole parfaitement l'internœud ($R_m \uparrow \uparrow, C_m \downarrow \downarrow$), \textbf{aucun courant ne fuit} latéralement. Tout le courant axial arrive intact au nœud suivant, le dépolarisant fortement jusqu'au seuil. Le PA « saute » ainsi d'un nœud à l'autre sans se propager continûment.
\par \textbf{Accélération :} Le temps de propagation est le temps de charge de la capacité du nœud suivant ($\tau_{\text{nœud}} = R_m^{\text{nœud}} C_m^{\text{nœud}}$), très court car surface nœudale petite. L'internœud est traversé quasi-instantanément par le courant axial (vitesse $\sim c/100$).
\item \textbf{Pathologie démyélinisante :} La sclérose en plaques (plus rare en zone tropicale mais existante) détruit la myéline du SNC. Prédire l'effet sur la vitesse de propagation, sur la fidélité de la conduction (bloc de conduction) et sur la fatigue musculaire. Relier à la sécurité factor.
\par \textbf{Effets de la démyélinisation :}
\begin{enumerate}
\item \textbf{Vitesse :} Chute drastique ($v \downarrow$). Perte de l'isolation $\to$ $C_m \uparrow$, $R_m \downarrow$ $\to$ courant local fuit, $\lambda \downarrow$, $\tau \uparrow$. Conduction redevient continue (lente) ou échoue.
\item \textbf{Fidélité / Bloc de conduction :} Si la démyélinisation est sévère, le courant axial qui arrive au nœud suivant est insuffisant pour atteindre le seuil (facteur de sécurité $< 1$) $\to$ \textbf{bloc de conduction} (le PA s'arrête). Phénomène dépendant de la fréquence (bloc à haute fréquence).
\item \textbf{Fatigue musculaire :} Lors d'une stimulation répétée, l'accumulation de $\text{K}^+$ extracellulaire et l'inactivation des canaux $\text{Na}^+$ réduisent davantage le facteur de sécurité déjà bas $\to$ \textbf{bloc de conduction par fatigue} (fatigabilité anormale, ex: ptosis, diplopie, faiblesse à l'effort en SEP).
\end{enumerate}
\par \textbf{Facteur de sécurité (Safety Factor) :} Rapport entre le courant disponible au nœud $n+1$ et le courant minimal nécessaire pour atteindre le seuil. Myéline $\uparrow$ Facteur de sécurité. Démyélinisation $\downarrow$ Facteur de sécurité.
\end{enumerate}
\end{exercice}

\begin{attention}[Pièges fréquents au Bac - Leçon 08]
\begin{itemize}
\item \textbf{Confusion Rheobase / Chronaxie :} Rheobase = \textbf{Intensité} (µA/mA). Chronaxie = \textbf{Durée} (ms/s). Ne jamais inverser.
\item \textbf{Signe du potentiel de repos :} Toujours \textbf{négatif} (intérieur négatif). $V_m = V_{\text{in}} - V_{\text{out}}$.
\item \textbf{Équation de Nernst pour $\text{Cl}^-$ :} $E_{\text{Cl}} = -58 \log \frac{[\text{Cl}^-]_o}{[\text{Cl}^-]_i}$ (ou $+58 \log \frac{[\text{Cl}^-]_i}{[\text{Cl}^-]_o}$). Ne pas oublier le signe moins pour les anions !
\item \textbf{Période réfractaire absolue vs relative :} Absolue = canaux $\text{Na}^+$ \textbf{inactivés} (impossible). Relative = canaux $\text{Na}^+$ \textbf{récupérés} mais $V_m$ hyperpolarisé (possible avec stimulus plus fort).
\item \textbf{Recrutement (nerf) vs Tout-ou-rien (fibre) :} Ne pas dire « le nerf obéit au tout-ou-rien ». Le nerf montre un \textbf{recrutement} (réponse graduée) car somme de fibres tout-ou-rien.
\item \textbf{Vitesse propagation :} Unités : $\text{m}\cdot\text{s}^{-1}$ (ou $\text{m/s}$). Distance en $\text{m}$, temps en $\text{s}$.
\item \textbf{TTX vs Lidoçaïne :} TTX = extracellulaire (pore externe). Lidoçaïne = intracellulaire (face interne, forme basique diffuse). TTX bloque \emph{tous} les canaux $\text{Na}_{\text{v}}$ (sauf $\text{Na}_{\text{v}}1.5$ cardiaque résistant, $\text{Na}_{\text{v}}1.8/1.9$ résistants). Lidoçaïne bloque préférentiellement canaux ouverts/inactivés (dépendance à l'usage).
\end{itemize}
\end{attention}

\begin{methode}[Méthode type Bac : Calculer la vitesse de propagation]
\begin{enumerate}
\item Identifier sur le tracé (oscilloscope) les deux pics de potentiel d'action enregistrés par les deux électrodes.
\item Mesurer le \textbf{décalage temporel} $\Delta t$ (en ms ou s) entre ces deux pics.
\item Connaître la \textbf{distance $d$} (en cm ou m) séparant les deux électrodes le long de l'axone.
\item Appliquer la formule : $v = \frac{d}{\Delta t}$.
\item Convertir les unités pour obtenir $v$ en $\text{m}\cdot\text{s}^{-1}$ ($1\,\text{cm} = 0,01\,\text{m}$ ; $1\,\text{ms} = 0,001\,\text{s}$).
\item \textbf{Vérification d'ordre de grandeur :} Amyélinisé $\sim 1-100\,\text{m/s}$ (selon diamètre). Myélinisé $\sim 10-120\,\text{m/s}$.
\end{enumerate}
\end{methode}

\begin{methode}[Méthode type Bac : Déterminer Rheobase et Chronaxie sur une courbe I(d)]
\begin{enumerate}
\item Tracer la courbe $I = f(d)$ en axes linéaires (ou papier millimétré).
\item \textbf{Rheobase ($I_{\infty}$) :} Prolonger la branche asymptotique horizontale (grands $d$) et lire l'ordonnée à l'infini (ou pour $d$ très grand). Unité : $\mu\text{A}$ ou $\text{mA}$.
\item \textbf{Chronaxie ($t_u$) :} Calculer $2 \times I_{\infty}$. Tracer l'horizontale $I = 2 I_{\infty}$. Lire l'abscisse $d$ correspondante sur la courbe. Cette abscisse est $t_u$. Unité : $\text{ms}$ ou $\text{s}$.
\item \textbf{Vérification (optionnel) :} Tester la loi de Lapicque $I = I_{\infty}(1 + t_u/d)$ sur 2-3 points expérimentaux.
\end{enumerate}
\end{methode}

\begin{savaistu}[Contexte camerounais \& Histoire des sciences]
\textbf{Alan Hodgkin \& Andrew Huxley (Prix Nobel 1963) :} Ont élucidé les mécanismes ioniques du PA sur l'\textbf{axone géant de calmar} (diamètre $\sim 500\,\mu\text{m}$, visible à l'œil nu, permet insertion d'électrodes). Leur modèle mathématique (équations différentielles) est la base de toute l'électrophysiologie moderne.
\textbf{Bernhard Katz (Prix Nobel 1970) :} Travaux sur la transmission synaptique et les potentiels de récepteur (corpuscules de Pacini, plaque motrice).
\textbf{Au Cameroun :} Les \textbf{neuropathies périphériques} sont fréquentes : diabète (HTA/Diabète = pandémie silencieuse), lèpre (atteinte fibres C/A$\delta$ $\to$ perte douleur/température $\to$ plantar ulcers), carences nutritionnelles (B1, B12), toxiques (alcool, médicaments neurotoxiques). La \textbf{sclérose en plaques} est rare mais diagnostiquée (IRM, potentiels évoqués). La \textbf{myasthénie} (auto-immune anti-récepteur ACh) affecte la plaque motrice, pas l'axone.
\textbf{Anesthésie locale :} La lidoçaïne (Xylocaïne\textsuperscript{\textregistered}) est le standard au bloc opératoire (CHU Yaoundé, Douala, hôpitaux de district). Comprendre sa sélectivité fibres C vs A$\alpha$ guide le choix de concentration/volume pour une analgésie sans bloc moteur complet (ex: péridurale accouchement).
\end{savaistu}

\newpage

\coursec{Corrigés des exercices}

\coursec{Exercices corrigés — Sensibilisation sur la nécessité de préserver le fonctionnement des neurones}

\begin{corrige}[Exercice 1 — QCM : Mécanismes ioniques]
\textbf{1. Réponse C.} Au repos, la membrane est bien plus perméable au \ce{K+} (canaux de fuite \ce{K+} ouverts) qu'au \ce{Na+} ; les gradients (${[\ce{K+}]}_{i}$ élevée, ${[\ce{Na+}]}_{o}$ élevée) sont maintenus par la pompe \ce{Na+}/\ce{K+}-ATPase. Le potentiel de repos s'établit donc près de $E_{\ce{K+}}$. A est faux (${[\ce{Na+}]}_{i}$ est \emph{faible}), B est faux (${[\ce{K+}]}_{o}$ est faible et la perméabilité au \ce{K+} est élevée), D est faux (la pompe fonctionne en permanence).

\textbf{2. Réponse C.} La dépolarisation résulte de l'ouverture massive et transitoire des canaux \ce{Na+} voltage-dépendants : entrée brutale de \ce{Na+} selon son gradient électrochimique, avec rétroaction positive. L'ouverture des canaux \ce{K+} (A) est responsable de la \emph{repolarisation} ; la pompe (D) n'intervient pas dans la cinétique du PA.

\textbf{3. Réponse B.} Pendant la période réfractaire absolue, les canaux \ce{Na+} sont inactivés : aucun stimulus, même supraliminaire, ne peut déclencher un nouveau PA. A décrit la période réfractaire \emph{relative} ; C décrit l'hyperpolarisation (qui coïncide avec la période relative) ; D décrit l'état de repos.

\textbf{4. Réponse B.} Sur les fibres myélinisées, la dépolarisation ne se régénère qu'aux nœuds de Ranvier (riches en canaux \ce{Na+}) : le PA « saute » de nœud en nœud, d'où une vitesse très supérieure à celle d'une fibre amyélinisée de même diamètre (C est donc faux). A décrit la conduction continue des fibres amyélinisées ; D est faux car c'est précisément aux nœuds que se produit la dépolarisation régénérative.

\begin{attention}[Points de vigilance]
Bien distinguer les \textbf{deux états de non-excitabilité} des canaux \ce{Na+} : \emph{fermé activable} (repos) et \emph{fermé inactivé} (période réfractaire absolue). Confondre ces états est l'erreur la plus fréquente.
\end{attention}
\end{corrige}

\begin{corrige}[Exercice 2 — Vrai ou Faux]
\textbf{1. Faux.} Au repos, $V_r \approx E_{\ce{K+}}$ donc la force électrochimique sur \ce{K+} est quasi nulle ; mais pour \ce{Na+}, $V_r$ est très éloigné de $E_{\ce{Na+}} \approx +55\,\text{mV}$ : la force motrice sur \ce{Na+} est forte et dirigée vers l'intérieur. La faible entrée permanente de \ce{Na+} est compensée par la pompe.

\textbf{2. Vrai.} Au seuil ($\approx -55\,\text{mV}$), le courant entrant \ce{Na+} égale puis dépasse les courants sortants (\ce{K+} et fuites) : la dépolarisation ouvre davantage de canaux \ce{Na+}, qui dépolarisent encore : c'est la rétroaction positive qui rend le PA auto-entretenu et tout-ou-rien.

\textbf{3. Faux.} Pour les fibres amyélinisées, $v \propto \sqrt{d}$ ; pour les fibres myélinisées, $v \propto d$ (approximativement). La myélinisation change donc la loi d'échelle, en plus d'augmenter fortement la vitesse à diamètre égal.

\textbf{4. Vrai.} Le potentiel de récepteur est local, gradué (amplitude proportionnelle à l'intensité du stimulus dans la plage de linéarité), sans seuil, non propagé sans décrément. Le passage au codage en fréquence de PA se fait au niveau de la zone de déclenchement (segment initial).

\textbf{5. Faux.} La rheobase est une \textbf{intensité} : l'intensité minimale efficace pour une durée de stimulation infinie (très longue). La durée minimale à intensité double de la rheobase est la \textbf{chronaxie} (temps utile).

\begin{attention}[Piège de rédaction]
Au Bac, une réponse « Vrai » ou « Faux » non justifiée ne rapporte aucun point : la justification (une phrase précise, avec les bons termes : gradient électrochimique, rétroaction positive, inactivation) est exigée.
\end{attention}
\end{corrige}

\begin{corrige}[Exercice 3 — Définitions et vocabulaire]
Les définitions données dans l'énoncé constituent le corrigé de référence. Points de rigueur exigés :
\begin{enumerate}
\item \textbf{Seuil d'excitation :} citer la valeur approchée ($\approx -55\,\text{mV}$) et le caractère \emph{autocatalytique} de l'ouverture des canaux \ce{Na+}.
\item \textbf{Période réfractaire absolue :} la cause moléculaire (inactivation des canaux \ce{Na+}) doit apparaître.
\item \textbf{Conduction saltatoire :} mentionner les nœuds de Ranvier, le rôle isolant de la myéline et le double avantage (vitesse accrue, économie d'énergie car moins d'ions échangés à repomper).
\item \textbf{Potentiel de récepteur :} les trois adjectifs clés : \emph{gradué, local, non propagé}.
\item \textbf{Rheobase / Chronaxie :} rheobase = intensité ($\mu\text{A}$), chronaxie = durée (ms) mesurée à $2 \times$ rheobase.
\item \textbf{Codage de l'intensité :} l'amplitude du PA ne code rien (tout-ou-rien) ; seule la \textbf{fréquence} des PA code l'intensité, bornée par la période réfractaire absolue ($f_{\max} \approx 1/\text{PRA}$, soit quelques centaines de Hz).
\end{enumerate}
\end{corrige}

\begin{corrige}[Exercice 4 — Calculs directs : vitesse et Nernst]
\textbf{1.} $v = \dfrac{d}{\Delta t} = \dfrac{0,10\,\text{m}}{2,5\times 10^{-3}\,\text{s}} = \dfrac{0,10}{0,0025} = \mathbf{40\,\text{m}\cdot\text{s}^{-1}}$.
Vérification : $0,10/0,0025 = 100/2,5 = 40$. Correct.

\textbf{2.} La fibre myélinisée ($50\,\text{m/s}$) conduit plus vite que l'amyélinisée ($40\,\text{m/s}$) à diamètre égal. La myéline augmente $R_m$ et diminue $C_m$ : la constante de longueur $\lambda$ augmente (le courant local atteint le nœud suivant sans fuite) et la constante de temps diminue (charge capacitive plus rapide) : c'est la conduction saltatoire.

\textbf{3.} $E_{\ce{K+}} = 61,5 \log_{10}(5/150) = 61,5 \log_{10}(1/30)$.
$\log_{10}(30) = 1{,}477$, donc $E_{\ce{K+}} = 61,5 \times (-1{,}477) \approx \mathbf{-90,8\,\text{mV}}$.
Vérification : $61,5 \times 1{,}477 = 90,8$. Correct.
$V_r = -70\,\text{mV}$ est moins négatif que $E_{\ce{K+}}$ : la membrane au repos possède une perméabilité résiduelle au \ce{Na+} (et au \ce{Cl-}) qui « tire » $V_m$ vers $E_{\ce{Na+}} > 0$. Le potentiel de repos résulte d'une moyenne pondérée par les perméabilités (équation de GHK), non d'un équilibre de Nernst pur pour \ce{K+}.

\begin{methode}[Vérification rapide d'un calcul de Nernst]
Pour un cation : si $[X]_o > [X]_i$, alors $E_X > 0$ ; si $[X]_o < [X]_i$, alors $E_X < 0$. Ici ${[\ce{K+}]}_{o} \ll {[\ce{K+}]}_{i} \Rightarrow E_{\ce{K+}} < 0$ : cohérent.
\end{methode}
\end{corrige}

\begin{corrige}[Exercice 5 — Concentrations ioniques et potentiel de repos]
\textbf{1. Calculs des potentiels d'équilibre à $20^\circ\text{C}$ :}
\begin{align*}
E_{\ce{Na+}} &= 58 \log_{10}\left(\tfrac{440}{50}\right) = 58 \times \log_{10}(8{,}8) = 58 \times 0{,}944 \approx \mathbf{+54,8\,\text{mV}} \\
E_{\ce{K+}} &= 58 \log_{10}\left(\tfrac{20}{400}\right) = 58 \times \log_{10}(0{,}05) = 58 \times (-1{,}301) \approx \mathbf{-75,5\,\text{mV}} \\
E_{\ce{Cl-}} &= 58 \log_{10}\left(\tfrac{{[\ce{Cl-}]}_{i}}{{[\ce{Cl-}]}_{o}}\right) = 58 \log_{10}\left(\tfrac{40}{560}\right) = 58 \times (-1{,}146) \approx \mathbf{-66,5\,\text{mV}}
\end{align*}
Vérifications : $\log_{10}(8{,}8) = 0{,}944$ ; $\log_{10}(0{,}05) = -1{,}301$ ; $\log_{10}(0{,}0714) = -1{,}146$. Pour l'anion \ce{Cl-}, on inverse le rapport des concentrations (ou on change le signe) : piège classique évité ici.

\textbf{2. Comparaison et déduction :} $V_r = -70\,\text{mV}$ est proche de $E_{\ce{K+}}$ et de $E_{\ce{Cl-}}$, très éloigné de $E_{\ce{Na+}}$. Dans l'équation GHK, chaque ion « tire » $V_m$ vers son potentiel d'équilibre avec un poids égal à sa perméabilité. La proximité de $V_r$ avec $E_{\ce{K+}}$ impose donc $P_{\ce{K+}} \gg P_{\ce{Na+}}$ ; la proximité avec $E_{\ce{Cl-}}$ indique une perméabilité au \ce{Cl-} non négligeable. Le léger écart $V_r - E_{\ce{K+}} = +5,5\,\text{mV}$ traduit l'effet dépolarisant du terme $P_{\ce{Na+}} \cdot {[\ce{Na+}]}_{o}$ au numérateur.

\textbf{3. Hyperkaliémie expérimentale :} $E_{\ce{K+}}' = 58 \log_{10}(100/400) = 58 \log_{10}(0{,}25) = 58 \times (-0{,}602) \approx \mathbf{-34,9\,\text{mV}}$.
Vérification : $\log_{10}(0{,}25) = -0{,}602$ ; $58 \times 0{,}602 = 34{,}9$. Correct.
$V_r$ se dépolarise fortement vers $-35\,\text{mV}$. Conséquence : les canaux \ce{Na+} voltage-dépendants passent durablement à l'état \emph{inactivé} (l'inactivation dépend du voltage) : les cellules deviennent \textbf{inexcitables} malgré un potentiel proche du seuil. C'est le mécanisme des paralysies et des troubles du rythme cardiaque de l'hyperkaliémie sévère (urgence vitale).

\begin{attention}[Point de vigilance Bac]
Ne pas écrire « la dépolarisation rend la cellule plus excitable car plus proche du seuil » : c'est vrai seulement pour une dépolarisation \emph{modérée et transitoire}. Une dépolarisation \emph{maintenue} inactive les canaux \ce{Na+} et produit l'effet inverse (inexcitabilité). Cette nuance est souvent sanctionnée.
\end{attention}
\end{corrige}

\begin{corrige}[Exercice 6 — Courbes de conductances $g_{\ce{Na+}}$ et $g_{\ce{K+}}$]
\textbf{1. Phases :} dépolarisation : $t \approx 1{,}5$ à $2{,}0\,\text{ms}$ ; repolarisation : $t \approx 2{,}0$ à $3{,}5\,\text{ms}$ ; hyperpolarisation : $t \approx 3{,}5$ à $5{,}0\,\text{ms}$ (le $V_m$ passe sous $V_r$ car $g_{\ce{K+}}$ est encore élevée et rapproche $V_m$ de $E_{\ce{K+}}$).

\textbf{2. $g_{\ce{Na+}}$ et dépolarisation :} la dépolarisation initiale ouvre des canaux \ce{Na+} voltage-dépendants ; l'entrée de \ce{Na+} dépolarise davantage, ce qui ouvre plus de canaux : \textbf{rétroaction positive} (emballement), d'où la montée quasi verticale de $g_{\ce{Na+}}$ et de $V_m$. La chute brutale de $g_{\ce{Na+}}$ après le pic est due à l'\textbf{inactivation} des canaux \ce{Na+} (mécanisme de verrouillage rapide, dépendant du temps et du voltage), relayée par la désactivation lors de la repolarisation.

\textbf{3. $g_{\ce{K+}}$ et repolarisation :} les canaux \ce{K+} voltage-dépendants s'ouvrent avec un \textbf{retard} (d'où leur nom de « courant retardé ») : le courant sortant de \ce{K+} ramène $V_m$ vers $E_{\ce{K+}}$ (repolarisation). Comme leur fermeture est lente, $g_{\ce{K+}}$ reste élevée après le retour à $V_r$ : $V_m$ continue vers $E_{\ce{K+}}$, d'où l'\textbf{hyperpolarisation}. L'absence de mécanisme d'inactivation rapide des canaux \ce{K+} explique le retour lent de $g_{\ce{K+}}$ à sa valeur de repos.

\textbf{4. Périodes réfractaires :}
\begin{itemize}
\item \textbf{Absolue} ($\approx 1{,}5$ à $3{,}5\,\text{ms}$) : les canaux \ce{Na+} sont ouverts puis inactivés ; aucun nouveau PA n'est possible, quelle que soit l'intensité du stimulus.
\item \textbf{Relative} ($\approx 3{,}5$ à $5{,}0\,\text{ms}$) : les canaux \ce{Na+} ont récupéré (état fermé activable) mais $V_m$ est hyperpolarisé et $g_{\ce{K+}}$ encore élevée : un PA est possible uniquement pour un stimulus \emph{supraliminaire}.
\end{itemize}
Conséquence fonctionnelle majeure : la période réfractaire impose le \textbf{sens unidirectionnel} de la propagation et borne la fréquence maximale des PA.
\end{corrige}

\begin{corrige}[Exercice 7 — Courbe Intensité-Durée]
\textbf{1. Tracé :} axes linéaires ; abscisse $d$ de 0 à $6\,\text{ms}$ ; ordonnée $I$ de 0 à $55\,\mu\text{A}$. Points : $(0{,}1\,;\,50)$, $(0{,}2\,;\,30)$, $(0{,}5\,;\,15)$, $(1\,;\,10)$, $(2\,;\,7{,}5)$, $(5\,;\,6)$. Courbe décroissante de type hyperbolique, asymptote horizontale vers $6\,\mu\text{A}$.

\textbf{2. Déterminations graphiques :}
\begin{itemize}
\item \textbf{Rheobase} : ordonnée de l'asymptote : $I_{\infty} \approx \mathbf{6\,\mu\text{A}}$.
\item \textbf{Chronaxie} : abscisse du point d'ordonnée $2 I_{\infty} = 12\,\mu\text{A}$ : lecture entre $0{,}5$ et $1{,}0\,\text{ms}$, soit $t_u \approx \mathbf{0{,}75\,\text{ms}}$.
\end{itemize}

\textbf{3. Loi de Lapicque :} $I = I_{\infty}\left(1 + \dfrac{t_u}{d}\right)$.
Vérifications :
\begin{itemize}
\item $d = 1\,\text{ms}$ : $I_{\text{théo}} = 6(1 + 0{,}75) = 10{,}5\,\mu\text{A}$ ; expérimental $10\,\mu\text{A}$ : bon accord (écart $5\%$).
\item $d = 0{,}5\,\text{ms}$ : $I_{\text{théo}} = 6(1 + 1{,}5) = 15\,\mu\text{A}$ ; expérimental $15\,\mu\text{A}$ : accord exact.
\item $d = 2\,\text{ms}$ : $I_{\text{théo}} = 6(1 + 0{,}375) = 8{,}25\,\mu\text{A}$ ; expérimental $7{,}5\,\mu\text{A}$ : accord correct.
\end{itemize}
La loi de Lapicque décrit bien les données.

\textbf{4. Comparaison d'excitabilité :} pour les stimuli \textbf{longs} ($d \gg t_u$), $I \approx I_{\infty}$ : les deux fibres ont la même rheobase, donc la même excitabilité. Pour les stimuli \textbf{brefs} ($d \ll t_u$), $I \approx I_{\infty}\, t_u / d$ : l'intensité seuil est proportionnelle à la chronaxie ; la fibre de chronaxie double exige une intensité double : elle est \textbf{moins excitable} pour les stimuli brefs. La chronaxie mesure donc l'excitabilité « rapide » de la fibre.

\begin{methode}[Lecture graphique sans erreur]
Toujours tracer l'horizontale $I = 2 I_{\infty}$ \emph{après} avoir déterminé la rheobase, puis projeter l'intersection avec la courbe sur l'axe des durées. Ne jamais confondre l'unité de la rheobase ($\mu\text{A}$) et celle de la chronaxie (ms).
\end{methode}
\end{corrige}

\begin{corrige}[Exercice 8 — Fibre isolée vs nerf entier]
\textbf{1. Courbe (A) :} réponse nulle sous le seuil $I_s$, puis amplitude maximale constante dès le seuil atteint. Le PA d'une fibre unique est \textbf{tout-ou-rien} : déclenché ou non ; s'il l'est, son amplitude est fixée par les gradients ioniques ($E_{\ce{Na+}}$, $E_{\ce{K+}}$) et non par l'intensité du stimulus. C'est la \textbf{loi du tout-ou-rien}.

\textbf{2. Courbe (B) :} réponse graduée (sigmoïde) : l'amplitude du \emph{potentiel d'action composé} croît avec $I$ jusqu'à un plateau. Explication : le nerf est un faisceau de fibres de seuils différents ; l'augmentation de l'intensité excite un nombre croissant de fibres : c'est le \textbf{recrutement spatial}. Le plateau correspond à l'excitation de toutes les fibres du nerf.

\textbf{3. Populations de fibres :} $I_{s1}$ (seuil bas) correspond aux \textbf{grosses fibres myélinisées} (A$\alpha$, A$\beta$) : leur faible résistance axoplasmique ($R_a \propto 1/d^2$) et leur conduction saltatoire leur confèrent le seuil d'excitation le plus bas. $I_{s2}$ correspond aux fibres de plus petit diamètre (A$\delta$, puis fibres C amyélinisées, de seuil encore plus élevé). Conséquence clinique : lors d'une stimulation électrique croissante de la peau, le toucher (A$\beta$) est perçu avant la douleur (A$\delta$, C).

\textbf{4. Effet de la TTX :} la TTX bloque les canaux \ce{Na+} voltage-dépendants de toutes les fibres : aucun PA ne peut naître ; la courbe (B) devient une \textbf{droite horizontale d'amplitude nulle} quelle que soit l'intensité du stimulus (le potentiel de repos, lui, est conservé).

\begin{attention}[Piège rédactionnel]
Ne jamais écrire « le nerf obéit à la loi du tout-ou-rien » : la loi du tout-ou-rien s'applique à la \textbf{fibre} (et à la cellule) ; le \textbf{nerf} présente une réponse graduée par recrutement. Cette distinction est un classique des sujets de Bac.
\end{attention}
\end{corrige}

\begin{corrige}[Exercice 9 — Sujet type Bac : potentiel de repos (axone de calmar)]
\textbf{Barème indicatif (sur 10) :} Q1 : 2,5 pts ; Q2 : 3 pts ; Q3 : 2 pts ; Q4 : 2,5 pts.

\textbf{1. Montage et observation.} Le circuit équivalent comprend : les résistances d'électrodes $R_{e1}$ (bain) et $R_{e2}$ (intracellulaire), la résistance du bain $R_s$, puis la membrane modélisée par la capacité $C_m$ en parallèle avec les trois branches ioniques ($R_{\ce{K+}}$, $R_{\ce{Na+}}$, $R_{\ce{Cl-}}$, chacune en série avec la pile de Nernst correspondante). $\text{E}_2$ étant intracellulaire et $\text{E}_1$ dans le bain (référence), le voltmètre différentiel affiche $V = V_{\text{in}} - V_{\text{out}} < 0$ : une tension \textbf{négative}, $V_r \approx -70\,\text{mV}$.

\textbf{2. Interprétation ionique.} $E_{\ce{K+}} = 58 \log_{10}(20/400) \approx -75{,}5\,\text{mV}$ ; $E_{\ce{Na+}} = 58 \log_{10}(440/50) \approx +54{,}8\,\text{mV}$ ; $E_{\ce{Cl-}} \approx -66{,}5\,\text{mV}$. $V_r$ est proche de $E_{\ce{K+}}$ et $E_{\ce{Cl-}}$, très loin de $E_{\ce{Na+}}$. Dans l'équation GHK, si $P_{\ce{Na+}}$ était comparable à $P_{\ce{K+}}$, le terme $P_{\ce{Na+}} \cdot {[\ce{Na+}]}_{o}$ (avec ${[\ce{Na+}]}_{o} = 440$) tirerait fortement $V_m$ vers les valeurs positives. La valeur mesurée ($-70\,\text{mV}$) impose donc que $P_{\ce{K+}}$ domine : $P_{\ce{K+}} \gg P_{\ce{Na+}}$ (rapport typique de 20 à 100 au repos).

\textbf{3. Rôle de la pompe.} La pompe \ce{Na+}/\ce{K+}-ATPase maintient les gradients de concentration (elle expulse 3 \ce{Na+} et fait entrer 2 \ce{K+} par ATP hydrolysé). Sans elle, les fuites passives dissipent progressivement les gradients : $E_{\ce{K+}}$ et $E_{\ce{Na+}}$ tendent vers 0, donc $V_r$ s'effondre. Sa contribution électrogène directe (courant net de 1 charge positive sortante par cycle, quelques mV) est secondaire : son rôle essentiel est \emph{indirect} (maintenance des gradients, source d'énergie potentielle électrochimique).

\textbf{4. Hyperkaliémie.} $\uparrow {[\ce{K+}]}_{o} \Rightarrow$ rapport ${[\ce{K+}]}_{o}/{[\ce{K+}]}_{i}$ augmente $\Rightarrow E_{\ce{K+}}$ devient moins négatif $\Rightarrow V_r$ se dépolarise. La dépolarisation maintenue inactive les canaux \ce{Na+} voltage-dépendants : inexcitabilité des fibres nerveuses (paralysies) et des cellules myocardiques (troubles de conduction, arythmies pouvant aller jusqu'à l'arrêt cardiaque). D'où l'urgence thérapeutique (calcium intraveineux, insuline-glucose, etc.).

\begin{attention}[Points de vigilance]
Citer les unités à chaque étape ; écrire explicitement la convention de signe ($V_m = V_{\text{in}} - V_{\text{out}}$) ; pour le \ce{Cl-}, inverser le rapport des concentrations dans la formule de Nernst ; ne pas attribuer le potentiel de repos à la seule pompe (elle est nécessaire mais indirecte).
\end{attention}
\end{corrige}

\begin{corrige}[Exercice 10 — Sujet type Bac : corpuscule de Pacini]
\textbf{Barème indicatif (sur 10) :} Q1 : 2 pts ; Q2 : 3 pts ; Q3 : 2,5 pts ; Q4 : 2,5 pts.

\textbf{1. Nature du potentiel de récepteur.} Gradué (amplitude proportionnelle à l'intensité du stimulus : $V_{\text{rec}}(P_2) > V_{\text{rec}}(P_1)$), local et non propagé (décroît avec la distance, conduction électrotonique), sans seuil, de durée calquée sur le stimulus. Ce n'est \emph{pas} un phénomène tout-ou-rien, contrairement au PA (amplitude fixe, propagé sans décrément, fréquence variable).

\textbf{2. Transduction et codage.} La pression déforme les lamelles et la membrane terminale, ouvrant des canaux mécanosensibles perméables aux cations (\ce{Na+}, \ce{Ca^{2+}}) : entrée nette de charges positives $\Rightarrow$ dépolarisation graduée ($V_{\text{rec}}$). Si $V_{\text{rec}}$ atteint le seuil au niveau de la zone de déclenchement (premier nœud), des PA sont émis. L'intensité est codée par : la \textbf{fréquence instantanée} des PA (plus élevée pour $P_2$), le \textbf{nombre total} de PA et la \textbf{latence} (plus courte pour $P_2$). C'est la loi de la fréquence : conversion d'un codage analogique (amplitude de $V_{\text{rec}}$) en codage fréquentiel (PA tout-ou-rien).

\textbf{3. Adaptation.} Les lamelles conjonctives entourant la terminaison contiennent un fluide visqueux : sous pression maintenue, le fluide se redistribue et la déformation de la membrane disparaît ; les canaux mécanosensibles se referment, $V_{\text{rec}}$ retombe, les PA cessent. Le corpuscule de Pacini est donc un détecteur de \textbf{variations} : début et fin du stimulus, vibrations (jusqu'à 200-300 Hz), et non un capteur de pression statique.

\textbf{4. Neuropathie diabétique.} La destruction des grosses fibres myélinisées A$\beta$ entraîne : perte de la sensibilité vibratoire (test au diapason 128 Hz), du toucher fin et de la proprioception, avec ataxie sensitive (troubles de l'équilibre, signe de Romberg). Les sensibilités douloureuse et thermique (fibres A$\delta$ et C) sont relativement préservées au début, ce qui expose aux plaies ignorées (pied diabétique, complication majeure au Cameroun).

\begin{attention}[Points de vigilance]
Bien distinguer les deux codages : \emph{amplitude} pour le potentiel de récepteur, \emph{fréquence} pour les PA. Ne jamais écrire que « l'amplitude des PA augmente avec l'intensité du stimulus » : faute grave (violation de la loi du tout-ou-rien).
\end{attention}
\end{corrige}

\begin{corrige}[Exercice 11 — Sujet type Bac : pharmacologie des canaux et conduction saltatoire]
\textbf{Barème indicatif (sur 12) :} Partie A : Q1 : 2 pts ; Q2 : 2 pts ; Q3 : 3 pts. Partie B : Q4 : 2 pts ; Q5 : 2 pts ; Q6 : 1 pt.

\textbf{Partie A.}
\textbf{1.} TTX : se fixe sur la \textbf{bouche extracellulaire du pore} des canaux \ce{Na+} voltage-dépendants et obstrue physiquement le passage des ions ; blocage non compétitif, réversible par lavage. Lidoçaïne : se fixe sur un \textbf{site intracellulaire} du canal (accessible lorsque le canal est ouvert ou inactivé) ; blocage non compétitif, réversible, et renforcé par l'activité (use-dependence).

\textbf{2.} La TTX, chargée positivement, ne traverse pas la bicouche lipidique : elle n'agit que depuis l'extérieur, où se trouve son site. La lidoçaïne, base faible (pKa $\approx 7{,}9$), diffuse à travers la membrane sous sa forme non ionisée lipophile, puis se reprotone dans le cytoplasme et se lie à la face interne du canal : d'où son efficacité par voie intracellulaire ou après diffusion transmembranaire.

\textbf{3. Bloc différentiel.} Trois mécanismes concourent : (i) la lidoçaïne se lie préférentiellement aux canaux ouverts/inactivés : les fibres nociceptives, qui déchargent à haute fréquence, accumulent plus vite le blocage ; (ii) les petites fibres (C, A$\delta$) ont un rapport surface/volume élevé et pas de myéline : grande densité de canaux accessibles ; (iii) les grosses fibres myélinisées ont un facteur de sécurité élevé : il faut bloquer une fraction importante de leurs canaux nodaux pour interrompre la conduction. Résultat clinique : analgésie (perte de la douleur) obtenue avant le bloc du toucher et de la motricité — principe des anesthésies locales et péridurales pratiquées dans les hôpitaux camerounais.

\textbf{Partie B.}
\textbf{4. Gain de vitesse :} $v_2/v_1 = 50/2 = \mathbf{25}$. La myéline : (i) \textbf{diminue fortement $C_m$} (épaisseur diélectrique multipliée, $C_m \propto 1/e$) : moins de charges à déplacer pour dépolariser, réponse plus rapide ; (ii) \textbf{augmente fortement $R_m$} : les fuites latérales sont supprimées, la constante de longueur $\lambda = \sqrt{R_m/R_i}$ augmente, le courant local atteint le nœud suivant quasiment intact.

\textbf{5. Conduction saltatoire.} Au nœud actif $n$ ($V_m \approx +30\,\text{mV}$), l'entrée de \ce{Na+} crée un courant axial dans l'axoplasme vers le nœud $n+1$ au repos ($-70\,\text{mV}$) ; la boucle se referme par le milieu extracellulaire. Sous la myéline, aucun courant ne fuit : toute l'intensité sert à charger la capacité du nœud $n+1$, qui atteint le seuil et régénère le PA. Le PA « saute » ainsi de nœud en nœud : la vitesse est bien supérieure à la conduction continue, et l'économie est double (moins de canaux, moins d'ions à repomper).

\textbf{6. Démyélinisation (SEP).} Perte de myéline $\Rightarrow C_m \uparrow$, $R_m \downarrow$ $\Rightarrow$ le courant local fuit et s'éteint avant le nœud suivant : vitesse effondrée, puis \textbf{bloc de conduction} lorsque le facteur de sécurité (courant disponible / courant seuil) devient inférieur à 1. La fatigue et l'hyperthermie aggravent le bloc (inactivation des canaux \ce{Na+}), d'où la fatigabilité et la fluctuation des signes (troubles moteurs, sensitifs, visuels).

\begin{attention}[Points de vigilance]
Pour le calcul de gain, exprimer le rapport sans unité (25, « fois plus rapide »). Dans le schéma de conduction saltatoire, flécher le courant local \emph{dans les deux compartiments} (axoplasme et milieu extracellulaire) et situer les canaux \ce{Na+} uniquement aux nœuds. Ne pas écrire que « la myéline conduit le courant » : elle \emph{isole}, c'est l'axoplasme qui conduit.
\end{attention}
\end{corrige}

\newpage

\coursec{Fiche bilan}

\courssub{Fiche Bilan : Le fonctionnement des neurones et sa préservation}

\begin{popfigure}
\begin{center}\tcbox[colback=popink,colframe=popink,coltext=white,arc=9pt,boxsep=4pt,left=6pt,right=6pt]{\bfseries\small Fonctionnement des neurones}\end{center}
\vspace{-2pt}
\begin{tcbraster}[raster columns=3,raster equal height=rows,raster column skip=6pt,raster row skip=6pt,enhanced,blankest]
\begin{tcolorbox}[enhanced,colback=white,colframe=popgreen,boxrule=0.9pt,arc=3pt,title={1. Potentiel de repos},fonttitle=\bfseries\scriptsize,coltitle=white,colbacktitle=popgreen,left=4pt,right=4pt,top=2pt,bottom=2pt,boxsep=1pt,toptitle=1.5pt,bottomtitle=1.5pt,halign=left]
\begin{itemize}[nosep,leftmargin=*,itemsep=1pt,font=\scriptsize]
\item Pompe Na$^+$/K$^+$ (3/2)
\item Canaux K$^+$ fuyards
\item Intérieur négatif ($\approx -70$ mV)
\item Équilibre électrochimique
\end{itemize}
\end{tcolorbox}
\begin{tcolorbox}[enhanced,colback=white,colframe=poporange,boxrule=0.9pt,arc=3pt,title={2. Potentiel d'action},fonttitle=\bfseries\scriptsize,coltitle=white,colbacktitle=poporange,left=4pt,right=4pt,top=2pt,bottom=2pt,boxsep=1pt,toptitle=1.5pt,bottomtitle=1.5pt,halign=left]
\begin{itemize}[nosep,leftmargin=*,itemsep=1pt,font=\scriptsize]
\item Dépolarisation (Na$^+$ entrant)
\item Repolarisation (K$^+$ sortant)
\item Hyperpolarisation
\item Tout ou rien
\end{itemize}
\end{tcolorbox}
\begin{tcolorbox}[enhanced,colback=white,colframe=poppurple,boxrule=0.9pt,arc=3pt,title={3. Naissance du message},fonttitle=\bfseries\scriptsize,coltitle=white,colbacktitle=poppurple,left=4pt,right=4pt,top=2pt,bottom=2pt,boxsep=1pt,toptitle=1.5pt,bottomtitle=1.5pt,halign=left]
\begin{itemize}[nosep,leftmargin=*,itemsep=1pt,font=\scriptsize]
\item Récepteur : énergie $\to$ PG
\item Potentiel générateur (gradué)
\item Seuil d'excitation
\item Codage : fréquence des PA
\end{itemize}
\end{tcolorbox}
\begin{tcolorbox}[enhanced,colback=white,colframe=poppink,boxrule=0.9pt,arc=3pt,title={4. Propagation},fonttitle=\bfseries\scriptsize,coltitle=white,colbacktitle=poppink,left=4pt,right=4pt,top=2pt,bottom=2pt,boxsep=1pt,toptitle=1.5pt,bottomtitle=1.5pt,halign=left]
\begin{itemize}[nosep,leftmargin=*,itemsep=1pt,font=\scriptsize]
\item Sens unique (période réfractaire)
\item Mode saltatoire (myéline)
\item Vitesse : $v = d/\Delta t$
\item Facteurs : diamètre, T°, myéline
\end{itemize}
\end{tcolorbox}
\begin{tcolorbox}[enhanced,colback=white,colframe=popgold,boxrule=0.9pt,arc=3pt,title={5. Relation I-D},fonttitle=\bfseries\scriptsize,coltitle=white,colbacktitle=popgold,left=4pt,right=4pt,top=2pt,bottom=2pt,boxsep=1pt,toptitle=1.5pt,bottomtitle=1.5pt,halign=left]
\begin{itemize}[nosep,leftmargin=*,itemsep=1pt,font=\scriptsize]
\item Rhéobase ($I_{rh}$)
\item Chronaxie ($t_{ch}$)
\item Courbe hyperbolique
\item $I = I_{rh}(1 + t_{ch}/t)$
\end{itemize}
\end{tcolorbox}
\begin{tcolorbox}[enhanced,colback=white,colframe=popdark,boxrule=0.9pt,arc=3pt,title={6. Préservation},fonttitle=\bfseries\scriptsize,coltitle=white,colbacktitle=popdark,left=4pt,right=4pt,top=2pt,bottom=2pt,boxsep=1pt,toptitle=1.5pt,bottomtitle=1.5pt,halign=left]
\begin{itemize}[nosep,leftmargin=*,itemsep=1pt,font=\scriptsize]
\item Toxiques : alcool, drogues, métaux lourds
\item Pathologies : épilepsie, SLA, myasthénie
\item Traumatismes (crâniens, rachidiens)
\item Hygiène : sommeil, nutrition, prévention
\end{itemize}
\end{tcolorbox}
\end{tcbraster}

\legende{Carte mentale récapitulative de la leçon sur le fonctionnement des neurones (Terminale D). Les six branches correspondent aux six compétences clés : potentiel de repos, potentiel d'action, naissance du message nerveux, propagation, relation intensité-durée et préservation des neurones.}
\end{popfigure}

\begin{aretenir}[Définitions clés à connaître par cœur]
\begin{itemize}
    \item \cle{Potentiel de repos ($V_r$)} : différence de potentiel transmembranaire au repos ($\approx -70$ mV, intérieur négatif par rapport à l'extérieur), due à la pompe Na$^+$/K$^+$ ATPase (3 Na$^+$ sortent pour 2 K$^+$ qui entrent) et à la perméabilité sélective de la membrane au repos (canaux K$^+$ fuyards, ouverts en permanence).
    \item \cle{Potentiel d'action (PA)} : inversion brève et propagée de la polarité membranaire (dépolarisation $\to$ pic $\approx +30$ mV $\to$ repolarisation $\to$ hyperpolarisation), phénomène obéissant à la loi du \emph{tout ou rien} : son amplitude est constante quelle que soit l'intensité du stimulus supraliminaire.
    \item \cle{Seuil d'excitation} : valeur critique de dépolarisation ($\approx -55$ mV) déclenchant l'ouverture massive des canaux Na$^+$ voltage-dépendants.
    \item \cle{Potentiel générateur (PG)} : variation graduée et locale (non propagée) du potentiel de membrane au niveau du récepteur sensoriel, dont l'amplitude croît avec l'intensité du stimulus.
    \item \cle{Période réfractaire absolue (PRA)} : temps pendant lequel un second PA est impossible, quelle que soit l'intensité du stimulus (canaux Na$^+$ inactivés). Elle garantit le sens unique de propagation.
    \item \cle{Période réfractaire relative (PRR)} : temps pendant lequel un PA est possible mais nécessite un stimulus plus intense que le seuil (canaux K$^+$ encore ouverts, hyperpolarisation).
    \item \cle{Vitesse de propagation} : $v = \dfrac{d}{\Delta t}$ (en m/s). Elle dépend du diamètre de l'axone, de la température et de la myélinisation (propagation saltatoire de nœud de Ranvier en nœud de Ranvier).
\end{itemize}
\end{aretenir}

\begin{aretenir}[Mécanismes ioniques et formules indispensables]
\begin{center}
\begin{tabularx}{\linewidth}{|l|X|X|}
\hline
\rowcolor{popblueL}
\textbf{Phénomène} & \textbf{Mécanisme ionique principal} & \textbf{Équation / Relation clé} \\
\hline
Potentiel de repos & Pompe Na$^+$/K$^+$ ATPase + canaux K$^+$ fuyards & Équation de Nernst : $E_K = \dfrac{RT}{F}\ln\dfrac{[K^+]_o}{[K^+]_i}$ \\
\hline
Montée du PA (dépolarisation) & Ouverture des canaux Na$^+$ voltage-dépendants $\Rightarrow$ entrée massive de Na$^+$ & $V_m$ tend vers $E_{Na}$ ($\approx +60$ mV) \\
\hline
Descente du PA (repolarisation) & Inactivation des canaux Na$^+$ + ouverture des canaux K$^+$ voltage-dépendants $\Rightarrow$ sortie de K$^+$ & $V_m$ tend vers $E_K$ \\
\hline
Rhéobase / Chronaxie & Intensité minimale d'excitation (durée infinie) / durée minimale d'excitation pour $I = 2\,I_{rh}$ & $I = I_{rh}\left(1 + \dfrac{t_{ch}}{t}\right)$ \\
\hline
Vitesse de propagation (fibre myélinisée) & Propagation saltatoire (saut d'un nœud de Ranvier au suivant) & $v$ croît avec le diamètre, la myélinisation et la température \\
\hline
\end{tabularx}
\end{center}
\end{aretenir}

\begin{attention}[Précisions scientifiques à ne pas négliger]
\begin{itemize}
    \item Le pic du PA ($\approx +30$ mV) n'atteint jamais $E_{Na}$ ($\approx +60$ mV) : les canaux Na$^+$ s'inactivent avant l'équilibre et les canaux K$^+$ commencent à s'ouvrir. Dire que $V_m$ \emph{tend vers} $E_{Na}$ est correct ; dire qu'il l'\emph{atteint} est faux.
    \item La pompe Na$^+$/K$^+$ ne crée pas directement l'essentiel du potentiel de repos : elle \emph{maintient les gradients de concentration} ; c'est la fuite de K$^+$ à travers les canaux fuyards qui établit l'essentiel de $V_r$.
    \item Le potentiel générateur est \emph{gradué} (amplitude proportionnelle au stimulus) et \emph{non propagé} ; le potentiel d'action est \emph{tout ou rien} et \emph{propagé sans décrément}. Ne jamais confondre les deux dans une copie.
    \item La chronaxie se lit à $I = 2\,I_{rh}$, pas à $I = I_{rh}$ : c'est une erreur fréquente au Bac.
    \item La période réfractaire absolue limite la fréquence maximale des PA (de l'ordre de 500 à 1000 PA/s selon les fibres) : c'est elle qui rend possible le codage en fréquence.
\end{itemize}
\end{attention}

\begin{methode}[Méthodes express (Savoir-faire Bac)]
\begin{enumerate}
    \item \textbf{Montage potentiel de repos :} deux électrodes (une microélectrode intracellulaire + une électrode de référence extracellulaire) reliées à un oscilloscope ou un voltmètre. Interprétation : déflexion stable $\approx -70$ mV $\Rightarrow$ rôle de la pompe Na$^+$/K$^+$ et de la perméabilité au K$^+$.
    \item \textbf{Montage potentiel d'action :} électrodes de stimulation + électrodes d'enregistrement. Description de la courbe : 4 phases (latence, dépolarisation, repolarisation, hyperpolarisation) avant retour au potentiel de repos.
    \item \textbf{Courbes de perméabilités $P_{Na}$ et $P_K$ :} le pic précoce de $P_{Na}$ correspond à la montée du PA ; le pic retardé et plus durable de $P_K$ correspond à la descente et à l'hyperpolarisation.
    \item \textbf{Récepteur sensoriel :} stimulus $\to$ potentiel générateur (gradué, local) $\to$ si PG $\ge$ seuil au niveau de la zone d'initiation $\to$ PA dont la \emph{fréquence} croît avec l'intensité du stimulus (codage en fréquence).
    \item \textbf{Courbe intensité-durée :} tracer $I = f(t)$. Lire $I_{rh}$ (asymptote horizontale) et $t_{ch}$ (abscisse du point d'ordonnée $2\,I_{rh}$).
    \item \textbf{Période réfractaire :} double stimulation (S1, S2) avec intervalle $\Delta t$ variable. PA2 absent quel que soit $I_2$ $\Rightarrow$ PRA ; PA2 présent seulement si $I_2$ supérieur au seuil $\Rightarrow$ PRR.
    \item \textbf{Vitesse de propagation :} $v = \dfrac{\text{distance entre les deux points d'enregistrement}}{\text{délai entre les deux PA}}$, avec les unités converties (m et s).
\end{enumerate}
\end{methode}

\begin{exoresolu}[Calcul d'une vitesse de propagation]
Sur un nerf de grenouille, deux électrodes d'enregistrement $R_1$ et $R_2$ sont placées à $d = 4{,}0$ cm l'une de l'autre le long d'une fibre. Le potentiel d'action apparaît en $R_2$ avec un retard $\Delta t = 2{,}0$ ms par rapport à $R_1$. Calculer la vitesse de propagation du message nerveux dans cette fibre.
\tcblower
\textbf{Solution.} On applique la relation $v = \dfrac{d}{\Delta t}$ en convertissant les unités :
\begin{align*}
d &= 4{,}0~\text{cm} = 4{,}0 \times 10^{-2}~\text{m} \\
\Delta t &= 2{,}0~\text{ms} = 2{,}0 \times 10^{-3}~\text{s} \\
v &= \frac{4{,}0 \times 10^{-2}}{2{,}0 \times 10^{-3}} = 20~\text{m/s}
\end{align*}
La vitesse de propagation est $v = 20$ m/s, valeur typique d'une fibre myélinisée de diamètre moyen chez un amphibien (la température, plus basse que chez les mammifères, limite la vitesse).
\end{exoresolu}

\begin{aretenir}[Checklist avant le Bac]
\begin{itemize}
    \item $\square$ Je sais schématiser et légender le montage de mise en évidence du potentiel de repos (électrodes, oscilloscope, signe de la tension).
    \item $\square$ Je sais décrire les 4 phases du potentiel d'action et lier chaque phase à la variation de perméabilité ($P_{Na}$, $P_K$).
    \item $\square$ J'explique le rôle de la pompe Na$^+$/K$^+$ (maintien des gradients) et des canaux fuyards K$^+$ (établissement de $V_r$).
    \item $\square$ Je distingue potentiel générateur (gradué, local) et potentiel d'action (tout ou rien, propagé) au niveau du récepteur.
    \item $\square$ Je sais lire une courbe intensité-durée : identifier $I_{rh}$ (rhéobase) et $t_{ch}$ (chronaxie) et appliquer la formule $I = I_{rh}\left(1 + \dfrac{t_{ch}}{t}\right)$.
    \item $\square$ J'interprète un tracé de période réfractaire (absolue vs relative) et j'en déduis le sens unique de propagation.
    \item $\square$ Je calcule la vitesse de propagation $v = d/\Delta t$ à partir d'un tracé à double enregistrement, en convertissant correctement les unités.
    \item $\square$ Je cite 3 facteurs augmentant la vitesse : grand diamètre de l'axone, myélinisation, température élevée.
    \item $\square$ Je donne 3 exemples de perturbations (toxiques, traumatismes, pathologies) et 3 mesures de prévention (hygiène de vie, protection, dépistage et suivi médical).
    \item $\square$ J'utilise le vocabulaire précis : \emph{dépolarisation} (et non « augmentation »), \emph{repolarisation}, \emph{hyperpolarisation}, \emph{inactivation} (canaux Na$^+$), \emph{propagation saltatoire}.
    \item $\square$ Je relie la structure (nœuds de Ranvier, gaine de myéline, canaux voltage-dépendants concentrés aux nœuds) à la fonction (vitesse accrue, économie d'énergie).
\end{itemize}
\end{aretenir}

\begin{savaistu}[Préserver ses neurones au quotidien]
Au Cameroun, les traumatismes crâniens et rachidiens dus aux accidents de moto-taxi sont une cause majeure de lésions nerveuses : le port du casque est une mesure de prévention directe du fonctionnement des neurones. De même, la consommation d'alcool et de drogues perturbe la membrane neuronale et les canaux ioniques, tandis qu'une alimentation équilibrée (apports en vitamines du groupe B, notamment B1 et B12, indispensables à la myélinisation) et un sommeil suffisant favorisent le maintien des gaines de myéline et donc la vitesse de propagation des messages nerveux.
\end{savaistu}

\findecours

\end{document}
